% Generated mechanically from frozen input inputs/p03/ja/tex/sites-cohomology.tex.
% Do not edit this generated file; edit the bound locale source instead.

% OK, start here.
%

\setcounter{chapter}{20}
\chapter{サイト上のコホモロジー}



\phantomsection
\label{sites-cohomology-section-phantom}


\section{はじめに}
\label{sites-cohomology-section-introduction}

\noindent
この文書では層のコホモロジーに関するいくつかの話題を扱う。
サイト上の層について何が起こるかを調べるが、
多くの場合、位相空間の場合の議論、
構成、および証明をそのまま
この場合にも繰り返すことになる。
基本的な参考文献は \cite{SGA4}, \cite{Godement}, \cite{Iversen} である。




\section{層のコホモロジー}
\label{sites-cohomology-section-cohomology-sheaves}

\noindent
サイト $\mathcal{C}$ をとる（
Sites, Definition \ref{sites-definition-site} 参照）。
$\mathcal{C}$ 上のアーベル層を $\mathcal{F}$ とする。
$\mathcal{C}$ 上のアーベル層の圏は十分な入射対象をもつことが分かっている（
Injectives, Theorem \ref{injectives-theorem-sheaves-injectives} 参照）。
したがって入射分解
$\mathcal{F}[0] \to \mathcal{I}^\bullet$.
を選ぶことができる。サイト $\mathcal{C}$ の任意の対象 $U$ に対して
\begin{equation}
\label{sites-cohomology-equation-cohomology-object-site}
H^i(U, \mathcal{F}) = H^i(\Gamma(U, \mathcal{I}^\bullet))
\end{equation}
を「対象 $U$ 上のアーベル層 $\mathcal{F}$ の {\it $i$ 次コホモロジー群}」
と定める。別の言い方をすれば、これは関手
$\mathcal{F} \mapsto \mathcal{F}(U)$ の右導来関手である。
関手族 $H^i(U, -)$ は普遍 $\delta$-関手
$\textit{Ab}(\mathcal{C}) \to \textit{Ab}$.

\medskip\noindent
サイト $\mathcal{C}$ が終対象をもたないこともある。この
場合、$\mathcal{C}$ 上の集合の前層 $\mathcal{F}$ の {\it 大域切断} を
次の集合として定める。
\begin{equation}
\label{sites-cohomology-equation-global-sections}
\Gamma(\mathcal{C}, \mathcal{F}) =
\Mor_{\textit{PSh}(\mathcal{C})}(e, \mathcal{F})
\end{equation}
ここで $e$ は $\mathcal{C}$ 上の前層の圏における終対象である。
この場合、$\mathcal{C}$ 上のアーベル層 $\mathcal{F}$ の
{\it $\mathcal{C}$ 上の $i$ 次コホモロジー群} を次で定める。
\begin{equation}
\label{sites-cohomology-equation-cohomology}
H^i(\mathcal{C}, \mathcal{F}) = H^i(\Gamma(\mathcal{C}, \mathcal{I}^\bullet))
\end{equation}
すなわち、これは大域切断関手の $i$ 次右導来関手である。
関手族 $H^i(\mathcal{C}, -)$ は普遍 $\delta$-関手
$\textit{Ab}(\mathcal{C}) \to \textit{Ab}$.

\medskip\noindent
トポスの射 $f : \Sh(\mathcal{C}) \to \Sh(\mathcal{D})$ をとる（
Sites, Definition \ref{sites-definition-topos} 参照）。
上と同じく $\mathcal{F}[0] \to \mathcal{I}^\bullet$ を用いて
次を定める。
\begin{equation}
\label{sites-cohomology-equation-higher-direct-image}
R^if_*\mathcal{F} = H^i(f_*\mathcal{I}^\bullet)
\end{equation}
を $\mathcal{F}$ の {\it $i$ 次高次順像} と呼ぶ。
これは $f_*$ の右導来関手である。
関手族 $R^if_*$ は
$\textit{Ab}(\mathcal{C}) \to \textit{Ab}(\mathcal{D})$ からの
普遍 $\delta$-関手をなす。

\medskip\noindent
環付きサイト $(\mathcal{C}, \mathcal{O})$ をとる（
Modules on Sites, Definition \ref{sites-modules-definition-ringed-site} 参照）。
$\mathcal{O}$-加群を $\mathcal{F}$ とする。
$\mathcal{O}$-加群の圏は十分な入射対象をもつことが分かっている（
Injectives, Theorem \ref{injectives-theorem-sheaves-modules-injectives} 参照）。
したがって入射分解
$\mathcal{F}[0] \to \mathcal{I}^\bullet$.
を選べる。サイト $\mathcal{C}$ の任意の対象 $U$ に対して
\begin{equation}
\label{sites-cohomology-equation-cohomology-object-site-modules}
H^i(U, \mathcal{F}) = H^i(\Gamma(U, \mathcal{I}^\bullet))
\end{equation}
を $\mathcal{F}$ の {\it $U$ 上の $i$ 次コホモロジー群} と定める。
関手族 $H^i(U, -)$ は普遍 $\delta$-関手
$\textit{Mod}(\mathcal{O}) \to \text{Mod}_{\mathcal{O}(U)}$. 同様に
\begin{equation}
\label{sites-cohomology-equation-cohomology-modules}
H^i(\mathcal{C}, \mathcal{F}) = H^i(\Gamma(\mathcal{C}, \mathcal{I}^\bullet))
\end{equation}
これは $\mathcal{F}$ の {\it $\mathcal{C}$ 上の $i$ 次コホモロジー群} である。
関手族 $H^i(\mathcal{C}, -)$ は普遍
$\delta$-関手
$\textit{Mod}(\mathcal{C}) \to \text{Mod}_{\Gamma(\mathcal{C}, \mathcal{O})}$.

\medskip\noindent
環付きトポスの射
$f : (\Sh(\mathcal{C}), \mathcal{O}) \to (\Sh(\mathcal{D}), \mathcal{O}')$
をとる（Modules on Sites, Definition \ref{sites-modules-definition-ringed-topos} 参照）。
上と同じ $\mathcal{F}[0] \to \mathcal{I}^\bullet$ に対して
次を定める。
\begin{equation}
\label{sites-cohomology-equation-higher-direct-image-modules}
R^if_*\mathcal{F} = H^i(f_*\mathcal{I}^\bullet)
\end{equation}
を $\mathcal{F}$ の {\it $i$ 次高次順像} と呼ぶ。
これは $f_*$ の右導来関手である。
関手族 $R^if_*$ は
$\textit{Mod}(\mathcal{O}) \to \textit{Mod}(\mathcal{O}')$ からの
普遍 $\delta$-関手をなす。








\section{導来関手}
\label{sites-cohomology-section-derived-functors}

\noindent
分解関手を用いて右導来関手を構成する方法を簡単に説明する。
すなわち $(\mathcal{C}, \mathcal{O})$ を環付きサイトとする。
この章では
$$
K(\mathcal{O}) = K(\textit{Mod}(\mathcal{O}))
\quad
\text{および}
\quad
D(\mathcal{O}) = D(\textit{Mod}(\mathcal{O}))
$$
また、次で導入された三角圏について、その有界版も同様に書く。
Derived Categories, Definition \ref{derived-definition-complexes-notation} および
Definition \ref{derived-definition-unbounded-derived-category}.
Derived Categories, Remark \ref{derived-remark-big-abelian-category} により
分解関手
$$
j = j_{(\mathcal{C}, \mathcal{O})} :
K^{+}(\textit{Mod}(\mathcal{O}))
\longrightarrow
K^{+}(\mathcal{I})
$$
ここで $\mathcal{I}$ は入射 $\mathcal{O}$-加群からなる
$\textit{Mod}(\mathcal{O})$ の充満加法部分圏である。
任意のアーベル圏 $\mathcal{B}$ への左完全関手
$F : \textit{Mod}(\mathcal{O}) \to \mathcal{B}$ に対し、
$RF$ で次の右導来関手を表す。
これは今述べた分解関手 $j$ を用いて
Derived Categories, Section \ref{derived-section-right-derived-functor} に
構成される。
\begin{equation}
\label{sites-cohomology-equation-RF}
RF = F \circ j' : D^{+}(\mathcal{O}) \longrightarrow D^{+}(\mathcal{B})
\end{equation}
記法については
Derived Categories, Lemma \ref{derived-lemma-right-derived-functor} を参照。
$RF$ は状況に応じて
$\textit{Mod}(\mathcal{O})$, $\text{Comp}^{+}(\textit{Mod}(\mathcal{O}))$, or
$K^{+}(\mathcal{O})$ 上で定義されていると考えてよい。
Derived Categories, Definition \ref{derived-definition-higher-derived-functors} により
$i$ 次右導来関手
\begin{equation}
\label{sites-cohomology-equation-RFi}
R^iF = H^i \circ RF : \textit{Mod}(\mathcal{O}) \longrightarrow \mathcal{B}
\end{equation}
を得る。したがって $R^0F = F$ であり、
$\{R^iF, \delta\}_{i \geq 0}$ は普遍 $\delta$-関手である。
Derived Categories, Lemma \ref{derived-lemma-higher-derived-functors} 参照。

\medskip\noindent
この構成の特殊な場合を二つ挙げる。環 $R$ に対して
$K(R) = K(\text{Mod}_R)$ および $D(R) = D(\text{Mod}_R)$ と書き、
有界版についても同様に書く。$\mathcal{C}$ の任意の対象 $U$ に対して左完全関手
$
\Gamma(U, -) :
\textit{Mod}(\mathcal{O})
\longrightarrow
\text{Mod}_{\mathcal{O}(U)}
$
が得られ、上の議論により
$$
R\Gamma(U, -) :
D^{+}(\mathcal{O})
\longrightarrow
D^{+}(\mathcal{O}(U))
$$
が定まる。$H^i(U, -) = R^i\Gamma(U, -)$ は
上の (\ref{sites-cohomology-equation-cohomology-object-site-modules}) と整合することに注意する。
同様に
$$
R\Gamma(\mathcal{C}, -) :
D^{+}(\mathcal{O})
\longrightarrow
D^{+}(\Gamma(\mathcal{C}, \mathcal{O}))
$$
は (\ref{sites-cohomology-equation-cohomology-modules}) と整合する。もし
$f : (\Sh(\mathcal{C}), \mathcal{O}) \to (\Sh(\mathcal{D}), \mathcal{O}')$
が環付きトポスの射なら、左完全関手
$f_* : \textit{Mod}(\mathcal{O}) \to \textit{Mod}(\mathcal{O}')$
が得られ、これが {\it 導来順像}
$$
Rf_* : D^{+}(\mathcal{O}) \to D^+(\mathcal{O}')
$$
$Rf_*\mathcal{F}^\bullet$ の $i$ 次コホモロジー層を
$R^if_*\mathcal{F}^\bullet$ と書き、(\ref{sites-cohomology-equation-higher-direct-image-modules}) に従って
$i$ 次の {\it 高次順像} と呼ぶ。
上に表示した関手は導来圏の完全関手である。







\section{第一コホモロジーとトーサー}
\label{sites-cohomology-section-h1-torsors}

\begin{definition}
\label{sites-cohomology-definition-torsor}
サイト $\mathcal{C}$ をとる。
$\mathcal{C}$ 上の（可換とは限らない）群の層 $\mathcal{G}$ をとる。
{\it 擬トーサー}、より正確には
{\it 擬 $\mathcal{G}$-トーサー} とは、$\mathcal{C}$ 上の集合の層
$\mathcal{F}$ に作用
$\mathcal{G} \times \mathcal{F} \to \mathcal{F}$ が備わり、次を満たすものをいう。
\begin{enumerate}
\item $\mathcal{F}(U)$ が空でないときはいつでも作用
$\mathcal{G}(U) \times \mathcal{F}(U) \to \mathcal{F}(U)$
は単純推移的である。
\end{enumerate}
{\it 擬 $\mathcal{G}$-トーサーの射}
$\mathcal{F} \to \mathcal{F}'$
とは、$\mathcal{G}$-作用と両立する集合の層の射である。
{\it トーサー}、より正確には
{\it $\mathcal{G}$-トーサー} とは、さらに次を満たす擬 $\mathcal{G}$-トーサーである。
\begin{enumerate}
\item[(2)] 任意の $U \in \Ob(\mathcal{C})$ に対し、
$\mathcal{F}(U_i)$ がすべての $i \in I$ で空でなくなるような
$U$ の被覆 $\{U_i \to U\}_{i \in I}$ が存在する。
\end{enumerate}
{\it $\mathcal{G}$-トーサーの射} は擬 $\mathcal{G}$-トーサーの射である。
{\it 自明な $\mathcal{G}$-トーサー} とは、明らかな左
$\mathcal{G}$-作用を備えた層 $\mathcal{G}$ である。
\end{definition}

\noindent
トーサーの射は自動的に同型であることは明らかである。

\begin{lemma}
\label{sites-cohomology-lemma-trivial-torsor}
サイト $\mathcal{C}$ をとり、$\mathcal{C}$ 上の（可換とは限らない）
群の層 $\mathcal{G}$ をとる。
$\mathcal{G}$-トーサー $\mathcal{F}$ が自明であることと
$\Gamma(\mathcal{C}, \mathcal{F}) \not = \emptyset$ は同値である。
\end{lemma}

\begin{proof}
省略する。
\end{proof}

\begin{lemma}
\label{sites-cohomology-lemma-torsors-h1}
サイト $\mathcal{C}$ をとる。
$\mathcal{C}$ 上のアーベル層を $\mathcal{H}$ とする。
$\mathcal{H}$-トーサーの同型類の集合と
$H^1(\mathcal{C}, \mathcal{H})$ の間には標準的な全単射が存在する。
\end{lemma}

\begin{proof}
$\mathcal{H}$-トーサー $\mathcal{F}$ をとる。
$\mathcal{F}$ 上の自由アーベル層 $\mathbf{Z}[\mathcal{F}]$ を考える。
これは各 $U \in \Ob(\mathcal{C})$ に有限形式和
$\sum n_i[s_i]$（$n_i \in \mathbf{Z}$、$s_i \in \mathcal{F}(U)$）
を対応させる規則の層化である。自然な射
$$
\sigma : \mathbf{Z}[\mathcal{F}] \longrightarrow \underline{\mathbf{Z}}
$$
は局所切断 $\sum n_i[s_i]$ を $\sum n_i$ に対応させる。
$\sigma$ の核は $[s]-[s']$ の形の切断で生成される。
標準的な射
$a : \Ker(\sigma) \to \mathcal{H}$
$[s]-[s'] \mapsto h$ を写すものがある。ここで $h$ は
$h\cdot s'=s$ を満たす $\mathcal{H}$ の局所切断である。
押し出し図式を考える。
$$
\xymatrix{
0 \ar[r] &
\Ker(\sigma) \ar[r] \ar[d]^a &
\mathbf{Z}[\mathcal{F}] \ar[r] \ar[d] &
\underline{\mathbf{Z}} \ar[r] \ar[d] &
0 \\
0 \ar[r] &
\mathcal{H} \ar[r] &
\mathcal{E} \ar[r] &
\underline{\mathbf{Z}} \ar[r] &
0
}
$$
ここで $\mathcal{E}$ は押し出しで得られる拡張である。
下段の短完全列に付随する長完全コホモロジー列から、
$1 \in H^0(\mathcal{C}, \underline{\mathbf{Z}})$ に境界作用素を適用して
$\xi = \xi_\mathcal{F} \in H^1(\mathcal{C}, \mathcal{H})$
を得る。

\medskip\noindent
逆に、$\xi \in H^1(\mathcal{C}, \mathcal{H})$ が与えられたとき、
次のように $\xi$ にトーサーを対応させる。
$\mathcal{H}$ を入射アーベル層 $\mathcal{I}$ に埋め込む射
$\mathcal{H} \to \mathcal{I}$ を選び、
$\mathcal{Q} = \mathcal{I}/\mathcal{H}$ とおくと、短完全列
$$
\xymatrix{
0 \ar[r] &
\mathcal{H} \ar[r] &
\mathcal{I} \ar[r] &
\mathcal{Q} \ar[r] &
0
}
$$
$\xi$ は大域切断
$q \in H^0(\mathcal{C}, \mathcal{Q})$ の像である。
これは $H^1(\mathcal{C}, \mathcal{I})=0$ による（
Derived Categories, Lemma \ref{derived-lemma-higher-derived-functors} 参照）。
$\mathcal{Q}$ の中で $q$ に写る切断全体からなる
$\mathcal{I}$ の（集合の）部分層を $\mathcal{F} \subset \mathcal{I}$ とする。
$\mathcal{F}$ が $\mathcal{H}$-トーサーであることは容易に確かめられる。

\medskip\noindent
上の二つの構成が互いに逆であることの確認は省略する。
\end{proof}






\section{第一コホモロジーと拡張}
\label{sites-cohomology-section-h1-extensions}

\begin{lemma}
\label{sites-cohomology-lemma-h1-extensions}
環付きサイト $(\mathcal{C}, \mathcal{O})$ をとる。
$\mathcal{C}$ 上の $\mathcal{O}$-加群の層を $\mathcal{F}$ とする。
標準的な全単射
$$
\Ext^1_{\textit{Mod}(\mathcal{O})}(\mathcal{O}, \mathcal{F})
\longrightarrow
H^1(\mathcal{C}, \mathcal{F})
$$
は拡張
$$
0 \to \mathcal{F} \to \mathcal{E} \to \mathcal{O} \to 0
$$
を $1 \in \Gamma(\mathcal{C}, \mathcal{O})$ の
$H^1(\mathcal{C}, \mathcal{F})$ における像へ対応させる。
\end{lemma}

\begin{proof}
補題の射の逆を構成する。
$\xi \in H^1(\mathcal{C}, \mathcal{F})$ とする。
$\mathcal{F} \subset \mathcal{I}$ であって
$\mathcal{I}$ が $\textit{Mod}(\mathcal{O})$ で入射となるものを選ぶ。
$\mathcal{Q} = \mathcal{I}/\mathcal{F}$ とおく。
コホモロジーの長完全列から、$\xi$ は切断
$\tilde \xi \in \Gamma(\mathcal{C}, \mathcal{Q}) =
\Hom_\mathcal{O}(\mathcal{O}, \mathcal{Q})$.
ここで引き戻しを形成する。
$$
\xymatrix{
0 \ar[r] &
\mathcal{F} \ar[r] \ar@{=}[d] &
\mathcal{E} \ar[r] \ar[d] &
\mathcal{O} \ar[r] \ar[d]^{\tilde \xi} &
0 \\
0 \ar[r] &
\mathcal{F} \ar[r] &
\mathcal{I} \ar[r] &
\mathcal{Q} \ar[r] &
0
}
$$
Homology, Section \ref{homology-section-extensions} 参照。
\end{proof}

\noindent
次の補題は、より一般的な
Lemma \ref{sites-cohomology-lemma-cohomology-modules-abelian-agree} によって置き換えられる。

\begin{lemma}
\label{sites-cohomology-lemma-h1-mod-ab-agree}
環付きサイト $(\mathcal{C}, \mathcal{O})$ をとる。
$\mathcal{C}$ 上の $\mathcal{O}$-加群の層を $\mathcal{F}$ とする。
$\mathcal{F}_{ab}$ で台となるアーベル群の層を表す。
すると関手的な同型
$$
H^1(\mathcal{C}, \mathcal{F}_{ab})
=
H^1(\mathcal{C}, \mathcal{F})
$$
が存在する。左辺のコホモロジーは
$\textit{Ab}(\mathcal{C})$ で計算し、右辺は
$\textit{Mod}(\mathcal{O})$ で計算する。
\end{lemma}

\begin{proof}
定数層 $\mathbf{Z}$ を $\underline{\mathbf{Z}}$ と書く。
$\textit{Ab}(\mathcal{C}) = \textit{Mod}(\underline{\mathbf{Z}})$ なので、
Lemma \ref{sites-cohomology-lemma-h1-extensions} を二度適用できるので、
示すべきことは
$$
\Ext^1_{\textit{Mod}(\mathcal{O})}(\mathcal{O}, \mathcal{F})
=
\Ext^1_{\textit{Mod}(\underline{\mathbf{Z}})}(
\underline{\mathbf{Z}}, \mathcal{F}_{ab}).
$$
まず
$0 \to \mathcal{F} \to \mathcal{E} \to \mathcal{O} \to 0$
が $\textit{Mod}(\mathcal{O})$ における拡張であるとする。
アーベル層の明らかな射
$1 : \underline{\mathbf{Z}} \to \mathcal{O}$
を用いて引き戻しを行うと、次の拡張 $\mathcal{E}_{ab}$ を得る。
$$
\xymatrix{
0 \ar[r] &
\mathcal{F}_{ab} \ar[r] \ar@{=}[d] &
\mathcal{E}_{ab} \ar[r] \ar[d] &
\underline{\mathbf{Z}} \ar[r] \ar[d]^{1} &
0 \\
0 \ar[r] &
\mathcal{F} \ar[r] &
\mathcal{E} \ar[r] &
\mathcal{O} \ar[r] &
0
}
$$
逆向きは少しだけ手間がかかる。
$0 \to \mathcal{F}_{ab} \to \mathcal{E}_{ab} \to \underline{\mathbf{Z}} \to 0$
が $\textit{Mod}(\underline{\mathbf{Z}})$ における拡張であるとする。
$\underline{\mathbf{Z}}$ は平坦な $\underline{\mathbf{Z}}$-加群なので、
次の列
$$
0 \to \mathcal{F}_{ab} \otimes_{\underline{\mathbf{Z}}} \mathcal{O}
\to \mathcal{E}_{ab} \otimes_{\underline{\mathbf{Z}}} \mathcal{O}
\to \underline{\mathbf{Z}} \otimes_{\underline{\mathbf{Z}}} \mathcal{O}
\to 0
$$
は完全である（
Modules on Sites, Lemma \ref{sites-modules-lemma-flat-tor-zero}.
もちろん
$\underline{\mathbf{Z}} \otimes_{\underline{\mathbf{Z}}} \mathcal{O}
= \mathcal{O}$.
したがって、（$\mathcal{O}$-線型）乗法写像
$\mu : \mathcal{F} \otimes_{\underline{\mathbf{Z}}} \mathcal{O}
\to \mathcal{F}$ による押し出しを形成すれば、次のように
$\mathcal{F}$ による $\mathcal{O}$ の拡張を得る。
$$
\xymatrix{
0 \ar[r] &
\mathcal{F}_{ab} \otimes_{\underline{\mathbf{Z}}} \mathcal{O}
\ar[r] \ar[d]^\mu &
\mathcal{E}_{ab} \otimes_{\underline{\mathbf{Z}}} \mathcal{O}
\ar[r] \ar[d] &
\mathcal{O} \ar[r] \ar@{=}[d] &
0 \\
0 \ar[r] &
\mathcal{F} \ar[r] &
\mathcal{E} \ar[r] &
\mathcal{O} \ar[r] &
0
}
$$
これが求める拡張である。これらの構成が互いに逆であることの確認は省略する。
\end{proof}





\section{第一コホモロジーと可逆層}
\label{sites-cohomology-section-invertible-sheaves}

\noindent
環付きサイトのピカール群は
Modules on Sites, Section \ref{sites-modules-section-invertible} で定義する。

\begin{lemma}
\label{sites-cohomology-lemma-h1-invertible}
局所環付きサイト $(\mathcal{C}, \mathcal{O})$ をとる。
標準的な同型
$$
H^1(\mathcal{C}, \mathcal{O}^*) = \Pic(\mathcal{O}).
$$
がアーベル群の同型として存在する。
\end{lemma}

\begin{proof}
$\mathcal{O}$-可逆加群を $\mathcal{L}$ とする。
次の規則で定まる前層 $\mathcal{L}^*$ を考える。
$$
U \longmapsto \{s \in \mathcal{L}(U)
\text{ かつ } \mathcal{O}_U \xrightarrow{s \cdot -} \mathcal{L}_U
\text{ は同型である}\}
$$
この前層は層条件を満たす。さらに
$f \in \mathcal{O}^*(U)$、$s \in \mathcal{L}^*(U)$ なら明らかに
$fs \in \mathcal{L}^*(U)$ である。同様に $s,s' \in \mathcal{L}^*(U)$
なら $fs=s'$ を満たす一意な $f \in \mathcal{O}^*(U)$ が存在する。
また、Modules on Sites, Lemma
\ref{sites-modules-lemma-invertible-is-locally-free-rank-1} により、
層 $\mathcal{L}^*$ は局所的に切断をもつ。
したがって $\mathcal{L}^*$ は $\mathcal{O}^*$-トーサーである。
ゆえに次の射を得る。
$$
\begin{matrix}
\text{可逆層の集合 }(\mathcal{C}, \mathcal{O})\text{ 上}\\
\text{同型まで}
\end{matrix}
\longrightarrow
\begin{matrix}
\mathcal{O}^*\text{-トーサーの集合}\\
\text{同型まで}
\end{matrix}
$$
これがアーベル群の準同型であることの確認は省略する。
Lemma \ref{sites-cohomology-lemma-torsors-h1} により、
右辺は標準的に $H^1(\mathcal{C}, \mathcal{O}^*)$ と全単射である。
したがって、この射が単射かつ全射であることを示せばよい。

\medskip\noindent
単射性。トーサー $\mathcal{L}^*$ が自明なら、
Lemma \ref{sites-cohomology-lemma-trivial-torsor}
により $\mathcal{L}^*$ は大域切断をもつ。
これはまさに $\mathcal{L} \cong \mathcal{O}$ が
$\Pic(\mathcal{O})$ の単位元であることを意味する。

\medskip\noindent
全射性。$\mathcal{O}^*$-トーサー $\mathcal{F}$ をとる。
集合の前層
$$
\mathcal{L}_1 : U \longmapsto
(\mathcal{F}(U) \times \mathcal{O}(U))/\mathcal{O}^*(U)
$$
ここで $f \in \mathcal{O}^*(U)$ の $(s,g)$ への作用は
$(fs,f^{-1}g)$ である。$\mathcal{L}_1$ は次で
$\mathcal{O}$-加群の前層となる。
$(s,g)+(s',g')=(s,g+(s'/s)g')$ とおく。ここで $s'/s$ は
$fs=s'$ を満たす $\mathcal{O}^*$ の局所切断 $f$ であり、
$h(s,g)=(s,hg)$（$h$ は $\mathcal{O}$ の局所切断）とする。
層化 $\mathcal{L}=\mathcal{L}_1^\#$ が可逆 $\mathcal{O}$-加群で、
それに付随する $\mathcal{O}^*$-トーサー $\mathcal{L}^*$ が
$\mathcal{F}$ と同型であることの確認は省略する。
\end{proof}









\section{コホモロジーの局所性}
\label{sites-cohomology-section-locality}

\noindent
次の補題は、サイトの対象上の層 $\mathcal{F}$ のコホモロジーを
定義する際に曖昧さがないことを述べる。

\begin{lemma}
\label{sites-cohomology-lemma-cohomology-of-open}
環付きサイト $(\mathcal{C}, \mathcal{O})$ をとる。
$U$ を $\mathcal{C}$ の対象とする。
\begin{enumerate}
\item $\mathcal{I}$ が入射 $\mathcal{O}$-加群なら、
$\mathcal{I}|_U$ は入射 $\mathcal{O}_U$-加群である。
\item 任意の $\mathcal{O}$-加群の層 $\mathcal{F}$ に対して
$H^p(U, \mathcal{F}) = H^p(\mathcal{C}/U, \mathcal{F}|_U)$.
\end{enumerate}
\end{lemma}

\begin{proof}
制限関手 $j_U^{-1}$ は零による延長関手 $j_{U!}$ の右随伴であることを思い出す（
Modules on Sites, Section
\ref{sites-modules-section-localize} 参照）。
さらに $j_{U!}$ は完全である。したがって (1) は
Homology, Lemma \ref{homology-lemma-adjoint-preserve-injectives} による。

\medskip\noindent
定義により $H^p(U, \mathcal{F}) = H^p(\mathcal{I}^\bullet(U))$ である。
ここで $\mathcal{F} \to \mathcal{I}^\bullet$ は
$\textit{Mod}(\mathcal{O})$ における入射分解である。
上の結果から $\mathcal{F}|_U \to \mathcal{I}^\bullet|_U$ は
$\textit{Mod}(\mathcal{O}_U)$ における入射分解である。
したがって $H^p(U, \mathcal{F}|_U)$ は
$H^p(\mathcal{I}^\bullet|_U(U))$.
もちろん $\mathcal{C}$ 上の任意の層 $\mathcal{F}$ について
$\mathcal{F}(U)=\mathcal{F}|_U(U)$ である。これで (2) の等式が従う。
\end{proof}

\noindent
次の補題は、部分宇宙を変更したときに何が起こるかを調べるため、
また小さいエタールサイトと大きいエタールサイトのコホモロジーを
比較するために用いる。

\begin{lemma}
\label{sites-cohomology-lemma-cohomology-bigger-site}
サイト $\mathcal{C},\mathcal{D}$ をとる。
関手 $u : \mathcal{C} \to \mathcal{D}$ をとる。
$u$ が
Sites, Lemma \ref{sites-lemma-bigger-site}.
に述べる仮定を満たすとする。
対応するトポスの射を
$g : \Sh(\mathcal{C}) \to \Sh(\mathcal{D})$ とする。
$\mathcal{D}$ 上の任意のアーベル層 $\mathcal{F}$ に対して同型
$$
R\Gamma(\mathcal{C}, g^{-1}\mathcal{F}) = R\Gamma(\mathcal{D}, \mathcal{F}),
$$
特に
$H^p(\mathcal{C}, g^{-1}\mathcal{F}) = H^p(\mathcal{D}, \mathcal{F})$
さらに任意の $U \in \Ob(\mathcal{C})$ に対して同型
$$
R\Gamma(U, g^{-1}\mathcal{F}) = R\Gamma(u(U), \mathcal{F}),
$$
特に
$H^p(U, g^{-1}\mathcal{F}) = H^p(u(U), \mathcal{F})$ である。
これらの同型はすべて $\mathcal{F}$ に関して関手的である。
\end{lemma}

\begin{proof}
仮定 (e) により
$\Gamma(\mathcal{C}, g^{-1}\mathcal{F}) = \Gamma(\mathcal{D}, \mathcal{F})$
が明らかなので、$g^{-1}$ が入射アーベル層を入射アーベル層へ移すことを
示せば十分である。いつものように
Homology, Lemma \ref{homology-lemma-adjoint-preserve-injectives}
を用いる。$g^{-1}$ の左随伴は、記法
Sites, Lemma \ref{sites-lemma-bigger-site}
で $g_! = f^{-1}$ であり、これは完全関手である。したがってこの補題を適用できる。
\end{proof}

\noindent
環付きサイト $(\mathcal{C}, \mathcal{O})$ をとる。
$\mathcal{O}$-加群の層を $\mathcal{F}$ とする。
$\mathcal{O}$ の射 $\varphi : U \to V$ をとる。
すると標準的な {\it 制限写像}
\begin{equation}
\label{sites-cohomology-equation-restriction-mapping}
H^n(V, \mathcal{F})
\longrightarrow
H^n(U, \mathcal{F}), \quad
\xi \longmapsto \xi|_U
\end{equation}
$\mathcal{F}$ に関して関手的である。実際、任意の入射分解
$\mathcal{F} \to \mathcal{I}^\bullet$ を選ぶ。
層 $\mathcal{I}^p$ の制限写像は複体の射
$$
\Gamma(V, \mathcal{I}^\bullet)
\longrightarrow
\Gamma(U, \mathcal{I}^\bullet)
$$
左辺は $R\Gamma(V, \mathcal{F})$ を表す複体であり、
右辺は $R\Gamma(U, \mathcal{F})$ を表す複体である。
関手 $H^n$ を適用してコホモロジー群の写像を得る。
上で示したように、この写像を表す記法として
$\xi \mapsto \xi|_U$ を用いる。
したがって規則 $U \mapsto H^n(U, \mathcal{F})$ は
$\mathcal{O}$-加群の前層である。この前層を通常
$\underline{H}^n(\mathcal{F})$ と書く。
この前層の別の解釈を Lemma \ref{sites-cohomology-lemma-include} で与える。

\medskip\noindent
次の補題は、被覆へ移ることで高次コホモロジー類を消せることを述べる。

\begin{lemma}
\label{sites-cohomology-lemma-kill-cohomology-class-on-covering}
環付きサイト $(\mathcal{C}, \mathcal{O})$ をとる。
$\mathcal{O}$-加群の層を $\mathcal{F}$ とする。
$U$ を $\mathcal{C}$ の対象とする。
$n>0$ および $\xi \in H^n(U,\mathcal{F})$ とする。
すべての $i \in I$ について $\xi|_{U_i}=0$ となる
$\mathcal{C}$ の被覆 $\{U_i \to U\}$ が存在する。
\end{lemma}

\begin{proof}
入射分解 $\mathcal{F} \to \mathcal{I}^\bullet$ をとる。すると
$$
H^n(U, \mathcal{F}) =
\frac{\Ker(\mathcal{I}^n(U) \to \mathcal{I}^{n + 1}(U))}
{\Im(\mathcal{I}^{n - 1}(U) \to \mathcal{I}^n(U))}.
$$
上の表示でコホモロジー類を表す元
$\tilde \xi \in \mathcal{I}^n(U)$ を選ぶ。
$\mathcal{I}^\bullet$ は $\mathcal{F}$ の入射分解であり $n>0$ なので、
$\mathcal{I}^\bullet$ は次数 $n$ で完全である。したがって
$\Im(\mathcal{I}^{n - 1} \to \mathcal{I}^n) =
\Ker(\mathcal{I}^n \to \mathcal{I}^{n + 1})$ as sheaves.
$\tilde \xi$ は $U$ 上の核層の切断なので、
$\tilde \xi|_{U_i}$ がある切断
$\xi_i \in \mathcal{I}^{n-1}(U_i)$ の $d$ による像となる
サイトの被覆 $\{U_i \to U\}$ が存在する。
$\xi|_{U_i}$ を $\tilde \xi|_{U_i}$ の類に対応させるという
制限の定義から結論を得る。
\end{proof}

\begin{lemma}
\label{sites-cohomology-lemma-higher-direct-images}
連続関手 $u : \mathcal{D} \to \mathcal{C}$ に対応する環付きサイトの射
$f : (\mathcal{C}, \mathcal{O}_\mathcal{C}) \to
(\mathcal{D}, \mathcal{O}_\mathcal{D})$ をとる。
$\mathcal{F} \in \Ob(\textit{Mod}(\mathcal{O}_\mathcal{C}))$ に対し、
層 $R^if_*\mathcal{F}$ は前層
$$
V \longmapsto H^i(u(V), \mathcal{F})
$$
\end{lemma}

\begin{proof}
入射分解 $\mathcal{F} \to \mathcal{I}^\bullet$ をとる。
定義により $R^if_*\mathcal{F}$ は次の複体の
$i$ 次コホモロジー層である。
$$
f_*\mathcal{I}^0 \to f_*\mathcal{I}^1 \to f_*\mathcal{I}^2 \to \ldots
$$
 $\mathcal{O}_\mathcal{D}$-加群のアーベル圏構造の定義により、
このコホモロジー層は前層
$$
V
\longmapsto
\frac{\Ker(f_*\mathcal{I}^i(V) \to f_*\mathcal{I}^{i + 1}(V))}
{\Im(f_*\mathcal{I}^{i - 1}(V) \to f_*\mathcal{I}^i(V))}
$$
であり、これは明らかに
$$
\frac{\Ker(\mathcal{I}^i(u(V)) \to \mathcal{I}^{i + 1}(u(V)))}
{\Im(\mathcal{I}^{i - 1}(u(V)) \to \mathcal{I}^i(u(V)))}
$$
に等しい。これは $H^i(u(V), \mathcal{F})$ に等しいので証明された。
\end{proof}






\section{{\v C}ech 複体と {\v C}ech コホモロジー}
\label{sites-cohomology-section-cech}

\noindent
圏 $\mathcal{C}$ をとる。固定された終域をもつ射の族
$\mathcal{U} = \{U_i \to U\}_{i \in I}$ をとる（
Sites, Definition \ref{sites-definition-family-morphisms-fixed-target} 参照）。
すべてのファイバー積 $U_{i_0} \times_U \ldots \times_U U_{i_p}$ が
$\mathcal{C}$ に存在すると仮定する。$\mathcal{C}$ 上のアーベル前層
$\mathcal{F}$ に対して
$$
\check{\mathcal{C}}^p(\mathcal{U}, \mathcal{F})
=
\prod\nolimits_{(i_0, \ldots, i_p) \in I^{p + 1}}
\mathcal{F}(U_{i_0} \times_U \ldots \times_U U_{i_p}).
$$
これはアーベル群である。
$s \in \check{\mathcal{C}}^p(\mathcal{U}, \mathcal{F})$ に対し、
$s_{i_0\ldots i_p}$ で因子
$\mathcal{F}(U_{i_0} \times_U \ldots \times_U U_{i_p})$ における値を表す。
次の射を定める。
$$
d : \check{\mathcal{C}}^p(\mathcal{U}, \mathcal{F})
\longrightarrow
\check{\mathcal{C}}^{p + 1}(\mathcal{U}, \mathcal{F})
$$
を次の式で定める。
\begin{equation}
\label{sites-cohomology-equation-d-cech}
d(s)_{i_0\ldots i_{p + 1}} =
\sum\nolimits_{j = 0}^{p + 1}
(-1)^j s_{i_0\ldots \hat i_j \ldots i_{p + 1}}
|_{U_{i_0} \times_U \ldots \times_U U_{i_{p + 1}}}
\end{equation}
ここで制限は射影写像
$$
U_{i_0} \times_U \ldots \times_U U_{i_{p + 1}} \longrightarrow
U_{i_0} \times_U \ldots \times_U \widehat{U_{i_j}} \times_U
\ldots \times_U U_{i_{p + 1}}.
$$
$d \circ d = 0$ は直接確かめられる。つまり
$\check{\mathcal{C}}^\bullet(\mathcal{U}, \mathcal{F})$ は複体である。

\begin{definition}
\label{sites-cohomology-definition-cech-complex}
圏 $\mathcal{C}$ をとる。固定された終域をもつ射の族
$\mathcal{U} = \{U_i \to U\}_{i \in I}$ をとり、すべてのファイバー積
$U_{i_0} \times_U \ldots \times_U U_{i_p}$ が $\mathcal{C}$ に存在するとする。
$\mathcal{C}$ 上のアーベル前層を $\mathcal{F}$ とする。
複体 $\check{\mathcal{C}}^\bullet(\mathcal{U}, \mathcal{F})$ を
$\mathcal{F}$ と族 $\mathcal{U}$ に付随する {\it {\v C}ech 複体} と呼ぶ。
そのコホモロジー群
$H^i(\check{\mathcal{C}}^\bullet(\mathcal{U}, \mathcal{F}))$
を $\mathcal{U}$ に関する $\mathcal{F}$ の
{\it {\v C}ech コホモロジー群} と呼び、
$\check H^i(\mathcal{U}, \mathcal{F})$ と書く。
\end{definition}

\noindent
サイト $\mathcal{C}$ の任意の被覆 $\{U_i \to U\}$ は、
この定義を適用できる固定終域付き射の族であることに注意する。

\begin{lemma}
\label{sites-cohomology-lemma-cech-h0}
サイト $\mathcal{C}$ をとる。
$\mathcal{C}$ 上のアーベル前層を $\mathcal{F}$ とする。
次は同値である。
\begin{enumerate}
\item $\mathcal{F}$ は $\mathcal{C}$ 上のアーベル層である。
\item $\mathcal{C}$ の任意の被覆
$\mathcal{U} = \{U_i \to U\}_{i \in I}$ に対して自然な写像
$$
\mathcal{F}(U) \to \check{H}^0(\mathcal{U}, \mathcal{F})
$$
(Sites, Section \ref{sites-section-sheafification} 参照）は全単射である。
\end{enumerate}
\end{lemma}

\begin{proof}
これは層条件がまさに、$\mathcal{C}$ の任意の被覆に対して
$\mathcal{F}(U) \to \check{H}^0(\mathcal{U}, \mathcal{F})$
が全単射であることだからである。
\end{proof}

\noindent
圏 $\mathcal{C}$ をとる。$\mathcal{C}$ の固定終域付き射の族
$\mathcal{U} = \{U_i \to U\}_{i\in I}$ で、すべてのファイバー積
$U_{i_0} \times_U \ldots \times_U U_{i_p}$ が $\mathcal{C}$ に存在するものをとる。
別の族 $\mathcal{V} = \{V_j \to V\}_{j\in J}$ もとる。
$f : U \to V$、$\alpha : I \to J$、および
$f_i : U_i \to V_{\alpha(i)}$ を固定終域付き射の族の射とする（
Sites, Section \ref{sites-section-refinements} 参照）。
このとき {\v C}ech 複体の写像
\begin{equation}
\label{sites-cohomology-equation-map-cech-complexes}
\varphi : \check{\mathcal{C}}^\bullet(\mathcal{V}, \mathcal{F})
\longrightarrow
\check{\mathcal{C}}^\bullet(\mathcal{U}, \mathcal{F})
\end{equation}
を得る。次数 $p$ では
$$
\varphi(s)_{i_0 \ldots i_p} =
(f_{i_0} \times \ldots \times f_{i_p})^*s_{\alpha(i_0) \ldots \alpha(i_p)}
$$


\section{前層上の関手としての {\v C}ech コホモロジー}
\label{sites-cohomology-section-cech-functor}

\noindent
警告：この節では圏上のアーベル前層だけを扱う。
層および層の圏という設定では、ここでの結果は完全に誤りとなる。

\medskip\noindent
圏 $\mathcal{C}$ をとる。固定終域付き射の族
$\mathcal{U} = \{U_i \to U\}_{i \in I}$ をとり、すべてのファイバー積
$U_{i_0} \times_U \ldots \times_U U_{i_p}$ が $\mathcal{C}$ に存在するとする。
$\mathcal{C}$ 上のアーベル前層を $\mathcal{F}$ とする。構成
$$
\mathcal{F} \longmapsto \check{\mathcal{C}}^\bullet(\mathcal{U}, \mathcal{F})
$$
は $\mathcal{F}$ に関して関手的である。実際、これは関手
\begin{equation}
\label{sites-cohomology-equation-cech-functor}
\check{\mathcal{C}}^\bullet(\mathcal{U}, -) :
\textit{PAb}(\mathcal{C})
\longrightarrow
\text{Comp}^{+}(\textit{Ab})
\end{equation}
記法については Derived Categories, Definition
\ref{derived-definition-complexes-notation} を参照。
アーベル圏における下に有界な複体の圏はアーベル圏であることを思い出す（
Homology, Lemma \ref{homology-lemma-cat-cochain-abelian} 参照）。

\begin{lemma}
\label{sites-cohomology-lemma-cech-exact-presheaves}
式 (\ref{sites-cohomology-equation-cech-functor}) で与えられる関手は完全関手である（
Homology, Lemma \ref{homology-lemma-exact-functor} 参照）。
\end{lemma}

\begin{proof}
任意の $\mathcal{C}$ の対象 $W$ に対して、関手
$\mathcal{F} \mapsto \mathcal{F}(W)$ は
$\textit{PAb}(\mathcal{C})$ から $\textit{Ab}$ への加法的完全関手である。
複体の項 $\check{\mathcal{C}}^p(\mathcal{U}, \mathcal{F})$ は
これらの完全関手の積なので完全である。
さらに複体の列が完全であることと、各次数の項の列が完全であることは同値である。
よって補題が従う。
\end{proof}

\begin{lemma}
\label{sites-cohomology-lemma-cech-cohomology-delta-functor-presheaves}
圏 $\mathcal{C}$ をとる。
固定終域付き射の族 $\mathcal{U} = \{U_i \to U\}_{i \in I}$ をとり、
すべてのファイバー積
$U_{i_0} \times_U \ldots \times_U U_{i_p}$ が $\mathcal{C}$ に存在するとする。
関手 $\mathcal{F} \mapsto \check{H}^n(\mathcal{U}, \mathcal{F})$ は、
アーベル圏 $\textit{PAb}(\mathcal{C})$ から $\mathbf{Z}$-加群の圏への
$\delta$-関手をなす（
Homology, Definition \ref{homology-definition-cohomological-delta-functor} 参照）。
\end{lemma}

\begin{proof}
Lemma \ref{sites-cohomology-lemma-cech-exact-presheaves} により、
アーベル前層の短完全列
$0 \to \mathcal{F}_1 \to \mathcal{F}_2 \to \mathcal{F}_3 \to 0$
は $\mathbf{Z}$-加群の複体の短完全列に移る。
したがって
Homology, Lemma \ref{homology-lemma-long-exact-sequence-cochain}
を用いて境界写像
$\delta_{\mathcal{F}_1 \to \mathcal{F}_2 \to \mathcal{F}_3} :
\check{H}^n(\mathcal{U}, \mathcal{F}_3) \to
\check{H}^{n + 1}(\mathcal{U}, \mathcal{F}_1)$
と対応する長完全列を得る。
これらの写像が前層の短完全列の射と両立することの確認は省略する。
\end{proof}

\begin{lemma}
\label{sites-cohomology-lemma-cech-map-into}
圏 $\mathcal{C}$ をとる。固定終域付き射の族
$\mathcal{U} = \{U_i \to U\}_{i \in I}$ をとり、すべてのファイバー積
$U_{i_0} \times_U \ldots \times_U U_{i_p}$ が $\mathcal{C}$ に存在するとする。
アーベル前層の鎖複体 $\mathbf{Z}_{\mathcal{U}, \bullet}$ を考える。
$$
\ldots
\to
\bigoplus_{i_0i_1i_2} \mathbf{Z}_{U_{i_0} \times_U U_{i_1} \times_U U_{i_2}}
\to
\bigoplus_{i_0i_1} \mathbf{Z}_{U_{i_0} \times_U U_{i_1}}
\to
\bigoplus_{i_0} \mathbf{Z}_{U_{i_0}}
\to 0 \to \ldots
$$
ここで最後の非零項は次数 $0$ に置き、写像
$$
\mathbf{Z}_{U_{i_0} \times_U \ldots \times_u U_{i_{p + 1}}}
\longrightarrow
\mathbf{Z}_{U_{i_0} \times_U
\ldots \widehat{U_{i_j}} \ldots \times_U U_{i_{p + 1}}}
$$
$(-1)^j$ 倍の標準写像で与える。すると同型
$$
\Hom_{\textit{PAb}(\mathcal{C})}(\mathbf{Z}_{\mathcal{U}, \bullet}, \mathcal{F})
=
\check{\mathcal{C}}^\bullet(\mathcal{U}, \mathcal{F})
$$
$\mathcal{F} \in \Ob(\textit{PAb}(\mathcal{C}))$ に関して関手的である。
\end{lemma}

\begin{proof}
これは次の事実から直ちに従う。
\begin{align*}
\Hom_{\textit{PAb}(\mathcal{C})}(
\bigoplus_{i_0 \ldots i_p}
\mathbf{Z}_{U_{i_0} \times_U \ldots \times_U U_{i_p}},
\mathcal{F})
& =
\prod_{i_0 \ldots i_p}
\Hom_{\textit{PAb}(\mathcal{C})}(
\mathbf{Z}_{U_{i_0} \times_U \ldots \times_U U_{i_p}},
\mathcal{F}) \\
& =
\prod_{i_0 \ldots i_p}
\mathcal{F}(U_{i_0} \times_U \ldots \times_U U_{i_p})
\end{align*}
Modules on Sites, Lemma \ref{sites-modules-lemma-obvious-adjointness} 参照。
\end{proof}

\begin{lemma}
\label{sites-cohomology-lemma-homology-complex}
圏 $\mathcal{C}$ をとる。固定終域付き射の族
$\mathcal{U} = \{f_i : U_i \to U\}_{i \in I}$ をとり、すべてのファイバー積
$U_{i_0} \times_U \ldots \times_U U_{i_p}$ が $\mathcal{C}$ に存在するとする。
上の Lemma \ref{sites-cohomology-lemma-cech-map-into} の前層の鎖複体
$\mathbf{Z}_{\mathcal{U}, \bullet}$ は正次数で完全である。
すなわちホモロジー前層
$H_i(\mathbf{Z}_{\mathcal{U}, \bullet})$ は $i>0$ で零である。
\end{lemma}

\begin{proof}
 $V$ を $\mathcal{C}$ の対象とする。アーベル群の鎖複体
$\mathbf{Z}_{\mathcal{U}, \bullet}(V)$ が次数 $>0$ で完全であることを示せばよい。
これは次の複体である。
$$
\xymatrix{
\ldots \ar[d] \\
\bigoplus_{i_0i_1i_2}
\mathbf{Z}[
\Mor_\mathcal{C}(V, U_{i_0} \times_U U_{i_1} \times_U U_{i_2})
]
\ar[d] \\
\bigoplus_{i_0i_1}
\mathbf{Z}[
\Mor_\mathcal{C}(V, U_{i_0} \times_U U_{i_1})
]
\ar[d] \\
\bigoplus_{i_0}
\mathbf{Z}[
\Mor_\mathcal{C}(V, U_{i_0})
] \ar[d] \\
0
}
$$
任意の射 $\varphi : V \to U$ に対して
$\Mor_\varphi(V, U_i) = \{\varphi_i : V \to U_i \mid
f_i \circ \varphi_i = \varphi\}$ とおく。同様の記法を
$\Mor_\varphi(V, U_{i_0} \times_U \ldots \times_U U_{i_p})$ にも用いる。
ファイバー積 $U_{i_0} \times_U \ldots \times_U U_{i_p}$ の間の
各射影写像との合成は、これらの射の集合を保つ。
したがって上の複体は次の複体と同じである。
$$
\xymatrix{
\ldots \ar[d] \\
\bigoplus_\varphi
\bigoplus_{i_0i_1i_2}
\mathbf{Z}[
\Mor_\varphi(V, U_{i_0} \times_U U_{i_1} \times_U U_{i_2})
]
\ar[d] \\
\bigoplus_\varphi
\bigoplus_{i_0i_1}
\mathbf{Z}[
\Mor_\varphi(V, U_{i_0} \times_U U_{i_1})
]
\ar[d] \\
\bigoplus_\varphi
\bigoplus_{i_0}
\mathbf{Z}[
\Mor_\varphi(V, U_{i_0})
] \ar[d] \\
0
}
$$
次の等式にも注意する。
$$
\Mor_\varphi(V, U_{i_0} \times_U \ldots \times_U U_{i_p})
=
\Mor_\varphi(V, U_{i_0}) \times \ldots
\times \Mor_\varphi(V, U_{i_p})
$$
これと $\mathbf{Z}[A] \oplus \mathbf{Z}[B] =
\mathbf{Z}[A \amalg B]$ という事実から、複体は
$$
\xymatrix{
\ldots \ar[d] \\
\bigoplus_\varphi
\mathbf{Z}\left[
\coprod_{i_0i_1i_2}
\Mor_\varphi(V, U_{i_0}) \times \Mor_\varphi(V, U_{i_1}) \times
\Mor_\varphi(V, U_{i_2})
\right]
\ar[d] \\
\bigoplus_\varphi
\mathbf{Z}\left[
\coprod_{i_0i_1}
\Mor_\varphi(V, U_{i_0}) \times \Mor_\varphi(V, U_{i_1})
\right]
\ar[d] \\
\bigoplus_\varphi
\mathbf{Z}\left[
\coprod_{i_0}
\Mor_\varphi(V, U_{i_0})
\right] \ar[d] \\
0
}
$$
最後に $S_\varphi = \coprod_{i \in I} \Mor_\varphi(V, U_i)$ とおくと、
次を得る。
$$
\bigoplus\nolimits_\varphi \left(\ldots \to
\mathbf{Z}[S_\varphi \times S_\varphi \times S_\varphi] \to
\mathbf{Z}[S_\varphi \times S_\varphi] \to
\mathbf{Z}[S_\varphi] \to 0 \to \ldots
\right)
$$
これで問題は簡単になった。すなわち、空でない集合 $S$ ごとに
自由アーベル群の（拡張された）複体
$$
\ldots \to
\mathbf{Z}[S \times S \times S] \to
\mathbf{Z}[S \times S] \to
\mathbf{Z}[S] \xrightarrow{\Sigma} \mathbf{Z} \to 0 \to \ldots
$$
がすべての次数で完全であることを示せば十分である。
これを見るため $s \in S$ を一つ固定し、ホモトピー
$$
n_{(s_0, \ldots, s_p)} \longmapsto n_{(s, s_0, \ldots, s_p)}
$$
を用いる（記法は明らかである）。
\end{proof}

\begin{lemma}
\label{sites-cohomology-lemma-complex-tensored-still-exact}
\begin{slogan}
整数係数の前層 Cech 複体は、整数の定数前層の平坦分解である。
\end{slogan}
$\mathcal{C}$ を圏とする。$\mathcal{U} = \{f_i : U_i \to U\}_{i \in I}$ を固定終域を
もつ射の族で、すべてのファイバー積
$U_{i_0} \times_U \ldots \times_U U_{i_p}$ が $\mathcal{C}$ に存在するものとする。
$\mathcal{O}$ を $\mathcal{C}$ 上の環の前層とする。鎖複体
$$
\mathbf{Z}_{\mathcal{U}, \bullet}
\otimes_{p, \mathbf{Z}}
\mathcal{O}
$$
は正次数で完全である。ここで $\mathbf{Z}_{\mathcal{U}, \bullet}$ は Lemma
\ref{sites-cohomology-lemma-cech-map-into} の鎖複体であり、テンソル積は値 $\mathbf{Z}$ をもつ
環の定数前層上でとる。
\end{lemma}

\begin{proof}
$V$ を $\mathcal{C}$ の対象とする。Lemma \ref{sites-cohomology-lemma-homology-complex} の証明で、
$\mathbf{Z}_{\mathcal{U}, \bullet}(V)$ は、次数 0 に置かれた $\mathbf{Z}$ とホモトピー
同値な複体の直和に複体として同型であることを見た。従って
$\mathbf{Z}_{\mathcal{U}, \bullet}(V) \otimes_\mathbf{Z} \mathcal{O}(V)$
も、次数 0 に置かれた $\mathcal{O}(V)$ とホモトピー同値な複体の直和に複体として
同型である。あるいは Modules on Sites, Lemma
\ref{sites-modules-lemma-flat-resolution-of-flat} を用いてもよい。実際、前層
$\mathbf{Z}_{\mathcal{U}, i}$ は平坦であり、Lemma \ref{sites-cohomology-lemma-homology-complex} の証明は
$H_0(\mathbf{Z}_{\mathcal{U}, \bullet})$ も平坦な前層であることを示している。
\end{proof}

\begin{lemma}
\label{sites-cohomology-lemma-cech-cohomology-derived-presheaves}
$\mathcal{C}$ を圏とする。$\mathcal{U} = \{f_i : U_i \to U\}_{i \in I}$ を固定終域を
もつ射の族で、すべてのファイバー積
$U_{i_0} \times_U \ldots \times_U U_{i_p}$ が $\mathcal{C}$ に存在するものとする。
Cech コホモロジー関手 $\check{H}^p(\mathcal{U}, -)$ は、$\delta$-関手として
次の関手の右導来関手と標準的に同型である。
$$
\check{H}^0(\mathcal{U}, -) :
\textit{PAb}(\mathcal{C})
\longrightarrow
\textit{Ab}.
$$
さらに、関手的な擬同型
$$
\check{\mathcal{C}}^\bullet(\mathcal{U}, \mathcal{F})
\longrightarrow
R\check{H}^0(\mathcal{U}, \mathcal{F})
$$
が存在する。ここで右辺は、左完全関手 $\check{H}^0(\mathcal{U}, -)$ の導来関手
$$
R\check{H}^0(\mathcal{U}, -) :
D^{+}(\textit{PAb}(\mathcal{C}))
\longrightarrow
D^{+}(\mathbf{Z})
$$
を表す。
\end{lemma}

\begin{proof}
アーベル前層の圏は十分多くの入射対象をもつことに注意する（Injectives, Proposition
\ref{injectives-proposition-presheaves-injectives} 参照）。また
$\check{H}^0(\mathcal{U}, -)$ はアーベル前層の圏から $\mathbf{Z}$-加群の圏への
左完全関手である。従って導来関手と右導来関手が存在する（Derived Categories,
Section \ref{derived-section-right-derived-functor} 参照）。

\medskip\noindent
$\mathcal{I}$ を入射アーベル前層とする。この場合、関手
$\Hom_{\textit{PAb}(\mathcal{C})}(-, \mathcal{I})$
は $\textit{PAb}(\mathcal{C})$ 上で完全である。Lemma
\ref{sites-cohomology-lemma-cech-map-into} により
$$
\Hom_{\textit{PAb}(\mathcal{C})}(
\mathbf{Z}_{\mathcal{U}, \bullet}, \mathcal{I})
=
\check{\mathcal{C}}^\bullet(\mathcal{U}, \mathcal{I}).
$$
Lemma \ref{sites-cohomology-lemma-homology-complex} により
$\mathbf{Z}_{\mathcal{U}, \bullet}$ は正次数で完全である。従って上で述べた
$\mathcal{I}$ への Hom の完全性により、すべての $i > 0$ に対して
$\check{H}^i(\mathcal{U}, \mathcal{I}) = 0$ である。従って $\delta$-関手
$(\check{H}^n, \delta)$（Lemma
\ref{sites-cohomology-lemma-cech-cohomology-delta-functor-presheaves} 参照）は Homology, Lemma
\ref{homology-lemma-efface-implies-universal} の仮定を満たし、したがって普遍
$\delta$-関手である。

\medskip\noindent
Derived Categories, Lemma \ref{derived-lemma-higher-derived-functors} により、列
$R^i\check{H}^0(\mathcal{U}, -)$ も普遍 $\delta$-関手をなす。普遍
$\delta$-関手の一意性（Homology, Lemma
\ref{homology-lemma-uniqueness-universal-delta-functor} 参照）により
$R^i\check{H}^0(\mathcal{U}, -) = \check{H}^i(\mathcal{U}, -)$
を得る。これは大部分の応用には十分であり、読者には証明の残りを読み飛ばすことを
勧める。

\medskip\noindent
$\mathcal{F}$ を $\mathcal{C}$ 上の任意のアーベル前層とする。圏
$\textit{PAb}(\mathcal{C})$ における入射分解
$\mathcal{F} \to \mathcal{I}^\bullet$ を選ぶ。二重複体
$\check{\mathcal{C}}^\bullet(\mathcal{U}, \mathcal{I}^\bullet)$
を考える。その項は $\check{\mathcal{C}}^p(\mathcal{U}, \mathcal{I}^q)$ である。
次に、この二重複体に付随する全複体
$\text{Tot}(\check{\mathcal{C}}^\bullet(\mathcal{U}, \mathcal{I}^\bullet))$
を考える（Homology, Section \ref{homology-section-double-complexes} 参照）。
複体の写像
$$
\check{\mathcal{C}}^\bullet(\mathcal{U}, \mathcal{F})
\longrightarrow
\text{Tot}(\check{\mathcal{C}}^\bullet(\mathcal{U}, \mathcal{I}^\bullet))
$$
があり、これは写像
$\check{\mathcal{C}}^p(\mathcal{U}, \mathcal{F})
\to \check{\mathcal{C}}^p(\mathcal{U}, \mathcal{I}^0)$
から生じる。また、複体の写像
$$
\check{H}^0(\mathcal{U}, \mathcal{I}^\bullet)
\longrightarrow
\text{Tot}(\check{\mathcal{C}}^\bullet(\mathcal{U}, \mathcal{I}^\bullet))
$$
があり、これは写像
$\check{H}^0(\mathcal{U}, \mathcal{I}^q) \to
\check{\mathcal{C}}^0(\mathcal{U}, \mathcal{I}^q)$
から生じる。Homology, Lemma \ref{homology-lemma-double-complex-gives-resolution} を
適用すると、これら二つの写像はいずれも擬同型である。実際、Cech 複体は関手として
完全なので（Lemma \ref{sites-cohomology-lemma-cech-exact-presheaves}）、二重複体の列は正次数で完全で
ある。また、先ほど見たように入射前層 $\mathcal{I}^q$ の高次 Cech コホモロジー群は
0 なので、二重複体の行も正次数で完全である。擬同型は $D^{+}(\mathbf{Z})$ で可逆に
なるから、これにより補題の最後に表示した射が得られる。この射が関手的であることの
確認は省略する。
\end{proof}





\section{{\v C}ech コホモロジーとコホモロジー}
\label{sites-cohomology-section-cech-cohomology-cohomology}

\noindent
コホモロジーと {\v C}ech コホモロジーの関係は、入射アーベル層の
{\v C}ech コホモロジーが消えることから従う。これを見るために、入射
アーベル層は入射アーベル前層でもあることに注意し、前節の {\v C}ech
コホモロジーの結果を適用する。

\begin{lemma}
\label{sites-cohomology-lemma-injective-abelian-sheaf-injective-presheaf}
$\mathcal{C}$ をサイトとする。入射アーベル層は圏
$\textit{PAb}(\mathcal{C})$ の対象としても入射的である。
\end{lemma}

\begin{proof}
Homology, Lemma \ref{homology-lemma-adjoint-preserve-injectives} を、圏
$\mathcal{A} = \textit{Ab}(\mathcal{C})$、
$\mathcal{B} = \textit{PAb}(\mathcal{C})$、包含関手、および層化に適用する。
この補題の仮定がすべて満たされることは、Modules on Sites, Section
\ref{sites-modules-section-abelian-sheaves} を参照すれば分かる。
\end{proof}

\begin{lemma}
\label{sites-cohomology-lemma-injective-trivial-cech}
$\mathcal{C}$ をサイトとする。
$\mathcal{U} = \{U_i \to U\}_{i \in I}$ を $\mathcal{C}$ の被覆とする。
$\mathcal{I}$ を入射アーベル層、すなわち $\textit{Ab}(\mathcal{C})$ の
入射対象とする。このとき
$$
\check{H}^p(\mathcal{U}, \mathcal{I}) =
\left\{
\begin{matrix}
\mathcal{I}(U) & \text{if} & p = 0 \\
0 & \text{if} & p > 0
\end{matrix}
\right.
$$
\end{lemma}

\begin{proof}
Lemma \ref{sites-cohomology-lemma-injective-abelian-sheaf-injective-presheaf} により、
$\mathcal{I}$ は $\textit{PAb}(\mathcal{C})$ の入射対象でもある。
したがって Lemma \ref{sites-cohomology-lemma-cech-cohomology-derived-presheaves}（または
その証明）を適用すれば、高次の {\v C}ech コホモロジー群が消えることが
分かる。0 次の場合は Lemma \ref{sites-cohomology-lemma-cech-h0} を参照。
\end{proof}

\begin{lemma}
\label{sites-cohomology-lemma-cech-cohomology}
$\mathcal{C}$ をサイトとする。
$\mathcal{U} = \{U_i \to U\}_{i \in I}$ を $\mathcal{C}$ の被覆とする。
変換
$$
\check{\mathcal{C}}^\bullet(\mathcal{U}, -)
\longrightarrow
R\Gamma(U, -)
$$
は関手 $\textit{Ab}(\mathcal{C}) \to D^{+}(\mathbf{Z})$ の間に存在する。
特に、$\textit{Ab}(\mathcal{C})$ を動く $\mathcal{F}$ に対して、関手の変換
$\check{H}^p(\mathcal{U}, \mathcal{F}) \to H^p(U, \mathcal{F})$ を得る。
\end{lemma}

\begin{proof}
$\mathcal{F}$ をアーベル層とする。入射分解
$\mathcal{F} \to \mathcal{I}^\bullet$ をとる。二重複体
$\check{\mathcal{C}}^\bullet(\mathcal{U}, \mathcal{I}^\bullet)$
の項は $\check{\mathcal{C}}^p(\mathcal{U}, \mathcal{I}^q)$ である。次に、
これに付随する全複体
$\text{Tot}(\check{\mathcal{C}}^\bullet(\mathcal{U}, \mathcal{I}^\bullet))$,
を考える（Homology, Definition
\ref{homology-definition-associated-simple-complex} 参照）。複体の射
$$
\alpha :
\Gamma(U, \mathcal{I}^\bullet)
\longrightarrow
\text{Tot}(\check{\mathcal{C}}^\bullet(\mathcal{U}, \mathcal{I}^\bullet))
$$
は写像 $\mathcal{I}^q(U) \to \check{H}^0(\mathcal{U}, \mathcal{I}^q)$
から得られる。また、複体の射
$$
\beta :
\check{\mathcal{C}}^\bullet(\mathcal{U}, \mathcal{F})
\longrightarrow
\text{Tot}(\check{\mathcal{C}}^\bullet(\mathcal{U}, \mathcal{I}^\bullet))
$$
は写像 $\mathcal{F} \to \mathcal{I}^0$ から得られる。
Homology, Lemma
\ref{homology-lemma-double-complex-gives-resolution} を適用すると、
$\alpha$ は準同型同値であることが分かる。実際、Lemma
\ref{sites-cohomology-lemma-injective-trivial-cech} により、二重複体
$\check{\mathcal{C}}^\bullet(\mathcal{U}, \mathcal{I}^\bullet)$ の第 $q$
行は $\Gamma(U, \mathcal{I}^q)$ の分解である。したがって $\alpha$ は
$D^{+}(\mathbf{Z})$ で可逆となり、補題の変換は $\beta$ と $\alpha$ の逆射の
合成である。これが関手的であることの確認は省略する。
\end{proof}

\begin{lemma}
\label{sites-cohomology-lemma-cech-h1}
$\mathcal{C}$ をサイトとし、$\mathcal{G}$ を $\mathcal{C}$ 上のアーベル層と
する。$\mathcal{U} = \{U_i \to U\}_{i \in I}$ を $\mathcal{C}$ の被覆とする。
写像
$$
\check{H}^1(\mathcal{U}, \mathcal{G})
\longrightarrow
H^1(U, \mathcal{G})
$$
は単射であり、Lemma \ref{sites-cohomology-lemma-torsors-h1} の全単射を通じて
$\check{H}^1(\mathcal{U}, \mathcal{G})$ を、各 $U_i$ 上で自明なトーサーに
制限される $\mathcal{G}|_U$-トーサーの同型類の集合と同一視する。
\end{lemma}

\begin{proof}
これを見るため逆写像を構成する。すなわち、$\mathcal{C}/U$ 上の
$\mathcal{G}|_U$-トーサー $\mathcal{F}$ で、その $\mathcal{C}/U_i$ への制限が
自明であるものをとる。Lemma \ref{sites-cohomology-lemma-trivial-torsor} により、これは
$s_i \in \mathcal{F}(U_i)$ なる切断が存在することを意味する。
$U_{i_0} \times_U U_{i_1}$ 上には、次を満たす $\mathcal{G}$ の一意な切断
$s_{i_0i_1}$ が存在する。
$s_{i_0i_1} \cdot s_{i_0}|_{U_{i_0} \times_U U_{i_1}} =
s_{i_1}|_{U_{i_0} \times_U U_{i_1}}$。
容易な計算により、$s_{i_0i_1}$ は {\v C}ech コサイクルであり、その類は
（切断 $s_i$ の選択によらず）well-defined であることが分かる。逆写像は
$\mathcal{F}$ の同型類をコサイクル $(s_{i_0i_1})$ のコホモロジー類へ写す。
この写像が実際に逆写像であることの確認は省略する。
\end{proof}

\begin{lemma}
\label{sites-cohomology-lemma-include}
$\mathcal{C}$ をサイトとする。関手
$i : \textit{Ab}(\mathcal{C}) \to \textit{PAb}(\mathcal{C})$。
は左完全関手であり、その右導来関手は
$$
R^pi(\mathcal{F}) = \underline{H}^p(\mathcal{F}) :
U \longmapsto H^p(U, \mathcal{F})
$$
で与えられる（Section \ref{sites-cohomology-section-locality} の議論参照）。
\end{lemma}

\begin{proof}
$i$ が左完全であることは明らかである。入射分解
$\mathcal{F} \to \mathcal{I}^\bullet$ をとる。定義により $R^pi$ は複体
$\mathcal{I}^\bullet$ の $p$ 次コホモロジー前層である。言い換えると、
$\mathcal{C}$ の対象 $U$ 上の $R^pi(\mathcal{F})$ の切断は
$$
\frac{\Ker(\mathcal{I}^n(U) \to \mathcal{I}^{n + 1}(U))}
{\Im(\mathcal{I}^{n - 1}(U) \to \mathcal{I}^n(U))}.
$$
で与えられ、これは $H^p(U, \mathcal{F})$ の定義である。
\end{proof}

\begin{lemma}
\label{sites-cohomology-lemma-cech-spectral-sequence}
$\mathcal{C}$ をサイトとし、$\mathcal{U} = \{U_i \to U\}_{i \in I}$ を
$\mathcal{C}$ の被覆とする。任意のアーベル層 $\mathcal{F}$ に対して、
次を満たすスペクトル系列 $(E_r, d_r)_{r \geq 0}$ が存在する。
$$
E_2^{p, q} = \check{H}^p(\mathcal{U}, \underline{H}^q(\mathcal{F}))
$$
は $H^{p + q}(U, \mathcal{F})$ に収束する。このスペクトル系列は
$\mathcal{F}$ について関手的である。
\end{lemma}

\begin{proof}
これは関手
$$
i :  \textit{Ab}(\mathcal{C}) \to \textit{PAb}(\mathcal{C})
\quad\text{および}\quad
\check{H}^0(\mathcal{U}, - ) : \textit{PAb}(\mathcal{C})
\to \textit{Ab}.
$$
に対する Grothendieck スペクトル系列である（Derived Categories, Lemma
\ref{derived-lemma-grothendieck-spectral-sequence} 参照）。実際、Lemma
\ref{sites-cohomology-lemma-cech-h0} により $\check{H}^0(\mathcal{U}, i(\mathcal{F})) =
\mathcal{F}(U)$ である。また Lemma \ref{sites-cohomology-lemma-injective-trivial-cech} により
$i(\mathcal{I})$ は {\v C}ech 非輪状であり、Lemma
\ref{sites-cohomology-lemma-cech-cohomology-derived-presheaves} により
$\textit{PAb}(\mathcal{C})$ 上の関手として
$\check{H}^p(\mathcal{U}, -) = R^p\check{H}^0(\mathcal{U}, -)$ である。
以上を合わせれば補題が従う。
\end{proof}

\begin{lemma}
\label{sites-cohomology-lemma-cech-spectral-sequence-application}
$\mathcal{C}$ をサイトとする。
$\mathcal{U} = \{U_i \to U\}_{i \in I}$ を被覆とする。
$\mathcal{F} \in \Ob(\textit{Ab}(\mathcal{C}))$ とする。
すべての $i > 0$、$p \geq 0$、$i_0, \ldots, i_p \in I$ に対して
$H^i(U_{i_0} \times_U \ldots \times_U U_{i_p}, \mathcal{F}) = 0$ と仮定する。
このとき $\check{H}^p(\mathcal{U}, \mathcal{F}) = H^p(U, \mathcal{F})$ である。
\end{lemma}

\begin{proof}
Lemma \ref{sites-cohomology-lemma-cech-spectral-sequence} のスペクトル系列を用いる。
仮定により、$q \not = 0$ を満たすすべての $(p, q)$ について
$E_2^{p, q} = 0$ である。したがってスペクトル系列は $E_2$ で退化し、
主張が従う。
\end{proof}

\begin{lemma}
\label{sites-cohomology-lemma-ses-cech-h1}
$\mathcal{C}$ をサイトとする。
次の短完全列
$$
0 \to \mathcal{F} \to \mathcal{G} \to \mathcal{H} \to 0
$$
は $\mathcal{C}$ 上のアーベル層の短完全列とする。$U$ を $\mathcal{C}$ の
対象とする。$U$ の被覆の共終系 $\mathcal{U}$ で
$\check{H}^1(\mathcal{U}, \mathcal{F}) = 0$,
を満たすものが存在するなら、写像 $\mathcal{G}(U) \to \mathcal{H}(U)$ は
全射である。
\end{lemma}

\begin{proof}
 $s \in \mathcal{H}(U)$ をとる。次を満たす被覆
$\mathcal{U} = \{U_i \to U\}_{i \in I}$ を選ぶ。
(a) $\check{H}^1(\mathcal{U}, \mathcal{F}) = 0$、および (b)
$s|_{U_i}$ が切断 $s_i \in \mathcal{G}(U_i)$ の像である。
 (b) を満たす被覆は確実に存在するので、補題の仮定から (a)、(b) の
両方を満たす被覆を選べる。切断
$$
s_{i_0i_1} =
s_{i_1}|_{U_{i_0} \times_U U_{i_1}} - s_{i_0}|_{U_{i_0} \times_U U_{i_1}}.
$$
 $s_i$ は $s$ の持ち上げなので、
$s_{i_0i_1} \in \mathcal{F}(U_{i_0} \times_U U_{i_1})$.
 $\check{H}^1(\mathcal{U}, \mathcal{F})$ の消滅により、次を満たす切断
$t_i \in \mathcal{F}(U_i)$ をとれる。
$$
s_{i_0i_1} =
t_{i_1}|_{U_{i_0} \times_U U_{i_1}} - t_{i_0}|_{U_{i_0} \times_U U_{i_1}}.
$$
すると切断 $s_i - t_i$ は明らかに層条件を満たし、$s$ に写る $U$ 上の
$\mathcal{G}$ の切断へ貼り合わされる。これで示された。
\end{proof}

\begin{lemma}
\label{sites-cohomology-lemma-cech-vanish-collection}
(Cohomology, Lemma \ref{cohomology-lemma-cech-vanish} の変形。)
$\mathcal{C}$ をサイトとする。$\text{Cov}_\mathcal{C}$ を $\mathcal{C}$ の
被覆全体の集合とする（Sites, Definition \ref{sites-definition-site} 参照）。
次をおく。
$\mathcal{B} \subset \Ob(\mathcal{C})$ および
$\text{Cov} \subset \text{Cov}_\mathcal{C}$
を部分集合とする。$\mathcal{F}$ を $\mathcal{C}$ 上のアーベル層とする。
次を仮定する。
\begin{enumerate}
\item 任意の $\mathcal{U} \in \text{Cov}$、$\mathcal{U} = \{U_i \to U\}_{i \in I}$
に対して $U, U_i \in \mathcal{B}$ であり、すべての
$U_{i_0} \times_U \ldots \times_U U_{i_p}$ は $\mathcal{B}$ に属する。
\item 任意の $U \in \mathcal{B}$ に対し、$\text{Cov}$ に現れる $U$ の被覆は
$U$ の被覆の共終系をなす。
\item 任意の $\mathcal{U} \in \text{Cov}$ に対して、すべての $p > 0$ で
$\check{H}^p(\mathcal{U}, \mathcal{F}) = 0$ である。
\end{enumerate}
すると任意の $U \in \mathcal{B}$ とすべての $p > 0$ に対して
$H^p(U, \mathcal{F}) = 0$ である。
\end{lemma}

\begin{proof}
$\mathcal{F}$ と $\text{Cov}$ は補題のとおりとする。これを
「任意の $\mathcal{U} \in \text{Cov}$ に対して $\mathcal{F}$ は高次
{\v C}ech コホモロジーが消える」と呼ぶことにする。$\mathcal{F}$ を入射
アーベル層 $\mathcal{I}$ へ埋め込む。Lemma
\ref{sites-cohomology-lemma-injective-trivial-cech} により、任意の $\mathcal{U} \in \text{Cov}$
で $\mathcal{I}$ の高次 {\v C}ech コホモロジーも消える。$\mathcal{Q} =
\mathcal{I}/\mathcal{F}$ とおくと、短完全列
$$
0 \to \mathcal{F} \to \mathcal{I} \to \mathcal{Q} \to 0.
$$
Lemma \ref{sites-cohomology-lemma-ses-cech-h1} と仮定 (2) により、この列は各
$U \in \mathcal{B}$ に対して完全列
$$
0 \to \mathcal{F}(U) \to \mathcal{I}(U) \to \mathcal{Q}(U) \to 0.
$$
を与える。したがって任意の $\mathcal{U} \in \text{Cov}$ に対し、{\v C}ech 複体の
短完全列
$$
0 \to
\check{\mathcal{C}}^\bullet(\mathcal{U}, \mathcal{F}) \to
\check{\mathcal{C}}^\bullet(\mathcal{U}, \mathcal{I}) \to
\check{\mathcal{C}}^\bullet(\mathcal{U}, \mathcal{Q}) \to 0
$$
を得る。これは仮定 (1) により、{\v C}ech 複体の各項が $\mathcal{B}$ の元上の
値の積からなるためである。特に任意の被覆 $\mathcal{U} \in \text{Cov}$ について
{\v C}ech コホモロジー群の長完全列を得る。従って $\mathcal{Q}$ もアーベル層で
あり、すべての $\mathcal{U} \in \text{Cov}$ に対して高次 {\v C}ech
コホモロジーが消える。

\medskip\noindent
次に、コホモロジーの長完全列
$$
\xymatrix{
0 \ar[r] &
H^0(U, \mathcal{F}) \ar[r] &
H^0(U, \mathcal{I}) \ar[r] &
H^0(U, \mathcal{Q}) \ar[lld] \\
&
H^1(U, \mathcal{F}) \ar[r] &
H^1(U, \mathcal{I}) \ar[r] &
H^1(U, \mathcal{Q}) \ar[lld] \\
&
\ldots & \ldots & \ldots \\
}
$$
を考える。これは任意の $U \in \mathcal{B}$ に対して成り立つ。
$\mathcal{I}$ は入射的なので、$n > 0$ なら
$H^n(U, \mathcal{I}) = 0$ である（Derived Categories, Lemma
\ref{derived-lemma-higher-derived-functors} 参照）。上の結果から
$H^0(U, \mathcal{I}) \to H^0(U, \mathcal{Q})$ は全射であり、従って
$H^1(U, \mathcal{F}) = 0$ となる。$\mathcal{F}$ は任意の、すべての
$\mathcal{U} \in \text{Cov}$ に対して高次 {\v C}ech コホモロジーが消える
アーベル層だったので、$\mathcal{Q}$ もそのような層であることから
$H^1(U, \mathcal{Q}) = 0$ も従う（上記参照）。長完全列により、さらに
$H^2(U, \mathcal{F}) = 0$ が従い、以下同様である。
\end{proof}














\section{第二コホモロジーと gerbe}
\label{sites-cohomology-section-gerbes}

\noindent
$p : \mathcal{S} \to \mathcal{C}$ を、すべての自己同型群が可換であるサイト上の
gerbe とする。この状況では、自己同型の層の第一および第二コホモロジー群
（Stacks, Lemma \ref{stacks-lemma-gerbe-abelian-auts}）が対象の存在を支配する。

\medskip\noindent
次の補題は、この関係について将来追加するより完全な議論により不要となる予定である。

\begin{lemma}
\label{sites-cohomology-lemma-existence}
$\mathcal{C}$ をサイトとする。$p : \mathcal{S} \to \mathcal{C}$ を、自己同型層が
アーベルであるサイト上の gerbe とする。$\mathcal{G}$ を Stacks, Lemma
\ref{stacks-lemma-gerbe-abelian-auts} で構成されるアーベル群の層とする。
$U$ を次を満たす $\mathcal{C}$ の対象とする。
\begin{enumerate}
\item $\mathcal{C}$ における $U$ の被覆 $\{U_i \to U\}$ の共終系で、
$H^1(U_i, \mathcal{G}) = 0$ およびすべての $i,j$ に対して
$H^1(U_i \times_U U_j, \mathcal{G}) = 0$ となるものが存在する。
\item $H^2(U, \mathcal{G}) = 0$。
\end{enumerate}
このとき $U$ 上にある $\mathcal{S}$ の対象が存在する。
\end{lemma}

\begin{proof}
Stacks, Definition \ref{stacks-definition-gerbe} により、被覆
$\mathcal{U} = \{U_i \to U\}$ と、各 $U_i$ 上にある $\mathcal{S}$ の対象
$x_i$ が存在する。$U_{ij} = U_i \times_U U_j$ と書く。被覆を細分すれば
(1) により $H^1(U_i, \mathcal{G}) = 0$ および
$H^1(U_{ij}, \mathcal{G}) = 0$ と仮定できる。層
$$
\mathcal{F}_{ij} =
\mathit{Isom}(x_i|_{U_{ij}}, x_j|_{U_{ij}})
$$
を $\mathcal{C}/U_{ij}$ 上で考える。
$\mathcal{G}|_{U_{ij}} = \mathit{Aut}(x_i|_{U_{ij}})$ なので、次の作用がある。
$$
\mathcal{G}|_{U_{ij}} \times \mathcal{F}_{ij} \to \mathcal{F}_{ij}
$$
これは前合成による作用である。$\mathcal{F}_{ij}$ は明らかに擬
$\mathcal{G}|_{U_{ij}}$-トーサーであり、gerbe の任意の二つの対象は局所的に
同型なので実際にトーサーである。被覆の選択と Lemma
\ref{sites-cohomology-lemma-torsors-h1} により、これらのトーサーは自明である（従って Lemma
\ref{sites-cohomology-lemma-trivial-torsor} により大域切断をもつ）。すなわち同型
$$
\varphi_{ij} : x_i|_{U_{ij}} \longrightarrow x_j|_{U_{ij}}
$$
 $U$ 上の対象 $x$ を得るため、これらの $\varphi_{ij}$ の選択を調整して
降下データを得る（$p : \mathcal{S} \to \mathcal{C}$ はスタックなので、これは
必ず実効的である）。降下データとなる妨げは、コサイクル条件が成立しない
可能性である。$U_{ijk} = U_i \times_U U_j \times_U U_k$ とおくと、次を考えられる。
$$
g_{ijk} = \varphi_{ik}^{-1}|_{U_{ijk}} \circ \varphi_{jk}|_{U_{ijk}} \circ
\varphi_{ij}|_{U_{ijk}}
$$
これは $U_{ijk}$ 上の $x_i$ の自己同型である。従って $g_{ijk}$ を
$U_{ijk}$ 上の $\mathcal{G}$ の切断とみなす。計算（省略）により、
$(g_{i_0i_1i_2})$ は被覆 $\mathcal{U}$ に関する $\mathcal{G}$ の
{\v C}ech 複体 ${\check C}^\bullet(\mathcal{U}, \mathcal{G})$ における
$2$-コサイクルである。
Lemma \ref{sites-cohomology-lemma-cech-spectral-sequence} のスペクトル系列と、すべての $i$ で
$H^1(U_i, \mathcal{G}) = 0$ であることから、
${\check H}^2(\mathcal{U}, \mathcal{G}) \to H^2(U, \mathcal{G})$ は単射である。
従って $H^2(U, \mathcal{G}) = 0$ という仮定により
$(g_{i_0i_1i_2})$ はコバウンダリである。よって
$g_{ij} \in \mathcal{G}(U_{ij})$ なる切断で
$g_{ik}^{-1}|_{U_{ijk}} g_{jk}|_{U_{ijk}}  g_{ij}|_{U_{ijk}} = g_{ijk}$
すべての $i,j,k$ に対して成り立つものをとれる。
$\varphi_{ij}$ を $\varphi_{ij}g_{ij}^{-1}$ で置き換えると、$\varphi_{ij}$ は
$U_i$ 上の対象 $x_i$ の降下データを与える。これで証明が完了する。
\end{proof}












\section{加群のコホモロジー}
\label{sites-cohomology-section-cohomology-modules}

\noindent
アーベル層のコホモロジーについて述べたことはすべて加群のコホモロジーにも
成り立つ。両者は一致するからである。

\begin{lemma}
\label{sites-cohomology-lemma-injective-module-injective-presheaf}
$(\mathcal{C}, \mathcal{O})$ を環付きサイトとする。入射加群層は圏
$\textit{PMod}(\mathcal{O})$ の対象としても入射的である。
\end{lemma}

\begin{proof}
Homology, Lemma \ref{homology-lemma-adjoint-preserve-injectives} を、圏
$\mathcal{A} = \textit{Mod}(\mathcal{O})$、
$\mathcal{B} = \textit{PMod}(\mathcal{O})$、包含関手、および層化に適用する。
この補題の仮定がすべて満たされることは、Modules on Sites,
Section \ref{sites-modules-section-sheafification-presheaves-modules} を
参照すれば分かる。
\end{proof}

\begin{lemma}
\label{sites-cohomology-lemma-include-modules}
$(\mathcal{C}, \mathcal{O})$ を環付きサイトとする。関手
$i : \textit{Mod}(\mathcal{C}) \to \textit{PMod}(\mathcal{C})$.
は左完全関手であり、その右導来関手は
$$
R^pi(\mathcal{F}) = \underline{H}^p(\mathcal{F}) :
U \longmapsto H^p(U, \mathcal{F})
$$
で与えられる（Section \ref{sites-cohomology-section-locality} の議論参照）。
\end{lemma}

\begin{proof}
$i$ が左完全であることは明らかである。$\textit{Mod}(\mathcal{O})$ における
入射分解 $\mathcal{F} \to \mathcal{I}^\bullet$ をとる。定義により $R^pi$ は
複体 $\mathcal{I}^\bullet$ の $p$ 次コホモロジー前層である。言い換えると、
$\mathcal{C}$ の対象 $U$ 上の $R^pi(\mathcal{F})$ の切断は
$$
\frac{\Ker(\mathcal{I}^n(U) \to \mathcal{I}^{n + 1}(U))}
{\Im(\mathcal{I}^{n - 1}(U) \to \mathcal{I}^n(U))}.
$$
で与えられ、これは $H^p(U, \mathcal{F})$ の定義である。
\end{proof}

\begin{lemma}
\label{sites-cohomology-lemma-injective-module-trivial-cech}
$(\mathcal{C}, \mathcal{O})$ を環付きサイトとする。
$\mathcal{U} = \{U_i \to U\}_{i \in I}$ を $\mathcal{C}$ の被覆とする。
$\mathcal{I}$ を入射 $\mathcal{O}$-加群、すなわち $\textit{Mod}(\mathcal{O})$ の
入射対象とする。このとき
$$
\check{H}^p(\mathcal{U}, \mathcal{I}) =
\left\{
\begin{matrix}
\mathcal{I}(U) & \text{if} & p = 0 \\
0 & \text{if} & p > 0
\end{matrix}
\right.
$$
\end{lemma}

\begin{proof}
Lemma \ref{sites-cohomology-lemma-cech-map-into} により、次の等式列の第一の等号を得る。
\begin{align*}
\check{\mathcal{C}}^\bullet(\mathcal{U}, \mathcal{I})
& =
\Mor_{\textit{PAb}(\mathcal{C})}(
\mathbf{Z}_{\mathcal{U}, \bullet}, \mathcal{I}) \\
& =
\Mor_{\textit{PMod}(\mathbf{Z})}(
\mathbf{Z}_{\mathcal{U}, \bullet}, \mathcal{I}) \\
& =
\Mor_{\textit{PMod}(\mathcal{O})}(
\mathbf{Z}_{\mathcal{U}, \bullet} \otimes_{p, \mathbf{Z}} \mathcal{O},
\mathcal{I})
\end{align*}
第三の等号は Modules on Sites,
Lemma \ref{sites-modules-lemma-adjointness-tensor-restrict-presheaves} による。
Lemma \ref{sites-cohomology-lemma-injective-module-injective-presheaf} により、$\mathcal{I}$ は
$\textit{PMod}(\mathcal{O})$ の入射対象である。従って
$\Hom_{\textit{PMod}(\mathcal{O})}(-, \mathcal{I})$ は完全関手である。
Lemma \ref{sites-cohomology-lemma-complex-tensored-still-exact} により、高次の {\v C}ech
コホモロジー群は消える。0 次の場合は Lemma \ref{sites-cohomology-lemma-cech-h0} を参照。
\end{proof}

\begin{lemma}
\label{sites-cohomology-lemma-cohomology-modules-abelian-agree}
$\mathcal{C}$ をサイトとし、$\mathcal{O}$ を $\mathcal{C}$ 上の環の層とする。
$\mathcal{F}$ を $\mathcal{O}$-加群とし、台となるアーベル群の層を
$\mathcal{F}_{ab}$ と書く。このとき
$$
H^i(\mathcal{C}, \mathcal{F}_{ab})
=
H^i(\mathcal{C}, \mathcal{F})
$$
また、$\mathcal{C}$ の任意の対象 $U$ に対して
$$
H^i(U, \mathcal{F}_{ab})
=
H^i(U, \mathcal{F}).
$$
左辺は $\textit{Ab}(\mathcal{C})$ で計算したコホモロジーであり、右辺は
$\textit{Mod}(\mathcal{O})$ で計算したコホモロジーである。
\end{lemma}

\begin{proof}
Derived Categories, Lemma \ref{derived-lemma-higher-derived-functors} により、
$\delta$-関手 $(\mathcal{F} \mapsto H^p(U, \mathcal{F}))_{p \geq 0}$ は普遍的である。
関手
$\textit{Mod}(\mathcal{O}) \to \textit{Ab}(\mathcal{C})$、
$\mathcal{F} \mapsto \mathcal{F}_{ab}$ は完全である。従って
$(\mathcal{F} \mapsto H^p(U, \mathcal{F}_{ab}))_{p \geq 0}$
も $\delta$-関手である。これが
$(\mathcal{F} \mapsto H^p(U, \mathcal{F}_{ab}))_{p \geq 0}$
普遍的であることを示せばよい。そうすれば普遍 $\delta$-関手の一意性により
補題の第二の等式が従う（Homology, Lemma
\ref{homology-lemma-uniqueness-universal-delta-functor} 参照）。
$\textit{Mod}(\mathcal{O})$ は十分な入射対象をもつので、任意の
$\textit{Mod}(\mathcal{O})$ の入射対象 $\mathcal{I}$ に対して
$H^i(U, \mathcal{I}_{ab}) = 0$ を示せば十分である（Homology, Lemma
\ref{homology-lemma-efface-implies-universal} 参照）。

\medskip\noindent
$\textit{Mod}(\mathcal{O})$ の入射対象 $\mathcal{I}$ をとる。Lemma
\ref{sites-cohomology-lemma-cech-vanish-collection} を $\mathcal{F} = \mathcal{I}$、
$\mathcal{B} = \mathcal{C}$、$\text{Cov} = \text{Cov}_\mathcal{C}$ として適用する。
この補題の仮定 (3) は Lemma \ref{sites-cohomology-lemma-injective-module-trivial-cech} により
満たされる。従って $\mathcal{C}$ のすべての対象 $U$ について
$H^i(U, \mathcal{I}_{ab}) = 0$ である。

\medskip\noindent
 $\mathcal{C}$ が終対象をもつなら、これから第一の等式も従う。もたない場合は、
Sites, Lemma \ref{sites-lemma-topos-good-site} により、環付きトポス
$(\Sh(\mathcal{C}), \mathcal{O})$ は、基礎サイトが終対象をもつ環付きトポスと
同値である。従って補題が従う。
\end{proof}

\begin{lemma}
\label{sites-cohomology-lemma-cohomology-products}
$\mathcal{C}$ をサイトとする。$I$ を集合とし、各 $i \in I$ に対して
$\mathcal{F}_i$ を $\mathcal{C}$ 上のアーベル層とする。$U \in \Ob(\mathcal{C})$
とする。標準写像
$$
H^p(U, \prod\nolimits_{i \in I} \mathcal{F}_i)
\longrightarrow
\prod\nolimits_{i \in I} H^p(U, \mathcal{F}_i)
$$
$p = 0$ では同型であり、$p = 1$ では単射である。
\end{lemma}

\begin{proof}
 $p = 0$ の主張は、層の積が基礎前層の積に等しいことから従う（Sites,
Lemma \ref{sites-lemma-limit-sheaf} 参照）。$p = 1$ を示す。
$\mathcal{F} = \prod \mathcal{F}_i$ とおく。$\xi \in H^1(U, \mathcal{F})$ が
$\prod H^1(U, \mathcal{F}_i)$ で 0 に写るとする。コホモロジーの局所性
（Lemma \ref{sites-cohomology-lemma-kill-cohomology-class-on-covering} 参照）により、すべての
$j$ で $\xi|_{U_j} = 0$ となる被覆 $\mathcal{U} = \{U_j \to U\}$ が存在する。
Lemma \ref{sites-cohomology-lemma-cech-h1} により、これは $\xi$ が次の元から来ることを意味する。
$\check \xi \in \check H^1(\mathcal{U}, \mathcal{F})$.
写像
$\check H^1(\mathcal{U}, \mathcal{F}_i) \to H^1(U, \mathcal{F}_i)$
はすべての $i$ で単射であり（Lemma \ref{sites-cohomology-lemma-cech-h1}）、また $\xi$ の像は
$\prod H^1(U, \mathcal{F}_i)$ で 0 なので、その像
$\check H^1(\mathcal{U}, \mathcal{F}_i)$ において $\check \xi_i = 0$ である。
しかし $\mathcal{F} = \prod \mathcal{F}_i$ なので、
$\check{\mathcal{C}}^\bullet(\mathcal{U}, \mathcal{F})$ は複体
$\check{\mathcal{C}}^\bullet(\mathcal{U}, \mathcal{F}_i)$,
の積である。従って Homology, Lemma
\ref{homology-lemma-product-abelian-groups-exact} により
$\check \xi = 0$ を得る。これは所望の結論である。
\end{proof}

\begin{lemma}
\label{sites-cohomology-lemma-restriction-along-monomorphism-surjective}
$(\mathcal{C}, \mathcal{O})$ を環付きサイトとし、$a : U' \to U$ を
$\mathcal{C}$ のモノ射とする。このとき任意の入射 $\mathcal{O}$-加群
$\mathcal{I}$ に対して制限写像 $\mathcal{I}(U) \to \mathcal{I}(U')$ は全射である。
\end{lemma}

\begin{proof}
$j : \mathcal{C}/U \to \mathcal{C}$ および
$j' : \mathcal{C}/U' \to \mathcal{C}$ を局所化射とする（Modules on Sites,
Section \ref{sites-modules-section-localize} 参照）。$j_!$ は制限の左随伴なので、
任意の $\mathcal{O}$-加群層 $\mathcal{F}$ に対して
$$
\Hom_\mathcal{O}(j_!\mathcal{O}_U, \mathcal{F})
=
\Hom_{\mathcal{O}_U}(\mathcal{O}_U, \mathcal{F}|_U)
=
\mathcal{F}(U)
$$
同様に、層 $j'_!\mathcal{O}_{U'}$ は関手 $\mathcal{F} \mapsto \mathcal{F}(U')$
を表現する。さらに次で、$\mathcal{O}$-加群の標準写像
$$
j'_!\mathcal{O}_{U'} \longrightarrow j_!\mathcal{O}_U
$$
を構成する。この写像は Yoneda の補題（Categories, Lemma
\ref{categories-lemma-yoneda}）を通じて制限写像
$\mathcal{F}(U) \to \mathcal{F}(U')$ に対応する。表示した加群の写像が単射で
あることを示せば十分である（Homology, Lemma
\ref{homology-lemma-characterize-injectives} 参照）。

\medskip\noindent
写像を構成するには、$\mathcal{C}$ の対象 $V$ に次の $\mathcal{O}(V)$-加群を
割り当てる前層の間の写像を構成すれば十分である。
$$
\bigoplus\nolimits_{\varphi' \in \Mor_\mathcal{C}(V, U')} \mathcal{O}(V)
  \quad\text{および}\quad
\bigoplus\nolimits_{\varphi \in \Mor_\mathcal{C}(V, U)} \mathcal{O}(V)
$$
（Modules on Sites, Lemma \ref{sites-modules-lemma-extension-by-zero} 参照）。
$\varphi'$ に対応する直和因子を $\varphi = a \circ \varphi'$ に対応する因子へ
$\mathcal{O}(V)$ 上の恒等写像で送る写像をとる。$a$ はモノ射なので、この写像は
単射である。層化は完全なので、主張が従う。
\end{proof}






\section{完全非輪状層}
\label{sites-cohomology-section-limp}

\noindent
$(\mathcal{C}, \mathcal{O})$ を環付きサイトとする。$K$ を $\mathcal{C}$ 上の集合の
前層とする（アーベル層と区別するため、ここでは意図的にローマン体の大文字を
用いる）。$\mathcal{O}$-加群層 $\mathcal{F}$ に対し、次のようにおく。
$$
\mathcal{F}(K) =
\Mor_{\textit{PSh}(\mathcal{C})}(K, \mathcal{F}) =
\Mor_{\Sh(\mathcal{C})}(K^\#, \mathcal{F})
$$
関手 $\mathcal{F} \mapsto \mathcal{F}(K)$ は左完全関手
$\textit{Mod}(\mathcal{O}) \to \textit{Ab}$ である。従ってその右導来関手が存在する。
これらを $H^p(K, \mathcal{F})$ と書き、$H^0(K, \mathcal{F}) = \mathcal{F}(K)$
とする。

\medskip\noindent
いくつか注意を述べる。
\begin{enumerate}
\item $\mathcal{F}(K) = \mathcal{F}(K^\#)$ なので、
$H^p(K, \mathcal{F}) = H^p(K^\#, \mathcal{F})$ である。上の定義で $K$ を
前層として許したのは、純粋に記法上の便宜である。
\item $K = h_U$ または $K = h_U^\#$（$U$ は $\mathcal{C}$ の対象）とする。
このとき $\Mor_{\Sh(\mathcal{C})}(h_U^\#, \mathcal{F}) = \mathcal{F}(U)$ なので、
$H^p(K, \mathcal{F}) = H^p(U, \mathcal{F})$ である（Sites, Section
\ref{sites-section-representable-sheaves} 参照）。
\item $\mathcal{O} = \mathbf{Z}$（定数層）なら、コホモロジー群は関手
$H^p(K, - ) : \textit{Ab}(\mathcal{C}) \to \textit{Ab}$
となる。この場合 $\textit{Mod}(\mathcal{O}) = \textit{Ab}(\mathcal{C})$ だからである。
\end{enumerate}
既に証明した結果のいくつかを、この言葉で言い換えられる。

\begin{lemma}
\label{sites-cohomology-lemma-compute-cohomology-on-sheaf-sets}
$(\mathcal{C}, \mathcal{O})$ を環付きサイトとする。$K$ を $\mathcal{C}$ 上の集合の
前層とし、$\mathcal{F}$ を $\mathcal{O}$-加群とする。台となるアーベル群の層を
$\mathcal{F}_{ab}$ と書く。このとき $H^p(K, \mathcal{F}) = H^p(K, \mathcal{F}_{ab})$
である。
\end{lemma}

\begin{proof}
 $K$ を層化して $K$ が層であると仮定してよい。$H^p(K, \mathcal{F})$ と
$H^p(K, \mathcal{F}_{ab})$ は基礎サイトではなくトポスだけに依存することに注意する。
従って Sites, Lemma \ref{sites-lemma-topos-good-site} により、$K = h_U$（$U$ は
$\mathcal{C}$ の対象）となる「より大きい」サイトへ $\mathcal{C}$ を取り替えてよい。
この場合は Lemma \ref{sites-cohomology-lemma-cohomology-modules-abelian-agree} から従う。
\end{proof}

\begin{lemma}
\label{sites-cohomology-lemma-cech-to-cohomology-sheaf-sets}
$\mathcal{C}$ をサイトとする。$K' \to K$ を $\mathcal{C}$ 上の集合の前層の射で、
その層化が全射となるものとする。$K'_p = K' \times_K \ldots \times_K K'$（因子は
$p + 1$ 個）とおく。任意のアーベル層 $\mathcal{F}$ に対し、
$E_1^{p, q} = H^q(K'_p, \mathcal{F})$ で $H^{p + q}(K, \mathcal{F})$ に収束する
スペクトル系列が存在する。
\end{lemma}

\begin{proof}
層化は完全なので、
$(K_p')^\#$ は
$(K')^\# \times_{K^\#} \ldots \times_{K^\#} (K')^\#$
（因子は $p + 1$ 個）。従って $K,K'$ を層化で置き換え、$K' \to K$ が層の
全射であると仮定してよい。Sites, Lemma \ref{sites-lemma-topos-good-site} のように
$\mathcal{C}$ を「より大きい」サイトへ取り替えれば、$K,K'$ は $\mathcal{C}$ の
対象であり、$\mathcal{U} = \{K' \to K\}$ が被覆であると仮定できる。このとき
Lemma \ref{sites-cohomology-lemma-cech-spectral-sequence} の {\v C}ech からコホモロジーへの
スペクトル系列が得られ、その $E_1$ ページは補題の記述どおりである。
\end{proof}

\begin{lemma}
\label{sites-cohomology-lemma-cohomology-on-sheaf-sets}
$\mathcal{C}$ をサイトとし、$K$ を $\mathcal{C}$ 上の集合の層とする。トポスの射
$j : \Sh(\mathcal{C}/K) \to \Sh(\mathcal{C})$（Sites, Lemma
\ref{sites-lemma-localize-topos-site} 参照）を考える。
このとき $j^{-1}$ は入射対象を保ち、$\mathcal{C}$ 上の任意のアーベル層
$\mathcal{F}$ に対して $H^p(K, \mathcal{F}) = H^p(\mathcal{C}/K, j^{-1}\mathcal{F})$
である。
\end{lemma}

\begin{proof}
Sites, Lemmas \ref{sites-lemma-localize-topos} および
\ref{sites-lemma-localize-topos-site} により、トポスの射 $j$ は局所化と同値である。
従って Lemma \ref{sites-cohomology-lemma-cohomology-of-open} から従う。
\end{proof}

\noindent
Lemma \ref{sites-cohomology-lemma-compute-cohomology-on-sheaf-sets} を念頭に置けば、次の定義は
環付きサイト上の加群層に対しても「正しい定義」であることが分かる。

\begin{definition}
\label{sites-cohomology-definition-limp}
$\mathcal{C}$ をサイトとする。
アーベル層 $\mathcal{F}$ が
{\it 完全非輪状}\footnote{この用語は \cite[Vbis, Proposition 1.3.10]{SGA4}
で用いられているが、標準的な記法ではないと思われる。
\cite[V, Definition 4.1]{SGA4} ではこの性質を ``flasque'' と呼んでいるが、
位相空間上の flasque 層に対する本書の定義と衝突するため、この語は使えない。
よりよい提案があれば
\href{mailto:stacks.project@gmail.com}{stacks.project@gmail.com}
にメールしてほしい。}
とは、任意の集合の層 $K$ に対して
$H^p(K, \mathcal{F}) = 0$ がすべての $p \geq 1$ で成り立つことをいう。
\end{definition}

\noindent
完全非輪状であることは内在的な性質、すなわちトポスの同値で保たれる性質である
ことは明らかである。完全非輪状層はサイトのすべての対象上で高次コホモロジーが
消えるが、一般には完全非輪状という条件の方が真に強い。次は、ときに有用な
完全非輪状層の特徴付けである。

\begin{lemma}
\label{sites-cohomology-lemma-characterize-limp}
$\mathcal{C}$ をサイトとし、$\mathcal{F}$ をアーベル層とする。次を仮定する。
\begin{enumerate}
\item $p > 0$ および $U \in \Ob(\mathcal{C})$ に対して
$H^p(U, \mathcal{F}) = 0$ である。
\item 集合の層の任意の全射 $K' \to K$ に対し、拡張 {\v C}ech 複体
$$
0 \to H^0(K, \mathcal{F}) \to H^0(K', \mathcal{F}) \to
H^0(K' \times_K K', \mathcal{F}) \to \ldots
$$
は完全である。
\end{enumerate}
このとき $\mathcal{F}$ は完全非輪状である（逆も成り立つ）。
\end{lemma}

\begin{proof}
仮定 (1) により、すべての $p > 0$ および $\mathcal{C}$ の対象 $U$ に対して
$H^p(h_U^\#, g^{-1}\mathcal{I}) = 0$ である。$K = \coprod K_i$ が
$\mathcal{C}$ 上の集合の層の余積ならば
$H^p(K, g^{-1}\mathcal{I}) = \prod H^p(K_i, g^{-1}\mathcal{I})$ であることに
注意する。任意の集合の層 $K$ に対して全射
$$
K' = \coprod h_{U_i}^\# \longrightarrow K
$$
が存在する（Sites, Lemma \ref{sites-lemma-sheaf-coequalizer-representable} 参照）。
従って次を得る。(*) 任意の集合の層 $K$ に対して、$p > 0$ で
$H^p(K', \mathcal{F}) = 0$ となる集合の層の全射 $K' \to K$ が存在する。
(*) と条件 (2) が $\mathcal{F}$ の完全非輪状性を導くことを示す。
(*) と (2) は基礎サイトではなくトポス $\Sh(\mathcal{C})$ の対象としての
$\mathcal{F}$ にのみ依存することに注意する（証明の残りでは性質 (1) を用いない）。

\medskip\noindent
 (*) と (2) から、$n \geq 0$ に関する帰納法で次の帰納法の仮定 $IH_n$ を
証明する。
$H^p(K, \mathcal{F}) = 0$ がすべての $0 < p \leq n$ および集合の層 $K$ に
対して成り立つ。$IH_0$ は成り立つことに注意する。$IH_n$ を仮定し、集合の層
$K$ をとる。$p > 0$ で $H^p(K', \mathcal{F}) = 0$ となる全射
$K' \to K$ をとる。次をもつスペクトル系列がある。
$$
E_1^{p, q} = H^q(K'_p, \mathcal{F})
$$
は $H^{p + q}(K, \mathcal{F})$ への収束を与える（Lemma
\ref{sites-cohomology-lemma-cech-to-cohomology-sheaf-sets} 参照）。$IH_n$ により
$0 < q \leq n$ なら $E_1^{p, q} = 0$ であり、仮定 (2) により $p > 0$ なら
$E_2^{p, 0} = 0$ である。さらに $K'$ の選択により
$H^q(K', \mathcal{F}) = 0$ なので、$q > 0$ なら $E_1^{0, q} = 0$ である。
従って $p + q = n + 1$ を満たすすべての $E_2^{p, q}$ が 0 となり、
$H^{n + 1}(K, \mathcal{F}) = 0$ が従う。
\end{proof}







\section{Leray スペクトル系列}
\label{sites-cohomology-section-leray}

\noindent
Leray スペクトル系列の存在を証明する鍵は次の補題である。

\begin{lemma}
\label{sites-cohomology-lemma-direct-image-injective-sheaf}
$f : (\Sh(\mathcal{C}), \mathcal{O}_\mathcal{C}) \to
(\Sh(\mathcal{D}), \mathcal{O}_\mathcal{D})$ を環付きトポスの射とする。
$\textit{Mod}(\mathcal{O}_\mathcal{C})$ の任意の入射対象 $\mathcal{I}$ に対して、
押し出し $f_*\mathcal{I}$ は完全非輪状である。
\end{lemma}

\begin{proof}
 $K$ を $\mathcal{D}$ 上の集合の層とする。
Modules on Sites, Lemma
\ref{sites-modules-lemma-morphism-ringed-topoi-comes-from-morphism-ringed-sites}
により、$\mathcal{C}, \mathcal{D}$ を「より大きい」サイトへ取り替え、$f$ が
連続関手 $u : \mathcal{D} \to \mathcal{C}$ により誘導される環付きサイトの射から
来るようにできる。また $K = h_V$（$V$ は $\mathcal{D}$ の対象）と仮定できる。

\medskip\noindent
従って $f$ が環付きサイトの射で与えられる場合、$\mathcal{D}$ のすべての対象
$V$ と $q > 0$ に対して $H^q(V, f_*\mathcal{I}) = 0$ を示せばよい。
$\mathcal{V} = \{V_j \to V\}$ を $\mathcal{D}$ の任意の被覆とする。
$u$ は連続なので $\mathcal{U} = \{u(V_j) \to u(V)\}$ は $\mathcal{C}$ の被覆である。
このとき {\v C}ech 複体について等式
$$
\check{\mathcal{C}}^\bullet(\mathcal{V}, f_*\mathcal{I})
=
\check{\mathcal{C}}^\bullet(\mathcal{U}, \mathcal{I})
$$
を得る（$f_*$ の定義による）。Lemma
\ref{sites-cohomology-lemma-injective-module-trivial-cech} により、この複体の正次数コホモロジーは
消える。Lemma \ref{sites-cohomology-lemma-cech-vanish-collection} を適用すれば主張が従う。
\end{proof}

\noindent
平坦射では関手 $f_*$ は入射加群を保つ。特に関手
$f_* : \textit{Ab}(\mathcal{C}) \to \textit{Ab}(\mathcal{D})$ は常に
入射アーベル層を入射アーベル層へ写す。

\begin{lemma}
\label{sites-cohomology-lemma-pushforward-injective-flat}
$f : (\Sh(\mathcal{C}), \mathcal{O}_\mathcal{C}) \to
(\Sh(\mathcal{D}), \mathcal{O}_\mathcal{D})$ を環付きトポスの射とする。
$f$ が平坦なら、任意の入射 $\mathcal{O}_\mathcal{C}$-加群 $\mathcal{I}$ に対して
$f_*\mathcal{I}$ は入射 $\mathcal{O}_\mathcal{D}$-加群である。
\end{lemma}

\begin{proof}
この場合、関手 $f^*$ は完全である（Modules on Sites, Lemma
\ref{sites-modules-lemma-flat-pullback-exact} 参照）。従って Homology, Lemma
\ref{homology-lemma-adjoint-preserve-injectives} から従う。
\end{proof}

\begin{lemma}
\label{sites-cohomology-lemma-limp-acyclic}
$(\Sh(\mathcal{C}), \mathcal{O}_\mathcal{C})$ を環付きトポスとする。
完全非輪状層は次の関手に関して右非輪状である。
\begin{enumerate}
\item $\mathcal{C}$ の任意の対象 $U$ に対する関手 $H^0(U, -)$。
\item 任意の集合の前層 $K$ に対する関手 $\mathcal{F} \mapsto \mathcal{F}(K)$。
\item 大域切断の関手 $\Gamma(\mathcal{C}, -)$。
\item 環付きトポスの任意の射
$f : (\Sh(\mathcal{C}), \mathcal{O}_\mathcal{C}) \to
(\Sh(\mathcal{D}), \mathcal{O}_\mathcal{D})$ に対する関手 $f_*$。
\end{enumerate}
\end{lemma}

\begin{proof}
 (2) は完全非輪状層の定義である。(1) は完全非輪状層の定義に続く議論で
述べたように (2) の帰結である。(3) は $K = e$ を $\Sh(\mathcal{C})$ の終対象と
した (2) の特別な場合である。

\medskip\noindent
 (4) を証明する。Modules on Sites, Lemma
\ref{sites-modules-lemma-morphism-ringed-topoi-comes-from-morphism-ringed-sites}
により $f$ がサイトの射で与えられると仮定してよい。この場合、完全非輪状層の
$R^if_*$（$i > 0$）は、Lemma \ref{sites-cohomology-lemma-higher-direct-images} における
高次直接像の記述により 0 である。
\end{proof}

\begin{remark}
\label{sites-cohomology-remark-before-Leray}
以上の結果から、次が分かる。
Derived Categories, Lemma \ref{derived-lemma-compose-derived-functors}
は多くの状況に適用できる。例えば、環付きトポスの射
$f : (\Sh(\mathcal{C}), \mathcal{O}_\mathcal{C}) \to
(\Sh(\mathcal{D}), \mathcal{O}_\mathcal{D})$ に対して
$$
R\Gamma(\mathcal{D}, Rf_*\mathcal{F}) = R\Gamma(\mathcal{C}, \mathcal{F})
$$
任意の $\mathcal{O}_\mathcal{C}$-加群層 $\mathcal{F}$ に対して成り立つ。
実際、入射 $\mathcal{O}_\mathcal{C}$-加群 $\mathcal{I}$ に対し、
$\mathcal{O}_\mathcal{D}$-加群 $f_*\mathcal{I}$ は
Lemma \ref{sites-cohomology-lemma-direct-image-injective-sheaf}
で完全非輪状であり、完全非輪状層は Lemma \ref{sites-cohomology-lemma-limp-acyclic} により
$\Gamma(\mathcal{D}, -)$ に関して非輪状である。
\end{remark}

\begin{lemma}[Leray スペクトル系列]
\label{sites-cohomology-lemma-Leray}
$f : (\Sh(\mathcal{C}), \mathcal{O}_\mathcal{C}) \to
(\Sh(\mathcal{D}), \mathcal{O}_\mathcal{D})$ を環付きトポスの射とする。
$\mathcal{F}^\bullet$ を $\mathcal{O}_\mathcal{C}$-加群の下に有界な複体とする。
スペクトル系列
$$
E_2^{p, q} = H^p(\mathcal{D}, R^qf_*(\mathcal{F}^\bullet))
$$
は $H^{p + q}(\mathcal{C}, \mathcal{F}^\bullet)$ に収束する。
\end{lemma}

\begin{proof}
これは関手の合成
$\Gamma(\mathcal{C}, -) = \Gamma(\mathcal{D}, -) \circ f_*$
から得られる Grothendieck スペクトル系列である（Derived Categories, Lemma
\ref{derived-lemma-grothendieck-spectral-sequence} 参照）。同補題の仮定が満たされることは
Derived Categories, Lemma \ref{derived-lemma-grothendieck-spectral-sequence} および
Lemmas \ref{sites-cohomology-lemma-direct-image-injective-sheaf} および
\ref{sites-cohomology-lemma-limp-acyclic} を参照。
\end{proof}

\begin{lemma}
\label{sites-cohomology-lemma-apply-Leray}
$f : (\Sh(\mathcal{C}), \mathcal{O}_\mathcal{C}) \to
(\Sh(\mathcal{D}), \mathcal{O}_\mathcal{D})$ を環付きトポスの射とする。
$\mathcal{F}$ を $\mathcal{O}_\mathcal{C}$-加群とする。
\begin{enumerate}
\item $q > 0$ に対して $R^qf_*\mathcal{F} = 0$ ならば、すべての $p$ に対して
$H^p(\mathcal{C}, \mathcal{F}) = H^p(\mathcal{D}, f_*\mathcal{F})$ である。
\item すべての $q$ および $p > 0$ に対して
$H^p(\mathcal{D}, R^qf_*\mathcal{F}) = 0$ ならば、すべての $q$ に対して
$H^q(\mathcal{C}, \mathcal{F}) = H^0(\mathcal{D}, R^qf_*\mathcal{F})$ である。
\end{enumerate}
\end{lemma}

\begin{proof}
これは Leray スペクトル系列を収束させる二つの単純な条件である。
これらの事実は（スペクトル系列を用いずに）直接証明することもでき、
層のコホモロジーにおけるよい演習となる。
\end{proof}

\begin{lemma}[相対 Leray スペクトル系列]
\label{sites-cohomology-lemma-relative-Leray}
\noindent
$f : (\Sh(\mathcal{C}), \mathcal{O}_\mathcal{C}) \to
(\Sh(\mathcal{D}), \mathcal{O}_\mathcal{D})$
および
$g : (\Sh(\mathcal{D}), \mathcal{O}_\mathcal{D}) \to
(\Sh(\mathcal{E}), \mathcal{O}_\mathcal{E})$
を環付きトポスの射とする。$\mathcal{F}$ を $\mathcal{O}_\mathcal{C}$-加群とする。
次のスペクトル系列が存在する。
$$
E_2^{p, q} = R^pg_*(R^qf_*\mathcal{F})
$$
は $R^{p + q}(g \circ f)_*\mathcal{F}$ に収束する。
このスペクトル系列は $\mathcal{F}$ に関して関手的であり、
$\mathcal{O}_\mathcal{C}$-加群の下に有界な複体に対する版も存在する。
\end{lemma}

\begin{proof}
これは関手の合成に対する Grothendieck スペクトル系列である（Derived Categories,
Lemma \ref{derived-lemma-grothendieck-spectral-sequence} および
Lemmas \ref{sites-cohomology-lemma-direct-image-injective-sheaf},
\ref{sites-cohomology-lemma-limp-acyclic} 参照）。
\end{proof}







\section{基底変換写像}
\label{sites-cohomology-section-base-change-map}

\noindent
この節では基底変換写像をいくつかの場合に構成する。一般の場合は
Remark \ref{sites-cohomology-remark-base-change} で扱う。この節では、環付きサイトの平坦射による
基底変換の場合に制限することで、導来引き戻しの使用を避ける。
結果を述べる前に、導来圏上の平坦引き戻しについて説明する。次を仮定する。
$g : (\Sh(\mathcal{C}), \mathcal{O}_\mathcal{C})
\to (\Sh(\mathcal{D}), \mathcal{O}_\mathcal{D})$
を環付きトポスの平坦射とする。Modules on Sites, Lemma
\ref{sites-modules-lemma-flat-pullback-exact} により、関手
$g^* : \textit{Mod}(\mathcal{O}_\mathcal{D}) \to
\textit{Mod}(\mathcal{O}_\mathcal{C})$ は完全である。従って導来関手
$$
g^* : D(\mathcal{O}_\mathcal{D}) \to D(\mathcal{O}_\mathcal{C})
$$
これは $D(\mathcal{O}_\mathcal{D})$ の対象の代表を単に引き戻すことで計算できる
（Derived Categories, Lemma \ref{derived-lemma-right-derived-exact-functor} 参照）。
この関手は上に有界および下に有界な部分圏を保つ。従って以下では、$Lg^*$ ではなく
$g^*$ でこの関手を表す。

\begin{lemma}
\label{sites-cohomology-lemma-base-change-map-flat-case}
次を考える。
$$
\xymatrix{
(\Sh(\mathcal{C}'), \mathcal{O}_{\mathcal{C}'})
\ar[r]_{g'} \ar[d]_{f'} &
(\Sh(\mathcal{C}), \mathcal{O}_\mathcal{C}) \ar[d]^f \\
(\Sh(\mathcal{D}'), \mathcal{O}_{\mathcal{D}'})
\ar[r]^g &
(\Sh(\mathcal{D}), \mathcal{O}_\mathcal{D})
}
$$
を環付きトポスの可換図式とする。$\mathcal{F}^\bullet$ を
$\mathcal{O}_\mathcal{C}$-加群の下に有界な複体とする。$g$ と $g'$ はともに平坦と
仮定する。このとき、標準的な基底変換写像
$$
g^*Rf_*\mathcal{F}^\bullet
\longrightarrow
R(f')_*(g')^*\mathcal{F}^\bullet
$$
が $D^{+}(\mathcal{O}_{\mathcal{D}'})$ に存在する。
\end{lemma}

\begin{proof}
入射分解 $\mathcal{F}^\bullet \to \mathcal{I}^\bullet$ と
$(g')^*\mathcal{F}^\bullet \to \mathcal{J}^\bullet$ を選ぶ。
Lemma \ref{sites-cohomology-lemma-pushforward-injective-flat} により、
$(g')_*\mathcal{J}^\bullet$ は $R(g')_*(g')^*\mathcal{F}^\bullet$ を表す入射対象の
複体である。従って Derived Categories, Lemmas
\ref{derived-lemma-morphisms-lift} および
\ref{derived-lemma-morphisms-equal-up-to-homotopy} により、図式の射 $\beta$ が
$$
\xymatrix{
(g')_*(g')^*\mathcal{F}^\bullet \ar[r] &
(g')_*\mathcal{J}^\bullet \\
\mathcal{F}^\bullet \ar[u]^{adjunction} \ar[r] &
\mathcal{I}^\bullet \ar[u]_\beta
}
$$
存在し、ホモトピーを除いて一意である。$\mathcal{D}$ へ押し出すと
$$
f_*\beta :
f_*\mathcal{I}^\bullet
\longrightarrow
f_*(g')_*\mathcal{J}^\bullet
=
g_*(f')_*\mathcal{J}^\bullet
$$
 $g^*$ と $g_*$ の随伴により、複体の射
$g^*f_*\mathcal{I}^\bullet \to (f')_*\mathcal{J}^\bullet$
を得る。この射はホモトピーを除いて一意であることに注意する。というのも、全過程
で選択したのは射 $\beta$ だけであり、すべて複体のレベルで行ったからである。
\end{proof}








\section{コホモロジーと余極限}
\label{sites-cohomology-section-limits}

\noindent
$(\mathcal{C}, \mathcal{O})$ を環付きサイトとする。
$\mathcal{I} \to \textit{Mod}(\mathcal{O})$, $i \mapsto \mathcal{F}_i$ を添字圏
$\mathcal{I}$ 上の図式とする（Categories, Section \ref{categories-section-limits} 参照）。
各 $i$ に対して標準写像 $\mathcal{F}_i \to \colim_i \mathcal{F}_i$ があり、
これはコホモロジー上の写像を誘導する。従って標準写像
$$
\colim_i H^p(U, \mathcal{F}_i)
\longrightarrow
H^p(U, \colim_i \mathcal{F}_i)
$$
がすべての $p \geq 0$ および $\mathcal{C}$ のすべての対象 $U$ に対して得られる。
一般には、$p = 0$ であってもこれらの写像は同型ではない。

\medskip\noindent
次の補題はコホモロジーに対する Sites, Lemma
\ref{sites-lemma-directed-colimits-sections} の類似である。

\begin{lemma}
\label{sites-cohomology-lemma-colim-works-over-collection}
$\mathcal{C}$ をサイトとする。$\text{Cov}_\mathcal{C}$ を $\mathcal{C}$ の被覆全体の集合
とする（Sites, Definition \ref{sites-definition-site} 参照）。
$\mathcal{B} \subset \Ob(\mathcal{C})$ および $\text{Cov} \subset \text{Cov}_\mathcal{C}$
を部分集合とする。次を仮定する。
\begin{enumerate}
\item すべての $\mathcal{U} \in \text{Cov}$ に対し、$I$ が有限で
$\mathcal{U} = \{U_i \to U\}_{i \in I}$、$U, U_i \in \mathcal{B}$、さらにすべての
$U_{i_0} \times_U \ldots \times_U U_{i_p} \in \mathcal{B}$ が成り立つ。
\item すべての $U \in \mathcal{B}$ に対し、$\text{Cov}$ に現れる $U$ の被覆は
$U$ の被覆の共終系である。
\end{enumerate}
このとき写像
$$
\colim_i H^p(U, \mathcal{F}_i)
\longrightarrow
H^p(U, \colim_i \mathcal{F}_i)
$$
は、すべての $p \geq 0$、すべての $U \in \mathcal{B}$、およびすべての
$\textit{Ab}(\mathcal{C})$ のフィルター付き図式 $\mathcal{I} \to \textit{Ab}(\mathcal{C})$
に対して同型である。
\end{lemma}

\begin{proof}
この補題は $p$ に関する帰納法で証明する。(1) では $\mathcal{U} \in \text{Cov}$ の
被覆が有限であることを要求しているので、$\mathcal{B}$ のすべての元は準コンパクトで
あることに注意する。(2) と (1) により、任意の $U \in \mathcal{B}$ は Sites, Lemma
\ref{sites-lemma-directed-colimits-sections} (4) の仮定を満たす。従って $p = 0$ では
結果が成り立つ。ここで補題が $p$ について成り立つと仮定し、$p + 1$ について証明する。

\medskip\noindent
フィルター付き図式
$\mathcal{F} : \mathcal{I} \to \textit{Ab}(\mathcal{C})$,
$i \mapsto \mathcal{F}_i$
を選ぶ。$\textit{Ab}(\mathcal{C})$ は関手的な入射埋め込みをもつので、Injectives,
Theorem \ref{injectives-theorem-sheaves-injectives} により、フィルター付き図式の射
$\mathcal{F} \to \mathcal{I}$
で、各 $\mathcal{F}_i \to \mathcal{I}_i$ がアーベル層から入射アーベル層への入射写像
となるものを取れる。余核を $\mathcal{Q}_i$ と書くと、短完全列
$$
0 \to
\mathcal{F}_i \to
\mathcal{I}_i \to
\mathcal{Q}_i \to 0.
$$
層の余極限は前層レベルの余極限の層化であり、層化は完全で、アーベル群のフィルター
付き余極限も完全である（Algebra, Lemma \ref{algebra-lemma-directed-colimit-exact} 参照）。
従って次の列
$$
0 \to
\colim_i \mathcal{F}_i \to
\colim_i \mathcal{I}_i \to
\colim_i \mathcal{Q}_i \to 0.
$$
も短完全列である。ここで、すべての $U \in \mathcal{B}$ および $q \geq 1$ に対して
$H^q(U, \colim_i \mathcal{I}_i) = 0$ であると主張する。この主張をひとまず認め、次の
図式を考える。
$$
\xymatrix{
\colim_i H^p(U, \mathcal{I}_i) \ar[d] \ar[r] &
\colim_i H^p(U, \mathcal{Q}_i) \ar[d] \ar[r] &
\colim_i H^{p + 1}(U, \mathcal{F}_i) \ar[d] \ar[r] &
0 \ar[d] \\
H^p(U, \colim_i \mathcal{I}_i) \ar[r] &
H^p(U, \colim_i \mathcal{Q}_i) \ar[r] &
H^{p + 1}(U, \colim_i \mathcal{F}_i) \ar[r] &
0
}
$$
右下隅の 0 は主張から従い、右上隅の 0 は層 $\mathcal{I}_i$ が入射であることから従う。
上段は Algebra, Lemma \ref{algebra-lemma-directed-colimit-exact} の適用により完全で
ある。従って snake lemma により $p + 1$ の結果が得られる。

\medskip\noindent
残りは主張を示すことである。Lemma \ref{sites-cohomology-lemma-cech-vanish-collection} を用いる。
$p = 0$ の結果により、$\mathcal{U} \in \text{Cov}$ に対して
$$
\check{\mathcal{C}}^\bullet(\mathcal{U}, \colim_i \mathcal{I}_i)
=
\colim_i \check{\mathcal{C}}^\bullet(\mathcal{U}, \mathcal{I}_i)
$$
すべての $U_{j_0} \times_U \ldots \times_U U_{j_p}$ が $\mathcal{B}$ に属するからである。
Lemma \ref{sites-cohomology-lemma-injective-trivial-cech} により、Cech 複体の余極限に現れる各複体は
次数 $\geq 1$ で非輪状である。従って Algebra, Lemma
\ref{algebra-lemma-directed-colimit-exact} により、次の Cech 複体も
$\check{\mathcal{C}}^\bullet(\mathcal{U}, \colim_i \mathcal{I}_i)$
は次数 $\geq 1$ で非輪状である。すなわち
$\check{H}^p(\mathcal{U}, \colim_i \mathcal{I}_i) = 0$
がすべての $p \geq 1$ で成り立つ。従って Lemma
\ref{sites-cohomology-lemma-cech-vanish-collection} の仮定が満たされ、主張が従う。
\end{proof}

\begin{lemma}
\label{sites-cohomology-lemma-colim-global}
$\mathcal{C}$ をサイトとする。$S \subset \Ob(\Sh(\mathcal{C}))$ を部分集合とし、
$*$ を $\Sh(\mathcal{C})$ の終対象と書く。次を仮定する。
\begin{enumerate}
\item ある $K \in S$ に対して $K \to *$ が全射である。
\item $K \in S$ への層の全射 $\mathcal{F} \to K$ が与えられたとき、$K' \in S$ と
$K' \to \mathcal{F}$ なる射で、合成 $K' \to K$ が全射となるものが存在する。
\item $K, K' \in S$ に対して、$K'' \in S$ かつ全射 $K'' \to K \times K'$ が存在する。
\item $K, K' \in S$ に対する $a, b : K \to K'$ が与えられたとき、$K'' \in S$ かつ
全射 $K'' \to \text{Equalizer}(a, b)$ が存在する。
\item すべての $K \in S$ は準コンパクトである
（Sites, Definition \ref{sites-definition-quasi-compact-topos}）。
\end{enumerate}
このとき、すべての $p \geq 0$ に対して写像
$$
\colim_\lambda H^p(\mathcal{C}, \mathcal{F}_\lambda)
\longrightarrow
H^p(\mathcal{C}, \colim_\lambda \mathcal{F}_\lambda)
$$
は、すべてのフィルター付き図式
$\Lambda \to \textit{Ab}(\mathcal{C})$, $\lambda \mapsto \mathcal{F}_\lambda$ に対して同型である。
\end{lemma}

\begin{proof}
これを $p$ に関する帰納法で証明する。$p = 0$ の基底の場合は、Sites, Lemma
\ref{sites-lemma-directed-colimits-global-sections} part (4) から従う。仮定が成り立つ
ことを確認するが、この部分は読み飛ばしてもよい。$\mathcal{F} \to *$ が全射とする。
(1) のように $K \in S$ および全射 $K \to *$ を選ぶ。このとき
$\mathcal{F} \times K \to K$ は全射である。(2) のように $K' \in S$ かつ全射
$K' \to K$ となる $K' \to \mathcal{F} \times K$ を選ぶ。すると $K' \to \mathcal{F}$
が存在し、$K' \to *$ は全射である。従って Sites, Lemma
\ref{sites-lemma-directed-colimits-global-sections} の仮定 (4)(a) が満たされる。
Sites, Lemma \ref{sites-lemma-image-of-quasi-compact} の仮定 (3), (5) により、
すべての $K \in S$ に対して $K \times K$ は準コンパクトである。従って Sites, Lemma
\ref{sites-lemma-directed-colimits-global-sections} の仮定 (4)(b) も満たされる。
これで基底の場合の証明が終わる。

\medskip\noindent
帰納段階。$p \leq p_0$ の $H^p$ について、(1)--(5) を満たす集合
$S \subset \Ob(\Sh(\mathcal{C}))$ が存在するすべてのトポス $\Sh(\mathcal{C})$ で結果が
成り立つと仮定する。Lemma \ref{sites-cohomology-lemma-colim-works-over-collection} の証明と全く同様に
議論すると、入射アーベル層 $\mathcal{I}_\lambda$ の余極限
$\mathcal{I} = \colim \mathcal{I}_\lambda$ に対して
$H^{p_0 + 1}(\mathcal{C}, \mathcal{I}) = 0$ を示せば十分である。(1) のように
$K \in S$ かつ全射 $K \to *$ を選ぶ。$K^n$ を $K$ の $n$ 回自己積と書く。次の
スペクトル系列を考える。
$$
E_1^{p, q} = H^q(K^{p + 1}, \mathcal{I}) \Rightarrow
H^{p + q}(*, \mathcal{I}) = H^{p + q}(\mathcal{C}, \mathcal{I})
$$
これは Lemma \ref{sites-cohomology-lemma-cech-to-cohomology-sheaf-sets} のスペクトル系列である。任意の
$\mathcal{C}$ 上のアーベル層について
$H^q(K^{p + 1}, \mathcal{F}) = H^q(\mathcal{C}/K^{p + 1}, j^{-1}\mathcal{F})$ である
（Lemma \ref{sites-cohomology-lemma-cohomology-on-sheaf-sets} 参照）。$j^{-1}$ は余極限と可換なので
$j^{-1}\mathcal{I} = \colim j^{-1}\mathcal{I}_\lambda$ である。Lemma
\ref{sites-cohomology-lemma-cohomology-of-open} により、制限 $j^{-1}\mathcal{I}_\lambda$ は
$\mathcal{C}/K^{p + 1}$ 上の入射アーベル層である。以下で帰納法の仮定が
$\mathcal{C}/K^{p + 1}$ に適用できることを示すので、$q < p_0 + 1$ に対して
$H^q(K^{p + 1}, \mathcal{I}) = \colim H^q(K^{p + 1}, \mathcal{I}_\lambda) = 0$
が従う（$\mathcal{I}_\lambda$ が入射であることによる消滅）。従って
$$
H^{p_0 + 1}(\mathcal{C}, \mathcal{I}) =
H^{p_0 + 1}\left(\ldots \to
H^0(K^{p_0}, \mathcal{I}) \to
H^0(K^{p_0 + 1}, \mathcal{I}) \to
H^0(K^{p_0 + 2}, \mathcal{I}) \to \ldots\right)
$$
さらに帰納法の仮定を用いると、右辺に描かれた複体は $\Lambda$ に関する次の複体の
余極限である。
$$
\ldots \to
H^0(K^{p_0}, \mathcal{I}_\lambda) \to
H^0(K^{p_0 + 1}, \mathcal{I}_\lambda) \to
H^0(K^{p_0 + 2}, \mathcal{I}_\lambda) \to \ldots
$$
これらの複体は $\mathcal{I}_\lambda$ が入射アーベル層であるため完全である（例えば
スペクトル系列から従う）。フィルター付き余極限はアーベル群の圏で完全なので、求める
消滅が得られる。

\medskip\noindent
すべての $n \geq 1$ に対してサイト $\mathcal{C}/K^n$ に帰納法の仮定が適用できることを
示す必要がある。$\Sh(\mathcal{C}/K^n) = \Sh(\mathcal{C})/K^n$ であることを思い出す
（Sites, Lemma \ref{sites-lemma-localize-topos-site} 参照）。従って
$\Sh(\mathcal{C})/K^n$ で作業してよい。$S_n \subset \Ob(\Sh(\mathcal{C}/K^n)$ を
$K' \to K^n$ という形の対象の集合と書く。各性質を順に確認する。
\begin{enumerate}
\item (3) と帰納法により、$K' \in S$ で全射 $K' \to K^n$ が存在する。このとき
$(K' \to K^n) \to (K^n \to K^n)$ は $\Sh(\mathcal{C})/K^n$ における全射であり、
$K^n \to K^n$ は $\Sh(\mathcal{C})/K^n$ の終対象である。従って $S_n$ について (1) が
成り立つ。
\item $S_n$ に対する性質 (2) は、$S$ に対する (2) から直ちに従う。
\item $S_n$ の対象 $a : K_1 \to K^n$ と $b : K_2 \to K^n$ をとる。このとき
$(K_1 \to K^n) \times (K_2 \to K^n)$ は $\Sh(\mathcal{C})/K^n$ の対象
$K_1 \times_{K^n} K_2 \to K^n$ である。部分層
$K_1 \times_{K^n} K_2 \subset K_1 \times K_2$ は $a \circ \text{pr}_1$ と
$b \circ \text{pr}_2$ の等化子である。$a = (a_1, \ldots, a_n)$ および
$b = (b_1, \ldots, b_n)$ と書く。$K_3 \in S$ かつ全射となる
$K_3 \to K_1 \times K_2$ を選ぶ。これは $\mathcal{C}$ に対する仮定 (3) により可能で
ある。さらに
$$
K_4
\longrightarrow
\text{Equalizer}(K_3 \to K_1 \times K_2 \xrightarrow{a_1, b_1} K)
$$
で $K_4 \in S$ かつ全射となるものを選ぶ。これは $\mathcal{C}$ に対する仮定 (4) により
可能である。次に
$$
K_5
\longrightarrow
\text{Equalizer}(K_4 \to K_1 \times K_2 \xrightarrow{a_2, b_2} K)
$$
で $K_5 \in S$ かつ全射となるものを選ぶ。これも可能である。この手順を続けると
$$
K_{3 + n}
\longrightarrow
\text{Equalizer}(K_{2 + n} \to K_1 \times K_2 \xrightarrow{a_n, b_n} K)
$$
で $K_{3 + n} \in S$ かつ全射となるものを得る。構成により
$K_{3 + n} \to K_1 \times_{K^n} K_2$ は全射である。従って
$(K_{3 + n} \to K^n)$ は $S_n$ に属し、積
$(K_1 \to K^n) \times (K_2 \to K^n)$ へ全射となる。従って $S_n$ について (3) が
成り立つ。
\item $S_n$ に対する性質 (4) は、$S$ に対する性質 (4) から直ちに従う。
\item $S_n$ に対する性質 (5) は、$S$ に対する性質 (5) から従う。実際、
$\Sh(\mathcal{C})/K^n$ の対象 $\mathcal{F} \to K^n$ が
$\Sh(\mathcal{C}/K^n)$ の準コンパクト対象に対応することと、$\mathcal{F}$ が
$\Sh(\mathcal{C})$ の準コンパクト対象であることは同値である。
\end{enumerate}
これで補題の証明が終わる。
\end{proof}

\begin{remark}
\label{sites-cohomology-remark-colim-global}
$\mathcal{C}$ をサイトとし、$\mathcal{B} \subset \Ob(\mathcal{C})$ を部分集合とする。
$S \subset \Ob(\Sh(\mathcal{C}))$ を、次の形をもつ層 $K$ の集合とする。
$$
K = \coprod\nolimits_{i = 1, \ldots, n} h_{U_i}^\#
$$
ここで $U_1, \ldots, U_n \in \mathcal{B}$ である。この集合がいつ Lemma
\ref{sites-cohomology-lemma-colim-global} の仮定を満たすかを問うことができる。一つの十分条件は次である。
\begin{enumerate}
\item ある $n \geq 0$ と $U_1, \ldots, U_n \in \mathcal{B}$ に対して、写像
$\coprod_{i = 1, \ldots, n} h_{U_i}^\# \to *$ が全射である。
\item $U \in \mathcal{B}$ のすべての被覆は、$U_i \in \mathcal{B}$ である形
$\{U_i \to U\}_{i = 1, \ldots, n}$ の被覆により細分される。
\item $U, U' \in \mathcal{B}$ が与えられたとき、ある $n \geq 0$、
$U_1, \ldots, U_n \in \mathcal{B}$、および射 $U_i \to U$, $U_i \to U'$ が存在し、
$\coprod_{i = 1, \ldots, n} h_{U_i}^\# \to h_U^\# \times h_{U'}^\#$ が全射である。
\item $U, U' \in \mathcal{B}$ である $\mathcal{C}$ の射 $a, b : U \to U'$ が与えられた
とき、$U_1, \ldots, U_n \in \mathcal{B}$ と $a,b$ を等化する射 $U_i \to U$ が存在し、
$\coprod_{i = 1, \ldots, n} h_{U_i}^\# \to
\text{Equalizer}(h_a^\#, h_b^\# : h_U^\# \to h_{U'}^\#)$
が全射である。
\end{enumerate}
詳細な確認は省略する。ただし上の (2) は $\mathcal{B}$ のすべての元が準コンパクトで
あることを保証し、従って Sites, Lemma
\ref{sites-lemma-coproduct-quasi-compact} により、すべての $K \in S$ も
準コンパクトであることを述べておく。
\end{remark}

\begin{lemma}
\label{sites-cohomology-lemma-higher-direct-image-colimit}
$f : \mathcal{C} \to \mathcal{D}$ を連続関手 $u : \mathcal{D} \to \mathcal{C}$ に対応する
サイトの射とする。$(\mathcal{F}_i, \varphi_{ii'})$ を $\mathcal{C}$ 上のアーベル層の
系とし、$\mathcal{F} = \colim \mathcal{F}_i$ とおく。$p \geq 0$ を整数とする。
$\mathcal{B}$ を
$H^p(u(V), \mathcal{F}) = \colim H^p(u(V), \mathcal{F}_i)$ を満たす
$V \in \Ob(\mathcal{D})$ の集合とする。$\mathcal{D}$ のすべての対象が
$\mathcal{B}$ の元による被覆をもつならば、
$R^pf_*\mathcal{F} = \colim R^pf_*\mathcal{F}_i$ である。
\end{lemma}

\begin{proof}
$R^pf_*\mathcal{F}$ は $V$ を $H^p(u(V), \mathcal{F})$ へ送る前層
$\mathcal{G}$ の層化である（Lemma \ref{sites-cohomology-lemma-higher-direct-images} 参照）。同様に、
$R^pf_*\mathcal{F}_i$ は $V$ を $H^p(u(V), \mathcal{F}_i)$ へ送る前層
$\mathcal{G}_i$ の層化である。層化は層から前層への包含の左随伴であることを思い出す
（Sites, Section \ref{sites-section-sheafification} 参照）。従って層化は余極限と可換する
（Categories, Lemma \ref{categories-lemma-adjoint-exact} 参照）。従って、前層の圏で
余極限をとった前層の写像
$$
\colim \mathcal{G}_i \longrightarrow \mathcal{G}
$$
が層化上の同型を誘導することを示せば十分である。これは Sites, Lemma
\ref{sites-lemma-isomorphism-sheafifications} と $\mathcal{B}$ に関する仮定から従う。
\end{proof}

\begin{lemma}
\label{sites-cohomology-lemma-colim-sites-injective}
$\mathcal{I}$ を余フィルター付き添字圏とし、$(\mathcal{C}_i, f_a)$ を Sites,
Situation \ref{sites-situation-inverse-limit-sites} のような $\mathcal{I}$ 上のサイトの
逆系とする。Sites, Lemmas \ref{sites-lemma-colimit-sites} および
\ref{sites-lemma-compute-pullback-to-limit} のように
$\mathcal{C} = \colim \mathcal{C}_i$ とおく。さらに次が与えられていると仮定する。
\begin{enumerate}
\item すべての $i \in \Ob(\mathcal{I})$ に対する $\mathcal{C}_i$ 上のアーベル層
$\mathcal{F}_i$。
\item $a : j \to i$ に対する $\mathcal{C}_j$ 上のアーベル層の写像
$\varphi_a : f_a^{-1}\mathcal{F}_i \to \mathcal{F}_j$
\end{enumerate}
で、$c = a \circ b$ のとき常に
$\varphi_c = \varphi_b \circ f_b^{-1}\varphi_a$ を満たすもの。このとき系の射
$(\mathcal{F}_i, \varphi_a) \to (\mathcal{G}_i, \psi_a)$
で、$\mathcal{F}_i \to \mathcal{G}_i$ が単射であり、$\mathcal{G}_i$ が入射アーベル層
となるものが存在する。
\end{lemma}

\begin{proof}
各 $i$ に対し、$\mathcal{A}_i$ が $\mathcal{C}_i$ 上の入射アーベル層となる単射
$\mathcal{F}_i \to \mathcal{A}_i$ を選ぶ。このとき写像族
$$
\gamma_i :
\mathcal{F}_i
\longrightarrow
\prod\nolimits_{b : k \to i} f_{b, *}\mathcal{A}_k = \mathcal{G}_i
$$
を考えられる。その成分写像は
$f_b^{-1}\mathcal{F}_i \to \mathcal{F}_k \to \mathcal{A}_k$ に随伴する写像である。
$\mathcal{I}$ の $a : j \to i$ に対し、標準写像
$$
\psi_a : f_a^{-1}\mathcal{G}_i \to \mathcal{G}_j
$$
があり、その $b : k \to j$ に対応する成分は標準写像
$f_b^{-1}f_{a \circ b, *}\mathcal{A}_k \to f_{b, *}\mathcal{A}_k$ である。従って系の単射
$(\gamma_i) : (\mathcal{F}_i, \varphi_a) \to (\mathcal{G}_i, \psi_a)$
を得る。$\mathcal{G}_i$ は $\mathcal{C}_i$ 上のアーベル群の入射層であることに注意する
（Lemma \ref{sites-cohomology-lemma-pushforward-injective-flat} および Homology, Lemma
\ref{homology-lemma-product-injectives} 参照）。これで構成が完了する。
\end{proof}

\begin{lemma}
\label{sites-cohomology-lemma-colimit}
Lemma \ref{sites-cohomology-lemma-colim-sites-injective} の状況で
$\mathcal{F} = \colim f_i^{-1}\mathcal{F}_i$ とおく。
$i \in \Ob(\mathcal{I})$, $X_i \in \text{Ob}(\mathcal{C}_i)$ とする。このとき
$$
\colim_{a : j \to i} H^p(u_a(X_i), \mathcal{F}_j) =
H^p(u_i(X_i), \mathcal{F})
$$
がすべての $p \geq 0$ に対して成り立つ。
\end{lemma}

\begin{proof}
$p = 0$ の場合は Sites, Lemma \ref{sites-lemma-colimit} である。

\medskip\noindent
Lemma \ref{sites-cohomology-lemma-colim-sites-injective} のように
$(\mathcal{F}_i, \varphi_a) \to (\mathcal{G}_i, \psi_a)$ を選ぶ。Lemma
\ref{sites-cohomology-lemma-colim-works-over-collection} の証明と全く同様に議論すると、$p > 0$ に対して
$H^p(X, \colim f_i^{-1}\mathcal{G}_i) = 0$ を証明すれば十分である。

\medskip\noindent
$\mathcal{G} = \colim f_i^{-1}\mathcal{G}_i$ とおく。$\mathcal{C}$ のすべての対象上で
$\mathcal{G}$ のコホモロジーが消えることを示すため、$\mathcal{C}$ の任意の被覆
$\mathcal{U}$ に対して $\mathcal{G}$ の Cech コホモロジーが 0 であることを示す
（Lemma \ref{sites-cohomology-lemma-cech-vanish-collection}）。被覆 $\mathcal{U}$ は、ある $i$ に対する
$\mathcal{C}_i$ の被覆 $\mathcal{U}_i$ から来る。次が成り立つ。
$$
\check{\mathcal{C}}^\bullet(\mathcal{U}, \mathcal{G}) =
\colim_{a : j \to i}
\check{\mathcal{C}}^\bullet(u_a(\mathcal{U}_i), \mathcal{G}_j)
$$
$p = 0$ の場合による。右辺は Lemma \ref{sites-cohomology-lemma-injective-trivial-cech} により
非輪状複体のフィルター付き余極限であるから、正次数で非輪状である（Algebra, Lemma
\ref{algebra-lemma-directed-colimit-exact} 参照）。
\end{proof}















\section{平坦分解}
\label{sites-cohomology-section-flat}

\noindent
この節では Cohomology, Section \ref{cohomology-section-flat} の議論を、
環付きサイトおよび環付きトポスの枠組みで改めて行う。

\begin{lemma}
\label{sites-cohomology-lemma-derived-tor-exact}
$(\mathcal{C}, \mathcal{O})$ を環付きサイトとし、
$\mathcal{G}^\bullet$ を $\mathcal{O}$-加群の複体とする。関手
$$
K(\textit{Mod}(\mathcal{O}))
\longrightarrow
K(\textit{Mod}(\mathcal{O})),
\quad
\mathcal{F}^\bullet \longmapsto
\text{Tot}(\mathcal{G}^\bullet \otimes_\mathcal{O} \mathcal{F}^\bullet)
$$
および
$$
K(\textit{Mod}(\mathcal{O}))
\longrightarrow
K(\textit{Mod}(\mathcal{O})),
\quad
\mathcal{F}^\bullet \longmapsto
\text{Tot}(\mathcal{F}^\bullet \otimes_\mathcal{O} \mathcal{G}^\bullet)
$$
は三角圏の完全関手である。
\end{lemma}

\begin{proof}
これは Derived Categories, Remark
\ref{derived-remark-double-complex-as-tensor-product-of} から従う。
\end{proof}

\begin{definition}
\label{sites-cohomology-definition-K-flat}
$(\mathcal{C}, \mathcal{O})$ を環付きサイトとする。
$\mathcal{O}$-加群の複体 $\mathcal{K}^\bullet$ が {\it K-平坦} であるとは、
$\mathcal{O}$-加群の任意の非輪状複体 $\mathcal{F}^\bullet$ に対して複体
$$
\text{Tot}(\mathcal{F}^\bullet \otimes_\mathcal{O} \mathcal{K}^\bullet)
$$
が非輪状であることをいう。
\end{definition}

\begin{lemma}
\label{sites-cohomology-lemma-K-flat-quasi-isomorphism}
$(\mathcal{C}, \mathcal{O})$ を環付きサイトとし、
$\mathcal{K}^\bullet$ を K-平坦複体とする。このとき関手
$$
K(\textit{Mod}(\mathcal{O}))
\longrightarrow
K(\textit{Mod}(\mathcal{O})), \quad
\mathcal{F}^\bullet
\longmapsto
\text{Tot}(\mathcal{F}^\bullet \otimes_\mathcal{O} \mathcal{K}^\bullet)
$$
は擬同型を擬同型へ写す。
\end{lemma}

\begin{proof}
Lemma \ref{sites-cohomology-lemma-derived-tor-exact} と、擬同型がその錐の非輪状性によって
特徴づけられることから従う。
\end{proof}

\begin{lemma}
\label{sites-cohomology-lemma-restriction-K-flat}
$(\mathcal{C}, \mathcal{O})$ を環付きサイトとし、$U$ を $\mathcal{C}$ の対象とする。
$\mathcal{K}^\bullet$ が $\mathcal{O}$-加群の K-平坦複体ならば、
$\mathcal{K}^\bullet|_U$ は $\mathcal{O}_U$-加群の K-平坦複体である。
\end{lemma}

\begin{proof}
$\mathcal{G}^\bullet$ を $\mathcal{O}_U$-加群の完全複体とする。
$j_{U!}$ は完全であり
（Modules on Sites, Lemma \ref{sites-modules-lemma-extension-by-zero-exact}）、
$\mathcal{K}^\bullet$ は $\mathcal{O}$-加群の K-平坦複体であるから、複体
$$
j_{U!}(\text{Tot}(\mathcal{G}^\bullet \otimes_{\mathcal{O}_U}
\mathcal{K}^\bullet|_U)) =
\text{Tot}(j_{U!}\mathcal{G}^\bullet \otimes_\mathcal{O} \mathcal{K}^\bullet)
$$
は完全である。ここで等号は Modules on Sites, Lemma
\ref{sites-modules-lemma-j-shriek-and-tensor} と、$j_{U!}$ が左随伴として
直和と可換であることによる。Modules on Sites, Lemma
\ref{sites-modules-lemma-j-shriek-reflects-exactness} により $j_{U!}$ は
完全性を反映するので、結論を得る。
\end{proof}

\begin{lemma}
\label{sites-cohomology-lemma-tensor-product-K-flat}
$(\mathcal{C}, \mathcal{O})$ を環付きサイトとする。
$\mathcal{K}^\bullet$, $\mathcal{L}^\bullet$ が $\mathcal{O}$-加群の
K-平坦複体ならば、
$\text{Tot}(\mathcal{K}^\bullet \otimes_\mathcal{O} \mathcal{L}^\bullet)$
は $\mathcal{O}$-加群の K-平坦複体である。
\end{lemma}

\begin{proof}
同型
$$
\text{Tot}(\mathcal{M}^\bullet \otimes_\mathcal{O}
\text{Tot}(\mathcal{K}^\bullet \otimes_\mathcal{O} \mathcal{L}^\bullet))
=
\text{Tot}(\text{Tot}(\mathcal{M}^\bullet \otimes_\mathcal{O}
\mathcal{K}^\bullet) \otimes_\mathcal{O} \mathcal{L}^\bullet)
$$
と定義から従う。
\end{proof}

\begin{lemma}
\label{sites-cohomology-lemma-K-flat-two-out-of-three}
$(\mathcal{C}, \mathcal{O})$ を環付きサイトとする。
$(\mathcal{K}_1^\bullet, \mathcal{K}_2^\bullet, \mathcal{K}_3^\bullet)$ を
$K(\textit{Mod}(\mathcal{O}))$ の識別三角形とする。
$\mathcal{K}_i^\bullet$ のうち二つが K-平坦ならば、残りの一つも K-平坦である。
\end{lemma}

\begin{proof}
Lemma \ref{sites-cohomology-lemma-derived-tor-exact} と、$K(\textit{Mod}(\mathcal{O}))$ の
識別三角形において三項のうち二項が非輪状ならば残りの一項も非輪状であることから従う。
\end{proof}

\begin{lemma}
\label{sites-cohomology-lemma-K-flat-two-out-of-three-ses}
$(\mathcal{C}, \mathcal{O})$ を環付きサイトとする。
$0 \to \mathcal{K}_1^\bullet \to \mathcal{K}_2^\bullet \to
\mathcal{K}_3^\bullet \to 0$ を、$\mathcal{K}_3^\bullet$ の各項が
平坦 $\mathcal{O}$-加群であるような複体の短完全列とする。
$\mathcal{K}_i^\bullet$ のうち二つが K-平坦ならば、残りの一つも K-平坦である。
\end{lemma}

\begin{proof}
Modules on Sites, Lemma \ref{sites-modules-lemma-flat-tor-zero} により、
任意の複体 $\mathcal{L}^\bullet$ に対して複体の短完全列
$$
0 \to
\text{Tot}(\mathcal{L}^\bullet \otimes_\mathcal{O} \mathcal{K}_1^\bullet) \to
\text{Tot}(\mathcal{L}^\bullet \otimes_\mathcal{O} \mathcal{K}_2^\bullet) \to
\text{Tot}(\mathcal{L}^\bullet \otimes_\mathcal{O} \mathcal{K}_3^\bullet) \to 0
$$
を得る。したがって、コホモロジー層の長完全列と K-平坦複体の定義から補題が従う。
\end{proof}

\begin{lemma}
\label{sites-cohomology-lemma-bounded-flat-K-flat}
$(\mathcal{C}, \mathcal{O})$ を環付きサイトとする。
平坦 $\mathcal{O}$-加群からなる上に有界な複体は K-平坦である。
\end{lemma}

\begin{proof}
$\mathcal{K}^\bullet$ を平坦 $\mathcal{O}$-加群からなる上に有界な複体とし、
$\mathcal{L}^\bullet$ を $\mathcal{O}$-加群の非輪状複体とする。項ごとの余極限を取れば
$\mathcal{L}^\bullet = \colim_m \tau_{\leq m}\mathcal{L}^\bullet$
であることに注意する。したがって、項ごとに
$$
\text{Tot}(\mathcal{K}^\bullet \otimes_\mathcal{O} \mathcal{L}^\bullet)
=
\colim_m \text{Tot}(
\mathcal{K}^\bullet \otimes_\mathcal{O} \tau_{\leq m}\mathcal{L}^\bullet)
$$
でもある。ゆえに左辺の複体が非輪状であることを示すには、右辺の各複体が
非輪状であることを示せば十分である。$\tau_{\leq m}\mathcal{L}^\bullet$ は
非輪状なので、$\mathcal{L}^\bullet$ が上に有界な場合に帰着する。この場合、
Homology, Lemma \ref{homology-lemma-first-quadrant-ss} のスペクトル系列は
$$
{}'E_1^{p, q} = H^p(\mathcal{L}^\bullet \otimes_R \mathcal{K}^q)
$$
をもち、$\mathcal{K}^q$ が平坦で $\mathcal{L}^\bullet$ が非輪状であるため、
これは 0 である。したがって結論を得る。
\end{proof}

\begin{lemma}
\label{sites-cohomology-lemma-colimit-K-flat}
$(\mathcal{C}, \mathcal{O})$ を環付きサイトとする。
$\mathcal{K}_1^\bullet \to \mathcal{K}_2^\bullet \to \ldots$ を
K-平坦複体の系とする。このとき $\colim_i \mathcal{K}_i^\bullet$ は K-平坦である。
\end{lemma}

\begin{proof}
項ごとの余極限を取っているので、明らかに
$$
\colim_i \text{Tot}(
\mathcal{F}^\bullet \otimes_\mathcal{O} \mathcal{K}_i^\bullet)
=
\text{Tot}(\mathcal{F}^\bullet \otimes_\mathcal{O}
\colim_i \mathcal{K}_i^\bullet)
$$
である。したがって、フィルター付き余極限が完全であることから補題が従う。
\end{proof}

\begin{lemma}
\label{sites-cohomology-lemma-resolution-by-direct-sums-extensions-by-zero}
$(\mathcal{C}, \mathcal{O})$ を環付きサイトとする。
$\mathcal{O}$-加群の任意の複体 $\mathcal{G}^\bullet$ に対して、
$\mathcal{O}$-加群の複体からなる可換図式
$$
\xymatrix{
\mathcal{K}_1^\bullet \ar[d] \ar[r] &
\mathcal{K}_2^\bullet \ar[d] \ar[r] & \ldots \\
\tau_{\leq 1}\mathcal{G}^\bullet \ar[r] &
\tau_{\leq 2}\mathcal{G}^\bullet \ar[r] & \ldots
}
$$
が存在し、次を満たす。(1) 縦の射は擬同型であり、各項で全射である。
(2) 各 $\mathcal{K}_n^\bullet$ は上に有界な複体で、その各項は
$j_{U!}\mathcal{O}_U$ という形の $\mathcal{O}$-加群の直和である。
(3) 写像 $\mathcal{K}_n^\bullet \to \mathcal{K}_{n + 1}^\bullet$ は
各項で分裂する単射であり、その余核は $j_{U!}\mathcal{O}_U$ という形の
$\mathcal{O}$-加群の直和である。さらに、写像
$\colim \mathcal{K}_n^\bullet \to \mathcal{G}^\bullet$ は擬同型である。
\end{lemma}

\begin{proof}
図式の存在と性質 (1), (2), (3) は、
Modules on Sites, Lemma \ref{sites-modules-lemma-module-quotient-flat}
および Derived Categories, Lemma \ref{derived-lemma-special-direct-system}
から直ちに従う。誘導される写像
$\colim \mathcal{K}_n^\bullet \to \mathcal{G}^\bullet$
は、フィルター付き余極限が完全であるため擬同型である。
\end{proof}

\begin{lemma}
\label{sites-cohomology-lemma-K-flat-resolution}
$(\mathcal{C}, \mathcal{O})$ を環付きサイトとする。任意の複体
$\mathcal{G}^\bullet$ に対して、各項が平坦 $\mathcal{O}$-加群である
K-平坦複体 $\mathcal{K}^\bullet$ と、各項で全射な擬同型
$\mathcal{K}^\bullet \to \mathcal{G}^\bullet$ が存在する。
\end{lemma}

\begin{proof}
Lemma \ref{sites-cohomology-lemma-resolution-by-direct-sums-extensions-by-zero} の図式を選ぶ。
各複体 $\mathcal{K}_n^\bullet$ は平坦加群からなる上に有界な複体である
（Modules on Sites, Lemma \ref{sites-modules-lemma-j-shriek-flat} 参照）。
したがって Lemma \ref{sites-cohomology-lemma-bounded-flat-K-flat} により
$\mathcal{K}_n^\bullet$ は K-平坦である。ゆえに Lemma
\ref{sites-cohomology-lemma-colimit-K-flat} により $\colim \mathcal{K}_n^\bullet$ は
K-平坦である。誘導される写像
$\colim \mathcal{K}_n^\bullet \to \mathcal{G}^\bullet$
は構成により擬同型であり、各項で全射である。Lemma
\ref{sites-cohomology-lemma-resolution-by-direct-sums-extensions-by-zero} の性質 (3) から、
$\colim \mathcal{K}_n^m$ は平坦加群の直和、したがって平坦である。
これで最後の主張も示された。
\end{proof}

\begin{lemma}
\label{sites-cohomology-lemma-derived-tor-quasi-isomorphism-other-side}
$(\mathcal{C}, \mathcal{O})$ を環付きサイトとする。
$\alpha : \mathcal{P}^\bullet \to \mathcal{Q}^\bullet$ を
$\mathcal{O}$-加群の K-平坦複体の間の擬同型とする。
$\mathcal{O}$-加群の任意の複体 $\mathcal{F}^\bullet$ に対して、誘導される写像
$$
\text{Tot}(\text{id}_{\mathcal{F}^\bullet} \otimes \alpha) :
\text{Tot}(\mathcal{F}^\bullet \otimes_\mathcal{O} \mathcal{P}^\bullet)
\longrightarrow
\text{Tot}(\mathcal{F}^\bullet \otimes_\mathcal{O} \mathcal{Q}^\bullet)
$$
は擬同型である。
\end{lemma}

\begin{proof}
$\mathcal{K}^\bullet$ が K-平坦複体であるような擬同型
$\mathcal{K}^\bullet \to \mathcal{F}^\bullet$ を選ぶ
（Lemma \ref{sites-cohomology-lemma-K-flat-resolution} 参照）。可換図式
$$
\xymatrix{
\text{Tot}(\mathcal{K}^\bullet
\otimes_\mathcal{O} \mathcal{P}^\bullet) \ar[r] \ar[d] &
\text{Tot}(\mathcal{K}^\bullet
\otimes_\mathcal{O} \mathcal{Q}^\bullet) \ar[d] \\
\text{Tot}(\mathcal{F}^\bullet
\otimes_\mathcal{O} \mathcal{P}^\bullet) \ar[r] &
\text{Tot}(\mathcal{F}^\bullet
\otimes_\mathcal{O} \mathcal{Q}^\bullet)
}
$$
を考える。Lemma \ref{sites-cohomology-lemma-K-flat-quasi-isomorphism} により縦の射と
上の横の射は擬同型であるから、結論が従う。
\end{proof}

\noindent
$(\mathcal{C}, \mathcal{O})$ を環付きサイトとする。
$\mathcal{F}^\bullet$ を $D(\mathcal{O})$ の対象とする。
K-平坦分解 $\mathcal{K}^\bullet \to \mathcal{F}^\bullet$ を選ぶ
（Lemma \ref{sites-cohomology-lemma-K-flat-resolution} 参照）。Lemma
\ref{sites-cohomology-lemma-derived-tor-exact} により、三角圏の完全関手
$$
K(\mathcal{O})
\longrightarrow
K(\mathcal{O}),
\quad
\mathcal{G}^\bullet
\longmapsto
\text{Tot}(\mathcal{G}^\bullet \otimes_\mathcal{O} \mathcal{K}^\bullet)
$$
を得る。Lemma \ref{sites-cohomology-lemma-K-flat-quasi-isomorphism} により、この関手は関手
$D(\mathcal{O}) \to D(\mathcal{O})$ を誘導する。実際、$D(\mathcal{O})$ は
$K(\mathcal{O})$ を擬同型で局所化したものだからである。
$\mathcal{F}^\bullet$ の K-平坦分解の圏は余フィルター圏であり、Lemma
\ref{sites-cohomology-lemma-derived-tor-quasi-isomorphism-other-side} があるので、得られる関手は
（同型を除いて）$\mathcal{K}^\bullet$ の選択に依存しない。

\begin{definition}
\label{sites-cohomology-definition-derived-tor}
$(\mathcal{C}, \mathcal{O})$ を環付きサイトとする。
$\mathcal{F}^\bullet$ を $D(\mathcal{O})$ の対象とする。上で記述した
三角圏の完全関手
$$
- \otimes_\mathcal{O}^{\mathbf{L}} \mathcal{F}^\bullet :
D(\mathcal{O})
\longrightarrow
D(\mathcal{O})
$$
を {\it 導来テンソル積}という。
\end{definition}

\noindent
明示的な構成から、$D(\mathcal{O})$ の
$\mathcal{G}^\bullet$ と $\mathcal{F}^\bullet$ に対して標準同型
$$
\mathcal{F}^\bullet \otimes_\mathcal{O}^{\mathbf{L}} \mathcal{G}^\bullet
\cong
\mathcal{G}^\bullet \otimes_\mathcal{O}^{\mathbf{L}} \mathcal{F}^\bullet
$$
があることは明らかである。ここでは More on Algebra, Section
\ref{more-algebra-section-sign-rules} で与えられる符号規則を用いる。したがって
$\mathcal{F}^\bullet \otimes_\mathcal{O}^{\mathbf{L}} \mathcal{G}^\bullet$
と書くとき、通常はどちらの変数を用いて導来テンソル積を定義するかを区別しない。

\begin{definition}
\label{sites-cohomology-definition-tor}
$(\mathcal{C}, \mathcal{O})$ を環付きサイトとする。
$\mathcal{F}$、$\mathcal{G}$ を $\mathcal{O}$-加群とする。
$\mathcal{F}$ と $\mathcal{G}$ の {\it Tor} を、公式
$$
\text{Tor}_p^\mathcal{O}(\mathcal{F}, \mathcal{G}) =
H^{-p}(\mathcal{F} \otimes_\mathcal{O}^\mathbf{L} \mathcal{G})
$$
によって定義する。導来テンソル積は上で定義したものを用いる。
\end{definition}

\noindent
この定義から、任意の $\mathcal{O}$-加群の短完全系列
$0 \to \mathcal{F}_1 \to \mathcal{F}_2 \to \mathcal{F}_3 \to 0$
と任意の $\mathcal{O}$-加群 $\mathcal{G}$ に対して、コホモロジーの長完全系列
$$
\xymatrix{
\mathcal{F}_1 \otimes_\mathcal{O} \mathcal{G} \ar[r] &
\mathcal{F}_2 \otimes_\mathcal{O} \mathcal{G} \ar[r] &
\mathcal{F}_3 \otimes_\mathcal{O} \mathcal{G} \ar[r] & 0 \\
\text{Tor}_1^\mathcal{O}(\mathcal{F}_1, \mathcal{G}) \ar[r] &
\text{Tor}_1^\mathcal{O}(\mathcal{F}_2, \mathcal{G}) \ar[r] &
\text{Tor}_1^\mathcal{O}(\mathcal{F}_3, \mathcal{G}) \ar[ull]
}
$$
が得られる。これをこの状況に付随する $\text{Tor}$ の長完全系列という。

\begin{lemma}
\label{sites-cohomology-lemma-flat-tor-zero}
$(\mathcal{C}, \mathcal{O})$ を環付きサイトとし、
$\mathcal{F}$ を $\mathcal{O}$-加群とする。次は同値である。
\begin{enumerate}
\item $\mathcal{F}$ は平坦 $\mathcal{O}$-加群である。
\item $\text{Tor}_1^\mathcal{O}(\mathcal{F}, \mathcal{G}) = 0$
が任意の $\mathcal{O}$-加群 $\mathcal{G}$ に対して成り立つ。
\end{enumerate}
\end{lemma}

\begin{proof}
もし $\mathcal{F}$ が平坦ならば、$\mathcal{F} \otimes_\mathcal{O} -$ は
完全関手であり、衛星関手は消滅する。逆に (2) が成り立つと仮定する。
$\mathcal{G} \to \mathcal{H}$ が単射で余核が $\mathcal{Q}$ ならば、
$\text{Tor}$ の長完全系列から、
$\mathcal{F} \otimes_\mathcal{O} \mathcal{G} \to
\mathcal{F} \otimes_\mathcal{O} \mathcal{H}$
の核は $\text{Tor}_1^\mathcal{O}(\mathcal{F}, \mathcal{Q})$ の商であることが
分かる。これは仮定により 0 である。したがって $\mathcal{F}$ は平坦である。
\end{proof}

\begin{lemma}
\label{sites-cohomology-lemma-K-flat-flat-acyclic}
$(\mathcal{C}, \mathcal{O})$ を環付きサイトとする。$\mathcal{K}^\bullet$ を、
各項が平坦である K-平坦な非輪状複体とする。このとき
$\mathcal{F} = \Ker(\mathcal{K}^n \to \mathcal{K}^{n + 1})$
は平坦 $\mathcal{O}$-加群である。
\end{lemma}

\begin{proof}
列
$$
\ldots \to \mathcal{K}^{n - 2} \to \mathcal{K}^{n - 1} \to
\mathcal{F} \to 0
$$
は加群 $\mathcal{F}$ の平坦分解であることに注意する。平坦加群からなる上に有界な
複体は K-平坦なので（Lemma \ref{sites-cohomology-lemma-bounded-flat-K-flat}）、任意の
$\mathcal{O}$-加群 $\mathcal{G}$ に対する
$\text{Tor}_i(\mathcal{F}, \mathcal{G})$ の計算にこの分解を用いてよい。一方で
$\mathcal{K}^\bullet \otimes_\mathcal{O}^\mathbf{L} \mathcal{G}$
は $\mathcal{K}^\bullet$ が非輪状であるため $D(\mathcal{O})$ で 0 であり、
他方では $\mathcal{K}^\bullet \otimes_\mathcal{O} \mathcal{G}$ によって表される。
したがって
$$
\mathcal{K}^{n - 3} \otimes_\mathcal{O} \mathcal{G} \to
\mathcal{K}^{n - 2} \otimes_\mathcal{O} \mathcal{G} \to
\mathcal{K}^{n - 1} \otimes_\mathcal{O} \mathcal{G}
$$
は完全である。ゆえに $\text{Tor}_1(\mathcal{F}, \mathcal{G}) = 0$ であり、
Lemma \ref{sites-cohomology-lemma-flat-tor-zero} により結論を得る。
\end{proof}

\begin{lemma}
\label{sites-cohomology-lemma-factor-through-K-flat}
$(\mathcal{C}, \mathcal{O})$ を環付きサイトとする。
$a : \mathcal{K}^\bullet \to \mathcal{L}^\bullet$ を
$\mathcal{O}$-加群の複体の写像とする。$\mathcal{K}^\bullet$ が K-平坦ならば、
複体 $\mathcal{N}^\bullet$ と複体の写像
$b : \mathcal{K}^\bullet \to \mathcal{N}^\bullet$
および $c : \mathcal{N}^\bullet \to \mathcal{L}^\bullet$ で、次を満たすものが存在する。
\begin{enumerate}
\item $\mathcal{N}^\bullet$ は K-平坦である。
\item $c$ は擬同型である。
\item $a$ は $c \circ b$ とホモトピックである。
\end{enumerate}
$\mathcal{K}^\bullet$ の各項が平坦ならば、$\mathcal{N}^\bullet$ についても
同じことが成り立つように $\mathcal{N}^\bullet$, $b$, $c$ を選べる。
\end{lemma}

\begin{proof}
$K(\textit{Mod}(\mathcal{O}))$ が三角圏であることを用いる
（Derived Categories, Proposition
\ref{derived-proposition-homotopy-category-triangulated} 参照）。識別三角形
$\mathcal{K}^\bullet \to \mathcal{L}^\bullet \to
\mathcal{C}^\bullet \to \mathcal{K}^\bullet[1]$ を選ぶ。
$\mathcal{M}^\bullet$ が平坦な各項をもつ K-平坦複体であるような擬同型
$\mathcal{M}^\bullet \to \mathcal{C}^\bullet$ を選ぶ
（Lemma \ref{sites-cohomology-lemma-K-flat-resolution} 参照）。三角圏の公理により、合成
$\mathcal{M}^\bullet \to \mathcal{C}^\bullet \to \mathcal{K}^\bullet[1]$
を識別三角形
$\mathcal{K}^\bullet \to \mathcal{N}^\bullet \to
\mathcal{M}^\bullet \to \mathcal{K}^\bullet[1]$ に組み込める。
Lemma \ref{sites-cohomology-lemma-K-flat-two-out-of-three} により
$\mathcal{N}^\bullet$ は K-平坦である。再び三角圏の公理を用いると、
次の識別三角形の射に組み込まれる写像
$\mathcal{N}^\bullet \to \mathcal{L}^\bullet$ を選べる。
$$
\xymatrix{
\mathcal{K}^\bullet \ar[r] \ar[d] &
\mathcal{N}^\bullet \ar[r] \ar[d] &
\mathcal{M}^\bullet \ar[r] \ar[d] &
\mathcal{K}^\bullet[1] \ar[d] \\
\mathcal{K}^\bullet \ar[r] &
\mathcal{L}^\bullet \ar[r] &
\mathcal{C}^\bullet \ar[r] &
\mathcal{K}^\bullet[1]
}
$$
三つの矢印のうち二つが擬同型なので、これらの識別三角形に付随する
コホモロジーの長完全系列により、第三の矢印
$\mathcal{N}^\bullet \to \mathcal{L}^\bullet$ も擬同型である
（好みに応じて、この図式の $D(\mathcal{O})$ における像を見て
Derived Categories, Lemma \ref{derived-lemma-third-isomorphism-triangle} を
用いてもよい）。これで (1), (2), (3) の証明が完了する。最後の主張を示すには、
$\mathcal{N}^n \cong \mathcal{M}^n \oplus \mathcal{K}^n$ となるように
$\mathcal{N}^\bullet$ を選べばよい（Derived Categories, Lemma
\ref{derived-lemma-improve-distinguished-triangle-homotopy} 参照）。したがって、
$\mathcal{K}^\bullet$ の各項が平坦ならば望む平坦性を得る。
\end{proof}












\section{導来引き戻し}
\label{sites-cohomology-section-derived-pullback}

\noindent
$f : (\Sh(\mathcal{C}), \mathcal{O}) \to
(\Sh(\mathcal{C}'), \mathcal{O}')$ を環付きトポスの射とする。
K-平坦分解を用いて導来引き戻し関手
$$
Lf^* : D(\mathcal{O}') \to D(\mathcal{O})
$$
を定義できる。

\begin{lemma}
\label{sites-cohomology-lemma-pullback-K-flat}
$f : (\Sh(\mathcal{C}'), \mathcal{O}') \to (\Sh(\mathcal{C}), \mathcal{O})$
を環付きトポスの射とする。$\mathcal{K}^\bullet$ を、各項が平坦
$\mathcal{O}$-加群である $\mathcal{O}$-加群の K-平坦複体とする。このとき
$f^*\mathcal{K}^\bullet$ は、各項が平坦 $\mathcal{O}'$-加群である
$\mathcal{O}'$-加群の K-平坦複体である。
\end{lemma}

\begin{proof}
Modules on Sites, Lemma \ref{sites-modules-lemma-pullback-flat} により、
各項 $f^*\mathcal{K}^n$ は平坦 $\mathcal{O}'$-加群である。Lemma
\ref{sites-cohomology-lemma-resolution-by-direct-sums-extensions-by-zero} のような図式
$$
\xymatrix{
\mathcal{K}_1^\bullet \ar[d] \ar[r] &
\mathcal{K}_2^\bullet \ar[d] \ar[r] & \ldots \\
\tau_{\leq 1}\mathcal{K}^\bullet \ar[r] &
\tau_{\leq 2}\mathcal{K}^\bullet \ar[r] & \ldots
}
$$
を選ぶ。以下では補題に述べたすべての性質を断りなく用いる。各
$\mathcal{K}_n^\bullet$ は平坦加群からなる上に有界な複体である
（Modules on Sites, Lemma \ref{sites-modules-lemma-j-shriek-flat} 参照）。
$\mathcal{M}^\bullet$ を定める複体の短完全系列
$$
0 \to \mathcal{M}^\bullet \to
\colim \mathcal{K}_n^\bullet \to
\mathcal{K}^\bullet \to 0
$$
を考える。Lemmas \ref{sites-cohomology-lemma-bounded-flat-K-flat} と
\ref{sites-cohomology-lemma-colimit-K-flat} により複体 $\colim \mathcal{K}_n^\bullet$ は
K-平坦であり、Modules on Sites, Lemma
\ref{sites-modules-lemma-colimits-flat} により各項が平坦である。
Modules on Sites, Lemma \ref{sites-modules-lemma-flat-ses} により
$\mathcal{M}^\bullet$ の各項は平坦であり、Lemma
\ref{sites-cohomology-lemma-K-flat-two-out-of-three-ses} により $\mathcal{M}^\bullet$ は
K-平坦であり、さらにコホモロジーの長完全系列により
$\mathcal{M}^\bullet$ は非輪状である（第二の射が擬同型だからである）。引き戻し
$f^*(\colim \mathcal{K}_n^\bullet) = \colim f^*\mathcal{K}_n^\bullet$
は、平坦 $\mathcal{O}'$-加群からなる下に有界な複体の余極限であり、したがって
K-平坦である（上と同じ補題による）。上の短完全系列の引き戻し
$$
0 \to f^*\mathcal{M}^\bullet \to
f^*(\colim \mathcal{K}_n^\bullet) \to
f^*\mathcal{K}^\bullet \to 0
$$
は Modules on Sites, Lemma \ref{sites-modules-lemma-pullback-ses} により
複体の短完全系列である。したがって Lemma
\ref{sites-cohomology-lemma-K-flat-two-out-of-three-ses} により、$f^*\mathcal{M}^\bullet$ が
K-平坦であることを示せば十分である。これは次の段落で論じる場合に帰着する。

\medskip\noindent
$\mathcal{K}^\bullet$ は K-平坦かつ非輪状で、各項が平坦であると仮定する。
このとき Lemma \ref{sites-cohomology-lemma-K-flat-flat-acyclic} により、
$\tau_{\leq n}\mathcal{K}^\bullet$ のすべての項は平坦 $\mathcal{O}$-加群である。
上のような図式を選び、この図式について上で示したすべての性質を用いる。
複体の写像
$\mathcal{K}_n^\bullet \to \tau_{\leq n}\mathcal{K}^\bullet$
の核を $\mathcal{M}_n^\bullet$ と書く。すると複体の短完全系列
$$
0 \to \mathcal{M}_n^\bullet \to \mathcal{K}_n^\bullet \to
\tau_{\leq n}\mathcal{K}^\bullet \to 0
$$
をもつ。Modules on Sites, Lemma \ref{sites-modules-lemma-flat-ses} により、
複体 $\mathcal{M}_n^\bullet$ の各項は平坦である。したがって、この場合
$\mathcal{M} = \colim \mathcal{M}_n^\bullet$ は平坦加群からなる下に有界な
複体のフィルター付き余極限である。ゆえに $f^*\mathcal{M}^\bullet$ は
K-平坦であり（上と同じ議論）、結論を得る。
\end{proof}

\begin{lemma}
\label{sites-cohomology-lemma-derived-base-change}
$f : (\Sh(\mathcal{C}), \mathcal{O}) \to (\Sh(\mathcal{C}'), \mathcal{O}')$
を環付きトポスの射とする。三角圏の完全関手
$$
Lf^* : D(\mathcal{O}') \longrightarrow D(\mathcal{O})
$$
が存在し、各項が平坦である任意の K-平坦複体 $\mathcal{K}^\bullet$ に対して
$Lf^*\mathcal{K}^\bullet = f^*\mathcal{K}^\bullet$ となる。とくにこれは、
平坦 $\mathcal{O}'$-加群からなる任意の上に有界な複体に対して成り立つ。
\end{lemma}

\begin{proof}
これを見るため、Derived Categories, Section
\ref{derived-section-derived-functors} で展開された一般論を用いる。
$\mathcal{D} = K(\mathcal{O}')$ および
$\mathcal{D}' = D(\mathcal{O})$ と置く。規則
$F(\mathcal{G}^\bullet) = f^*\mathcal{G}^\bullet$ で定まる三角圏の完全関手を
$F : \mathcal{D} \to \mathcal{D}'$ と書く。$S$ を
$\mathcal{D} = K(\mathcal{O}')$ における擬同型の集合とする。
これにより Derived Categories, Situation
\ref{derived-situation-derived-functor} の状況を得るので、Derived Categories,
Definition \ref{derived-definition-right-derived-functor-defined} が適用できる。
$LF$ はいたるところ定義されると主張する。これは Derived Categories, Lemma
\ref{derived-lemma-find-existence-computes} から従う。ここで
$\mathcal{P} \subset \Ob(\mathcal{D})$ を、各項が平坦である K-平坦複体
$\mathcal{K}^\bullet$ の族とする。(1) は Lemma
\ref{sites-cohomology-lemma-K-flat-resolution} から従う。(2) を見るには、$\mathcal{P}$ の
対象の間の擬同型 $\mathcal{K}_1^\bullet \to \mathcal{K}_2^\bullet$ に対し、写像
$f^*\mathcal{K}_1^\bullet \to f^*\mathcal{K}_2^\bullet$ が擬同型であることを
示さなければならない。この写像を
$$
f^{-1}\mathcal{K}_1^\bullet \otimes_{f^{-1}\mathcal{O}'} \mathcal{O}
\longrightarrow
f^{-1}\mathcal{K}_2^\bullet \otimes_{f^{-1}\mathcal{O}'} \mathcal{O}
$$
と書く。関手 $f^{-1}$ は完全なので、写像
$f^{-1}\mathcal{K}_1^\bullet \to f^{-1}\mathcal{K}_2^\bullet$ は擬同型である。
環付きトポスの射
$(\Sh(\mathcal{C}), f^{-1}\mathcal{O}') \to
(\Sh(\mathcal{C}'), \mathcal{O}')$ を考えられるので、Lemma
\ref{sites-cohomology-lemma-pullback-K-flat} により、複体
$f^{-1}\mathcal{K}_1^\bullet$ と $f^{-1}\mathcal{K}_2^\bullet$ は
$f^{-1}\mathcal{O}'$-加群の K-平坦複体である。したがって Lemma
\ref{sites-cohomology-lemma-derived-tor-quasi-isomorphism-other-side} により、表示した写像は
擬同型である。ゆえに導来関手
$$
LF :
D(\mathcal{O}') = S^{-1}\mathcal{D}
\longrightarrow
\mathcal{D}' = D(\mathcal{O})
$$
を得る（Derived Categories, Equation
(\ref{derived-equation-everywhere}) 参照）。最後に、Derived Categories, Lemma
\ref{derived-lemma-find-existence-computes} は、$\mathcal{K}^\bullet$ が
$\mathcal{P}$ に属するとき
$LF(\mathcal{K}^\bullet) = F(\mathcal{K}^\bullet) = f^*\mathcal{K}^\bullet$
であることも保証する。Lemma \ref{sites-cohomology-lemma-bounded-flat-K-flat} により、
平坦加群からなる上に有界な複体は $\mathcal{P}$ に属することに注意すれば、
証明が完了する。
\end{proof}

\begin{lemma}
\label{sites-cohomology-lemma-derived-pullback-composition}
環付きトポスの射
$f : (\Sh(\mathcal{C}), \mathcal{O}_\mathcal{C}) \to
(\Sh(\mathcal{D}), \mathcal{O}_\mathcal{D})$
および
$g : (\Sh(\mathcal{D}), \mathcal{O}_\mathcal{D}) \to
(\Sh(\mathcal{E}), \mathcal{O}_\mathcal{E})$ を考える。このとき関手
$D(\mathcal{O}_\mathcal{E}) \to D(\mathcal{O}_\mathcal{C})$ として
$Lf^* \circ Lg^* = L(g \circ f)^*$ である。
\end{lemma}

\begin{proof}
$E$ を $D(\mathcal{O}_\mathcal{E})$ の対象とする。$E$ は、各項が平坦な
K-平坦複体 $\mathcal{K}^\bullet$ によって表せる
（Lemma \ref{sites-cohomology-lemma-K-flat-resolution} 参照）。構成により $Lg^*E$ は
$g^*\mathcal{K}^\bullet$ で計算される（Lemma
\ref{sites-cohomology-lemma-derived-base-change} 参照）。Lemma \ref{sites-cohomology-lemma-pullback-K-flat} により、
複体 $g^*\mathcal{K}^\bullet$ は各項が平坦な K-平坦複体である。したがって
$Lf^*Lg^*E$ は $f^*g^*\mathcal{K}^\bullet$ によって表される。一方、
$L(g \circ f)^*E$ も
$(g \circ f)^*\mathcal{K}^\bullet = f^*g^*\mathcal{K}^\bullet$ によって
表されるので、結論を得る。
\end{proof}

\begin{lemma}
\label{sites-cohomology-lemma-pullback-tensor-product}
$f : (\Sh(\mathcal{C}), \mathcal{O}) \to (\Sh(\mathcal{D}), \mathcal{O}')$
を環付きトポスの射とする。
$\mathcal{F}^\bullet, \mathcal{G}^\bullet \in \Ob(D(\mathcal{O}'))$ に対し、
二変数について関手的な標準同型
$$
Lf^*(
\mathcal{F}^\bullet \otimes_{\mathcal{O}'}^{\mathbf{L}} \mathcal{G}^\bullet
) =
Lf^*\mathcal{F}^\bullet 
\otimes_{\mathcal{O}}^{\mathbf{L}}
Lf^*\mathcal{G}^\bullet 
$$
がある。
\end{lemma}

\begin{proof}
Lemma \ref{sites-cohomology-lemma-derived-base-change} における導来引き戻しの構成と
Lemma \ref{sites-cohomology-lemma-K-flat-resolution} における分解の存在により、
$\mathcal{F}^\bullet$ と $\mathcal{G}^\bullet$ を、各項が平坦である
$\mathcal{O}'$-加群の K-平坦複体で置き換えてよい。この場合
$\mathcal{F}^\bullet \otimes_{\mathcal{O}'}^{\mathbf{L}} \mathcal{G}^\bullet$
は二重複体 $\mathcal{F}^\bullet \otimes_{\mathcal{O}'} \mathcal{G}^\bullet$
に付随する全複体にほかならない。複体
$\text{Tot}(\mathcal{F}^\bullet \otimes_{\mathcal{O}'} \mathcal{G}^\bullet)$
は Lemma \ref{sites-cohomology-lemma-tensor-product-K-flat} と Modules on Sites, Lemma
\ref{sites-modules-lemma-tensor-flats} により、各項が平坦な K-平坦複体である。
したがって補題の同型は、同型
$$
\text{Tot}(f^*\mathcal{F}^\bullet \otimes_{\mathcal{O}}
f^*\mathcal{G}^\bullet)
\longrightarrow
f^*\text{Tot}(\mathcal{F}^\bullet \otimes_{\mathcal{O}'} \mathcal{G}^\bullet)
$$
から生じる。その成分は Modules on Sites, Lemma
\ref{sites-modules-lemma-tensor-product-pullback} の同型
$f^*\mathcal{F}^p \otimes_{\mathcal{O}} f^*\mathcal{G}^q \to
f^*(\mathcal{F}^p \otimes_{\mathcal{O}'} \mathcal{G}^q)$ である。
\end{proof}

\begin{lemma}
\label{sites-cohomology-lemma-variant-derived-pullback}
$f : (\Sh(\mathcal{C}), \mathcal{O}) \to (\Sh(\mathcal{C}'), \mathcal{O}')$
を環付きトポスの射とする。$D(\mathcal{O})$ の $\mathcal{F}^\bullet$ と
$D(\mathcal{O}')$ の $\mathcal{G}^\bullet$ に対し、二変数について関手的な標準同型
$$
\mathcal{F}^\bullet
\otimes_\mathcal{O}^{\mathbf{L}}
Lf^*\mathcal{G}^\bullet
=
\mathcal{F}^\bullet 
\otimes_{f^{-1}\mathcal{O}'}^{\mathbf{L}}
f^{-1}\mathcal{G}^\bullet 
$$
がある。
\end{lemma}

\begin{proof}
$\mathcal{F}$ を $\mathcal{O}$-加群とし、$\mathcal{G}$ を
$\mathcal{O}'$-加群とする。このとき
$\mathcal{F} \otimes_{\mathcal{O}} f^*\mathcal{G} =
\mathcal{F} \otimes_{f^{-1}\mathcal{O}'} f^{-1}\mathcal{G}$
である。実際、
$f^*\mathcal{G} =
\mathcal{O} \otimes_{f^{-1}\mathcal{O}'} f^{-1}\mathcal{G}$ である。
補題はこれと定義から従う。
\end{proof}











\begin{lemma}
\label{sites-cohomology-lemma-check-K-flat-stalks}
$(\mathcal{C}, \mathcal{O})$ を環付きサイトとし、$\mathcal{K}^\bullet$ を
$\mathcal{O}$-加群の複体とする。
\begin{enumerate}
\item $\mathcal{K}^\bullet$ が K-平坦ならば、サイト $\mathcal{C}$ の任意の点
$p$ に対して、$\mathcal{O}_p$-加群の複体 $\mathcal{K}_p^\bullet$ は
More on Algebra, Definition \ref{more-algebra-definition-K-flat} の意味で
K-平坦である。
\item $\mathcal{C}$ が十分な点をもつならば、逆も成り立つ。
\end{enumerate}
\end{lemma}

\begin{proof}
(2) の証明。$\mathcal{C}$ が十分な点をもち、そのすべての点 $p$ に対して
$\mathcal{K}_p^\bullet$ が K-平坦であるとする。$\otimes$ と直和は茎を
取ることと可換であり、完全性は茎で確認できるので、$\mathcal{K}^\bullet$ は
K-平坦である（Modules on Sites, Lemma
\ref{sites-modules-lemma-check-exactness-stalks} 参照）。

\medskip\noindent
(1) の証明。$\mathcal{K}^\bullet$ が K-平坦であると仮定する。
$\mathcal{L}^\bullet$ が各項の平坦な K-平坦複体であるような擬同型
$a : \mathcal{L}^\bullet \to \mathcal{K}^\bullet$ を選ぶ
（Lemma \ref{sites-cohomology-lemma-K-flat-resolution} 参照）。$\mathcal{L}^\bullet$ の任意の
引き戻しは K-平坦である（Lemma \ref{sites-cohomology-lemma-pullback-K-flat} 参照）。とくに茎
$\mathcal{L}_p^\bullet$ は $\mathcal{O}_p$-加群の K-平坦複体である。したがって
$a$ の錐 $C(a)$ は $\mathcal{O}$-加群の K-平坦な
（Lemma \ref{sites-cohomology-lemma-K-flat-two-out-of-three}）非輪状複体であり、$C(a)$ の茎が
K-平坦であることを示せば十分である（More on Algebra, Lemma
\ref{more-algebra-lemma-K-flat-two-out-of-three} による）。ゆえに
$\mathcal{K}^\bullet$ が K-平坦かつ非輪状であると仮定してよい。

\medskip\noindent
$\mathcal{K}^\bullet$ は非輪状かつ K-平坦であると仮定する。先へ進む前に、
Sites, Lemma \ref{sites-lemma-topos-good-site} のようにサイト
$\mathcal{C}$ を別のサイトで置き換え、$\mathcal{C}$ がすべての有限極限を
もつようにする。すると $p$ の近傍の圏はフィルター圏となり
（Sites, Lemma \ref{sites-lemma-neighbourhoods-directed}）、層の茎を定義する
余極限はフィルター付きとなる。$M$ を有限表示 $\mathcal{O}_p$-加群とする。
More on Algebra, Lemma \ref{more-algebra-lemma-universally-acyclic-K-flat} により、
$\mathcal{K}^\bullet \otimes_{\mathcal{O}_p} M$ が非輪状であることを示せば
十分である。$\mathcal{O}_p$ は、$U$ が $p$ の近傍を走るときの
$\mathcal{O}(U)$ のフィルター付き余極限なので、$p$ の近傍 $(U, x)$ と、
$\mathcal{O}_p$ への基底変換が $M$ となる有限表示 $\mathcal{O}(U)$-加群
$M'$ を見いだせる（Algebra, Lemma
\ref{algebra-lemma-colimit-category-fp-modules} 参照）。Lemma
\ref{sites-cohomology-lemma-restriction-K-flat} により、$\mathcal{C}, \mathcal{O},
\mathcal{K}^\bullet$ を $\mathcal{C}/U, \mathcal{O}_U,
\mathcal{K}^\bullet|_U$ で置き換えてよい。したがって、
$M \cong \mathcal{F}_p$ を満たす $\mathcal{O}$-加群 $\mathcal{F}$ が存在すると
仮定してよい。$\mathcal{K}^\bullet$ は K-平坦かつ非輪状なので、
$\mathcal{K}^\bullet \otimes_\mathcal{O} \mathcal{F}$ は非輪状である
（定義により導来テンソル積を計算するため）。茎を取ることは完全関手なので、
望みどおり $\mathcal{K}^\bullet \otimes_{\mathcal{O}_p} M$ は非輪状である。
\end{proof}

\begin{lemma}
\label{sites-cohomology-lemma-pullback-K-flat-points}
$f : (\Sh(\mathcal{C}), \mathcal{O}) \to (\Sh(\mathcal{C}'), \mathcal{O}')$
を環付きトポスの射とする。$\mathcal{C}$ が十分な点をもつならば、
$\mathcal{O}'$-加群の K-平坦複体の引き戻しは
$\mathcal{O}$-加群の K-平坦複体である。
\end{lemma}

\begin{proof}
これは Lemma \ref{sites-cohomology-lemma-check-K-flat-stalks}, Modules on Sites, Lemma
\ref{sites-modules-lemma-pullback-stalk}, および More on Algebra, Lemma
\ref{more-algebra-lemma-base-change-K-flat} から従う。
\end{proof}

\begin{lemma}
\label{sites-cohomology-lemma-tensor-pull-compatibility}
$f : (\Sh(\mathcal{C}), \mathcal{O}_\mathcal{C}) \to
(\Sh(\mathcal{D}), \mathcal{O}_\mathcal{D})$ を環付きトポスの射とする。
$\mathcal{K}^\bullet$ と $\mathcal{M}^\bullet$ を
$\mathcal{O}_\mathcal{D}$-加群の複体とする。図式
$$
\xymatrix{
Lf^*(\mathcal{K}^\bullet
\otimes_{\mathcal{O}_\mathcal{D}}^\mathbf{L}
\mathcal{M}^\bullet) \ar[r] \ar[d] &
Lf^*\text{Tot}(\mathcal{K}^\bullet
\otimes_{\mathcal{O}_\mathcal{D}}
\mathcal{M}^\bullet) \ar[d] \\
Lf^*\mathcal{K}^\bullet \otimes_{\mathcal{O}_\mathcal{C}}^\mathbf{L}
Lf^*\mathcal{M}^\bullet \ar[d] &
f^*\text{Tot}(\mathcal{K}^\bullet
\otimes_{\mathcal{O}_\mathcal{D}}
\mathcal{M}^\bullet) \ar[d] \\
f^*\mathcal{K}^\bullet \otimes_{\mathcal{O}_\mathcal{C}}^\mathbf{L}
f^*\mathcal{M}^\bullet \ar[r] &
\text{Tot}(f^*\mathcal{K}^\bullet \otimes_{\mathcal{O}_\mathcal{C}}
f^*\mathcal{M}^\bullet)
}
$$
は可換である。
\end{lemma}

\begin{proof}
各項が平坦な K-平坦分解の存在（Lemma \ref{sites-cohomology-lemma-K-flat-resolution}）、
導来引き戻しがそのような複体で計算されること
（Lemma \ref{sites-cohomology-lemma-derived-base-change}）、および引き戻しがこれらの性質を
保つこと（Lemma \ref{sites-cohomology-lemma-pullback-K-flat}）を用いる。そのような分解
$\mathcal{P}^\bullet \to \mathcal{K}^\bullet$ および
$\mathcal{Q}^\bullet \to \mathcal{M}^\bullet$ を選ぶと、次の図式
$$
\xymatrix{
Lf^*\text{Tot}(\mathcal{P}^\bullet
\otimes_{\mathcal{O}_\mathcal{D}}
\mathcal{Q}^\bullet) \ar[r] \ar[d] &
Lf^*\text{Tot}(\mathcal{K}^\bullet
\otimes_{\mathcal{O}_\mathcal{D}}
\mathcal{M}^\bullet) \ar[d] \\
f^*\text{Tot}(\mathcal{P}^\bullet
\otimes_{\mathcal{O}_\mathcal{D}}
\mathcal{Q}^\bullet) \ar[d] \ar[r] &
f^*\text{Tot}(\mathcal{K}^\bullet
\otimes_{\mathcal{O}_\mathcal{D}}
\mathcal{M}^\bullet) \ar[d] \\
\text{Tot}(f^*\mathcal{P}^\bullet \otimes_{\mathcal{O}_\mathcal{C}}
f^*\mathcal{Q}^\bullet) \ar[r] &
\text{Tot}(f^*\mathcal{K}^\bullet \otimes_{\mathcal{O}_\mathcal{C}}
f^*\mathcal{M}^\bullet)
}
$$
が可換であることが分かる。しかし今や、$\mathcal{P}^\bullet$ と
$\mathcal{Q}^\bullet$ の選択および Lemma
\ref{sites-cohomology-lemma-tensor-product-K-flat} により、この図式の左辺は図式の左辺である。
\end{proof}









\section{非有界複体のコホモロジー}
\label{sites-cohomology-section-unbounded}

\noindent
$(\mathcal{C}, \mathcal{O})$ を環付きサイトとする。圏
$\textit{Mod}(\mathcal{O})$ は Grothendieck アーベル圏である。すなわち、
すべての余極限をもち、フィルター付き余極限は完全であり、さらに生成元
$$
\bigoplus\nolimits_{U \in \Ob(\mathcal{C})} j_{U!}\mathcal{O}_U,
$$
をもつ（Modules on Sites, Section \ref{sites-modules-section-kernels} および
Lemmas \ref{sites-modules-lemma-j-shriek-flat} と
\ref{sites-modules-lemma-module-quotient-flat} 参照）。Injectives, Theorem
\ref{injectives-theorem-K-injective-embedding-grothendieck}
により、$\mathcal{O}$-加群の任意の複体 $\mathcal{F}^\bullet$ に対して、
$\mathcal{O}$-加群の K-入射複体 $\mathcal{I}^\bullet$ への単射な擬同型
$\mathcal{F}^\bullet \to \mathcal{I}^\bullet$ が存在し、しかもこの埋め込みは
$\mathcal{F}^\bullet$ に関して関手的に選べる。Derived Categories, Lemma
\ref{derived-lemma-enough-K-injectives-implies} から次が従う。
\begin{enumerate}
\item 三角圏 $\mathcal{D}$ への任意の完全関手
$F : K(\textit{Mod}(\mathcal{O})) \to \mathcal{D}$ は右導来関手
$RF : D(\mathcal{O}) \to \mathcal{D}$ をもつ。
\item アーベル圏 $\mathcal{A}$ への任意の加法関手
$F : \textit{Mod}(\mathcal{O}) \to \mathcal{A}$
に対し、$F$ が誘導する完全関手
$F : K(\textit{Mod}(\mathcal{O})) \to K(\mathcal{A})$ を考えると、
右導来関手 $RF : D(\mathcal{O}) \to D(\mathcal{A})$ が得られる。
\end{enumerate}
構成により $RF(\mathcal{F}^\bullet) = F(\mathcal{I}^\bullet)$ である。
ここで $\mathcal{F}^\bullet \to \mathcal{I}^\bullet$ は上のものである。

\medskip\noindent
以上の例をいくつか挙げる。
\begin{enumerate}
\item 関手 $\Gamma(\mathcal{C}, -) : \textit{Mod}(\mathcal{O}) \to
\text{Mod}_{\Gamma(\mathcal{C}, \mathcal{O})}$ から
$$
R\Gamma(\mathcal{C}, -) :
D(\mathcal{O})
\longrightarrow
D(\Gamma(\mathcal{C}, \mathcal{O}))
$$
が得られる。コホモロジーには記法
$H^i(\mathcal{C}, K) = H^i(R\Gamma(\mathcal{C}, K))$ を用いる。
\item $\mathcal{C}$ の対象 $U$ に対し、関手
$\Gamma(U, -) : \textit{Mod}(\mathcal{O}) \to
\text{Mod}_{\Gamma(U, \mathcal{O})}$ を考える。これから
$$
R\Gamma(U, -) : D(\mathcal{O}) \to D(\Gamma(U, \mathcal{O}))
$$
が得られる。コホモロジーには記法
$H^i(U, K) = H^i(R\Gamma(U, K))$ を用いる。
\item 環付きトポスの射
$f : (\Sh(\mathcal{C}), \mathcal{O}) \to (\Sh(\mathcal{D}), \mathcal{O}')$
に対し、関手
$f_* : \textit{Mod}(\mathcal{O}) \to \textit{Mod}(\mathcal{O}')$
を考える。これから非有界導来圏上の全導来順像
$$
Rf_* : D(\mathcal{O}) \longrightarrow D(\mathcal{O}')
$$
が得られる。
\end{enumerate}

\begin{lemma}
\label{sites-cohomology-lemma-adjoint}
$f : (\Sh(\mathcal{C}), \mathcal{O}) \to (\Sh(\mathcal{D}), \mathcal{O}')$
を環付きトポスの射とする。上で定義した関手 $Rf_*$ と、Lemma
\ref{sites-cohomology-lemma-derived-base-change} で定義した関手 $Lf^*$ は随伴である。
$$
\Hom_{D(\mathcal{O})}(Lf^*\mathcal{G}^\bullet, \mathcal{F}^\bullet)
=
\Hom_{D(\mathcal{O}')}(\mathcal{G}^\bullet, Rf_*\mathcal{F}^\bullet)
$$
この同型は $\mathcal{F}^\bullet \in \Ob(D(\mathcal{O}))$ と
$\mathcal{G}^\bullet \in \Ob(D(\mathcal{O}'))$ に関して双関手的である。
\end{lemma}

\begin{proof}
これは $Rf_*$ と $Lf^*$ が存在することから形式的に従う
（Derived Categories, Lemma \ref{derived-lemma-derived-adjoint-functors} 参照）。
\end{proof}

\begin{lemma}
\label{sites-cohomology-lemma-derived-pushforward-composition}
環付きトポスの射
$f : (\Sh(\mathcal{C}), \mathcal{O}_\mathcal{C}) \to
(\Sh(\mathcal{D}), \mathcal{O}_\mathcal{D})$
および
$g : (\Sh(\mathcal{D}), \mathcal{O}_\mathcal{D}) \to
(\Sh(\mathcal{E}), \mathcal{O}_\mathcal{E})$
を考える。このとき関手 $D(\mathcal{O}_\mathcal{C}) \to
D(\mathcal{O}_\mathcal{E})$ として
$Rg_* \circ Rf_* = R(g \circ f)_*$ である。
\end{lemma}

\begin{proof}
Lemma \ref{sites-cohomology-lemma-adjoint} により、$Rg_* \circ Rf_*$ は
$Lf^* \circ Lg^*$ の随伴である。また Lemma
\ref{sites-cohomology-lemma-derived-pullback-composition} により
$Lf^* \circ Lg^* = L(g \circ f)^*$ であるから、随伴関手の一意性により
$Rg_* \circ Rf_* = R(g \circ f)_*$ を得る。
\end{proof}

\begin{remark}
\label{sites-cohomology-remark-base-change}
非有界導来関手 $Lf^*$ と $Rf_*$ の構成により、完全な一般性の下で
基底変換写像を構成できる。すなわち、
$$
\xymatrix{
(\Sh(\mathcal{C}'), \mathcal{O}_{\mathcal{C}'})
\ar[r]_{g'} \ar[d]_{f'} &
(\Sh(\mathcal{C}), \mathcal{O}_\mathcal{C}) \ar[d]^f \\
(\Sh(\mathcal{D}'), \mathcal{O}_{\mathcal{D}'})
\ar[r]^g &
(\Sh(\mathcal{D}), \mathcal{O}_\mathcal{D})
}
$$
が環付きトポスの可換図式であると仮定する。$K$ を
$D(\mathcal{O}_\mathcal{C})$ の対象とする。このとき標準的な基底変換写像
$$
Lg^*Rf_*K \longrightarrow R(f')_*L(g')^*K
$$
が $D(\mathcal{O}_{\mathcal{D}'})$ に存在する。実際、この写像は写像
$L(f')^*Lg^*Rf_*K \to L(g')^*K$ に随伴する。
$L(f')^* \circ Lg^* = L(g')^* \circ Lf^*$ であるから、これは写像
$L(g')^*Lf^*Rf_*K \to L(g')^*K$ と同じであり、後者として随伴写像
$Lf^*Rf_*K \to K$ に $L(g')^*$ を施したものを取れる。
\end{remark}

\begin{remark}
\label{sites-cohomology-remark-compose-base-change}
環付きトポスの可換図式
$$
\xymatrix{
(\Sh(\mathcal{B}'), \mathcal{O}_{\mathcal{B}'})
\ar[r]_k \ar[d]_{f'} &
(\Sh(\mathcal{B}), \mathcal{O}_\mathcal{B}) \ar[d]^f \\
(\Sh(\mathcal{C}'), \mathcal{O}_{\mathcal{C}'})
\ar[r]^l \ar[d]_{g'} &
(\Sh(\mathcal{C}), \mathcal{O}_\mathcal{C}) \ar[d]^g \\
(\Sh(\mathcal{D}'), \mathcal{O}_{\mathcal{D}'})
\ar[r]^m &
(\Sh(\mathcal{D}), \mathcal{O}_\mathcal{D}) \\
}
$$
を考える。このとき二つの正方形に対する Remark
\ref{sites-cohomology-remark-base-change} の基底変換写像を合成すると、外側の長方形に対する
基底変換写像が得られる。より正確には、合成
\begin{align*}
Lm^* \circ R(g \circ f)_*
& =
Lm^* \circ Rg_* \circ Rf_* \\
& \to Rg'_* \circ Ll^* \circ Rf_* \\
& \to Rg'_* \circ Rf'_* \circ Lk^* \\
& = R(g' \circ f')_* \circ Lk^*
\end{align*}
がこの長方形に対する基底変換写像である。検証は省略する。
\end{remark}

\begin{remark}
\label{sites-cohomology-remark-compose-base-change-horizontal}
環付きトポスの可換図式
$$
\xymatrix{
(\Sh(\mathcal{C}''), \mathcal{O}_{\mathcal{C}''})
\ar[r]_{g'} \ar[d]_{f''} &
(\Sh(\mathcal{C}'), \mathcal{O}_{\mathcal{C}'})
\ar[r]_g \ar[d]_{f'} &
(\Sh(\mathcal{C}), \mathcal{O}_\mathcal{C}) \ar[d]^f \\
(\Sh(\mathcal{D}''), \mathcal{O}_{\mathcal{D}''})
\ar[r]^{h'} &
(\Sh(\mathcal{D}'), \mathcal{O}_{\mathcal{D}'})
\ar[r]^h &
(\Sh(\mathcal{D}), \mathcal{O}_\mathcal{D})
}
$$
を考える。このとき二つの正方形に対する Remark
\ref{sites-cohomology-remark-base-change} の基底変換写像を合成すると、外側の長方形に対する
基底変換写像が得られる。より正確には、合成
\begin{align*}
L(h \circ h')^* \circ Rf_*
& =
L(h')^* \circ Lh^* \circ Rf_* \\
& \to L(h')^* \circ Rf'_* \circ Lg^* \\
& \to Rf''_* \circ L(g')^* \circ Lg^* \\
& = Rf''_* \circ L(g \circ g')^*
\end{align*}
がこの長方形に対する基底変換写像である。検証は省略する。
\end{remark}



\begin{lemma}
\label{sites-cohomology-lemma-adjoints-push-pull-compatibility}
$f : (\Sh(\mathcal{C}), \mathcal{O}_\mathcal{C}) \to
(\Sh(\mathcal{D}), \mathcal{O}_\mathcal{D})$ を環付きトポスの射とし、
$\mathcal{K}^\bullet$ を $\mathcal{O}_\mathcal{C}$-加群の複体とする。
複体上の $Lf^* \to f^*$、複体上の $f_* \to Rf_*$、および随伴
$Lf^* \circ Rf_* \to \text{id}$ から得られる図式
$$
\xymatrix{
Lf^*f_*\mathcal{K}^\bullet \ar[r] \ar[d] &
f^*f_*\mathcal{K}^\bullet \ar[d] \\
Lf^*Rf_*\mathcal{K}^\bullet \ar[r] &
\mathcal{K}^\bullet
}
$$
は $D(\mathcal{O}_\mathcal{C})$ で可換である。
\end{lemma}

\begin{proof}
K-平坦分解と K-入射分解の存在を用いる。Lemmas
\ref{sites-cohomology-lemma-K-flat-resolution}, \ref{sites-cohomology-lemma-derived-base-change},
\ref{sites-cohomology-lemma-pullback-K-flat} および上の議論を参照せよ。
$\mathcal{I}^\bullet$ が $\mathcal{O}_\mathcal{C}$-加群の複体として
K-入射的であるような擬同型
$\mathcal{K}^\bullet \to \mathcal{I}^\bullet$ を選ぶ。また、
$\mathcal{Q}^\bullet$ が各項の平坦な $\mathcal{O}_\mathcal{D}$-加群の
K-平坦複体であるような擬同型
$\mathcal{Q}^\bullet \to f_*\mathcal{I}^\bullet$ を選ぶ。
各項の平坦な $\mathcal{O}_\mathcal{D}$-加群の K-平坦複体
$\mathcal{P}^\bullet$ と、複体の射の図式
$$
\xymatrix{
\mathcal{P}^\bullet \ar[r] \ar[d] &
f_*\mathcal{K}^\bullet \ar[d] \\
\mathcal{Q}^\bullet \ar[r] & f_*\mathcal{I}^\bullet
}
$$
を、ホモトピーを除いて可換で、かつ上の水平射が擬同型となるように選べる。
実際、複体のホモトピー圏において擬同型は乗法系をなすので、まずある複体
$\mathcal{P}^\bullet$ に対してこのような図式を選び、次に
$\mathcal{P}^\bullet$ を各項の平坦な K-平坦複体で分解すればよい。
引き戻しを取ると、複体の射の図式
$$
\xymatrix{
f^*\mathcal{P}^\bullet \ar[r] \ar[d] &
f^*f_*\mathcal{K}^\bullet \ar[d] \ar[r] &
\mathcal{K}^\bullet \ar[d] \\
f^*\mathcal{Q}^\bullet \ar[r] &
f^*f_*\mathcal{I}^\bullet \ar[r] &
\mathcal{I}^\bullet
}
$$
を得る。これはホモトピーを除いて可換である。外側の長方形が、補題の主張が
成り立つことを示している。
\end{proof}



\begin{remark}
\label{sites-cohomology-remark-cup-product}
$f : (\Sh(\mathcal{C}), \mathcal{O}_\mathcal{C}) \to
(\Sh(\mathcal{D}), \mathcal{O}_\mathcal{D})$ を環付きトポスの射とする。
$Lf^*$ と $Rf_*$ の随伴性により、$D(\mathcal{O}_\mathcal{C})$ の任意の
$K, L$ に対して、$D(\mathcal{O}_\mathcal{D})$ における相対カップ積
$$
Rf_*K \otimes_{\mathcal{O}_\mathcal{D}}^\mathbf{L} Rf_*L
\longrightarrow
Rf_*(K \otimes_{\mathcal{O}_\mathcal{C}}^\mathbf{L} L)
$$
を構成できる。実際、この写像は写像
$Lf^*(Rf_*K \otimes_{\mathcal{O}_\mathcal{D}}^\mathbf{L} Rf_*L) \to
K \otimes_{\mathcal{O}_\mathcal{C}}^\mathbf{L} L$ に随伴する。後者としては、同型
$Lf^*(Rf_*K \otimes_{\mathcal{O}_\mathcal{D}}^\mathbf{L} Rf_*L) =
Lf^*Rf_*K \otimes_{\mathcal{O}_\mathcal{C}}^\mathbf{L} Lf^*Rf_*L$
（Lemma \ref{sites-cohomology-lemma-pullback-tensor-product}）と、余単位
$Lf^* \circ Rf_* \to \text{id}$ から得られる写像
$Lf^*Rf_*K \otimes_{\mathcal{O}_\mathcal{C}}^\mathbf{L} Lf^*Rf_*L
\to K \otimes_{\mathcal{O}_\mathcal{C}}^\mathbf{L} L$ との合成を取れる。
\end{remark}

\begin{lemma}
\label{sites-cohomology-lemma-torsion}
$\mathcal{C}$ をサイトとする。$\mathcal{A} \subset \textit{Ab}(\mathcal{C})$
を捩れアーベル層からなる Serre 部分圏とする。このとき関手
$D(\mathcal{A}) \to D_\mathcal{A}(\mathcal{C})$ は同値である。
\end{lemma}

\begin{proof}
鍵となる観察は、入射アーベル層 $\mathcal{I}$ が可除であることである。
実際、$s \in \mathcal{I}(U)$ が局所切断ならば、$s$ を写像
$s : j_{U!}\mathbf{Z} \to \mathcal{I}$ と解釈し、層の単射
$n : j_{U!}\mathbf{Z} \to j_{U!}\mathbf{Z}$ に入射対象の定義的性質を
適用する。すると $ns' = s$ を満たす $s' \in \mathcal{I}(U)$ が存在する。

\medskip\noindent
層 $\mathcal{F}$ に対し、その捩れ部分層を $\mathcal{F}_{tor}$ と書く。
$\mathcal{I}^\bullet$ が入射アーベル層の複体で、そのコホモロジー層が
捩れ層ならば、
$$
\mathcal{I}^\bullet_{tor} \to \mathcal{I}^\bullet
$$
は擬同型であると主張する。実際、$\mathbf{Q}$ は $\mathbf{Z}$ 上平坦なので
$$
H^p(\mathcal{I}^\bullet) \otimes_\mathbf{Z} \mathbf{Q} =
H^p(\mathcal{I}^\bullet \otimes_\mathbf{Z} \mathbf{Q})
$$
であり、これはコホモロジー層が捩れ層であるため 0 である。上で示した
可除性により、
$\mathcal{I}^\bullet \to \mathcal{I}^\bullet \otimes_\mathbf{Z} \mathbf{Q}$
は全射で、核は $\mathcal{I}^\bullet_{tor}$ である。得られる短完全系列に
付随するコホモロジー層の長完全系列から主張が従う。

\medskip\noindent
補題を示すため、右随伴 $T : D(\mathcal{C}) \to D(\mathcal{A})$ を構成する。
実際、$D(\mathcal{C})$ の $K$ が与えられたとき、$K$ をコホモロジー層が
入射的である K-入射複体 $\mathcal{I}^\bullet$ で表せる
（Injectives, Theorem
\ref{injectives-theorem-K-injective-embedding-grothendieck} 参照）。そこで
$T(K) = \mathcal{I}^\bullet_{tor}$ と置く。言い換えれば、$T$ は捩れ部分を
取る関手の右導来関手である。関手 $T$ は
$i : D(\mathcal{A}) \to D_\mathcal{A}(\mathcal{C})$ の右随伴である。
これは、$\mathcal{F}^\bullet$ が捩れ層の複体ならば
$$
\Hom_{K(\mathcal{A})}(\mathcal{F}^\bullet, I^\bullet_{tor}) =
\Hom_{K(\textit{Ab}(\mathcal{C}))}(\mathcal{F}^\bullet, I^\bullet)
$$
であるという観察から直ちに従う。とくに
$\mathcal{I}^\bullet_{tor}$ は $\mathcal{A}$ の K-入射複体である。いくつかの
細部は省略するが、疑義があれば、より一般的な Derived Categories, Lemma
\ref{derived-lemma-derived-adjoint-functors} からも従う。上の主張により、
$D(\mathcal{A})$ の $L$ に対して $L = T(i(L))$ であり、
$D_\mathcal{A}(\mathcal{C})$ の $K$ に対して $i(T(K)) = K$ である。
Categories, Lemma \ref{categories-lemma-adjoint-fully-faithful} を用いれば、
結論が従う。
\end{proof}




\section{K-入射複体のいくつかの性質}
\label{sites-cohomology-section-properties-K-injective}

\noindent
$(\mathcal{C}, \mathcal{O})$ を環付きサイトとし、$U$ を
$\mathcal{C}$ の対象とする。対応する局所化射を
$j : (\Sh(\mathcal{C}/U), \mathcal{O}_U) \to (\Sh(\mathcal{C}), \mathcal{O})$
と書く。引き戻し関手 $j^*$ は単なる制限関手なので完全である。
したがって導来引き戻し $Lj^*$ は、任意の複体を単に制限することで
計算される。対応する関手をしばしば単に
$$
D(\mathcal{O}) \to D(\mathcal{O}_U),
\quad
E \mapsto j^*E = E|_U
$$
と書く。同様に、零による延長
$j_! : \textit{Mod}(\mathcal{O}_U) \to \textit{Mod}(\mathcal{O})$
（Modules on Sites, Definition
\ref{sites-modules-definition-localize-ringed-site} 参照）は完全関手である
（Modules on Sites, Lemma \ref{sites-modules-lemma-extension-by-zero-exact}）。
したがって関手
$$
j_! : D(\mathcal{O}_U) \to D(\mathcal{O}), \quad
F \mapsto j_!F
$$
を誘導する。これは対象 $F$ を表す任意の複体に $j_!$ を施すだけで得られる。

\begin{lemma}
\label{sites-cohomology-lemma-restrict-K-injective-to-open}
$(\mathcal{C}, \mathcal{O})$ を環付きサイトとし、$U$ を
$\mathcal{C}$ の対象とする。$\mathcal{O}$-加群の K-入射複体を
$\mathcal{C}/U$ に制限したものは、$\mathcal{O}_U$-加群の
K-入射複体である。
\end{lemma}

\begin{proof}
これは Derived Categories, Lemma
\ref{derived-lemma-adjoint-preserve-K-injectives} と、制限関手が
完全な左随伴 $j_!$ をもつことから直ちに従う。上の議論を参照せよ。
\end{proof}

\begin{lemma}
\label{sites-cohomology-lemma-unbounded-cohomology-of-open}
$(\mathcal{C}, \mathcal{O})$ を環付きサイトとし、
$U \in \Ob(\mathcal{C})$ とする。$D(\mathcal{O})$ の $K$ に対して
$H^p(U, K) = H^p(\mathcal{C}/U, K|_{\mathcal{C}/U})$ である。
\end{lemma}

\begin{proof}
$\mathcal{I}^\bullet$ を $K$ を表す $\mathcal{O}$-加群の
K-入射複体とする。このとき
$$
H^q(U, K) = H^q(\Gamma(U, \mathcal{I}^\bullet)) =
H^q(\Gamma(\mathcal{C}/U, \mathcal{I}^\bullet|_{\mathcal{C}/U}))
$$
である。これはコホモロジーの構成による。Lemma
\ref{sites-cohomology-lemma-restrict-K-injective-to-open} により、
複体 $\mathcal{I}^\bullet|_{\mathcal{C}/U}$ は
$K|_{\mathcal{C}/U}$ を表す K-入射複体であり、補題が従う。
\end{proof}

\begin{lemma}
\label{sites-cohomology-lemma-sheafification-cohomology}
$(\mathcal{C}, \mathcal{O})$ を環付きサイトとし、$K$ を
$D(\mathcal{O})$ の対象とする。
$$
U \mapsto H^q(U, K) = H^q(\mathcal{C}/U, K|_{\mathcal{C}/U})
$$
の層化は、$K$ の第 $q$ コホモロジー層 $H^q(K)$ である。
\end{lemma}

\begin{proof}
等式 $H^q(U, K) = H^q(\mathcal{C}/U, K|_{\mathcal{C}/U})$ は
Lemma \ref{sites-cohomology-lemma-unbounded-cohomology-of-open} による。
$K$ を表す K-入射複体 $\mathcal{I}^\bullet$ を選ぶ。このとき
$$
H^q(U, K) =
\frac{\Ker(\mathcal{I}^q(U) \to \mathcal{I}^{q + 1}(U))}
{\Im(\mathcal{I}^{q - 1}(U) \to \mathcal{I}^q(U))}.
$$
これはコホモロジーの構成による。また、
$H^q(K) = \Ker(\mathcal{I}^q \to \mathcal{I}^{q + 1})/
\Im(\mathcal{I}^{q - 1} \to \mathcal{I}^q)$ であるから、結果は明らかである。
\end{proof}

\begin{lemma}
\label{sites-cohomology-lemma-restrict-direct-image-open}
$f : (\mathcal{C}, \mathcal{O}_\mathcal{C}) \to
(\mathcal{D}, \mathcal{O}_\mathcal{D})$ を、連続関手
$u : \mathcal{D} \to \mathcal{C}$ に対応する環付きサイトの射とする。
$V \in \mathcal{D}$ が与えられたとき $U = u(V)$ と置き、誘導される
環付きサイトの射
$g : (\mathcal{C}/U, \mathcal{O}_U) \to (\mathcal{D}/V, \mathcal{O}_V)$
を考える（Modules on Sites, Lemma
\ref{sites-modules-lemma-localize-morphism-ringed-sites}）。
このとき $D(\mathcal{O}_\mathcal{C})$ の $E$ に対して
$(Rf_*E)|_{\mathcal{D}/V} = Rg_*(E|_{\mathcal{C}/U})$ である。
\end{lemma}

\begin{proof}
$E$ を $\mathcal{O}_\mathcal{C}$-加群の K-入射複体
$\mathcal{I}^\bullet$ で表す。このとき
$Rf_*(E) = f_*\mathcal{I}^\bullet$ であり、Lemma
\ref{sites-cohomology-lemma-restrict-K-injective-to-open} により
$Rg_*(E|_{\mathcal{C}/U}) = g_*(\mathcal{I}^\bullet|_{\mathcal{C}/U})$
である。$\mathcal{C}$ 上の任意の層 $\mathcal{F}$ に対して
$(f_*\mathcal{F})|_{\mathcal{D}/V} = g_*(\mathcal{F}|_{\mathcal{C}/U})$
であることは明らかなので、結果が従う（Modules on Sites, Lemma
\ref{sites-modules-lemma-localize-morphism-ringed-sites} または、より基本的な
Sites, Lemma \ref{sites-lemma-localize-morphism} を参照せよ）。
\end{proof}


\begin{lemma}
\label{sites-cohomology-lemma-Leray-unbounded}
$f : (\mathcal{C}, \mathcal{O}_\mathcal{C}) \to
(\mathcal{D}, \mathcal{O}_\mathcal{D})$ を、連続関手
$u : \mathcal{D} \to \mathcal{C}$ に対応する環付きサイトの射とする。
このとき $D(\mathcal{O}_\mathcal{C}) \to
D(\Gamma(\mathcal{O}_\mathcal{D}))$ の関手として
$R\Gamma(\mathcal{D}, -) \circ Rf_* = R\Gamma(\mathcal{C}, -)$
である。より一般に、$V \in \mathcal{D}$ と $U = u(V)$ に対して
$R\Gamma(U, -) = R\Gamma(V, -) \circ Rf_*$ である。
\end{lemma}

\begin{proof}
一点トポス $pt$ に、環 $\Gamma(\mathcal{O}_\mathcal{D})$ で与えられる
$\mathcal{O}_{pt}$ を入れる。構造層の大域切断上で同一視を誘導する
環付きトポスの標準射
$(\mathcal{D}, \mathcal{O}_\mathcal{D}) \to (pt, \mathcal{O}_{pt})$
がある。このとき
$D(\mathcal{O}_{pt}) = D(\Gamma(\mathcal{O}_\mathcal{D}))$ である。
したがって主張
$R\Gamma(\mathcal{D}, -) \circ Rf_* = R\Gamma(\mathcal{C}, -)$ は、
Lemma \ref{sites-cohomology-lemma-derived-pushforward-composition} を
$$
(\mathcal{C}, \mathcal{O}_\mathcal{C}) \to
(\mathcal{D}, \mathcal{O}_\mathcal{D}) \to (pt, \mathcal{O}_{pt})
$$
に適用して得られる。第二の（より一般的な）主張は、第一の主張に
Lemma \ref{sites-cohomology-lemma-restrict-direct-image-open} を適用して従う。
\end{proof}

\begin{lemma}
\label{sites-cohomology-lemma-unbounded-describe-higher-direct-images}
$f : (\mathcal{C}, \mathcal{O}_\mathcal{C}) \to
(\mathcal{D}, \mathcal{O}_\mathcal{D})$ を、連続関手
$u : \mathcal{D} \to \mathcal{C}$ に対応する環付きサイトの射とする。
$K$ を $D(\mathcal{O}_\mathcal{C})$ の対象とする。このとき
$H^i(Rf_*K)$ は前層
$$
V \mapsto H^i(u(V), K) = H^i(V, Rf_*K)
$$
に付随する層である。
\end{lemma}

\begin{proof}
等式 $H^i(u(V), K) = H^i(V, Rf_*K)$ は、Lemma
\ref{sites-cohomology-lemma-Leray-unbounded} の第二の主張でコホモロジーを取って得られる。
層化に関する主張は Lemma \ref{sites-cohomology-lemma-sheafification-cohomology} から従う。
\end{proof}

\begin{lemma}
\label{sites-cohomology-lemma-modules-abelian-unbounded}
$(\mathcal{C}, \mathcal{O}_\mathcal{C})$ を環付きサイトとする。
$K$ を $D(\mathcal{O}_\mathcal{C})$ の対象とし、
$D(\underline{\mathbf{Z}}_\mathcal{C})$ におけるその像を
$K_{ab}$ と書く。
\begin{enumerate}
\item 標準写像
$R\Gamma(\mathcal{C}, K) \to R\Gamma(\mathcal{C}, K_{ab})$
があり、これは $D(\textit{Ab})$ における同型である。
\item 任意の $U \in \mathcal{C}$ に対して標準写像
$R\Gamma(U, K) \to R\Gamma(U, K_{ab})$
があり、これは $D(\textit{Ab})$ における同型である。
\item $f : (\mathcal{C}, \mathcal{O}_\mathcal{C}) \to
(\mathcal{D}, \mathcal{O}_\mathcal{D})$ を環付きサイトの射とする。
標準写像 $Rf_*K \to Rf_*(K_{ab})$ があり、これは
$D(\underline{\mathbf{Z}}_\mathcal{D})$ における同型である。
\end{enumerate}
\end{lemma}

\begin{proof}
写像を次のように構成する。$K$ を表す K-入射複体
$\mathcal{I}^\bullet$ を選ぶ。また、$\mathcal{J}^\bullet$ が
アーベル群の層の K-入射複体であるような擬同型
$\mathcal{I}^\bullet \to \mathcal{J}^\bullet$ を選ぶ。このとき (1) の写像は
$\Gamma(\mathcal{C}, \mathcal{I}^\bullet) \to
\Gamma(\mathcal{C}, \mathcal{J}^\bullet)$
で与えられ、(2) の写像は
$\Gamma(U, \mathcal{I}^\bullet) \to \Gamma(U, \mathcal{J}^\bullet)$
で与えられ、(3) の写像は
$f_*\mathcal{I}^\bullet \to f_*\mathcal{J}^\bullet$ で与えられる。
これらの写像が同型であることを示すには、コホモロジー群および
コホモロジー層上で同型を誘導することを証明すれば十分である。
Lemmas \ref{sites-cohomology-lemma-unbounded-cohomology-of-open} および
\ref{sites-cohomology-lemma-unbounded-describe-higher-direct-images} により、写像
$$
H^0(\mathcal{C}, K) \longrightarrow H^0(\mathcal{C}, K_{ab})
$$
が同型であることを示せば十分である。ここで
$$
H^0(\mathcal{C}, K) =
\Hom_{D(\mathcal{O}_\mathcal{C})}(\mathcal{O}_\mathcal{C}, K)
$$
であり、もう一方の群についても同様である。$K$ を表す任意の
$\mathcal{O}_\mathcal{C}$-加群の複体 $\mathcal{K}^\bullet$ を選ぶ。
導来圏を局所化として構成することにより、
$$
\Hom_{D(\mathcal{O}_\mathcal{C})}(\mathcal{O}_\mathcal{C}, K) =
\colim_{s : \mathcal{F}^\bullet \to \mathcal{O}_\mathcal{C}}
\Hom_{K(\mathcal{O}_\mathcal{C})}(\mathcal{F}^\bullet, \mathcal{K}^\bullet)
$$
を得る。ここで余極限は、$\mathcal{O}_\mathcal{C}$-加群の複体の
擬同型 $s$ 全体にわたって取る。同様に、
$$
\Hom_{D(\underline{\mathbf{Z}}_\mathcal{C})}
(\underline{\mathbf{Z}}_\mathcal{C}, K) =
\colim_{s : \mathcal{G}^\bullet \to \underline{\mathbf{Z}}_\mathcal{C}}
\Hom_{K(\underline{\mathbf{Z}}_\mathcal{C})}
(\mathcal{G}^\bullet, \mathcal{K}^\bullet)
$$
を得る。次に、$\mathcal{G}^\bullet$ が平坦な
$\underline{\mathbf{Z}}_\mathcal{C}$-加群からなる上に有界な複体であるような
擬同型 $s : \mathcal{G}^\bullet \to \underline{\mathbf{Z}}_\mathcal{C}$
はこの系で共終であることに注意する（これは Modules on Sites, Lemma
\ref{sites-modules-lemma-module-quotient-flat} および
Derived Categories, Lemma \ref{derived-lemma-subcategory-left-resolution}
から従う。Section \ref{sites-cohomology-section-flat} の議論を参照せよ）。
したがって写像
$H^0(\mathcal{C}, K) \longrightarrow H^0(\mathcal{C}, K_{ab})$
の逆を次のように構成できる。元
$\xi \in H^0(\mathcal{C}, K_{ab})$ を対
$$
(s : \mathcal{G}^\bullet \to \underline{\mathbf{Z}}_\mathcal{C},
a : \mathcal{G}^\bullet \to \mathcal{K}^\bullet)
$$
で表す。ここで $\mathcal{G}^\bullet$ は平坦な
$\underline{\mathbf{Z}}_\mathcal{C}$-加群からなる上に有界な複体である。
これを
$$
(\mathcal{G}^\bullet \otimes_{\underline{\mathbf{Z}}_\mathcal{C}}
\mathcal{O}_\mathcal{C}
\to \mathcal{O}_\mathcal{C},
\mathcal{G}^\bullet  \otimes_{\underline{\mathbf{Z}}_\mathcal{C}}
\mathcal{O}_\mathcal{C}
\to \mathcal{K}^\bullet)
$$
に送る。ここで注意すべき唯一の点は、Lemmas
\ref{sites-cohomology-lemma-derived-tor-quasi-isomorphism-other-side} および
\ref{sites-cohomology-lemma-bounded-flat-K-flat} により、第一の矢印が擬同型であることである。
この構成が実際に逆写像であることの詳細な検証は省略する。
\end{proof}


\begin{lemma}
\label{sites-cohomology-lemma-adjoint-lower-shriek-restrict}
$(\mathcal{C}, \mathcal{O})$ を環付きサイトとし、$U$ を
$\mathcal{C}$ の対象とする。対応する局所化射を
$j : (\Sh(\mathcal{C}/U), \mathcal{O}_U) \to (\Sh(\mathcal{C}), \mathcal{O})$
と書く。制限関手 $D(\mathcal{O}) \to D(\mathcal{O}_U)$ は、
零による延長
$j_! : D(\mathcal{O}_U) \to D(\mathcal{O})$ の右随伴である。
\end{lemma}

\begin{proof}
示すべきことは
$$
\Hom_{D(\mathcal{O})}(j_!E, F) = \Hom_{D(\mathcal{O}_U)}(E, F|_U)
$$
である。$E$ を表す $\mathcal{O}_U$-加群の複体
$\mathcal{E}^\bullet$ と、$F$ を表す K-入射複体
$\mathcal{I}^\bullet$ を選ぶ。Lemma
\ref{sites-cohomology-lemma-restrict-K-injective-to-open} により、複体
$\mathcal{I}^\bullet|_U$ も K-入射的である。したがって上の等式は
$$
\Hom_{D(\mathcal{O})}(j_!\mathcal{E}^\bullet, \mathcal{I}^\bullet) =
\Hom_{D(\mathcal{O}_U)}(\mathcal{E}^\bullet, \mathcal{I}^\bullet|_U)
$$
となる。これは $|_U$ と $j_!$ が随伴関手であること
（Modules on Sites, Lemma
\ref{sites-modules-lemma-extension-by-zero}）および
Derived Categories, Lemma \ref{derived-lemma-K-injective} から従う。
\end{proof}

\begin{lemma}
\label{sites-cohomology-lemma-j-shriek-and-tensor}
$(\mathcal{C}, \mathcal{O})$ を環付きサイトとし、
$U \in \Ob(\mathcal{C})$ とする。$D(\mathcal{O}_U)$ の $L$ と
$D(\mathcal{O})$ の $K$ に対して
$j_!L \otimes_\mathcal{O}^\mathbf{L} K =
j_!(L \otimes_{\mathcal{O}_U}^\mathbf{L} K|_U)$ である。
\end{lemma}

\begin{proof}
$L$ を $\mathcal{O}_U$-加群の複体で、$K$ を
$\mathcal{O}$-加群の K-平坦複体で表し、Modules on Sites, Lemma
\ref{sites-modules-lemma-j-shriek-and-tensor} を適用する。詳細は省略する。
\end{proof}

\begin{lemma}
\label{sites-cohomology-lemma-K-injective-flat}
$f : (\Sh(\mathcal{C}), \mathcal{O}_\mathcal{C}) \to
(\Sh(\mathcal{D}), \mathcal{O}_\mathcal{D})$ を環付きトポスの平坦射とする。
$\mathcal{I}^\bullet$ が $\mathcal{O}_\mathcal{C}$-加群の
K-入射複体ならば、$f_*\mathcal{I}^\bullet$ は
$\mathcal{O}_\mathcal{D}$-加群の複体として K-入射的である。
\end{lemma}

\begin{proof}
実際、
$$
\Hom_{K(\mathcal{O}_\mathcal{D})}(\mathcal{F}^\bullet, f_*\mathcal{I}^\bullet)
=
\Hom_{K(\mathcal{O}_\mathcal{C})}(f^*\mathcal{F}^\bullet, \mathcal{I}^\bullet)
$$
である。これは Modules on Sites, Lemma
\ref{sites-modules-lemma-adjoint-pullback-pushforward-modules} と、
$f$ が平坦であるため $f^*$ が完全であることによる。
\end{proof}

\begin{lemma}
\label{sites-cohomology-lemma-hom-K-injective}
$\mathcal{C}$ をサイトとし、$\mathcal{O} \to \mathcal{O}'$ を
環の層の写像とする。$\mathcal{I}^\bullet$ が
$\mathcal{O}$-加群の K-入射複体ならば、
$\SheafHom_\mathcal{O}(\mathcal{O}', \mathcal{I}^\bullet)$ は
$\mathcal{O}'$-加群の K-入射複体である。
\end{lemma}

\begin{proof}
実際、Modules on Sites, Lemma
\ref{sites-modules-lemma-adjoint-hom-restrict} により
$\Hom_{K(\mathcal{O}')}(\mathcal{G}^\bullet,
\Hom_\mathcal{O}(\mathcal{O}', \mathcal{I}^\bullet)) =
\Hom_{K(\mathcal{O})}(\mathcal{G}^\bullet, \mathcal{I}^\bullet)$
だからである。
\end{proof}






\section{局所化とコホモロジー}
\label{sites-cohomology-section-localization}

\noindent
$\mathcal{C}$ をサイトとし、$f : X \to Y$ を $\mathcal{C}$ の射とする。
このときトポスの射
$$
j_{X/Y} : \Sh(\mathcal{C}/X) \longrightarrow \Sh(\mathcal{C}/Y)
$$
が得られる。Sites, Sections \ref{sites-section-localize} および
\ref{sites-section-more-localize} を参照せよ。この型のトポスの射では、
コホモロジーに関するいくつかの問題がより容易になる。ここでは、
自明な型の基底変換定理が得られる例を挙げる。

\begin{lemma}
\label{sites-cohomology-lemma-localize-cartesian-square}
$\mathcal{C}$ をサイトとし、
$$
\xymatrix{
X' \ar[d] \ar[r] & X \ar[d] \\
Y' \ar[r] & Y
}
$$
を $\mathcal{C}$ の Cartesian 図式とする。このとき
$D(\mathcal{C}/X) \to D(\mathcal{C}/Y')$ の関手として
$j_{Y'/Y}^{-1} \circ Rj_{X/Y, *} =
Rj_{X'/Y', *} \circ j_{X'/X}^{-1}$ である。
\end{lemma}

\begin{proof}
$E \in D(\mathcal{C}/X)$ とする。$E$ を表す
$\mathcal{C}/X$ 上のアーベル層の K-入射複体
$\mathcal{I}^\bullet$ を選ぶ。Lemma
\ref{sites-cohomology-lemma-restrict-K-injective-to-open} により、
$j_{X'/X}^{-1}\mathcal{I}^\bullet$ も K-入射的である。
したがって $Rj_{X'/Y'}(j_{X'/X}^{-1}E)$ は
$j_{X'/Y', *}j_{X'/X}^{-1}\mathcal{I}^\bullet$ によって計算できる。
よって Sites, Lemma \ref{sites-lemma-localize-cartesian-square} から
この等式が従う。
\end{proof}

\noindent
環付きサイト $(\mathcal{C}, \mathcal{O})$ と $\mathcal{C}$ の射
$f : X \to Y$ が与えられると、$j_{X/Y}$ は環付きトポスの射
$$
j_{X/Y} :
(\Sh(\mathcal{C}/X), \mathcal{O}_X)
\longrightarrow
(\Sh(\mathcal{C}/Y), \mathcal{O}_Y)
$$
となる。Modules on Sites, Lemma
\ref{sites-modules-lemma-relocalize} を参照せよ。

\begin{lemma}
\label{sites-cohomology-lemma-localize-cartesian-square-modules}
$(\mathcal{C}, \mathcal{O})$ を環付きサイトとし、
$$
\xymatrix{
X' \ar[d] \ar[r] & X \ar[d] \\
Y' \ar[r] & Y
}
$$
を $\mathcal{C}$ の Cartesian 図式とする。このとき
$D(\mathcal{O}_X) \to D(\mathcal{O}_{Y'})$ の関手として
$j_{Y'/Y}^* \circ Rj_{X/Y, *} =
Rj_{X'/Y', *} \circ j_{X'/X}^*$ である。
\end{lemma}

\begin{proof}
$j_{Y'/Y}^{-1}\mathcal{O}_Y = \mathcal{O}_{Y'}$ なので
$j_{Y'/Y}^* = Lj_{Y'/Y}^* = j_{Y'/Y}^{-1}$ である。同様に
$j_{X'/X}^* = Lj_{X'/X}^* = j_{X'/X}^{-1}$ である。したがって
Lemma \ref{sites-cohomology-lemma-modules-abelian-unbounded} により、
アーベル層の導来圏上で結果を示せば十分であり、これは
Lemma \ref{sites-cohomology-lemma-localize-cartesian-square} で示した。
\end{proof}













\section{逆系とコホモロジー}
\label{sites-cohomology-section-inverse-systems}

\noindent
加群の層の逆系に関するいくつかの結果を証明する。

\begin{lemma}
\label{sites-cohomology-lemma-ML-general}
$I$ を環 $A$ のイデアル、$\mathcal{C}$ をサイトとする。
$$
\ldots \to \mathcal{F}_3 \to \mathcal{F}_2 \to \mathcal{F}_1
$$
を、$\mathcal{F}_n = \mathcal{F}_{n + 1}/I^n\mathcal{F}_{n + 1}$
を満たす $\mathcal{C}$ 上の $A$-加群の層の逆系とする。$p \geq 0$ とする。
$$
\bigoplus\nolimits_{n \geq 0} H^{p + 1}(\mathcal{C}, I^n\mathcal{F}_{n + 1})
$$
が次数付き $\bigoplus_{n \geq 0} I^n/I^{n + 1}$-加群として
昇鎖条件を満たすと仮定する。このとき逆系
$M_n = H^p(\mathcal{C}, \mathcal{F}_n)$ は Mittag--Leffler 条件を満たす。
\footnote{実際、ある $c \geq 0$ が存在して、すべての $n \geq c$ に対し
$\Im(M_n \to M_{n - c})$ が安定像となる。}
\end{lemma}

\begin{proof}
$N_n = H^{p + 1}(\mathcal{C}, I^n\mathcal{F}_{n + 1})$ と置き、
$\delta_n : M_n \to N_n$ を短完全列
$0 \to I^n\mathcal{F}_{n + 1} \to \mathcal{F}_{n + 1} \to \mathcal{F}_n \to 0$
から得られるコホモロジー上の境界写像とする。このとき
$\bigoplus \Im(\delta_n) \subset \bigoplus N_n$ は次数付き部分加群である。
実際、$s \in M_n$ かつ $f \in I^m$ ならば、可換図式
$$
\xymatrix{
0 \ar[r] &
I^n\mathcal{F}_{n + 1} \ar[d]_f \ar[r] &
\mathcal{F}_{n + 1} \ar[d]_f \ar[r] &
\mathcal{F}_n \ar[d]_f \ar[r] & 0 \\
0 \ar[r] &
I^{n + m}\mathcal{F}_{n + m + 1} \ar[r] &
\mathcal{F}_{n + m + 1} \ar[r] &
\mathcal{F}_{n + m} \ar[r] & 0
}
$$
を得る。中央の垂直写像は、$\mathcal{F}_{n + 1}$ の局所切断を
$\mathcal{F}_{n + m + 1}$ の切断へ持ち上げてから $f$ を掛けることで
与えられる。他の垂直射についても同様である。したがって
$\delta_{n + m}(fs) = f \delta_n(s)$ である。仮定により、
$\delta_{n_j}(s_j)$ が $\bigoplus \Im(\delta_n)$ を次数付き加群として
生成するような $s_j \in M_{n_j}$, $j = 1, \ldots, N$ を取れる。
$n > c = \max(n_j)$ とし、$s \in M_n$ とする。このとき
$\delta_n(s) = \sum f_j \delta_{n_j}(s_j)$ を満たす
$f_j \in I^{n - n_j}$ が存在する。したがって
$\delta(s - \sum f_j s_j) = 0$ であり、すなわち
$M_n$ で $s - \sum f_js_j$ へ写る $s' \in M_{n + 1}$ が存在する。
ゆえに
$$
\Im(M_{n + 1} \to M_{n - c}) = \Im(M_n \to M_{n - c})
$$
である。実際、元 $f_js_j$ は $M_{n - c}$ で零に写る。
これで補題が証明された。
\end{proof}

\begin{lemma}
\label{sites-cohomology-lemma-ML-general-better}
$I$ を環 $A$ のイデアル、$\mathcal{C}$ をサイトとする。
$$
\ldots \to \mathcal{F}_3 \to \mathcal{F}_2 \to \mathcal{F}_1
$$
を、$\mathcal{F}_n = \mathcal{F}_{n + 1}/I^n\mathcal{F}_{n + 1}$
を満たす $\mathcal{C}$ 上の $A$-加群の逆系とする。$p \geq 0$ とする。
$n$ に対して
$$
N_n =
\bigcap\nolimits_{m \geq n}
\Im\left(
H^{p + 1}(\mathcal{C}, I^n\mathcal{F}_{m + 1}) \to H^{p + 1}(\mathcal{C}, I^n\mathcal{F}_{n + 1})
\right)
$$
と定める。$\bigoplus N_n$ が次数付き
$\bigoplus_{n \geq 0} I^n/I^{n + 1}$-加群として昇鎖条件を満たすならば、
逆系 $M_n = H^p(\mathcal{C}, \mathcal{F}_n)$ は Mittag--Leffler 条件を満たす。
\footnote{実際、ある $c \geq 0$ が存在して、すべての $n \geq c$ に対し
$\Im(M_n \to M_{n - c})$ が安定像となる。}
\end{lemma}

\begin{proof}
証明は補題 \ref{sites-cohomology-lemma-ML-general} の証明とまったく同じである。
実際、$\bigoplus N_n$ が
$\bigoplus H^{p + 1}(\mathcal{C}, I^n\mathcal{F}_{n + 1})$ の次数付き
$\bigoplus_{n \geq 0} I^n/I^{n + 1}$-部分加群であること、および境界写像
$\delta_n : M_n \to H^{p + 1}(\mathcal{C}, I^n\mathcal{F}_{n + 1})$
の像が $N_n$ に含まれることを示せば、上の議論から結果が従う。

\medskip\noindent
$\xi \in N_n$ および $f \in I^k$ と仮定する。
$m \gg n + k$ を選び、$\xi$ の持ち上げ
$\xi' \in H^{p + 1}(\mathcal{C}, I^n\mathcal{F}_{m + 1})$ を選ぶ。
図式
$$
\xymatrix{
0 \ar[r] &
I^n\mathcal{F}_{m + 1} \ar[d]_f \ar[r] &
\mathcal{F}_{m + 1} \ar[d]_f \ar[r] &
\mathcal{F}_n \ar[d]_f \ar[r] & 0 \\
0 \ar[r] &
I^{n + k}\mathcal{F}_{m + 1} \ar[r] &
\mathcal{F}_{m + 1} \ar[r] &
\mathcal{F}_{n + k} \ar[r] & 0
}
$$
を補題 \ref{sites-cohomology-lemma-ML-general} の証明と同様に構成する。
コホモロジー上に誘導される写像を得て、
$f \xi' \in H^{p + 1}(\mathcal{C}, I^{n + k}\mathcal{F}_{m + 1})$ が
$f \xi$ に写ることが分かる。これはすべての $m \gg n + k$ に対して
成り立つので、望むとおり $f\xi$ は $N_{n + k}$ に属する。

\medskip\noindent
境界写像 $\delta_n$ の像が $N_n$ に含まれることを見るため、図式
$$
\xymatrix{
0 \ar[r] &
I^n\mathcal{F}_{m + 1} \ar[d] \ar[r] &
\mathcal{F}_{m + 1} \ar[d] \ar[r] &
\mathcal{F}_n \ar[d] \ar[r] & 0 \\
0 \ar[r] &
I^n\mathcal{F}_{n + 1} \ar[r] &
\mathcal{F}_{n + 1} \ar[r] &
\mathcal{F}_n \ar[r] & 0
}
$$
を $m \geq n$ に対して考える。コホモロジー上に誘導される写像を
調べれば結論を得る。
\end{proof}

\begin{lemma}
\label{sites-cohomology-lemma-topology-I-adic-general}
$I$ を環 $A$ のイデアル、$\mathcal{C}$ をサイトとする。
$$
\ldots \to \mathcal{F}_3 \to \mathcal{F}_2 \to \mathcal{F}_1
$$
を、$\mathcal{F}_n = \mathcal{F}_{n + 1}/I^n\mathcal{F}_{n + 1}$
を満たす $\mathcal{C}$ 上の $A$-加群の層の逆系とする。$p \geq 0$ とする。
$$
\bigoplus\nolimits_{n \geq 0} H^p(\mathcal{C}, I^n\mathcal{F}_{n + 1})
$$
が次数付き $\bigoplus_{n \geq 0} I^n/I^{n + 1}$-加群として
昇鎖条件を満たすと仮定する。このとき
$M = \lim H^p(\mathcal{C}, \mathcal{F}_n)$ 上の極限位相は $I$-進位相である。
\end{lemma}

\begin{proof}
$n \geq 1$ に対して
$F^n = \Ker(M \to H^p(\mathcal{C}, \mathcal{F}_n))$ と置き、$F^0 = M$ とする。
$I F^n \subset F^{n + 1}$ であり、特に $I^n M \subset F^n$ である。
したがって $I$-進位相は極限位相より細かい。逆を示すため、任意の $n$ に対し
$F^m \subset I^nM$ を満たす $m \geq n$ が存在することを示す。
\footnote{実際、ある $c \geq 0$ が存在して、すべての $n$ に対し
$F^{n + c} \subset I^nM$ となる。}
単射
$$
F^n/F^{n + 1} \longrightarrow H^p(\mathcal{C}, \mathcal{F}_{n + 1})
$$
があり、その像は
$H^p(\mathcal{C}, I^n\mathcal{F}_{n + 1}) \to H^p(\mathcal{C}, \mathcal{F}_{n + 1})$
の像に含まれる。この写像による $F^n/F^{n + 1}$ の逆像を
$$
E_n \subset H^p(\mathcal{C}, I^n\mathcal{F}_{n + 1})
$$
と書く。このとき $\bigoplus E_n$ は
$\bigoplus H^p(\mathcal{C}, I^n\mathcal{F}_{n + 1})$ の次数付き
$\bigoplus I^n/I^{n + 1}$-部分加群であり、
$\bigoplus E_n \to \bigoplus F^n/F^{n + 1}$ は次数付き加群の準同型である。
詳細は省略する。仮定により、$\bigoplus E_n$ は
$\bigoplus I^n/I^{n + 1}$ 上有限個の斉次な元で生成される。
$E_n \to F^n/F^{n + 1}$ は全射なので、
$\bigoplus F^n/F^{n + 1}$ についても同じことが成り立つ。
したがって、$r$、$c_1, \ldots, c_r \geq 0$、および
$\bigoplus F^n/F^{n + 1}$ における像がこれを生成する
$a_i \in F^{c_i}$ を取れる。$c = \max(c_i)$ と置く。

\medskip\noindent
$n \geq c$ に対して $I F^n = F^{n + 1}$ と主張する。この主張から、
望むように $F^{n + c} = I^nF^c \subset I^nM$ が従う。
主張を証明するため $a \in F^{n + 1}$ とする。
$F^{n + 1}/F^{n + 2}$ における $a$ の像は $a_i$ の線形結合である。
したがって、ある $f_i \in I^{n + 1 - c_i}$ に対して
$a  - \sum f_i a_i \in F^{n + 2}$ である。
$n \geq c_i$ なので
$I^{n + 1 - c_i} = I \cdot I^{n - c_i}$ であり、
$g_{i, j} \in I$ かつ $h_{i, j}a_i \in F^n$ となるように
$f_i = \sum g_{i, j} h_{i, j}$ と書ける。ゆえに
$F^{n + 1} = F^{n + 2} + IF^n$ である。
簡単な帰納法により、すべての $e > 0$ に対して
$F^{n + 1} = F^{n + e} + IF^n$ を得る。したがって $IF^n$ は
$F^{n + 1}$ で稠密である。$I$ の生成元
$k_1, \ldots, k_r$ を選び、連続写像
$$
u : (F^n)^{\oplus r} \longrightarrow F^{n + 1},\quad
(x_1, \ldots, x_r) \mapsto \sum k_i x_i
$$
を考える（極限位相に関して）。上の議論により、すべての $m \geq n$ に対し、
$u$ による $(F^m)^{\oplus r}$ の像は $F^{m + 1}$ で稠密である。
開写像補題（More on Algebra, Lemma
\ref{more-algebra-lemma-open-mapping}）により、$u$ は開写像である。
したがって $u$ は全射である。ゆえに $n \geq c$ に対して
$IF^n = F^{n + 1}$ であり、証明は完了する。
\end{proof}

\begin{lemma}
\label{sites-cohomology-lemma-topology-I-adic-general-better}
$I$ を環 $A$ のイデアル、$\mathcal{C}$ をサイトとする。
$$
\ldots \to \mathcal{F}_3 \to \mathcal{F}_2 \to \mathcal{F}_1
$$
を、$\mathcal{F}_n = \mathcal{F}_{n + 1}/I^n\mathcal{F}_{n + 1}$
を満たす $\mathcal{C}$ 上の $A$-加群の層の逆系とする。$p \geq 0$ とする。
$n$ に対して
$$
N_n =
\bigcap\nolimits_{m \geq n}
\Im\left(
H^p(\mathcal{C}, I^n\mathcal{F}_{m + 1}) \to H^p(\mathcal{C}, I^n\mathcal{F}_{n + 1})
\right)
$$
と定める。$\bigoplus N_n$ が次数付き
$\bigoplus_{n \geq 0} I^n/I^{n + 1}$-加群として昇鎖条件を満たすならば、
$M = \lim H^p(\mathcal{C}, \mathcal{F}_n)$ 上の極限位相は $I$-進位相である。
\end{lemma}

\begin{proof}
証明は補題 \ref{sites-cohomology-lemma-topology-I-adic-general} の証明とまったく同じである。
実際、$\bigoplus N_n$ が
$\bigoplus H^{p + 1}(\mathcal{C}, I^n\mathcal{F}_{n + 1})$ の次数付き
$\bigoplus_{n \geq 0} I^n/I^{n + 1}$-部分加群であること、および
$F^n/F^{n + 1} \subset H^p(\mathcal{C}, \mathcal{F}_{n + 1})$ が
$N_n \to H^p(\mathcal{C}, \mathcal{F}_{n + 1})$ の像に含まれることを示せば、
上の議論から結果が従う。加群構造に関する主張は補題
\ref{sites-cohomology-lemma-ML-general-better} の証明で見た。

\medskip\noindent
$t \in F^n$ とする。$H^p(\mathcal{C}, \mathcal{F}_{n + 1})$ における $t$ の像へ
写る元 $s \in H^p(\mathcal{C}, I^n\mathcal{F}_{n + 1})$ を選ぶ。
$s \in N_n$ を示す必要がある。$F^n$ は写像
$M \to H^p(\mathcal{C}, \mathcal{F}_n)$ の核なので、すべての $m \geq n$ に対し、
$t$ を $H^p(\mathcal{C}, \mathcal{F}_n)$ で零に写る元
$t_m \in H^p(\mathcal{C}, \mathcal{F}_{m + 1})$ へ写せる。短完全列
$0 \to I^n\mathcal{F}_{m + 1} \to \mathcal{F}_{m + 1} \to \mathcal{F}_n \to 0$
から得られるコホモロジー列
$$
H^{p - 1}(\mathcal{C}, \mathcal{F}_n) \to
H^p(\mathcal{C}, I^n\mathcal{F}_{m + 1}) \to
H^p(\mathcal{C}, \mathcal{F}_{m + 1}) \to
H^p(\mathcal{C}, \mathcal{F}_n)
$$
を考える。$t_m$ へ写る
$s_m \in H^p(\mathcal{C}, I^n\mathcal{F}_{m + 1})$ を選べる。
上の列を $m = n$ の場合の列と比較すると、$s_m$ は
$H^{p - 1}(\mathcal{C}, \mathcal{F}_n) \to H^p(\mathcal{C}, I^n\mathcal{F}_{n + 1})$
の像に属する元を除いて $s$ へ写ることが分かる。しかしこの写像は
$H^p(\mathcal{C}, I^n\mathcal{F}_{m + 1}) \to H^p(\mathcal{C}, I^n\mathcal{F}_{n + 1})$
を経由するので、望むとおり $s$ はその像に属する。
\end{proof}





\section{導来極限とホモトピー極限}
\label{sites-cohomology-section-derived-limits}

\noindent
$\mathcal{C}$ をサイトとする。圏 $\mathcal{C} \times \mathbf{N}$ を、
$n > m$ ならば
$\Mor((U, n), (V, m)) = \emptyset$、
そうでなければ
$\Mor((U, n), (V, m)) = \Mor(U, V)$
となるように考える。$\{U_i \to U\}$ が $\mathcal{C}$ の被覆であるような
族 $\{(U_i, n) \to (U, n)\}$ を被覆とすることで、これをサイトにする。
このとき $\mathcal{C} \times \mathbf{N}$ 上の層は
$\mathcal{C}$ 上の層の逆系と同じものであることを、読者は直ちに確認できる。
とくに $\textit{Ab}(\mathcal{C} \times \mathbf{N})$ は
$\mathcal{C}$ 上のアーベル層の逆系の圏である。ここで、逆系をその極限に送る関手
$$
\lim : \textit{Ab}(\mathcal{C} \times \mathbf{N}) \to \textit{Ab}(\mathcal{C})
$$
を考える。これは、連続かつ余連続な関手
$\mathcal{C} \times \mathbf{N} \to \mathcal{C}$ に付随するトポスの射
$g : \Sh(\mathcal{C} \times \mathbf{N}) \to \Sh(\mathcal{C})$
に対する $g_*$ にほかならない（$g^{-1}$ は $\mathcal{C}$ 上の層に
対応する定値逆系を割り当てることに注意せよ）。

\medskip\noindent
上で説明した一般論により、導来関手
$$
R\lim = Rg_* : D(\mathcal{C} \times \mathbf{N}) \to D(\mathcal{C})
$$
を得る。表示したように、この関手はしばしば $R\lim$ と書かれる。

\medskip\noindent
一方、連続かつ余連続な関手
$\mathcal{C} \to \mathcal{C} \times \mathbf{N}$,
$U \mapsto (U, n)$ はトポスの射
$i_n : \Sh(\mathcal{C}) \to \Sh(\mathcal{C} \times \mathbf{N})$
を定める。もちろん $i_n^{-1}$ は逆系の第 $n$ 項を取り出す関手である。
したがって関手の変換 $i_{n + 1}^{-1} \to i_n^{-1}$ がある。
ゆえに $K \in D(\mathcal{C} \times \mathbf{N})$ が与えられると、
$K_n = i_n^{-1}K \in D(\mathcal{C})$ と写像 $K_{n + 1} \to K_n$ を得る。
Derived Categories, Definition \ref{derived-definition-derived-limit} では
ホモトピー極限
$$
R\lim K_n \in D(\mathcal{C})
$$
の概念を定義した。この二つの概念は、意味をなす範囲で一致すると主張する。

\begin{lemma}
\label{sites-cohomology-lemma-derived-limit-is-ok}
$\mathcal{C}$ をサイトとし、$K$ を
$D(\mathcal{C} \times \mathbf{N})$ の対象とする。上のように
$K_n = i_n^{-1}K$ と置く。このとき $D(\mathcal{C})$ において
$$
R\lim K \cong R\lim K_n
$$
である。
\end{lemma}

\begin{proof}
$D(\mathcal{C} \times \mathbf{N})$ の対象 $K$ に対する
$R\lim$ を計算するため、各項が
$\textit{Ab}(\mathcal{C} \times \mathbf{N})$ の入射対象であるような
K-入射代表 $\mathcal{I}^\bullet$ を選ぶ。Injectives, Theorem
\ref{injectives-theorem-K-injective-embedding-grothendieck} を参照せよ。
$\mathcal{I}^\bullet$ を複体の逆系
$(\mathcal{I}_n^\bullet)$ とみなすと、
$$
R\lim K = \lim \mathcal{I}_n^\bullet
$$
であることが分かる。ここで右辺は各項ごとの逆極限である。

\medskip\noindent
$\mathcal{J} = (\mathcal{J}_n)$ を
$\textit{Ab}(\mathcal{C} \times \mathbf{N})$ の入射対象とする。
$(U, n) \to (U, n + 1)$ は
$\mathcal{C} \times \mathbf{N}$ の単射なので、Lemma
\ref{sites-cohomology-lemma-restriction-along-monomorphism-surjective} により
$\mathcal{J}(U, n + 1) \to \mathcal{J}(U, n)$ は全射である。
したがって $\mathcal{J}_{n + 1} \to \mathcal{J}_n$ は
前層の写像として全射である。

\medskip\noindent
関手 $i_n^{-1}$ は完全な左随伴 $i_{n, !}$ をもつことに注意する。
実際、$i_{n, !}\mathcal{F}$ は逆系
$\ldots 0 \to 0 \to \mathcal{F} \to \ldots \to \mathcal{F}$
である。したがって Derived Categories, Lemma
\ref{derived-lemma-adjoint-preserve-K-injectives} により、
複体 $i_n^{-1}\mathcal{I}^\bullet = \mathcal{I}_n^\bullet$ は
K-入射的である。

\medskip\noindent
K-入射複体を各項が入射的となるように選んだので、
$$
0 \to  \lim \mathcal{I}_n^\bullet \to \prod \mathcal{I}_n^\bullet
\to \prod \mathcal{I}_n^\bullet \to 0
$$
はアーベル層の複体の短完全列である。実際、これはアーベル群の前層の
複体の短完全列である。さらに、中央と右の積は
$D(\mathcal{C})$ における積を表す。Injectives, Lemma
\ref{injectives-lemma-derived-products} とその証明を参照せよ
（ここで $\mathcal{I}_n^\bullet$ が K-入射的であることを用いる）。
したがって、三角圏におけるホモトピー極限の定義により、
$R\lim K$ は逆系 $(K_n)$ のホモトピー極限である。
\end{proof}


\begin{lemma}
\label{sites-cohomology-lemma-RGamma-commutes-with-Rlim}
$(\mathcal{C}, \mathcal{O})$ を環付きサイトとする。
関手 $R\Gamma(\mathcal{C}, -)$ および
$U \in \Ob(\mathcal{C})$ に対する $R\Gamma(U, -)$ は
$R\lim$ と可換する。さらに、$D(\mathcal{O})$ の任意の逆系
$(K_n)$ と任意の $m \in \mathbf{Z}$ に対して短完全列
$$
0 \to
R^1\lim H^{m - 1}(U, K_n) \to H^m(U, R\lim K_n) \to
\lim H^m(U, K_n) \to 0
$$
が存在する。$H^m(\mathcal{C}, R\lim K_n)$ についても同様である。
\end{lemma}

\begin{proof}
最初の主張は Injectives, Lemma
\ref{injectives-lemma-RF-commutes-with-Rlim} から従う。次に
$R\lim R\Gamma(U, K_n) = R\Gamma(U, R\lim K_n)$ に
More on Algebra, Remark \ref{more-algebra-remark-compare-derived-limit}
を適用すれば、短完全列を得る。
\end{proof}

\begin{lemma}
\label{sites-cohomology-lemma-Rf-commutes-with-Rlim}
$f : (\Sh(\mathcal{C}), \mathcal{O}) \to
(\Sh(\mathcal{C}'), \mathcal{O}')$ を環付きトポスの射とする。
このとき $Rf_*$ は $R\lim$ と可換する。すなわち、
$Rf_*$ は導来極限と可換する。
\end{lemma}

\begin{proof}
$(K_n)$ を $D(\mathcal{O})$ の対象の逆系とする。$n$ に関する帰納法により、
$\mathcal{O}$-加群の実際の複体 $\mathcal{K}_n^\bullet$ と、
$D(\mathcal{O})$ における写像 $K_{n + 1} \to K_n$ を表す複体の写像
$\mathcal{K}_{n + 1}^\bullet \to \mathcal{K}_n^\bullet$ を選べる。
言い換えれば、付随する逆系が与えられたものとなる対象
$K \in D(\mathcal{C} \times \mathbf{N})$ が存在する。次に、
トポスの射の可換図式
$$
\xymatrix{
\Sh(\mathcal{C} \times \mathbf{N}) \ar[r]_g \ar[d]_{f \times 1} &
\Sh(\mathcal{C}) \ar[d]_f \\
\Sh(\mathcal{C}' \times \mathbf{N}) \ar[r]^{g'} &
\Sh(\mathcal{C}')
}
$$
を考える。このとき
$R\lim R(f \times 1)_*K = Rf_* R\lim K$ である。定義をたどり、
Lemma \ref{sites-cohomology-lemma-derived-limit-is-ok} を用いると、
$R\lim (Rf_*K_n) = Rf_*(R\lim K_n)$ を得る。

\medskip\noindent
$\mathcal{C}$ が十分な点をもつ場合の別証明。
$D(\mathcal{O})$ における定義完全三角
$$
R\lim K_n \to \prod K_n \to \prod K_n
$$
を考える。完全関手 $Rf_*$ を適用すると、
$D(\mathcal{O}')$ における完全三角
$$
Rf_*(R\lim K_n) \to Rf_*\left(\prod K_n\right)
\to Rf_*\left(\prod K_n\right)
$$
を得る。したがって $Rf_*$ が導来圏の積と可換することを示せば十分である
（この積は単に複体の積で与えられるわけではない。Injectives, Lemma
\ref{injectives-lemma-derived-products} を参照せよ）。
一方、Lemma \ref{sites-cohomology-lemma-adjoint} により $Rf_*$ は右随伴なので、
これは形式的に従う（Categories, Lemma
\ref{categories-lemma-adjoint-exact} 参照）。
注意：$R\lim K_n$ は $D(\mathcal{O})$ における極限ではないので、
Categories, Lemma \ref{categories-lemma-adjoint-exact}
を直接適用することはできない。
\end{proof}

\begin{remark}
\label{sites-cohomology-remark-discuss-derived-limit}
$(\mathcal{C}, \mathcal{O})$ を環付きサイトとし、
$(K_n)$ を $D(\mathcal{O})$ の逆系とする。$K = R\lim K_n$ と置く。
各 $n,m$ に対し、$\mathcal{H}^m_n = H^m(K_n)$ を
$K_n$ の第 $m$ コホモロジー層とし、同様に
$\mathcal{H}^m = H^m(K)$ と置く。前層
$$
U \longmapsto \underline{\mathcal{H}}^m_n(U) = H^m(U, K_n)
$$
を $\underline{\mathcal{H}}^m_n$ と書く。同様に
$\underline{\mathcal{H}}^m(U) = H^m(U, K)$ と置く。
Lemma \ref{sites-cohomology-lemma-sheafification-cohomology} により、
$\mathcal{H}^m_n$ は $\underline{\mathcal{H}}^m_n$ の層化であり、
$\mathcal{H}^m$ は $\underline{\mathcal{H}}^m$ の層化である。
次の図式がある。
$$
\xymatrix{
K \ar@{=}[d] &
\underline{\mathcal{H}}^m \ar[d] \ar[r] & 
\mathcal{H}^m \ar[d] \\
R\lim K_n &
\lim \underline{\mathcal{H}}^m_n \ar[r] & 
\lim \mathcal{H}^m_n
}
$$
一般には、$\lim \mathcal{H}^m_n$ が
$\lim \underline{\mathcal{H}}^m_n$ の層化であるとは限らない。
$U \in \mathcal{C}$ ならば、Lemma
\ref{sites-cohomology-lemma-RGamma-commutes-with-Rlim} により短完全列
\begin{equation}
\label{sites-cohomology-equation-ses-Rlim-over-U}
0 \to
R^1\lim \underline{\mathcal{H}}^{m - 1}_n(U) \to
\underline{\mathcal{H}}^m(U) \to
\lim \underline{\mathcal{H}}^m_n(U) \to 0
\end{equation}
を得る。
\end{remark}

\noindent
次の補題は、代数空間または代数スタック上の準連接加群で、
遷移写像が全射であるような逆系に適用できる。


\begin{lemma}
\label{sites-cohomology-lemma-inverse-limit-is-derived-limit}
$(\mathcal{C}, \mathcal{O})$ を環付きサイトとし、
$(\mathcal{F}_n)$ を $\mathcal{O}$-加群の逆系、
$\mathcal{B} \subset \Ob(\mathcal{C})$ を部分集合とする。次を仮定する。
\begin{enumerate}
\item $\mathcal{C}$ の各対象は、その要素が $\mathcal{B}$ に属する被覆をもつ。
\item $p > 0$ かつ $U \in \mathcal{B}$ ならば
$H^p(U, \mathcal{F}_n) = 0$ である。
\item $U \in \mathcal{B}$ に対し、逆系 $\mathcal{F}_n(U)$ の
$R^1\lim$ は零である。
\end{enumerate}
このとき $R\lim \mathcal{F}_n = \lim \mathcal{F}_n$ であり、
$p > 0$ かつ $U \in \mathcal{B}$ に対して
$H^p(U, \lim \mathcal{F}_n) = 0$ である。
\end{lemma}

\begin{proof}
$K_n = \mathcal{F}_n$、$K = R\lim \mathcal{F}_n$ と置く。
Remark \ref{sites-cohomology-remark-discuss-derived-limit} の記法と仮定 (2) により、
$U \in \mathcal{B}$ に対し、$m \not = 0$ ならば
$\underline{\mathcal{H}}_n^m(U) = 0$ であり、
$\underline{\mathcal{H}}_n^0(U) = \mathcal{F}_n(U)$ である。
Equation (\ref{sites-cohomology-equation-ses-Rlim-over-U}) と仮定 (3) から、
$m \not = 0$ ならば $\underline{\mathcal{H}}^m(U) = 0$ であり、
$m = 0$ ならばこれは $\lim \mathcal{F}_n(U)$ に等しい。
(1) を用いて層化すると、$m \not = 0$ ならば
$\mathcal{H}^m = 0$ であり、$m = 0$ ならばこれは
$\lim \mathcal{F}_n$ に等しい。ゆえに $K = \lim \mathcal{F}_n$ である。
$m > 0$ に対し
$H^m(U, K) = \underline{\mathcal{H}}^m(U) = 0$ であるから
（上を参照）、第二の主張も従う。
\end{proof}


\begin{lemma}
\label{sites-cohomology-lemma-cohomology-derived-limit-injective}
$(\mathcal{C}, \mathcal{O})$ を環付きサイトとし、
$(K_n)$ を $D(\mathcal{O})$ の逆系とする。
$V \in \Ob(\mathcal{C})$、$m \in \mathbf{Z}$ とする。
整数 $n(V)$ と $V$ の被覆の共終系 $\text{Cov}_V$ が存在し、
$\{V_i \to V\} \in \text{Cov}_V$ に対して次が成り立つと仮定する。
\begin{enumerate}
\item $R^1\lim H^{m - 1}(V_i, K_n) = 0$ である。
\item $n \geq n(V)$ に対し
$H^m(V_i, K_n) \to H^m(V_i, K_{n(V)})$ は単射である。
\end{enumerate}
このとき切断上の写像
$H^m(R\lim K_n)(V) \to H^m(K_{n(V)})(V)$ は単射である。
\end{lemma}

\begin{proof}
$\gamma \in H^m(R\lim K_n)(V)$ が
$H^m(K_{n(V)})(V)$ で零に写るとする。Lemma
\ref{sites-cohomology-lemma-sheafification-cohomology} により
$H^m(R\lim K_n)$ は $U \mapsto H^m(U, R\lim K_n)$ の層化なので、
被覆 $\{V_i \to V\} \in \text{Cov}_V$ と、
$\gamma|_{V_i}$ に写る元
$\tilde\gamma_i \in H^m(V_i, R\lim K_n)$ を選べる。
すると $\tilde\gamma_i$ は
$\tilde\gamma_{i, n(V)} \in H^m(V_i, K_{n(V)})$ に写る。
再び Lemma \ref{sites-cohomology-lemma-sheafification-cohomology} により
$H^m(K_{n(V)})$ は $U \mapsto H^m(U, K_{n(V)})$ の層化なので、
$\{V_i \to V\}$ を細分で置き換えることにより、すべての $i$ に対して
$\tilde\gamma_{i, n(V)} = 0$ と仮定できる。この被覆に対し、Lemma
\ref{sites-cohomology-lemma-RGamma-commutes-with-Rlim} の短完全列
$$
0 \to
R^1\lim H^{m - 1}(V_i, K_n) \to H^m(V_i, R\lim K_n) \to
\lim H^m(V_i, K_n) \to 0
$$
を考える。仮定 (1) により左の群は零であり、仮定 (2) により
右の群は $H^m(V_i, K_{n(V)})$ に単射で写る。
よって $\tilde\gamma_i = 0$、したがって望むとおり $\gamma = 0$ となる。
\end{proof}

\begin{lemma}
\label{sites-cohomology-lemma-is-limit-per-object}
$(\mathcal{C}, \mathcal{O})$ を環付きサイトとし、
$E \in D(\mathcal{O})$ とする。
$\mathcal{B} \subset \Ob(\mathcal{C})$ を部分集合とし、次を仮定する。
\begin{enumerate}
\item $\mathcal{C}$ の各対象は、その要素が $\mathcal{B}$ に属する被覆をもつ。
\item 各 $V \in \mathcal{B}$ に対し、関数
$p(V, -) : \mathbf{Z} \to \mathbf{Z}$ と $V$ の被覆の共終系
$\text{Cov}_V$ が存在し、
$$
H^p(V_i, H^{m - p}(E)) = 0
$$
が、$\{V_i \to V\} \in \text{Cov}_V$ および
$p > p(V, m)$ を満たすすべての整数 $p,m$ に対して成り立つ。
\end{enumerate}
このとき Derived Categories, Remark
\ref{derived-remark-map-into-derived-limit-truncations}
の写像 $E \to R\lim \tau_{\geq -n} E$ は
$D(\mathcal{O})$ における同型である。
\end{lemma}

\begin{proof}
$K_n = \tau_{\geq -n}E$、$K = R\lim K_n$ と置く。
標準写像 $E \to K$ は標準写像
$E \to K_n = \tau_{\geq -n}E$ から生じる。
$E \to K$ がコホモロジー層の同型
$H^m(E) \to H^m(K)$ を誘導することを示さなければならない。
以下の証明では $m$ を固定する。$n \geq -m$ ならば、写像
$E \to \tau_{\geq -n}E = K_n$ は同型
$H^m(E) \to H^m(K_n)$ を誘導する。
証明を終えるには、各 $V \in \mathcal{B}$ に対して、
$H^m(K)(V) \to H^m(K_{n(V)})(V)$ が単射となるような整数
$n(V) \geq -m$ が存在することを示せば十分である。実際、そのとき合成
$$
H^m(E)(V) \to H^m(K)(V) \to H^m(K_{n(V)})(V)
$$
は全単射で第二の矢印は単射だから、第一の矢印は全単射である。
性質 (1) により、これは $H^m(E) \to H^m(K)$ が同型であることを含意する。
ここで
$$
n(V) = 1 + \max\{-m, p(V, m - 1) - m, -1 + p(V, m) - m, -2 + p(V, m + 1) - m\}.
$$
と置けば、いずれの場合も $n(V) \geq -m$ である。主張：写像
$$
H^{m - 1}(V_i, K_{n + 1}) \to H^{m - 1}(V_i, K_n)
\quad\text{and}\quad
H^m(V_i, K_{n + 1}) \to H^m(V_i, K_n)
$$
は $n \geq n(V)$ かつ
$\{V_i \to V\} \in \text{Cov}_V$ に対して同型である。
この主張から Lemma
\ref{sites-cohomology-lemma-cohomology-derived-limit-injective} の条件 (1), (2) が満たされ、
したがって望む単射性が従う。完全三角
$$
H^{-n - 1}(E)[n + 1] \to
K_{n + 1} \to K_n \to H^{-n - 1}(E)[n + 2]
$$
が存在することを思い出そう（Derived Categories, Remark
\ref{derived-remark-truncation-distinguished-triangle}）。
対応するコホモロジー長完全列を見ると、$n \geq n(V)$ かつ
$\{V_i \to V\} \in \text{Cov}_V$ に対して
$$
H^{m + n}(V_i, H^{-n - 1}(E)),\quad
H^{m + n + 1}(V_i, H^{-n - 1}(E)),\quad
H^{m + n + 2}(V_i, H^{-n - 1}(E))
$$
が零ならば主張が従う。これは $n(V)$ の選び方と補題の仮定から従う。
\end{proof}


\begin{lemma}
\label{sites-cohomology-lemma-is-limit-spaltenstein}
$(\mathcal{C}, \mathcal{O})$ を環付きサイトとし、
$E \in D(\mathcal{O})$ とする。
$\mathcal{B} \subset \Ob(\mathcal{C})$ を部分集合とし、次を仮定する。
\begin{enumerate}
\item $\mathcal{C}$ の各対象は、その要素が $\mathcal{B}$ に属する被覆をもつ。
\item 各 $V \in \mathcal{B}$ に対し、整数 $d_V \geq 0$ と
$V$ の被覆の共終系 $\text{Cov}_V$ が存在し、
$$
H^p(V_i, H^q(E)) = 0 \text{ for }
\{V_i \to V\} \in \text{Cov}_V,\ p > d_V, \text{ and }q < 0
$$
が成り立つ。
\end{enumerate}
このとき Derived Categories, Remark
\ref{derived-remark-map-into-derived-limit-truncations}
の写像 $E \to R\lim \tau_{\geq -n} E$ は
$D(\mathcal{O})$ における同型である。
\end{lemma}

\begin{proof}
Lemma \ref{sites-cohomology-lemma-is-limit-per-object} で
$p(V, m) = d_V + \max(0, m)$ とすれば従う。
\end{proof}

\begin{lemma}
\label{sites-cohomology-lemma-is-limit}
$(\mathcal{C}, \mathcal{O})$ を環付きサイトとし、
$E \in D(\mathcal{O})$ とする。関数
$p(-) : \mathbf{Z} \to \mathbf{Z}$ と部分集合
$\mathcal{B} \subset \Ob(\mathcal{C})$ が存在し、次を満たすと仮定する。
\begin{enumerate}
\item $\mathcal{C}$ の各対象は、その要素が $\mathcal{B}$ に属する被覆をもつ。
\item $p > p(m)$ かつ $V \in \mathcal{B}$ ならば
$H^p(V, H^{m - p}(E)) = 0$ である。
\end{enumerate}
このとき Derived Categories, Remark
\ref{derived-remark-map-into-derived-limit-truncations}
の写像 $E \to R\lim \tau_{\geq -n} E$ は
$D(\mathcal{O})$ における同型である。
\end{lemma}

\begin{proof}
Lemma \ref{sites-cohomology-lemma-is-limit-per-object} を
$p(V, m) = p(m)$ とし、$\text{Cov}_V$ を
すべての $i$ に対して $V_i \in \mathcal{B}$ となる被覆
$\{V_i \to V\}$ の集合として適用する。
\end{proof}

\begin{lemma}
\label{sites-cohomology-lemma-is-limit-dimension}
$(\mathcal{C}, \mathcal{O})$ を環付きサイトとし、
$E \in D(\mathcal{O})$ とする。整数 $d \geq 0$ と部分集合
$\mathcal{B} \subset \Ob(\mathcal{C})$ が存在し、次を満たすと仮定する。
\begin{enumerate}
\item $\mathcal{C}$ の各対象は、その要素が $\mathcal{B}$ に属する被覆をもつ。
\item $p > d$、$q < 0$、$V \in \mathcal{B}$ ならば
$H^p(V, H^q(E)) = 0$ である。
\end{enumerate}
このとき Derived Categories, Remark
\ref{derived-remark-map-into-derived-limit-truncations}
の写像 $E \to R\lim \tau_{\geq -n} E$ は
$D(\mathcal{O})$ における同型である。
\end{lemma}

\begin{proof}
Lemma \ref{sites-cohomology-lemma-is-limit-spaltenstein} を $d_V = d$ とし、
$\text{Cov}_V$ を、すべての $i$ に対して $V_i \in \mathcal{B}$ となる被覆
$\{V_i \to V\}$ の集合として適用する。
\end{proof}

\noindent
上の補題を用いると、いくつかの場合にコホモロジーを計算できる。

\begin{lemma}
\label{sites-cohomology-lemma-cohomology-over-U-trivial}
$(\mathcal{C}, \mathcal{O})$ を環付きサイトとし、
$K$ を $D(\mathcal{O})$ の対象とする。
$\mathcal{B} \subset \Ob(\mathcal{C})$ を部分集合とし、次を仮定する。
\begin{enumerate}
\item $\mathcal{C}$ の各対象は、その要素が $\mathcal{B}$ に属する被覆をもつ。
\item すべての $p > 0$、$q \in \mathbf{Z}$、$U \in \mathcal{B}$ に対し
$H^p(U, H^q(K)) = 0$ である。
\end{enumerate}
このとき $q \in \mathbf{Z}$、$U \in \mathcal{B}$ に対し
$H^q(U, K) = H^0(U, H^q(K))$ である。
\end{lemma}

\begin{proof}
Lemma \ref{sites-cohomology-lemma-is-limit-dimension} で $d = 0$ とすれば、
$K = R\lim \tau_{\geq -n} K$ であることに注意する。
$U \in \mathcal{B}$ とする。Equation
(\ref{sites-cohomology-equation-ses-Rlim-over-U}) により短完全列
$$
0 \to R^1\lim H^{q - 1}(U, \tau_{\geq -n}K) \to
H^q(U, K) \to \lim H^q(U, \tau_{\geq -n}K) \to 0
$$
を得る。Derived Categories, Lemma
\ref{derived-lemma-two-ss-complex-functor} のスペクトル系列を用いると、
条件 (2) からすべての $q$ に対して
$H^q(U, \tau_{\geq -n} K) = H^0(U, H^q(\tau_{\geq -n} K))$
が従う。$\tau_{\geq -n}K$ は下に有界なので、このスペクトル系列は収束する。
$n > -q$ ならば $H^q(\tau_{\geq -n}K) = H^q(K)$ である。
したがって、表示した短完全列の左辺と右辺に現れる系は、それぞれ
$H^0(U, H^{q - 1}(K))$ と $H^0(U, H^q(K))$ を値として最終的に定常となる。
これで補題が従う。
\end{proof}

\noindent
導来極限を記述できる別の場合を挙げる。

\begin{lemma}
\label{sites-cohomology-lemma-derived-limit-suitable-system}
$(\mathcal{C}, \mathcal{O})$ を環付きサイトとし、
$(K_n)$ を $D(\mathcal{O})$ の対象の逆系とする。
$\mathcal{B} \subset \Ob(\mathcal{C})$ を部分集合とし、次を仮定する。
\begin{enumerate}
\item $\mathcal{C}$ の各対象は、その要素が $\mathcal{B}$ に属する被覆をもつ。
\item すべての $U \in \mathcal{B}$ とすべての
$q \in \mathbf{Z}$ に対し、次が成り立つ。
\begin{enumerate}
\item $p > 0$ に対し $H^p(U, H^q(K_n)) = 0$ である。
\item 逆系 $H^0(U, H^q(K_n))$ の $R^1\lim$ は零である。
\end{enumerate}
\end{enumerate}
このとき $q \in \mathbf{Z}$ に対し
$H^q(R\lim K_n) = \lim H^q(K_n)$ である。
\end{lemma}

\begin{proof}
$K = R\lim K_n$ と置く。Remark
\ref{sites-cohomology-remark-discuss-derived-limit} の記法を用いる。
$U \in \mathcal{B}$ とする。Lemma
\ref{sites-cohomology-lemma-cohomology-over-U-trivial} と (2)(a) により
$H^q(U, K_n) = H^0(U, H^q(K_n))$ である。
関手 $R\Gamma(U, -)$ が導来極限と可換することを用いると、
$$
H^q(U, K) = H^q(R\lim R\Gamma(U, K_n)) = \lim H^0(U, H^q(K_n))
$$
を得る。ここで最後の等式は More on Algebra, Remark
\ref{more-algebra-remark-compare-derived-limit} と仮定 (2)(b) による。
したがって $H^q(U, K)$ は、$U$ 上の層 $H^q(K_n)$ の切断の逆極限である。
$\lim H^q(K_n)$ は層なので、仮定 (1) を用いると、
前層 $U \mapsto H^q(U, K)$ の層化である $H^q(K)$ は
$\lim H^q(K_n)$ に等しい。これで補題が証明された。
\end{proof}








\section{K-入射分解の構成}
\label{sites-cohomology-section-K-injective}

\noindent
$(\mathcal{C}, \mathcal{O})$ を環付きサイトとし、
$\mathcal{F}^\bullet$ を $\mathcal{O}$-加群の複体とする。
圏 $\textit{Mod}(\mathcal{O})$ は十分な入射対象をもつので、
Derived Categories, Lemma \ref{derived-lemma-special-inverse-system}
を用いて、$\mathcal{O}$-加群の複体の圏における図式
$$
\xymatrix{
\ldots \ar[r] &
\tau_{\geq -2}\mathcal{F}^\bullet \ar[r] \ar[d] &
\tau_{\geq -1}\mathcal{F}^\bullet \ar[d] \\
\ldots \ar[r] & \mathcal{I}_2^\bullet \ar[r] & \mathcal{I}_1^\bullet
}
$$
で、次を満たすものを構成できる。
\begin{enumerate}
\item 垂直の矢印は擬同型である。
\item $\mathcal{I}_n^\bullet$ は入射対象からなる下に有界な複体である。
\item 矢印 $\mathcal{I}_{n + 1}^\bullet \to \mathcal{I}_n^\bullet$ は
各項で分裂する全射である。
\end{enumerate}
$\mathcal{O}$-加群の圏は極限をもち（これらは前層のレベルで計算される）、
したがって各項ごとの極限
$\mathcal{I}^\bullet = \lim_n \mathcal{I}_n^\bullet$ を作れる。
Derived Categories, Lemmas
\ref{derived-lemma-bounded-below-injectives-K-injective} および
\ref{derived-lemma-limit-K-injectives}
により、これは K-入射的複体である。一般には標準写像
\begin{equation}
\label{sites-cohomology-equation-into-candidate-K-injective}
\mathcal{F}^\bullet \to \mathcal{I}^\bullet
\end{equation}
は擬同型とは限らない。次の補題で、これが擬同型となるための
いくつかの条件を述べる。

\begin{lemma}
\label{sites-cohomology-lemma-K-injective}
上で述べた状況を考える。
$\mathcal{H}^m = H^m(\mathcal{F}^\bullet)$ を第 $m$ コホモロジー層とする。
$\mathcal{B} \subset \Ob(\mathcal{C})$ を部分集合とし、
$d \in \mathbf{N}$ とする。次を仮定する。
\begin{enumerate}
\item $\mathcal{C}$ の各対象は、その要素が $\mathcal{B}$ に属する被覆をもつ。
\item すべての $U \in \mathcal{B}$ に対し、$p > d$ かつ $q < 0$
ならば $H^p(U, \mathcal{H}^q) = 0$ である\footnote{条件
$\forall m$, $\exists p(m)$, $H^p(U. \mathcal{H}^{m - p}) = 0$ が
$p > p(m)$ のとき成り立てば十分である。Lemma
\ref{sites-cohomology-lemma-is-limit} を参照せよ。}。
\end{enumerate}
このとき (\ref{sites-cohomology-equation-into-candidate-K-injective}) は擬同型である。
\end{lemma}

\begin{proof}
Derived Categories, Lemma \ref{derived-lemma-difficulty-K-injectives}
により、写像
$\mathcal{F}^\bullet \to R\lim \tau_{\geq -n} \mathcal{F}^\bullet$
が同型であることを示せば十分である。これは Lemma
\ref{sites-cohomology-lemma-is-limit-dimension} から従う。
\end{proof}

\noindent
加群の複体の逆系を各項ごとに極限したもののコホモロジー層に関する
技術的な補題を挙げる。この補題の使用はできる限り避け、代わりに
導来逆極限を用いる議論を使うべきである。

\begin{lemma}
\label{sites-cohomology-lemma-inverse-limit-complexes}
$(\mathcal{C}, \mathcal{O})$ を環付きサイトとし、
$(\mathcal{F}_n^\bullet)$ を $\mathcal{O}$-加群の複体の逆系とする。
$m \in \mathbf{Z}$ とする。部分集合
$\mathcal{B} \subset \Ob(\mathcal{C})$ と整数 $n_0$ が与えられ、
次を満たすと仮定する。
\begin{enumerate}
\item $\mathcal{C}$ の各対象は、その要素が $\mathcal{B}$ に属する被覆をもつ。
\item すべての $U \in \mathcal{B}$ に対し、次が成り立つ。
\begin{enumerate}
\item アーベル群の系
$\mathcal{F}_n^{m - 2}(U)$ および $\mathcal{F}_n^{m - 1}(U)$
の $R^1\lim$ は零である（例えば、これらが Mittag--Leffler
条件を満たせばよい）。
\item アーベル群の系 $H^{m - 1}(\mathcal{F}_n^\bullet(U))$
の $R^1\lim$ は零である（例えば、この系が Mittag--Leffler
条件を満たせばよい）。
\item すべての $n \geq n_0$ に対し
$H^m(\mathcal{F}_n^\bullet(U)) = H^m(\mathcal{F}_{n_0}^\bullet(U))$
である。
\end{enumerate}
\end{enumerate}
このとき写像 $H^m(\mathcal{F}^\bullet) \to \lim H^m(\mathcal{F}_n^\bullet)
\to H^m(\mathcal{F}_{n_0}^\bullet)$ は層の同型である。ここで
$\mathcal{F}^\bullet = \lim \mathcal{F}_n^\bullet$ は各項ごとの
逆極限である。
\end{lemma}

\begin{proof}
$U \in \mathcal{B}$ とする。$H^m(\mathcal{F}^\bullet(U))$ は
$$
\lim_n \mathcal{F}_n^{m - 2}(U) \to
\lim_n \mathcal{F}_n^{m - 1}(U) \to
\lim_n \mathcal{F}_n^m(U) \to
\lim_n \mathcal{F}_n^{m + 1}(U)
$$
の左から三番目におけるコホモロジーであることに注意する。
仮定 (2)(a), (2)(b) により
More on Algebra, Lemma \ref{more-algebra-lemma-apply-Mittag-Leffler-again}
を適用でき、
$$
H^m(\mathcal{F}^\bullet(U)) = \lim H^m(\mathcal{F}_n^\bullet(U))
$$
と結論する。仮定 (2)(c) により、すべての $n \geq n_0$ に対して
$$
H^m(\mathcal{F}^\bullet(U)) = H^m(\mathcal{F}_n^\bullet(U))
$$
となる。仮定 (1) により、すべての $n \geq n_0$ に対して、
$U \mapsto H^m(\mathcal{F}^\bullet(U))$ の層化は
$U \mapsto H^m(\mathcal{F}_n^\bullet(U))$ の層化に等しい。
したがって層の逆系 $H^m(\mathcal{F}_n^\bullet)$ は
$n \geq n_0$ において値 $H^m(\mathcal{F}^\bullet)$ をもつ定値系である。
これで補題が証明された。
\end{proof}







\section{有界なコホモロジー次元}
\label{sites-cohomology-section-bounded}

\noindent
この節では、関手 $Rf_*$ がいつ有界なコホモロジー次元をもつかを問う。
非有界複体を考える場合、これはかなり微妙な問題である。

\begin{situation}
\label{sites-cohomology-situation-olsson-laszlo}
$\mathcal{C}$ をサイト、$\mathcal{O}$ を $\mathcal{C}$ 上の環の層とする。
$\mathcal{A} \subset \textit{Mod}(\mathcal{O})$ を弱 Serre 部分圏とする。
次を満たす部分集合 $\mathcal{B} \subset \Ob(\mathcal{C})$ が存在すると仮定する。
\begin{enumerate}
\item $\mathcal{C}$ の各対象は、その要素が $\mathcal{B}$ に属する被覆をもつ。
\item 各 $V \in \mathcal{B}$ に対し、整数 $d_V$ と $V$ の被覆の共終系
$\text{Cov}_V$ が存在し、
$$
H^p(V_i, \mathcal{F}) = 0 \text{ for }
\{V_i \to V\} \in \text{Cov}_V,\ p > d_V, \text{ and }
\mathcal{F} \in \Ob(\mathcal{A})
$$
が成り立つ。
\end{enumerate}
\end{situation}

\begin{lemma}
\label{sites-cohomology-lemma-olsson-laszlo}
\begin{reference}
これは \cite[Proposition 2.1.4]{six-I} の仮定をわずかに変更したものであり、
サイトに対する \cite[Proposition 3.13]{Spaltenstein} の類似である。
\end{reference}
Situation \ref{sites-cohomology-situation-olsson-laszlo} において、任意の
$E \in D_\mathcal{A}(\mathcal{O})$ に対し、Derived Categories, Remark
\ref{derived-remark-map-into-derived-limit-truncations} の写像
$E \to R\lim \tau_{\geq -n} E$
は $D(\mathcal{O})$ における同型である。
\end{lemma}

\begin{proof}
Lemma \ref{sites-cohomology-lemma-is-limit-spaltenstein} から直ちに従う。
\end{proof}

\begin{lemma}
\label{sites-cohomology-lemma-olsson-laszlo-modified}
Situation \ref{sites-cohomology-situation-olsson-laszlo} において、
$(K_n)$ を $D_\mathcal{A}^+(\mathcal{O})$ の逆系とする。
各 $j$ に対し、$\mathcal{A}$ における逆系 $(H^j(K_n))$ が
最終的に値 $\mathcal{H}^j$ の定値系になると仮定する。このとき、
すべての $j$ に対して $H^j(R\lim K_n) = \mathcal{H}^j$ である。
\end{lemma}

\begin{proof}
$V \in \mathcal{B}$ とし、$\{V_i \to V\}$ を Situation
\ref{sites-cohomology-situation-olsson-laszlo} の集合 $\text{Cov}_V$ に属する被覆とする。
$K_n$ は下に有界なので、スペクトル系列
$$
E_2^{p, q} = H^p(V_i, H^q(K_n))
$$
が存在し、$H^{p + q}(V_i, K_n)$ に収束する。Derived Categories, Lemma
\ref{derived-lemma-two-ss-complex-functor} を参照せよ。
仮定により $p > d_V$ に対して $E_2^{p, q} = 0$ であることに注意する。
すべての $n \geq n_0$ に対して
$$
\begin{matrix}
\mathcal{H}^{j + 1} & = & H^{j + 1}(K_n), \\
\mathcal{H}^j & = & H^j(K_n), \\
\ldots, \\
\mathcal{H}^{j - d_V - 2} & = & H^{j - d_V - 2}(K_n)
\end{matrix}
$$
となる $n_0$ を選ぶ。$K_n$ と $K_{n_0}$ に対する上のスペクトル系列を
比較すると、$n \geq n_0$ に対しコホモロジー群
$H^{j - 1}(V_i, K_n)$ および $H^j(V_i, K_n)$ は $n$ に依存しない。
したがって、十分大きな $n$（$V$ に依存する）に対し、切断上の写像
$H^j(R\lim K_n)(V) \to H^j(K_n)(V)$ は単射である。Lemma
\ref{sites-cohomology-lemma-cohomology-derived-limit-injective} を参照せよ。
$\mathcal{C}$ の各対象は $\mathcal{B}$ の要素で被覆できるので、写像
$H^j(R\lim K_n) \to \mathcal{H}^j$ は単射である。

\medskip\noindent
全射性も同様に示される。実際、$U \in \Ob(\mathcal{C})$ と
$\gamma \in \mathcal{H}^j(U)$ を取る。$U$ を被覆の要素で置き換えた後に、
$\gamma$ を $H^j(R\lim K_n)$ の切断へ持ち上げたい。したがって Situation
\ref{sites-cohomology-situation-olsson-laszlo} の性質 (1) により、
$U = V \in \mathcal{B}$ と仮定してよい。すべての $n \geq n_0$ に対して
$$
\begin{matrix}
\mathcal{H}^{j + 1} & = & H^{j + 1}(K_n), \\
\mathcal{H}^j & = & H^j(K_n), \\
\ldots, \\
\mathcal{H}^{j - d_V - 2} & = & H^{j - d_V - 2}(K_n)
\end{matrix}
$$
となる $n_0$ を選ぶ。$\gamma|_{V_i} \in
\mathcal{H}^j(V_i) = H^j(K_{n_0})(V_i)$ が元
$\gamma_{n_0, i} \in H^j(V_i, K_{n_0})$ に持ち上がるような
$\text{Cov}_V$ の要素 $\{V_i \to V\}$ を選ぶ。Lemma
\ref{sites-cohomology-lemma-sheafification-cohomology} により $H^j(K_{n_0})$ は
$U \mapsto H^j(U, K_{n_0})$ の層化なので、これは可能である。
証明の第一段落の議論により、$H^{j - 1}(V_i, K_n)$ と
$H^j(V_i, K_n)$ は $n \geq n_0$ に依存しない。したがって Lemma
\ref{sites-cohomology-lemma-RGamma-commutes-with-Rlim} により、$\gamma_{n_0, i}$ は元
$\gamma_i \in H^j(V_i, R\lim K_n)$ に持ち上がる。これで証明が終わる。
\end{proof}

\begin{lemma}
\label{sites-cohomology-lemma-olsson-laszlo-map-version-one}
\begin{reference}
これは \cite[Lemma 2.1.10]{six-I} の仮定をわずかに変更した版である。
\end{reference}
$f : (\Sh(\mathcal{C}), \mathcal{O}) \to
(\Sh(\mathcal{C}'), \mathcal{O}')$ を環付きトポスの射とする。
$\mathcal{A} \subset \textit{Mod}(\mathcal{O})$ および
$\mathcal{A}' \subset \textit{Mod}(\mathcal{O}')$ を弱 Serre 部分圏とする。
次を満たす整数 $N$ が存在すると仮定する。
\begin{enumerate}
\item $\mathcal{C}, \mathcal{O}, \mathcal{A}$ は Situation
\ref{sites-cohomology-situation-olsson-laszlo} の仮定を満たす。
\item $\mathcal{C}', \mathcal{O}', \mathcal{A}'$ は Situation
\ref{sites-cohomology-situation-olsson-laszlo} の仮定を満たす。
\item $p \geq 0$ かつ $\mathcal{F} \in \Ob(\mathcal{A})$ ならば
$R^pf_*\mathcal{F} \in \Ob(\mathcal{A}')$ である。
\item $p > N$ かつ $\mathcal{F} \in \Ob(\mathcal{A})$ ならば
$R^pf_*\mathcal{F} = 0$ である。
\end{enumerate}
$K \in D_\mathcal{A}(\mathcal{O})$ に対し、次が成り立つ。
\begin{enumerate}
\item[(a)] $Rf_*K$ は $D_{\mathcal{A}'}(\mathcal{O}')$ に属する。
\item[(b)] $j \geq N - n$ に対し、写像
$H^j(Rf_*K) \to H^j(Rf_*(\tau_{\geq -n}K))$ は同型である。
\end{enumerate}
\end{lemma}

\begin{proof}
Lemma \ref{sites-cohomology-lemma-olsson-laszlo} により
$K = R\lim \tau_{\geq -n}K$ である。Lemma
\ref{sites-cohomology-lemma-Rf-commutes-with-Rlim} により
$Rf_*K = R\lim Rf_*\tau_{\geq -n}K$ である。
複体 $Rf_*\tau_{\geq -n}K$ は下に有界である。スペクトル系列
$$
E_2^{p, q} = R^pf_*H^q(\tau_{\geq -n}K)
$$
は $H^{p + q}(Rf_*\tau_{\geq -n}K)$ に収束し
（Derived Categories, Lemma \ref{derived-lemma-two-ss-complex-functor}）、
仮定 (3) と合わせると $Rf_*\tau_{\geq -n}K$ が
$D^+_{\mathcal{A}'}(\mathcal{O}')$ に属することが分かる。
Homology, Lemma \ref{homology-lemma-biregular-ss-converges} を参照せよ。
$m \geq n$ に対し、写像
$$
Rf_*(\tau_{\geq -m}K) \longrightarrow Rf_*(\tau_{\geq -n}K)
$$
は、上のスペクトル系列により、次数 $j \geq -n + N$ の
コホモロジー層上で同型を誘導することに注意する。したがって Lemma
\ref{sites-cohomology-lemma-olsson-laszlo-modified} を適用すれば結論を得る。
\end{proof}

\noindent
仮定が少し異なる上の補題の変種が必要になることがある。

\begin{situation}
\label{sites-cohomology-situation-olsson-laszlo-prime}
$f : (\mathcal{C}, \mathcal{O}) \to (\mathcal{C}', \mathcal{O}')$
を環付きサイトの射とし、$u : \mathcal{C}' \to \mathcal{C}$ を
対応するサイトの連続関手とする。
$\mathcal{A} \subset \textit{Mod}(\mathcal{O})$ を弱 Serre 部分圏とする。
次を満たす部分集合 $\mathcal{B}' \subset \Ob(\mathcal{C}')$ が存在すると仮定する。
\begin{enumerate}
\item $\mathcal{C}'$ の各対象は、その要素が $\mathcal{B}'$ に属する被覆をもつ。
\item 各 $V' \in \mathcal{B}'$ に対し、整数 $d_{V'}$ と
$V'$ の被覆の共終系 $\text{Cov}_{V'}$ が存在し、
$$
H^p(u(V'_i), \mathcal{F}) = 0 \text{ for }
\{V'_i \to V'\} \in \text{Cov}_{V'},\ p > d_{V'}, \text{ and }
\mathcal{F} \in \Ob(\mathcal{A})
$$
が成り立つ。
\end{enumerate}
\end{situation}

\begin{lemma}
\label{sites-cohomology-lemma-olsson-laszlo-map-version-two}
\begin{reference}
これは \cite[Lemma 2.1.10]{six-I} の仮定をわずかに変更した版である。
\end{reference}
$f : (\mathcal{C}, \mathcal{O}) \to (\mathcal{C}', \mathcal{O}')$
を環付きサイトの射とする。さらに、次を満たす整数 $N$ が存在すると仮定する。
\begin{enumerate}
\item $\mathcal{C}, \mathcal{O}, \mathcal{A}$ は Situation
\ref{sites-cohomology-situation-olsson-laszlo} の仮定を満たす。
\item $f : (\mathcal{C}, \mathcal{O}) \to (\mathcal{C}', \mathcal{O}')$
と $\mathcal{A}$ は Situation \ref{sites-cohomology-situation-olsson-laszlo-prime}
の仮定を満たす。
\item $p > N$ かつ $\mathcal{F} \in \Ob(\mathcal{A})$ ならば
$R^pf_*\mathcal{F} = 0$ である。
\end{enumerate}
$K \in D_\mathcal{A}(\mathcal{O})$ に対し、$j \geq N - n$ ならば写像
$H^j(Rf_*K) \to H^j(Rf_*(\tau_{\geq -n}K))$ は同型である。
\end{lemma}

\begin{proof}
$K \in D_\mathcal{A}(\mathcal{O})$ とする。Lemma
\ref{sites-cohomology-lemma-olsson-laszlo} により $K = R\lim \tau_{\geq -n}K$ である。
Lemma \ref{sites-cohomology-lemma-Rf-commutes-with-Rlim} により
$Rf_*K = R\lim Rf_*(\tau_{\geq -n}K)$ である。
$V' \in \mathcal{B}'$ とし、$\{V'_i \to V'\}$ を
$\text{Cov}_{V'}$ の要素とする。このとき
$$
H^j(V'_i, Rf_*K) = H^j(u(V'_i), K)
\quad\text{and}\quad
H^j(V'_i, Rf_*(\tau_{\geq -n}K)) = H^j(u(V'_i), \tau_{\geq -n}K)
$$
を考える。Situation \ref{sites-cohomology-situation-olsson-laszlo-prime} の仮定により、
最後の群は $j$ と $d_{V'}$ に依存する十分大きな $n$ に対して
$n$ に依存しない。細部は省略する。これを $j$ と $j - 1$ に適用し、
Lemma \ref{sites-cohomology-lemma-RGamma-commutes-with-Rlim} を用いると、
$$
H^j(V'_i, Rf_*K) = \lim H^j(V'_i, Rf_*(\tau_{\geq -n} K))
$$
を得る。また右辺の系は、$d_{V'}$ と $j$ のみに依存する定数より
大きな $n$ に対して定値である。したがって Lemma
\ref{sites-cohomology-lemma-cohomology-derived-limit-injective} により、
$$
H^j(Rf_*K)(V') \longrightarrow
\left(\lim H^j(Rf_*(\tau_{\geq -n}K))\right)(V')
$$
は単射である。$\mathcal{B}'$ の要素 $V'$ が
$\mathcal{C}'$ の各対象を被覆するので、写像
$H^j(Rf_*K) \to \lim H^j(Rf_*(\tau_{\geq -n}K))$ は単射である。
スペクトル系列
$$
E_2^{p, q} = R^pf_*H^q(\tau_{\geq -n}K)
$$
は $H^{p + q}(Rf_*(\tau_{\geq -n}K))$ に収束し
（Derived Categories, Lemma \ref{derived-lemma-two-ss-complex-functor}）、
仮定 (3) と合わせると、$H^j(Rf_*(\tau_{\geq -n}K))$ は
$n \geq N - j$ に対して定値であることが分かる。
したがって $j \geq N - n$ に対し、
$H^j(Rf_*K) \to H^j(Rf_*(\tau_{\geq -n}K))$ は単射である。

\medskip\noindent
以上で、補題の最終行の「同型」を「単射」に置き換えた主張を証明した。
ここで $j \geq N - n$ を満たす $j,n$ を選び、完全三角
$$
\tau_{\leq -n - 1}K \to K \to \tau_{\geq -n}K \to (\tau_{\leq -n - 1}K)[1]
$$
を考える。Derived Categories, Remark
\ref{derived-remark-truncation-distinguished-triangle} を参照せよ。
$\tau_{\geq -n}\tau_{\leq -n -1}K = 0$ なので、
$\tau_{-n - 1}K$ に対してすでに証明した単射性から
$$
0 = H^j(Rf_*(\tau_{\leq -n - 1}K)) =
H^{j + 1}(Rf_*(\tau_{\leq -n - 1}K)) =
H^{j + 2}(Rf_*(\tau_{\leq -n - 1}K)) = \ldots
$$
が従う。完全三角
$$
Rf_*(\tau_{\leq -n - 1}K) \to Rf_*K \to Rf_*(\tau_{\geq -n}K) \to
Rf_*(\tau_{\leq -n - 1}K)[1]
$$
に付随するコホモロジー長完全列により、これは
$H^j(Rf_*K) \to H^j(Rf_*(\tau_{\geq -n}K))$ が同型であることを含意する。
\end{proof}








\section{Mayer-Vietoris}
\label{sites-cohomology-section-mayer-vietoris}

\noindent
通常の Mayer--Vietoris の主張と証明については、Cohomology, Section
\ref{cohomology-section-mayer-vietoris} を参照せよ。

\medskip\noindent
$(\mathcal{C}, \mathcal{O})$ を環付きサイトとし、圏 $\mathcal{C}$ における
可換図式
$$
\xymatrix{
E \ar[d] \ar[r] & Y \ar[d] \\
Z \ar[r] & X
}
$$
を考える。この状況で $D(\mathcal{O})$ の対象 $K$ が与えられると、
完全三角の初めの部分のように見えるもの
$$
R\Gamma(X, K) \to
R\Gamma(Z, K) \oplus
R\Gamma(Y, K) \to
R\Gamma(E, K)
$$
を得る。次の補題でこれをより正確に述べる。

\begin{lemma}
\label{sites-cohomology-lemma-c-square}
上の状況で、$K$ を表す $\mathcal{O}$-加群の K-入射複体
$\mathcal{I}^\bullet$ を選ぶ。四本の矢印の一つについて標準写像の
$-1$ 倍を用いると、複体の写像
$$
\mathcal{I}^\bullet(X) \xrightarrow{\alpha}
\mathcal{I}^\bullet(Z) \oplus
\mathcal{I}^\bullet(Y) \xrightarrow{\beta}
\mathcal{I}^\bullet(E)
$$
を得て、$\beta \circ \alpha = 0$ である。したがって標準写像
$$
c^K_{X, Z, Y, E} :
\mathcal{I}^\bullet(X)
\longrightarrow
C(\beta)^\bullet[-1]
$$
がある。この写像が標準的であるというのは、$K$ を表す別の
K-入射複体を選んでも、アーベル群の導来圏における同型な矢印が
定まるという意味である。$c^K_{X, Z, Y, E}$ が同型ならば、
その逆を用いて標準的な完全三角
$$
R\Gamma(X, K) \to
R\Gamma(Z, K) \oplus
R\Gamma(Y, K) \to
R\Gamma(E, K) \to
R\Gamma(X, K)[1]
$$
を得る。これらの構成はすべて $K$ に関して関手的である。
\end{lemma}

\begin{proof}
主張は構成から明らかである。例えば、$\mathcal{J}^\bullet$ を
$K$ を表すもう一つの K-入射複体とすると、擬同型
$\mathcal{I}^\bullet \to \mathcal{J}^\bullet$
を選べ、これは現れるすべての複体の間に擬同型を定める。細部は省略する。
錐の構成および完全三角との関係については Derived Categories, Sections
\ref{derived-section-cones} および
\ref{derived-section-homotopy-triangulated} を参照せよ。
\end{proof}

\begin{lemma}
\label{sites-cohomology-lemma-two-out-of-three-blow-up-square}
上の状況で、$K_1 \to K_2 \to K_3 \to K_1[1]$ を
$D(\mathcal{O})$ の完全三角とする。
$\{1, 2, 3\}$ のうち二つの $i$ に対して
$c^{K_i}_{X, Z, Y, E}$ が擬同型ならば、残る $i$ に対しても擬同型である。
\end{lemma}

\begin{proof}
完全三角を回転することにより、$c^{K_1}_{X, Z, Y, E}$ と
$c^{K_2}_{X, Z, Y, E}$ が擬同型であると仮定してよい。
$K_1 \to K_2$ を表す $\mathcal{O}$-加群の K-入射複体の写像
$f : \mathcal{I}^\bullet_1 \to \mathcal{I}^\bullet_2$ を選ぶ。
このとき $K_3$ は K-入射複体 $C(f)^\bullet$ で表される。
Derived Categories, Lemma \ref{derived-lemma-triangle-K-injective} を参照せよ。
このとき射 $c^{K_3}_{X, Z, Y, E}$ は同型である。実際、これは
$K(\textit{Ab})$ における完全三角の射で、ほかの二本の脚が擬同型であるものの
第三の脚である。細部は省略する。Derived Categories, Lemma
\ref{derived-lemma-third-isomorphism-triangle} を用いよ。
\end{proof}

\noindent
これが実際に完全三角を与えるための判定条件を述べる。

\begin{lemma}
\label{sites-cohomology-lemma-square-triangle}
上の状況で、次を仮定する。
\begin{enumerate}
\item $h_X^\# = h_Y^\# \amalg_{h_E^\#} h_Z^\#$ である。
\item $h_E^\# \to h_Y^\#$ は単射である。
\end{enumerate}
このとき Lemma \ref{sites-cohomology-lemma-c-square} の構成は完全三角
$$
R\Gamma(X, K) \to
R\Gamma(Z, K) \oplus
R\Gamma(Y, K) \to
R\Gamma(E, K) \to R\Gamma(X, K)[1]
$$
を与え、これは $D(\mathcal{C})$ の $K$ に関して関手的である。
\end{lemma}

\begin{proof}
$K$ は各項が入射アーベル層である K-入射複体で表せる。Section
\ref{sites-cohomology-section-unbounded} を参照せよ。したがって、$\mathcal{I}$ が
入射アーベル層ならば
$$
0 \to \mathcal{I}(X) \to
\mathcal{I}(Z) \oplus \mathcal{I}(Y) \to
\mathcal{I}(E) \to 0
$$
が短完全列であることを示せば十分である。第一の矢印は単射である。
実際、条件 (1) により写像 $h_Y \amalg h_Z \to h_X$ は層化すると全射になり、
これは $\{Y \to X, Z \to X\}$ が $X$ の被覆によって細分できることを意味する。
最後の矢印が全射なのは $\mathcal{I}(Y) \to \mathcal{I}(E)$ が全射だからである。
実際、
$\mathcal{I}(E) = \Hom(\mathbf{Z}_E^\#, \mathcal{I})$,
$\mathcal{I}(Y) = \Hom(\mathbf{Z}_Y^\#, \mathcal{I})$,
写像 $\mathbf{Z}_E^\# \to \mathbf{Z}_Y^\#$ は (2) により単射であり、
$\mathcal{I}$ は入射アーベル層である。Modules on Sites, Section
\ref{sites-modules-section-free-abelian-sheaf} と比較せよ。
最後に、$s \in \mathcal{I}(Y)$ と $t \in \mathcal{F}(Z)$ が
$\mathcal{I}(E)$ の同じ元に写ると仮定する。このとき $s,t$ は写像
$$
s \amalg t : h_Y^\# \amalg h_Z^\# \longrightarrow \mathcal{I}
$$
を定め、これは仮定により $h_Y^\# \amalg_{h_E^\#} h_Z^\#$ を経由する。
したがって仮定 (1) により一意な写像 $h_X^\# \to \mathcal{I}$ を得る。
これは $Y$ 上で $s$ に、$Z$ 上で $t$ に制限される
$\mathcal{I}(X)$ の元に対応する。
\end{proof}

\begin{lemma}
\label{sites-cohomology-lemma-square-triangle-general}
$\mathcal{C}$ をサイトとする。$\mathcal{C}$ 上の集合の前層の可換図式
$$
\xymatrix{
\mathcal{D} \ar[r] \ar[d] & \mathcal{F} \ar[d] \\
\mathcal{E} \ar[r] & \mathcal{G}
}
$$
を考え、次を仮定する。
\begin{enumerate}
\item $\mathcal{G}^\# =
\mathcal{E}^\# \amalg_{\mathcal{D}^\#} \mathcal{F}^\#$ である。
\item $\mathcal{D}^\# \to \mathcal{F}^\#$ は単射である。
\end{enumerate}
このとき標準的な完全三角
$$
R\Gamma(\mathcal{G}, K) \to
R\Gamma(\mathcal{E}, K) \oplus
R\Gamma(\mathcal{F}, K) \to
R\Gamma(\mathcal{D}, K) \to
R\Gamma(\mathcal{G}, K)[1]
$$
があり、$K \in D(\mathcal{C})$ に関して関手的である。ここで
$R\Gamma(\mathcal{G}, -)$ は Section \ref{sites-cohomology-section-limp} で論じた
コホモロジーである。
\end{lemma}

\begin{proof}
層化は完全であり、かつ
$R\Gamma(\mathcal{G}, -) = R\Gamma(\mathcal{G}^\#, -)$ なので、
$\mathcal{D}, \mathcal{E}, \mathcal{F}, \mathcal{G}$ は集合の層であると
仮定してよい。さらにコホモロジー $R\Gamma(\mathcal{G}, -)$ は
トポスのみに依存し、もとのサイトには依存しない。したがって Sites, Lemma
\ref{sites-lemma-topos-good-site} により、$\mathcal{C}$ を
準標準位相をもつ「より大きな」サイトで、
$\mathcal{G} = h_X$,
$\mathcal{F} = h_Y$,
$\mathcal{E} = h_Z$ および
$\mathcal{D} = h_E$
が $\mathcal{C}$ のある対象 $X,Y,Z,E$ に対して成り立つものに
置き換えてよい。この場合、結果は Lemma
\ref{sites-cohomology-lemma-square-triangle} から従う。
\end{proof}







\section{二つの位相の比較}
\label{sites-cohomology-section-compare}

\noindent
$\mathcal{C}$ を圏とする。
$\text{Cov}(\mathcal{C}) \supset \text{Cov}'(\mathcal{C})$
を $\mathcal{C}$ にサイト構造を与える二つの方法とする。
$\text{Cov}(\mathcal{C})$ に対応する位相を $\tau$、
$\text{Cov}'(\mathcal{C})$ に対応する位相を $\tau'$ と書く。
このとき $\mathcal{C}$ 上の恒等関手はサイトの射
$$
\epsilon : \mathcal{C}_\tau \longrightarrow \mathcal{C}_{\tau'}
$$
を定める。ここで $\epsilon_*$ は台となる前層上の恒等関手であり、
$\epsilon^{-1}$ は $\tau'$-層の $\tau$-層化である。Sites, Examples
\ref{sites-example-finer-topology} および
\ref{sites-example-finer-topology-bis} を参照せよ。
上の状況で次が成り立つ。
\begin{enumerate}
\item $\epsilon_* : \Sh(\mathcal{C}_\tau) \to \Sh(\mathcal{C}_{\tau'})$
は充満忠実であり、$\epsilon^{-1} \circ \epsilon_* = \text{id}$ である。
\item $\epsilon_* : \textit{Ab}(\mathcal{C}_\tau) \to
\textit{Ab}(\mathcal{C}_{\tau'})$
は充満忠実であり、$\epsilon^{-1} \circ \epsilon_* = \text{id}$ である。
\item $R\epsilon_* : D(\mathcal{C}_\tau) \to D(\mathcal{C}_{\tau'})$
は充満忠実であり、$\epsilon^{-1} \circ R\epsilon_* = \text{id}$ である。
\item $\mathcal{O}$ が $\tau$-位相に関する環の層ならば、
$\mathcal{O}$ は $\tau'$-位相に関しても層であり、$\epsilon$ は環付きサイトの平坦射
$$
\epsilon :
(\mathcal{C}_\tau, \mathcal{O}_\tau)
\longrightarrow
(\mathcal{C}_{\tau'}, \mathcal{O}_{\tau'})
$$
となる。
\item $\epsilon_* : \textit{Mod}(\mathcal{O}_\tau) \to
\textit{Mod}(\mathcal{O}_{\tau'})$ は充満忠実であり、
$\epsilon^* \circ \epsilon_* = \text{id}$ である。
\item $R\epsilon_* : D(\mathcal{O}_\tau) \to D(\mathcal{O}_{\tau'})$
は充満忠実であり、$\epsilon^* \circ R\epsilon_* = \text{id}$ である。
\end{enumerate}
以下、説明を加える。

\medskip\noindent
(1) について。$\mathcal{F}$ を $\tau$-位相における集合の層とする。
このとき $\epsilon_*\mathcal{F}$ は、$\mathcal{F}$ を
$\tau'$-位相における層とみなしたものにすぎない。
$\epsilon^{-1}$ を適用することは $\mathcal{F}$ の $\tau$-層化を取ることを
意味するが、$\mathcal{F}$ はすでに $\tau$-層なので何も変わらない。したがって
$\epsilon^{-1}(\epsilon_*\mathcal{F})) = \mathcal{F}$.
充満忠実性は Categories, Lemma
\ref{categories-lemma-adjoint-fully-faithful} から従う。

\medskip\noindent
(2) について。アーベル層の引き戻しと押し出しは、台となる集合の層に
それらの操作を行うことと同じなので、これは (1) の帰結である。

\medskip\noindent
(3) について。$K$ を $D(\mathcal{C}_\tau)$ の対象とする。
$R\epsilon_*K$ を計算するため、$K$ を表す K-入射複体
$\mathcal{I}^\bullet$ を選び、
$R\epsilon_*K = \epsilon_*\mathcal{I}^\bullet$ と置く。
$\epsilon^{-1} : D(\mathcal{C}_{\tau'}) \to D(\mathcal{C}_\tau)$ は、
対象 $L$ を表す任意のアーベル層の複体に完全関手 $\epsilon^{-1}$ を
適用することで計算される。したがって $\epsilon^{-1}R\epsilon_*K$ は
$\epsilon^{-1}\epsilon_*\mathcal{I}^\bullet$ で表される。(1) により
$\mathcal{I}^\bullet = \epsilon^{-1}\epsilon_*\mathcal{I}^\bullet$.
言い換えれば、$\epsilon^{-1} \circ R\epsilon_* = \text{id}$ であり、
前と同様に結論を得る。

\medskip\noindent
(4) について。(1) の議論を参照すると、
$\epsilon^{-1}\mathcal{O}_{\tau'} = \mathcal{O}_\tau$ である。
したがって $\epsilon$ は環付きサイトの平坦射である。Modules on Sites,
Definition \ref{sites-modules-definition-flat-morphism} を参照せよ。
さらに、$\mathcal{O}_{\tau'}$-加群上で
$\epsilon^* = \epsilon^{-1}$ であることも明らかである
（加群としての引き戻しの台となるアーベル層は、台となるアーベル層の
引き戻しと同じである）。

\medskip\noindent
(5) について。これは (2) と (4) で述べたことから明らかである。

\medskip\noindent
(6) について。これは (3) と同様である。細部は省略する。












\section{コホモロジー降下の形式論}
\label{sites-cohomology-section-formal-cohomological-descent}

\noindent
本節では、環付きトポスの射が、下側の導来圏から上側の導来圏への
埋め込みをどの程度定めるかだけを論じる。典型的な結果は次である。

\begin{lemma}
\label{sites-cohomology-lemma-downstairs}
$f : (\Sh(\mathcal{C}), \mathcal{O}_\mathcal{C}) \to
(\Sh(\mathcal{D}), \mathcal{O}_\mathcal{D})$ を環付きトポスの射とする。
次の標準写像が同型となる対象 $K$ からなる充満部分圏
$D' \subset D(\mathcal{O}_\mathcal{D})$ を考える：
$$
K \longrightarrow Rf_*Lf^*K
$$
このとき $D'$ は $D(\mathcal{O}_\mathcal{D})$ の厳密充満飽和三角部分圏であり、
関手 $Lf^* : D' \to D(\mathcal{O}_\mathcal{C})$ は充満忠実である。
\end{lemma}

\begin{proof}
この状況における飽和の定義については Derived Categories, Definition
\ref{derived-definition-saturated} を参照せよ。三角部分圏の議論については
Derived Categories, Lemma \ref{derived-lemma-triangulated-subcategory} を参照せよ。
補題の標準写像は関手の随伴対 $(Lf^*, Rf_*)$ の単位である；補題
\ref{sites-cohomology-lemma-adjoint} を参照せよ。以上を踏まえ、$D'$ が飽和三角部分圏であることの
証明は省略する。それは $Lf^*$ と $Rf_*$ が三角圏の間の完全関手であるという事実から
形式的に従う。最後の主張も、$Lf^*$ と $Rf_*$ が随伴であるという事実から
形式的に従う；Categories, Lemma
\ref{categories-lemma-adjoint-fully-faithful} と比較せよ。
\end{proof}

\begin{lemma}
\label{sites-cohomology-lemma-upstairs}
$f : (\Sh(\mathcal{C}), \mathcal{O}_\mathcal{C}) \to
(\Sh(\mathcal{D}), \mathcal{O}_\mathcal{D})$ を環付きトポスの射とする。
次の標準写像が同型となる対象 $K$ からなる充満部分圏
$D' \subset D(\mathcal{O}_\mathcal{C})$ を考える：
$$
Lf^*Rf_*K \longrightarrow K
$$
このとき $D'$ は $D(\mathcal{O}_\mathcal{C})$ の厳密充満飽和三角部分圏であり、
関手 $Rf_* : D' \to D(\mathcal{O}_\mathcal{D})$ は充満忠実である。
\end{lemma}

\begin{proof}
この状況における飽和の定義については Derived Categories, Definition
\ref{derived-definition-saturated} を参照せよ。三角部分圏の議論については
Derived Categories, Lemma \ref{derived-lemma-triangulated-subcategory} を参照せよ。
補題の標準写像は関手の随伴対 $(Lf^*, Rf_*)$ の余単位である；補題
\ref{sites-cohomology-lemma-adjoint} を参照せよ。以上を踏まえ、$D'$ が飽和三角部分圏であることの
証明は省略する。それは $Lf^*$ と $Rf_*$ が三角圏の間の完全関手であるという事実から
形式的に従う。最後の主張も、$Lf^*$ と $Rf_*$ が随伴であるという事実から
形式的に従う；Categories, Lemma
\ref{categories-lemma-adjoint-fully-faithful} と比較せよ。
\end{proof}

\begin{lemma}
\label{sites-cohomology-lemma-bounded-in-image-upstairs}
$f : (\Sh(\mathcal{C}), \mathcal{O}_\mathcal{C}) \to
(\Sh(\mathcal{D}), \mathcal{O}_\mathcal{D})$ を環付きトポスの射とする。
$K$ を $D(\mathcal{O}_\mathcal{C})$ の対象とする。次を仮定する。
\begin{enumerate}
\item $f$ は平坦である。
\item $K$ は下に有界である。
\item $f^*Rf_*H^q(K) \to H^q(K)$ は同型である。
\end{enumerate}
このとき $f^*Rf_*K \to K$ は同型である。
\end{lemma}

\begin{proof}
写像 $f^*Rf_*K \to K$ が同型であることは、各コホモロジー層 $H^j$ 上で
同型であることと同値である。また、
$H^j(f^*Rf_*K) = f^*H^j(Rf_*K) = f^*H^j(Rf_*\tau_{\leq j}K) =
H^j(f^*Rf_*\tau_{\leq j}K)$.
したがって $K$ は有界であると仮定してよい。このとき条件 (3) より、
コホモロジー層は補題 \ref{sites-cohomology-lemma-upstairs} の三角部分圏
$D' \subset D(\mathcal{O}_\mathcal{C})$ に属する。ゆえに $K$ もそこに属する。
\end{proof}

\begin{lemma}
\label{sites-cohomology-lemma-bounded-in-image-downstairs}
$f : (\Sh(\mathcal{C}), \mathcal{O}_\mathcal{C}) \to
(\Sh(\mathcal{D}), \mathcal{O}_\mathcal{D})$ を環付きトポスの射とする。
$K$ を $D(\mathcal{O}_\mathcal{D})$ の対象とする。次を仮定する。
\begin{enumerate}
\item $f$ は平坦である。
\item $K$ は下に有界である。
\item $H^q(K) \to Rf_*f^*H^q(K)$ は同型である。
\end{enumerate}
このとき $K \to Rf_*f^*K$ は同型である。
\end{lemma}

\begin{proof}
写像 $K \to Rf_*f^*K$ が同型であることは、各コホモロジー層 $H^j$ 上で
同型であることと同値である。また、
$H^j(Rf_*f^*K) = H^j(Rf_*\tau_{\leq j}f^*K) = H^j(Rf_*f^*\tau_{\leq j}K)$.
したがって $K$ は有界であると仮定してよい。このとき条件 (3) より、
コホモロジー層は補題 \ref{sites-cohomology-lemma-downstairs} の三角部分圏
$D' \subset D(\mathcal{O}_\mathcal{D})$ に属する。ゆえに $K$ もそこに属する。
\end{proof}

\begin{lemma}
\label{sites-cohomology-lemma-equivalence-bounded}
$f : (\Sh(\mathcal{C}), \mathcal{O}) \to (\Sh(\mathcal{C}'), \mathcal{O}')$
を環付きトポスの射とする。$\mathcal{A} \subset \textit{Mod}(\mathcal{O})$ と
$\mathcal{A}' \subset \textit{Mod}(\mathcal{O}')$ を弱 Serre 部分圏とする。次を仮定する。
\begin{enumerate}
\item $f$ は平坦である。
\item $f^*$ は圏同値 $\mathcal{A}' \to \mathcal{A}$ を誘導する。
\item $\mathcal{F}' \in \Ob(\mathcal{A}')$ に対して
$\mathcal{F}' \to Rf_*f^*\mathcal{F}'$ は同型である。
\end{enumerate}
このとき
$f^* : D_{\mathcal{A}'}^+(\mathcal{O}') \to D_\mathcal{A}^+(\mathcal{O})$
は圏同値であり、その擬逆は
$Rf_* : D_\mathcal{A}^+(\mathcal{O}) \to D_{\mathcal{A}'}^+(\mathcal{O}')$.
\end{lemma}

\begin{proof}
仮定 (2)、(3) と補題 \ref{sites-cohomology-lemma-bounded-in-image-downstairs}、
\ref{sites-cohomology-lemma-downstairs} により、
$f^* : D_{\mathcal{A}'}^+(\mathcal{O}') \to D_\mathcal{A}^+(\mathcal{O})$
は充満忠実である。$\mathcal{F} \in \Ob(\mathcal{A})$ とする。このとき
$\mathcal{F} = f^*\mathcal{F}'$ と書ける。したがって
$Rf_*\mathcal{F} = Rf_* f^*\mathcal{F}' = \mathcal{F}'$.
特に $p > 0$ に対して $R^pf_*\mathcal{F} = 0$ であり、
$f_*\mathcal{F} \in \Ob(\mathcal{A}')$ である。したがって任意の
$K \in D^+_\mathcal{A}(\mathcal{O})$ に対し、$R^{p + q}f_*K$ に収束する
スペクトル系列 $E_2^{p, q} = R^pf_*H^q(K)$ を用いると、
$Rf_*K$ は $D^+_{\mathcal{A}'}(\mathcal{O}')$ に属することが分かる。
もちろん、補題 \ref{sites-cohomology-lemma-bounded-in-image-upstairs} と
\ref{sites-cohomology-lemma-upstairs} から、
$Rf_* : D_\mathcal{A}^+(\mathcal{O}) \to D_{\mathcal{A}'}^+(\mathcal{O}')$
も充満忠実であることが従う。$f^*$ と $Rf_*$ は随伴であるから、例えば
Categories, Lemma \ref{categories-lemma-adjoint-fully-faithful} により
補題の結論を得る。
\end{proof}

\begin{lemma}
\label{sites-cohomology-lemma-equivalence-unbounded-one}
\begin{reference}
これは \cite[Theorem 2.2.3]{six-I} と同様である。
\end{reference}
$f : (\Sh(\mathcal{C}), \mathcal{O}) \to (\Sh(\mathcal{C}'), \mathcal{O}')$
を環付きトポスの射とする。$\mathcal{A} \subset \textit{Mod}(\mathcal{O})$ と
$\mathcal{A}' \subset \textit{Mod}(\mathcal{O}')$ を弱 Serre 部分圏とする。次を仮定する。
\begin{enumerate}
\item $f$ は平坦である。
\item $f^*$ は圏同値 $\mathcal{A}' \to \mathcal{A}$ を誘導する。
\item $\mathcal{F}' \in \Ob(\mathcal{A}')$ に対して
$\mathcal{F}' \to Rf_*f^*\mathcal{F}'$ は同型である。
\item $\mathcal{C}, \mathcal{O}, \mathcal{A}$ は状況
\ref{sites-cohomology-situation-olsson-laszlo} の仮定を満たす。
\item $\mathcal{C}', \mathcal{O}', \mathcal{A}'$ は状況
\ref{sites-cohomology-situation-olsson-laszlo} の仮定を満たす。
\end{enumerate}
このとき $f^* : D_{\mathcal{A}'}(\mathcal{O}') \to D_\mathcal{A}(\mathcal{O})$
は圏同値であり、その擬逆は
$Rf_* : D_\mathcal{A}(\mathcal{O}) \to D_{\mathcal{A}'}(\mathcal{O}')$.
\end{lemma}

\begin{proof}
関手 $f^*$ は完全なので、$f^*$ が補題の主張にある関手
$f^* : D_{\mathcal{A}'}(\mathcal{O}') \to D_\mathcal{A}(\mathcal{O})$
を定め、さらにこの関手が切断関手 $\tau_{\geq -n}$ と可換であることは明らかである。
対応する下に有界な圏上では $f^*$ と $Rf_*$ が互いに擬逆な圏同値であることを
すでに知っている；補題 \ref{sites-cohomology-lemma-equivalence-bounded} を参照せよ。
補題 \ref{sites-cohomology-lemma-olsson-laszlo-map-version-one} で $N = 0$ と取ると、
$Rf_*$ は実際に関手
$Rf_* : D_\mathcal{A}(\mathcal{O}) \to D_{\mathcal{A}'}(\mathcal{O}')$
を定め、さらにこの関手が切断関手 $\tau_{\geq -n}$ と可換であることが分かる。
したがって $D_\mathcal{A}(\mathcal{O})$ の $K$ に対する写像
$f^*Rf_*K \to K$ は同型である。実際、各切断上で同型である。
同様に、$D_{\mathcal{A}'}(\mathcal{O}')$ の $K'$ に対する写像
$K' \to Rf_*f^*K'$ は、各切断上で同型であるから同型である。
これで証明が完了する。
\end{proof}

\begin{lemma}
\label{sites-cohomology-lemma-equivalence-unbounded-two}
\begin{reference}
これは \cite[Theorem 2.2.3]{six-I} と同様である。
\end{reference}
$f : (\mathcal{C}, \mathcal{O}) \to (\mathcal{C}', \mathcal{O}')$
を環付きサイトの射とする。$\mathcal{A} \subset \textit{Mod}(\mathcal{O})$ と
$\mathcal{A}' \subset \textit{Mod}(\mathcal{O}')$ を弱 Serre 部分圏とする。次を仮定する。
\begin{enumerate}
\item $f$ は平坦である。
\item $f^*$ は圏同値 $\mathcal{A}' \to \mathcal{A}$ を誘導する。
\item $\mathcal{F}' \in \Ob(\mathcal{A}')$ に対して
$\mathcal{F}' \to Rf_*f^*\mathcal{F}'$ は同型である。
\item $\mathcal{C}, \mathcal{O}, \mathcal{A}$ は状況
\ref{sites-cohomology-situation-olsson-laszlo} の仮定を満たす。
\item $f : (\mathcal{C}, \mathcal{O}) \to (\mathcal{C}', \mathcal{O}')$ と
$\mathcal{A}$ は状況 \ref{sites-cohomology-situation-olsson-laszlo-prime} の仮定を満たす。
\end{enumerate}
このとき $f^* : D_{\mathcal{A}'}(\mathcal{O}') \to D_\mathcal{A}(\mathcal{O})$
は圏同値であり、その擬逆は
$Rf_* : D_\mathcal{A}(\mathcal{O}) \to D_{\mathcal{A}'}(\mathcal{O}')$.
\end{lemma}

\begin{proof}
この補題の証明は補題 \ref{sites-cohomology-lemma-equivalence-unbounded-one} の証明と
まったく同じである。ただし、補題 \ref{sites-cohomology-lemma-olsson-laszlo-map-version-one}
への参照を補題 \ref{sites-cohomology-lemma-olsson-laszlo-map-version-two} への参照で置き換える。
\end{proof}








\section{二つの位相の比較、II}
\label{sites-cohomology-section-compare-II}

\noindent
$\mathcal{C}$ を圏とする。
$\text{Cov}(\mathcal{C}) \supset \text{Cov}'(\mathcal{C})$
を $\mathcal{C}$ にサイト構造を与える二つの方法とする。
$\text{Cov}(\mathcal{C})$ に対応する位相を $\tau$、
$\text{Cov}'(\mathcal{C})$ に対応する位相を $\tau'$ と書く。
このとき $\mathcal{C}$ 上の恒等関手はサイトの射
$$
\epsilon : \mathcal{C}_\tau \longrightarrow \mathcal{C}_{\tau'}
$$
を定める。ここで $\epsilon_*$ は基礎となる前層上の恒等関手であり、
$\epsilon^{-1}$ は $\tau'$-層の $\tau$-層化である（したがって明らかに完全である）。
$\mathcal{O}$ を $\tau$-位相に関する環の層とする。このとき $\mathcal{O}$ は
$\tau'$-位相に関しても層であり、$\epsilon$ は環付きサイトの射
$$
\epsilon :
(\mathcal{C}_\tau, \mathcal{O}_\tau)
\longrightarrow
(\mathcal{C}_{\tau'}, \mathcal{O}_{\tau'})
$$
となる。詳しくは第 \ref{sites-cohomology-section-compare} 節を参照せよ。

\begin{lemma}
\label{sites-cohomology-lemma-compare-topologies-derived-adequate-modules}
$\epsilon : (\mathcal{C}_\tau, \mathcal{O}_\tau) \to
(\mathcal{C}_{\tau'}, \mathcal{O}_{\tau'})$ を上の射とする。
$\mathcal{B} \subset \Ob(\mathcal{C})$ を部分集合とし、
$\mathcal{A} \subset \textit{PMod}(\mathcal{O})$ を充満部分圏とする。次を仮定する。
\begin{enumerate}
\item $\mathcal{A}$ のすべての対象は $\tau$-位相に関する層である。
\item $\mathcal{A}$ は $\textit{Mod}(\mathcal{O}_\tau)$ の弱 Serre 部分圏である。
\item $\mathcal{C}$ の各対象は、すべての要素が $\mathcal{B}$ に属する
$\tau'$-被覆をもつ。
\item すべての $U \in \mathcal{B}$、$\mathcal{F} \in \mathcal{A}$、$p > 0$ に対して
$H^p_\tau(U, \mathcal{F}) = 0$ である。
\end{enumerate}
このとき $\mathcal{A}$ は $\textit{Mod}(\mathcal{O}_{\tau'})$ の
弱 Serre 部分圏であり、$\epsilon^*$ と $R\epsilon_*$ によって与えられる三角圏の同値
$D_\mathcal{A}(\mathcal{O}_\tau) = D_\mathcal{A}(\mathcal{O}_{\tau'})$
が存在する。
\end{lemma}

\begin{proof}
$\epsilon^{-1}\mathcal{O}_{\tau'} = \mathcal{O}_\tau$ であるから、
$\epsilon$ は環付きサイトの平坦射であり、実際、加群層上では
$\epsilon^{-1} = \epsilon^*$ である。条件 (1) により、$\mathcal{A}$ の各対象を
$\mathcal{O}_\tau$-加群の層とも $\mathcal{O}_{\tau'}$-加群の層ともみなせる。
言い換えると、充満忠実な包含関手
$$
\mathcal{A} \to \textit{Mod}(\mathcal{O}_\tau) \to
\textit{Mod}(\mathcal{O}_{\tau'})
$$
がある。混乱を避けるため、$\mathcal{A}$ の像を
$\mathcal{A}' \subset \textit{Mod}(\mathcal{O}_{\tau'})$ と書く。このとき
$\epsilon_* : \mathcal{A} \to \mathcal{A}'$ と
$\epsilon^* : \mathcal{A}' \to \mathcal{A}$ が互いに擬逆な同値であることは明らかである
（補題の前の議論を参照し、$\mathcal{A}'$ の対象が $\tau$-位相に関する層であることを用いよ）。

\medskip\noindent
条件 (3)、(4) より、$p > 0$ かつ $\mathcal{F} \in \Ob(\mathcal{A})$ に対して
$R^p\epsilon_*\mathcal{F} = 0$ である。実際、$R^p\epsilon_*$ は前層
$U \mapsto H^p_\tau(U, \mathcal{F})$ に付随する層である；補題
\ref{sites-cohomology-lemma-higher-direct-images} を参照せよ。したがって $\mathcal{A}$ 内の任意の
完全な複体（これは各項が $\mathcal{A}$ に属する
$\textit{Mod}(\mathcal{O}_\tau)$ 内の完全な複体と同じである；Homology, Lemma
\ref{homology-lemma-characterize-weak-serre-subcategory} を参照せよ）は、
関手 $\epsilon_*$ を適用しても完全なままである。

\medskip\noindent
次の完全列を考える：
$$
\mathcal{F}'_0 \to \mathcal{F}'_1 \to
\mathcal{F}'_2 \to \mathcal{F}'_3 \to
\mathcal{F}'_4
$$
これは $\textit{Mod}(\mathcal{O}_{\tau'})$ 内の列であり、
$\mathcal{F}'_0, \mathcal{F}'_1, \mathcal{F}'_3, \mathcal{F}'_4$ は
$\mathcal{A}'$ に属するとする。完全関手 $\epsilon^*$ を適用すると完全列
$$
\epsilon^*\mathcal{F}'_0 \to \epsilon^*\mathcal{F}'_1 \to
\epsilon^*\mathcal{F}'_2 \to \epsilon^*\mathcal{F}'_3 \to
\epsilon^*\mathcal{F}'_4
$$
を $\textit{Mod}(\mathcal{O}_\tau)$ 内に得る。$\mathcal{A}$ は弱 Serre 部分圏であり、
$\epsilon^*\mathcal{F}'_0, \epsilon^*\mathcal{F}'_1,
\epsilon^*\mathcal{F}'_3, \epsilon^*\mathcal{F}'_4$ は
$\mathcal{A}$ に属するので、Homology, Definition
\ref{homology-definition-serre-subcategory} により
$\epsilon^*\mathcal{F}_2$ は $\mathcal{A}$ に属する。次の列の射を考える：
$$
\xymatrix{
\mathcal{F}'_0 \ar[r] \ar[d] &
\mathcal{F}'_1 \ar[r] \ar[d] &
\mathcal{F}'_2 \ar[r] \ar[d] &
\mathcal{F}'_3 \ar[r] \ar[d] &
\mathcal{F}'_4 \ar[d] \\
\epsilon_*\epsilon^*\mathcal{F}'_0 \ar[r] &
\epsilon_*\epsilon^*\mathcal{F}'_1 \ar[r] &
\epsilon_*\epsilon^*\mathcal{F}'_2 \ar[r] &
\epsilon_*\epsilon^*\mathcal{F}'_3 \ar[r] &
\epsilon_*\epsilon^*\mathcal{F}'_4
}
$$
前段の議論により下段は完全である。添字 $0$、$1$、$3$、$4$ の縦矢印は、
最初の段落の議論により同型である。五項補題（Homology, Lemma
\ref{homology-lemma-five-lemma}）から
$\mathcal{F}'_2 \cong \epsilon_*\epsilon^*\mathcal{F}'_2$ を得るので、
$\mathcal{F}'_2$ は $\mathcal{A}'$ に属する。こうして $\mathcal{A}'$ は
$\textit{Mod}(\mathcal{O}_{\tau'})$ の弱 Serre 部分圏であることが分かる；
Homology, Definition \ref{homology-definition-serre-subcategory} を参照せよ。

\medskip\noindent
この時点で導来圏 $D_\mathcal{A}(\mathcal{O}_\tau)$ と
$D_{\mathcal{A}'}(\mathcal{O}_{\tau'})$ を論じることに意味がある；
Derived Categories, Section \ref{derived-section-triangulated-sub} を参照せよ。
証明を終えるため、補題 \ref{sites-cohomology-lemma-equivalence-unbounded-two} の条件
(1)--(5) が適用できることを示す。(1)、(2)、(3) は上ですでに確認した。
各対象は $\mathcal{B}$ の対象からなる $\tau'$-被覆をもつので、まして
$\mathcal{B}$ の対象からなる $\tau$-被覆をもつ。したがって補題
\ref{sites-cohomology-lemma-equivalence-unbounded-two} の条件 (4) は満たされる。
同様に条件 (5) も満たされる。
\end{proof}

\begin{lemma}
\label{sites-cohomology-lemma-descent-squares}
$\epsilon : (\mathcal{C}_\tau, \mathcal{O}_\tau) \to
(\mathcal{C}_{\tau'}, \mathcal{O}_{\tau'})$ を上の射とする。
$A$ を集合とし、各 $\alpha \in A$ に対して
$$
\xymatrix{
E_\alpha \ar[d] \ar[r] & Y_\alpha \ar[d] \\
Z_\alpha \ar[r] & X_\alpha
}
$$
を圏 $\mathcal{C}$ における可換図式とする。次を仮定する。
\begin{enumerate}
\item $\tau'$-層 $\mathcal{F}'$ がすべての $\alpha$ に対して
$\mathcal{F}'(X_\alpha) =
\mathcal{F}'(Z_\alpha) \times_{\mathcal{F}'(E_\alpha)} 
\mathcal{F}'(Y_\alpha)$ を満たすならば、$\mathcal{F}'$ は $\tau$-層である。
\item $R\epsilon_*$ の本質像に属する $D(\mathcal{O}_{\tau'})$ の $K'$ に対し、
補題 \ref{sites-cohomology-lemma-c-square} の写像
$c^{K'}_{X_\alpha, Z_\alpha, Y_\alpha, E_\alpha}$ は
すべての $\alpha$ について同型である。
\end{enumerate}
このとき $K' \in D^+(\mathcal{O}_{\tau'})$ が $R\epsilon_*$ の本質像に属することと、
写像 $c^{K'}_{X_\alpha, Z_\alpha, Y_\alpha, E_\alpha}$ が
すべての $\alpha$ について同型であることは同値である。
\end{lemma}

\begin{proof}
必要性は仮定 (2) から従う。一方、$K'$ の非零コホモロジー層が
ただ一つならば、十分性は仮定 (1) から従う。一般の場合には、
帰納法で十分性を証明する。$D^+(\mathcal{O}_{\tau'})$ の対象 $K'$ が
(P) を満たすとは、写像 $c^{K'}_{X_\alpha, Z_\alpha, Y_\alpha, E_\alpha}$ が
すべての $\alpha \in A$ について同型であることとする。

\medskip\noindent
具体的に、$K'$ を (P) を満たす $D^+(\mathcal{O}_{\tau'})$ の対象とする。
$D^+(\mathcal{O}_{\tau'})$ における完全三角
$$
K' \to R\epsilon_*\epsilon^{-1}K' \to M' \to K'[1]
$$
を選ぶ。ここで最初の矢印は随伴写像である。(2) と補題
\ref{sites-cohomology-lemma-two-out-of-three-blow-up-square} により $M'$ は (P) を満たす。
他方、$\epsilon^{-1}$ を適用し、第 \ref{sites-cohomology-section-compare} 節の
$\epsilon^{-1}R\epsilon_* = \text{id}$ を用いると、$\epsilon^{-1}M' = 0$ を得る。
次の段落で $M' = 0$ を示せば証明が完了する。

\medskip\noindent
$K'$ を (P) と $\epsilon^{-1}K' = 0$ を満たす
$D^+(\mathcal{O}_{\tau'})$ の対象とする。$K' = 0$ を示そう。
具体的には、$i < n$ に対して $H^i(K') = 0$ となる $n \in \mathbf{Z}$ が
与えられたとき、$H^n(K') = 0$ を示す。各 $\alpha \in A$ に対し、補題
\ref{sites-cohomology-lemma-c-square} により完全三角
$$
R\Gamma_{\tau'}(X_\alpha, K') \to
R\Gamma_{\tau'}(Z_\alpha, K') \oplus
R\Gamma_{\tau'}(Y_\alpha, K') \to
R\Gamma_{\tau'}(E_\alpha, K') \to
R\Gamma_{\tau'}(X_\alpha, K')[1]
$$
がある。次数 $n$ のコホモロジーを取り、$K'$ のコホモロジー層に仮定した
消滅を用いると、完全列
$$
0 \to
H^n_{\tau'}(X_\alpha, K') \to
H^n_{\tau'}(Z_\alpha, K') \oplus
H^n_{\tau'}(Y_\alpha, K') \to
H^n_{\tau'}(E_\alpha, K')
$$
を得る。これは完全列
$$
0 \to
\Gamma(X_\alpha, H^n(K')) \to
\Gamma(Z_\alpha, H^n(K')) \oplus
\Gamma(Y_\alpha, H^n(K')) \to
\Gamma(E_\alpha, H^n(K'))
$$
と同じである。仮定 (1) により、$H^n(K')$ は $\tau$-層である。
しかし、その $\tau$-層化
$\epsilon^{-1}H^n(K') = H^n(\epsilon^{-1}K')$ は、
$\epsilon^{-1}K' = 0$ なので $0$ である。したがって所望のとおり
$H^n(K') = 0$ である。
\end{proof}

\begin{lemma}
\label{sites-cohomology-lemma-descent-squares-helper}
$\epsilon : (\mathcal{C}_\tau, \mathcal{O}_\tau) \to
(\mathcal{C}_{\tau'}, \mathcal{O}_{\tau'})$ を上の射とする。
$$
\xymatrix{
E \ar[d] \ar[r] & Y \ar[d] \\
Z \ar[r] & X
}
$$
を圏 $\mathcal{C}$ における可換図式とし、次を仮定する。
\begin{enumerate}
\item $h_X^\# = h_Y^\# \amalg_{h_E^\#} h_Z^\#$ である。
\item $h_E^\# \to h_Y^\#$ は単射である。
\end{enumerate}
ここで ${}^\#$ は $\tau$-層化を表す。このとき $R\epsilon_*$ の本質像に属する
$K' \in D(\mathcal{O}_{\tau'})$ に対し、補題 \ref{sites-cohomology-lemma-c-square} の写像
$c^{K'}_{X, Z, Y, E}$（$\tau'$-位相を用いる）は同型である。
\end{lemma}

\begin{proof}
この補助補題は補題 \ref{sites-cohomology-lemma-square-triangle} のほぼ直ちに得られる帰結であり、
読者には証明を飛ばすことを強く勧める。$K' = R\epsilon_*K$ と書く。
$K$ を表す $\mathcal{O}_\tau$-加群の K-入射複体 $\mathcal{J}^\bullet$ を選ぶ。
このとき $\epsilon_*\mathcal{J}^\bullet$ は $K'$ を表す
$\mathcal{O}_{\tau'}$-加群の K-入射複体である；補題
\ref{sites-cohomology-lemma-K-injective-flat} を参照せよ。次に、
$$
0 \to
\mathcal{J}^\bullet(X) \xrightarrow{\alpha}
\mathcal{J}^\bullet(Z) \oplus
\mathcal{J}^\bullet(Y) \xrightarrow{\beta}
\mathcal{J}^\bullet(E)
\to 0
$$
はアーベル群の複体の短完全列である；補題
\ref{sites-cohomology-lemma-square-triangle} とその証明を参照せよ。これは
$c^{K'}_{X, Z, Y, E}$ の定義に用いるアーベル群の複体の列と同じなので、
結論が従う。
\end{proof}







\section{コホモロジーの比較}
\label{sites-cohomology-section-compare-general}

\noindent
異なる位相に関するコホモロジーを比較するための一般論を展開する。
第 \ref{sites-cohomology-section-compare} 節の $\mathcal{C}$、$\tau$、$\tau'$ と、
$\mathcal{C}$ の射 $f : X \to Y$ が与えられると、トポスの射の可換図式
\begin{equation}
\label{sites-cohomology-equation-commutative-epsilon}
\vcenter{
\xymatrix{
\Sh(\mathcal{C}_\tau/X) \ar[r]_{f_\tau} \ar[d]_{\epsilon_X} &
\Sh(\mathcal{C}_\tau/Y) \ar[d]^{\epsilon_Y} \\
\Sh(\mathcal{C}_{\tau'}/X) \ar[r]^{f_{\tau'}} &
\Sh(\mathcal{C}_{\tau'}/X)
}
}
\end{equation}
を得る。ここで射 $\epsilon_X$（それぞれ $\epsilon_Y$）は、二つの位相
$\tau$ と $\tau'$ を備えた圏 $\mathcal{C}/X$ に対する第
\ref{sites-cohomology-section-compare} 節の比較射である。射 $f_\tau$ と $f_{\tau'}$ は
「再局所化」射である（Sites, Lemma \ref{sites-lemma-relocalize}）。
この図式の可換性は Sites, Lemma \ref{sites-lemma-localize-morphism} の特別な場合である
（$\mathcal{C} = \mathcal{C}_\tau/Y$、
$\mathcal{D} = \mathcal{C}_{\tau'}/Y$、$u = \text{id}$、$U = X$、$V = X$
として適用する）。さらに
$\epsilon_{X, *} \circ f_\tau^{-1} = f_{\tau'}^{-1} \circ \epsilon_{Y, *}$
も得る。これは同補題から従い、また直接にも明らかである。

\begin{situation}
\label{sites-cohomology-situation-compare}
$\mathcal{C}$、$\tau$、$\tau'$ は第 \ref{sites-cohomology-section-compare} 節のものとする。
部分集合 $\mathcal{P} \subset \text{Arrows}(\mathcal{C})$ と、
$\mathcal{C}$ の各対象 $X$ に対して弱 Serre 部分圏
$\mathcal{A}'_X \subset \textit{Ab}(\mathcal{C}_{\tau'}/X)$ が与えられているとする。
次を仮定する。
\begin{enumerate}
\item
\label{sites-cohomology-item-base-change-P}
$f : X \to Y$ が $\mathcal{P}$ に属し、$Y' \to Y$ が任意の射ならば、
$X \times_Y Y'$ が存在し、$X \times_Y Y' \to Y'$ は $\mathcal{P}$ に属する。
\item
\label{sites-cohomology-item-restriction-A}
$\mathcal{C}$ の任意の射 $f : X \to Y$ に対し、$f_{\tau'}^{-1}$ は
$\mathcal{A}'_Y$ を $\mathcal{A}'_X$ に写す。
\item
\label{sites-cohomology-item-A-sheaf}
$X$ が $\mathcal{C}$ の対象で $\mathcal{F}'$ が $\mathcal{A}'_X$ に属するならば、
$\mathcal{F}'$ は $\tau$-被覆に関する層条件を満たす。すなわち
$\mathcal{F}' = \epsilon_{X, *}\epsilon_X^{-1}\mathcal{F}'$,
\item
\label{sites-cohomology-item-A-and-P}
$f : X \to Y$ が $\mathcal{P}$ に属し、
$\mathcal{F}' \in \Ob(\mathcal{A}'_X)$ ならば、
$R^if_{\tau', *}\mathcal{F}' \in \Ob(\mathcal{A}'_Y)$ が
すべての $i \geq 0$ について成り立つ。
\item
\label{sites-cohomology-item-refine-tau-by-P}
$\{U_i \to U\}_{i \in I}$ が $\tau$-被覆ならば、次が存在する。
\begin{enumerate}
\item $\tau'$-被覆 $\{V_j \to U\}_{j \in J}$。
\item 単一の $f_j \in \mathcal{P}$ からなる $\tau$-被覆
$\{f_j : W_j \to V_j\}$。
\item $\tau'$-被覆 $\{W_{jk} \to W_j\}_{k \in K_j}$。
\end{enumerate}
ここで $\{W_{jk} \to U\}_{j \in J, k \in K_j}$ は
$\{U_i \to U\}_{i \in I}$ の細分である。
\end{enumerate}
\end{situation}

\begin{lemma}
\label{sites-cohomology-lemma-A}
状況 \ref{sites-cohomology-situation-compare} において、$\mathcal{C}$ の対象 $X$ に対し、
$\mathcal{F}' \in \mathcal{A}'_X$ を用いて $\epsilon_X^{-1}\mathcal{F}'$ と書ける
$\textit{Ab}(\mathcal{C}_\tau/X)$ の対象全体を $\mathcal{A}_X$ と書く。このとき
\begin{enumerate}
\item $\textit{Ab}(\mathcal{C}_\tau/X)$ の $\mathcal{F}$ に対し
$\mathcal{F} \in \mathcal{A}_X \Leftrightarrow
\epsilon_{X, *}\mathcal{F} \in \mathcal{A}'_X$ である。
\item $\mathcal{C}$ の任意の射 $f : X \to Y$ に対し、$f_\tau^{-1}$ は
$\mathcal{A}_Y$ を $\mathcal{A}_X$ に写す。
\end{enumerate}
\end{lemma}

\begin{proof}
(1) は (\ref{sites-cohomology-item-A-sheaf}) から従う。(2) は
(\ref{sites-cohomology-item-restriction-A}) と、次を与える
(\ref{sites-cohomology-equation-commutative-epsilon}) の可換性から従う：
$\epsilon_X^{-1} \circ f_{\tau'}^{-1} = f_\tau^{-1} \circ \epsilon_Y^{-1}$.
\end{proof}

\noindent
次の目標は補題 \ref{sites-cohomology-lemma-compare-cohomology-general} と
\ref{sites-cohomology-lemma-cohomological-descent-general} を証明することである。
次の帰納法の仮定を用いて証明する。

\medskip\noindent
$(V_n)$ $X$ が $\mathcal{C}$ の対象で $\mathcal{F}$ が $\mathcal{A}_X$ に属するとき、
$1 \leq i \leq n$ に対して $R^i\epsilon_{X, *}\mathcal{F} = 0$ である。

\begin{lemma}
\label{sites-cohomology-lemma-V-implies-C-general}
状況 \ref{sites-cohomology-situation-compare} において $(V_n)$ が成り立つと仮定する。
$\mathcal{P}$ に属する $f : X \to Y$ と $\mathcal{A}_X$ に属する
$\mathcal{F}$ に対し、$i \leq n$ ならば
$R^if_{\tau', *}\epsilon_{X, *}\mathcal{F} =
\epsilon_{Y, *}R^if_{\tau, *}\mathcal{F}$ である。
\end{lemma}

\begin{proof}
以下、可換図式 (\ref{sites-cohomology-equation-commutative-epsilon}) を断りなく用いる。特に
$$
Rf_{\tau', *}R\epsilon_{X, *}\mathcal{F} =
R\epsilon_{Y, *}Rf_{\tau, *}\mathcal{F}
$$
である。仮定 $(V_n)$ により、写像
$\epsilon_{X, *}\mathcal{F} \to R\epsilon_{X, *}\mathcal{F}$ は
次数 $\leq n$ で同型である。したがって
$Rf_{\tau', *}\epsilon_{X, *}\mathcal{F} \to
Rf_{\tau', *}R\epsilon_{X, *}\mathcal{F}$ は次数 $\leq n$ で同型である。ゆえに
$$
R^if_{\tau', *}\epsilon_{X, *}\mathcal{F} \to
H^i(R\epsilon_{Y, *}Rf_{\tau, *}\mathcal{F})
$$
は $i \leq n$ に対して同型である。補題 \ref{sites-cohomology-lemma-relative-Leray} の
$R\epsilon_{Y, *}Rf_{\tau, *}\mathcal{F}$ に対するスペクトル系列の
第 2 頁を調べて補題を証明する。その図は次のとおりである：
$$
\begin{matrix}
\ldots &
\ldots &
\ldots &
\ldots \\
\epsilon_{Y, *}R^2f_{\tau, *}\mathcal{F} &
R^1\epsilon_{Y, *}R^2f_{\tau, *}\mathcal{F} &
R^2\epsilon_{Y, *}R^2f_{\tau, *}\mathcal{F} &
\ldots \\
\epsilon_{Y, *}R^1f_{\tau, *}\mathcal{F} &
R^1\epsilon_{Y, *}R^1f_{\tau, *}\mathcal{F} &
R^2\epsilon_{Y, *}R^1f_{\tau, *}\mathcal{F} &
\ldots \\
\epsilon_{Y, *}f_{\tau, *}\mathcal{F} &
R^1\epsilon_{Y, *}f_{\tau, *}\mathcal{F} &
R^2\epsilon_{Y, *}f_{\tau, *}\mathcal{F} &
\ldots
\end{matrix}
$$
$(C_m)$ を、$i \leq m$ に対して
$R^if_{\tau', *}\epsilon_{X, *}\mathcal{F} =
\epsilon_{Y, *}R^if_{\tau, *}\mathcal{F}$ が成り立つという仮定とする。
$(C_0)$ は成り立つ。$m < n$ に対して
$(C_{m - 1}) \Rightarrow (C_m)$ を示そう。実際、$(C_{m - 1})$ が成り立つならば、
$n \geq p > 0$ かつ $q \leq m - 1$ に対して
\begin{align*}
R^p\epsilon_{Y, *}R^qf_{\tau, *}\mathcal{F}
& =
R^p\epsilon_{Y, *}
\epsilon_Y^{-1} \epsilon_{Y, *} R^qf_{\tau, *}\mathcal{F} \\
& =
R^p\epsilon_{Y, *}
\epsilon_Y^{-1}R^qf_{\tau', *}\epsilon_{X, *}\mathcal{F} = 0
\end{align*}
となる。最初の等号は $\epsilon_Y^{-1}\epsilon_{Y, *} = \text{id}$ により、
二番目は $(C_{m - 1})$ により、最後は $(V_n)$ による。実際、
(\ref{sites-cohomology-item-A-and-P}) により
$\epsilon_Y^{-1}R^qf_{\tau', *}\epsilon_{X, *}\mathcal{F}$ は
$\mathcal{A}_Y$ に属する。スペクトル系列を見ると、
$E_2^{0, m} = \epsilon_{Y, *}R^mf_{\tau, *}\mathcal{F}$
は $p + q = m$ を満たす唯一の非零項 $E_2^{p, q}$ である。
$\text{d}_r^{p, q} : E_r^{p, q} \to E_r^{p + r, q - r + 1}$
であったことを思い出す。したがって $r \geq 2$ に対し、この位置から出る、または
この位置に入る非零微分 $\text{d}_r^{p, q}$ はない。ゆえに
$H^m(R\epsilon_{Y, *}Rf_{\tau, *}\mathcal{F}) =
\epsilon_{Y, *}R^mf_{\tau, *}\mathcal{F}$ であり、上の議論から $(C_m)$ が従う。

\medskip\noindent
最後に $(C_{n - 1})$ を仮定する。同じ解析により
$E_2^{0, n} = \epsilon_{Y, *}R^nf_{\tau, *}\mathcal{F}$
は $p + q = n$ を満たす唯一の非零項 $E_2^{p, q}$ である。
この位置に入る非零微分は依然としてないが、ここから出る非零微分はあり得る。
それは写像
$d_{n + 1}^{0, n} : \epsilon_{Y, *}R^nf_{\tau, *}\mathcal{F} \to
R^{n + 1}\epsilon_{Y, *}f_{\tau, *}\mathcal{F}$ である。したがって完全列
$$
0 \to R^nf_{\tau', *}\epsilon_{X, *}\mathcal{F} \to 
\epsilon_{Y, *}R^nf_{\tau, *}\mathcal{F} \to
R^{n + 1}\epsilon_{Y, *}f_{\tau, *}\mathcal{F}
$$
を得る。(\ref{sites-cohomology-item-A-and-P}) と (\ref{sites-cohomology-item-A-sheaf}) により、層
$R^nf_{\tau', *}\epsilon_{X, *}\mathcal{F}$
は、$\epsilon_{Y, *}R^nf_{\tau, *}\mathcal{F}$ と同様、$\tau$-被覆に関する
層条件を満たす（第 \ref{sites-cohomology-section-compare} 節の $\epsilon_*$ の記述を用いよ）。
しかし $\tau'$-層
$R^{n + 1}\epsilon_{Y, *}f_{\tau, *}\mathcal{F}$
の $\tau$-層化は $0$ である（コホモロジーの局所性による；補題
\ref{sites-cohomology-lemma-kill-cohomology-class-on-covering} と
\ref{sites-cohomology-lemma-higher-direct-images} を用いよ）。したがって
$R^nf_{\tau', *}\epsilon_{X, *}\mathcal{F} \to
\epsilon_{Y, *}R^nf_{\tau, *}\mathcal{F}$
は同型でなければならず、証明が完了する。
\end{proof}

\noindent
$E'$（それぞれ $E$）が $D(\mathcal{C}_{\tau'}/X)$（それぞれ
$D(\mathcal{C}_\tau/X)$）の対象ならば、$\mathcal{C}/X$ の対象 $U$ 上の
$E'$（それぞれ $E$）のコホモロジーを
$H^n_{\tau'}(U, E')$（それぞれ $H^n_\tau(U, E)$）
と書く。

\begin{lemma}
\label{sites-cohomology-lemma-V-implies-cohomology-general}
状況 \ref{sites-cohomology-situation-compare} において $(V_n)$ が成り立つとする。
$X$ が $\mathcal{C}$ の対象で、$L \in D(\mathcal{C}_{\tau'}/X)$ が
$i < 0$ に対して $H^i(L) = 0$、$0 \leq i \leq n$ に対して
$H^i(L) \in \mathcal{A}'_X$ を満たすならば、
$H^n_{\tau'}(X, L) = H^n_\tau(X, \epsilon_X^{-1}L)$.
\end{lemma}

\begin{proof}
補題 \ref{sites-cohomology-lemma-Leray-unbounded} により
$H^n_\tau(X, \epsilon_X^{-1}L) =
H^n_{\tau'}(X, R\epsilon_{X, *}\epsilon_X^{-1}L)$.
また、次のスペクトル系列がある：
$$
E_2^{p, q} = R^p\epsilon_{X, *}\epsilon_X^{-1}H^q(L)
$$
$H^{p + q}(R\epsilon_{X, *}\epsilon_X^{-1}L)$ に収束する。
$(V_n)$ により、$0 < p \leq n$ かつ $0 \leq q \leq n$ に対して
$E_2^{p, q}$ は消滅する。したがって
$E_2^{0, q} = \epsilon_{X, *}\epsilon_X^{-1}H^q(L) = H^q(L)$
が、$p + q \leq n$ を満たす唯一の非零項 $E_2^{p, q}$ である。よって写像
$$
L \longrightarrow R\epsilon_{X, *}\epsilon_X^{-1}L
$$
は次数 $\leq n$ で同型である（細部は省略する）。したがって
$H^i_{\tau'}(X, L) = H^i_{\tau'}(X, R\epsilon_{X, *}\epsilon_X^{-1}L)$
が $i \leq n$ に対して成り立つ。これで補題が証明された。
\end{proof}

\begin{lemma}
\label{sites-cohomology-lemma-V-implies-cohomology-extra-general}
状況 \ref{sites-cohomology-situation-compare} において $(V_n)$ が成り立つとする。
$X$ が $\mathcal{C}$ の対象で $\mathcal{F}$ が $\mathcal{A}_X$ に属するとき、写像
$H^{n + 1}_{\tau'}(X, \epsilon_{X, *}\mathcal{F}) \to
H^{n + 1}_\tau(X, \mathcal{F})$
は単射であり、その像は $X$ のある $\tau'$-被覆上で零となる類全体である。
\end{lemma}

\begin{proof}
$\epsilon_X^{-1}\epsilon_{X, *}\mathcal{F} = \mathcal{F}$ であったことを思い出す。
したがってこの写像は、$\epsilon_X$ によるコホモロジー類の引き戻しで与えられる。
Leray スペクトル系列（補題 \ref{sites-cohomology-lemma-Leray}）
$$
E_2^{p, q} = H^p_{\tau'}(X, R^q\epsilon_{X, *}\mathcal{F})
\Rightarrow
H^{p + q}_\tau(X, \mathcal{F})
$$
と仮定した消滅を組み合わせると、完全列
$$
0 \to
H^{n + 1}_{\tau'}(X, \epsilon_{X, *}\mathcal{F}) \to
H^{n + 1}_\tau(X, \mathcal{F}) \to
H^0_{\tau'}(X, R^{n + 1}\epsilon_{X, *}\mathcal{F})
$$
を得る。これは補題の言い換えである。
\end{proof}

\begin{lemma}
\label{sites-cohomology-lemma-make-class-zero-general}
状況 \ref{sites-cohomology-situation-compare} において $f : X \to Y$ は $\mathcal{P}$ に属し、
$\{X \to Y\}$ は $\tau$-被覆であるとする。
$\mathcal{F}' \in \mathcal{A}'_Y$ とする。$n \geq 0$ かつ
$$
\theta \in
\text{Equalizer}\left(
\xymatrix{
H^{n + 1}_{\tau'}(X, \mathcal{F}')
\ar@<1ex>[r] \ar@<-1ex>[r] &
H^{n + 1}_{\tau'}(X \times_Y X, \mathcal{F}')
}
\right)
$$
ならば、$\theta$ の $H^{n + 1}_{\tau'}(Y_i \times_Y X, \mathcal{F}')$
への制限が零となるような $\tau'$-被覆 $\{Y_i \to Y\}$ が存在する。
\end{lemma}

\begin{proof}
(\ref{sites-cohomology-item-base-change-P}) により $X \times_Y X$ が存在することに注意する。
$\mathcal{C}/Y$ の対象 $Z$ に対し、$\mathcal{F}'$ の
$\mathcal{C}_{\tau'}/Z$ への制限を $\mathcal{F}'|_Z$ と書く。補題
\ref{sites-cohomology-lemma-cohomology-of-open} により
$H^{n + 1}_{\tau'}(X, \mathcal{F}') =
H^{n + 1}(\mathcal{C}_{\tau'}/X, \mathcal{F}'|_X)$
であったことを思い出す。補題の主張は、$\theta$ の像
$\overline{\theta} \in H^0(Y, R^{n + 1}f_{\tau', *}\mathcal{F}'|_X)$
が零であるということである。Cartesian 図式
$$
\xymatrix{
X \times_Y X \ar[d]_{\text{pr}_1} \ar[r]_{\text{pr}_2} &
X \ar[d]^f \\
X \ar[r]^f & Y
}
$$
を考える。自明な基底変換（補題 \ref{sites-cohomology-lemma-localize-cartesian-square}）により
$$
f_{\tau'}^{-1}R^{n + 1}f_{\tau', *}(\mathcal{F}'|_X) =
R^{n + 1}\text{pr}_{1, \tau', *}(\mathcal{F}'|_{X \times_Y X})
$$
である。$\text{pr}_1^{-1}\theta = \text{pr}_2^{-1}\theta$ ならば、
$f_{\tau'}^{-1}R^{n + 1}f_{\tau', *}(\mathcal{F}'|_X)$ の切断
$f_{\tau'}^{-1}\overline{\theta}$ は零である。実際、
$\text{pr}_1^{-1}\theta$ が
$H^0(X, R^{n + 1}\text{pr}_{1, \tau', *}(\mathcal{F}'|_{X \times_Y X}))$
の零元に写ることは明らかである。(\ref{sites-cohomology-item-restriction-A}) により
$\mathcal{F}'|_X \in \mathcal{A}'_X$ である。したがって
$\mathcal{G}' = R^{n + 1}f_{\tau', *}(\mathcal{F}'|_X)$ は
(\ref{sites-cohomology-item-A-and-P}) により $\mathcal{A}'_Y$ の対象である。ゆえに
(\ref{sites-cohomology-item-A-sheaf}) により $\mathcal{G}'$ は $\tau$-被覆に関する層条件を満たす。
$\{X \to Y\}$ は $\tau$-被覆なので、制限写像
$\mathcal{G}'(Y) \to \mathcal{G}'(X)$ は単射である。したがって
$\overline{\theta}$ は零である。
\end{proof}

\begin{lemma}
\label{sites-cohomology-lemma-induction-step-V-C-general}
状況 \ref{sites-cohomology-situation-compare} において $(V_n) \Rightarrow (V_{n + 1})$ である。
\end{lemma}

\begin{proof}
$X$ を $\mathcal{C}$ の対象、$\mathcal{F}$ を $\mathcal{A}_X$ の対象とする。
ある $U/X$ に対して $\xi \in H^{n + 1}_\tau(U, \mathcal{F})$ とする。
$U$ のある $\tau'$-被覆の各要素への $\xi$ の制限が零であることを示す必要がある；
補題 \ref{sites-cohomology-lemma-higher-direct-images} を参照せよ。このことから、$U$ を
$U$ の $\tau'$-被覆の各要素で置き換えてよい。

\medskip\noindent
コホモロジーの局所性（補題 \ref{sites-cohomology-lemma-kill-cohomology-class-on-covering}）により、
$\xi$ の各 $U_i$ への制限が零となる $\tau$-被覆
$\{U_i \to U\}$ を選べる。(\ref{sites-cohomology-item-refine-tau-by-P}) のように
$\{V_j \to V\}$、$\{f_j : W_j \to V_j\}$、
$\{W_{jk} \to W_j\}$ を選ぶ。$U$ を $V_j$ で、$\mathcal{F}$ を
$\mathcal{C}_\tau/V_j$ への制限で置き換えることが
(\ref{sites-cohomology-item-base-change-P}) により許されるので、次の段落で論じる場合に帰着する。

\medskip\noindent
ここで $f : X \to Y$ は $\mathcal{P}$ の要素で、$\{X \to Y\}$ は
$\tau$-被覆、$\mathcal{F}$ は $\mathcal{A}_Y$ の対象であるとする。また、
$\xi \in H^{n + 1}_\tau(Y, \mathcal{F})$ に対し、すべての $i \in I$ について
$\xi$ の $X_i$ への制限が零となる $\tau'$-被覆
$\{X_i \to X\}_{i \in I}$ が存在するとする。
課題は、$Y$ のある $\tau'$-被覆上で $\xi$ が零となることを示すことである。

\medskip\noindent
補題 \ref{sites-cohomology-lemma-V-implies-cohomology-extra-general} により、一意な
$\tau'$-コホモロジー類
$\theta \in H^{n + 1}_{\tau'}(X, \epsilon_{X, *}\mathcal{F})$
で像が $\xi|_X$ となるものが存在する。$\xi|_X$ は二つの射影によって
$X \times_Y X$ 上の同じ類へ引き戻されるので、$\theta$ についても同じことが
成り立つ（一意性による）。補題 \ref{sites-cohomology-lemma-make-class-zero-general} により、
$Y$ をある $\tau'$-被覆の各要素で置き換えれば $\theta = 0$ と仮定してよい。
したがって $\xi|_X$ が零であると仮定してよい。

\medskip\noindent
$f : X \to Y$ を $\mathcal{P}$ の要素で $\{X \to Y\}$ が $\tau$-被覆となるもの、
$\mathcal{F}$ を $\mathcal{A}_Y$ の対象とする。また、
$\xi \in H^{n + 1}_\tau(Y, \mathcal{F})$ の
$H^{n + 1}_\tau(X, \mathcal{F})$ における像は零であるとする。
課題は、$Y$ のある $\tau'$-被覆上で $\xi$ が零となることを示すことである。

\medskip\noindent
仮定により、$\xi$ は写像
$$
\mathcal{F} \longrightarrow Rf_{\tau, *}f_\tau^{-1}\mathcal{F}
$$
によって零に写る。補題 \ref{sites-cohomology-lemma-Leray-unbounded} を用いよ。
切断の完全三角（Derived Categories, Remark
\ref{derived-remark-truncation-distinguished-triangle}）を用いる簡単な議論により、
$\xi$ は写像
$$
\mathcal{F} \longrightarrow \tau_{\leq n}Rf_{\tau, *}f_\tau^{-1}\mathcal{F}
$$
によって零に写る。これを写像 $\epsilon_{Y, *}\mathcal{F} \to K$ と比較する。ここで
$$
K = \tau_{\leq n}Rf_{\tau', *}f_{\tau'}^{-1}\epsilon_{Y, *}\mathcal{F} =
\tau_{\leq n}Rf_{\tau', *}\epsilon_{X, *}f_{\tau}^{-1}\mathcal{F}
$$
である。等式
$\epsilon_{X, *} f_\tau^{-1} = f_{\tau'}^{-1} \epsilon_{Y, *}$
は (\ref{sites-cohomology-equation-commutative-epsilon}) の性質である。補題
\ref{sites-cohomology-lemma-V-implies-C-general} の証明で用いた写像
$$
Rf_{\tau', *}\epsilon_{X, *}f_{\tau}^{-1}\mathcal{F} \longrightarrow
Rf_{\tau', *}R\epsilon_{X, *}f_{\tau}^{-1}\mathcal{F} =
R\epsilon_{Y, *}Rf_{\tau, *}f_\tau^{-1}\mathcal{F}
$$
を考える。これは随伴により写像
$$
\epsilon_Y^{-1} Rf_{\tau', *}\epsilon_{X, *}f_{\tau}^{-1}\mathcal{F} \to
Rf_{\tau, *}f_\tau^{-1}\mathcal{F}
$$
を誘導する。切断を取ると写像
$$
\epsilon_Y^{-1}K
\longrightarrow
\tau_{\leq n}Rf_{\tau, *}f_\tau^{-1}\mathcal{F}
$$
を得る。これは補題 \ref{sites-cohomology-lemma-V-implies-C-general} により同型である。
実際、補題 \ref{sites-cohomology-lemma-A} により $f_\tau^{-1}\mathcal{F}$ は
$\mathcal{A}_X$ に属するので、同補題を適用できる。完全三角
$$
\epsilon_{Y, *}\mathcal{F} \to K \to L \to \epsilon_{Y, *}\mathcal{F}[1]
$$
を選ぶ。$\{X \to Y\}$ は $\tau$-被覆なので、写像
$\mathcal{F} \to f_{\tau, *}f_\tau^{-1}\mathcal{F}$ は単射である。したがって
$\epsilon_{Y, *}\mathcal{F} \to
\epsilon_{Y, *}f_{\tau, *}f_\tau^{-1}\mathcal{F} =
f_{\tau', *}f_{\tau'}^{-1}\epsilon_{Y, *}\mathcal{F}$
も単射である。ゆえに $L$ の非零コホモロジー層は次数 $0, \ldots, n$ にしかない。
(\ref{sites-cohomology-item-restriction-A}) と (\ref{sites-cohomology-item-A-and-P}) により
$f_{\tau', *}f_{\tau'}^{-1}\epsilon_{Y, *}\mathcal{F}$ は
$\mathcal{A}'_Y$ に属するので、
$$
H^0(L) = \Coker(\epsilon_{Y, *}\mathcal{F} \to
f_{\tau', *}f_{\tau'}^{-1}\epsilon_{Y, *}\mathcal{F})
$$
は弱 Serre 部分圏 $\mathcal{A}'_Y$ に属する。$1 \leq i \leq n$ に対しても、
(\ref{sites-cohomology-item-restriction-A}) と (\ref{sites-cohomology-item-A-and-P}) により
$H^i(L) = R^if_{\tau', *}f_{\tau'}^{-1}\epsilon_{Y, *}\mathcal{F}$ は
$\mathcal{A}'_Y$ に属する。上の完全三角を $\epsilon_Y$ で引き戻すと、完全三角
$$
\mathcal{F} \to \tau_{\leq n}Rf_{\tau, *}f_\tau^{-1}\mathcal{F}
\to \epsilon_Y^{-1}L \to \mathcal{F}[1]
$$
を得る。$\xi$ は中間項で零に写るので、ある元
$\xi' \in H^n_\tau(Y, \epsilon_Y^{-1}L)$ の像である。補題
\ref{sites-cohomology-lemma-V-implies-cohomology-general} により
$$
H^n_{\tau'}(Y, L) = H^n_\tau(Y, \epsilon_Y^{-1}L),
$$
である。したがって $\xi'$ を $H^n_{\tau'}(Y, L)$ の元に持ち上げ、
$H^{n + 1}_{\tau'}(Y, \epsilon_{Y, *}\mathcal{F})$ への境界を取ると、
$\xi$ が標準写像
$H^{n + 1}_{\tau'}(Y, \epsilon_{Y, *}\mathcal{F}) \to
H^{n + 1}_\tau(Y, \mathcal{F})$
の像に属することが分かる。
$H^{n + 1}_{\tau'}(Y,\epsilon_{Y, *}\mathcal{F})$ に対するコホモロジーの局所性
（補題 \ref{sites-cohomology-lemma-kill-cohomology-class-on-covering}）により結論を得る。
\end{proof}

\begin{lemma}
\label{sites-cohomology-lemma-V-C-all-n-general}
状況 \ref{sites-cohomology-situation-compare} において、$(V_n)$ は任意の $n$ に対して成り立つ。
さらに、次が成り立つ。
\begin{enumerate}
\item $X$ を $\mathcal{C}$ の対象とし、
$K' \in D^+_{\mathcal{A}'_X}(\mathcal{C}_{\tau'}/X)$ とする。このとき写像
$K' \to R\epsilon_{X, *}(\epsilon_X^{-1}K')$ は同型である。
\item $f : X \to Y$ を $\mathcal{P}$ の射とし、
$K' \in D^+_{\mathcal{A}'_X}(\mathcal{C}_{\tau'}/X)$ とする。このとき
$Rf_{\tau', *}K' \in D^+_{\mathcal{A}'_X}(\mathcal{C}_{\tau'}/Y)$ であり、
$\epsilon_Y^{-1}(Rf_{\tau', *}K') = Rf_{\tau, *}(\epsilon_X^{-1}K')$。
\end{enumerate}
\end{lemma}

\begin{proof}
$(V_0)$ は空条件なので成り立つことに注意する。
したがって補題 \ref{sites-cohomology-lemma-induction-step-V-C-general} により、
任意の $n$ に対して $(V_n)$ が成り立つ。

\medskip\noindent
（1）の証明。対象 $K = \epsilon_X^{-1}K'$ のコホモロジー層
$H^i(K) = \epsilon_X^{-1}H^i(K')$ は $\mathcal{A}_X$ に属する。
したがってスペクトル系列
$$
E_2^{p, q} = R^p\epsilon_{X, *} H^q(K) \Rightarrow
H^{p + q}(R\epsilon_{X, *}K)
$$
は、任意の $n$ に対する $(V_n)$ により退化し、
$$
H^n(R\epsilon_{X, *}K) = \epsilon_{X, *}H^n(K) =
\epsilon_{X, *}\epsilon_X^{-1}H^i(K') = H^i(K').
$$
を得る。最後の等号でも $H^i(K')$ が $\mathcal{A}'_X$ に属することを用いた。
ゆえに標準写像 $K' \to R\epsilon_{X, *}(\epsilon_X^{-1}K')$
は同型である。

\medskip\noindent
（2）の証明。スペクトル系列
$$
E_2^{p, q} = R^pf_{\tau', *}H^q(K') \Rightarrow R^{p + q}f_{\tau', *}K'
$$
と、（\ref{sites-cohomology-item-A-and-P}）により $R^pf_{\tau', *}H^q(K')$ が
$\mathcal{A}'_Y$ に属すること、$\mathcal{A}'_Y$ が
$\textit{Ab}(\mathcal{C}_{\tau'}/Y)$ の弱 Serre 部分圏であること、
および Homology, Lemma \ref{homology-lemma-biregular-ss-converges} を用いると、
$Rf_{\tau', *}K' \in D^+_{\mathcal{A}'_X}(\mathcal{C}_{\tau'}/X)$
を得る。証明を終えるには、基底変換写像
$$
\epsilon_Y^{-1}(Rf_{\tau', *}K')
\longrightarrow
Rf_{\tau, *}(\epsilon_X^{-1}K')
$$
が同型であることを示せばよい。上のスペクトル系列をスペクトル系列
$$
E_2^{p, q} = R^pf_{\tau, *}H^q(\epsilon_X^{-1}K')
\Rightarrow R^{p + q}f_{\tau, *}\epsilon_X^{-1}K'
$$
と比較すれば、$K'$ が $\mathcal{A}'_X$ に属するただ一つの非零コホモロジー層
$\mathcal{F}'$ をもつ場合に帰着する；詳細は省略する。
この場合、補題 \ref{sites-cohomology-lemma-V-implies-C-general} により、任意の $i$ に対して
$\epsilon_Y^{-1}R^if_{\tau', *}\mathcal{F}' =
R^if_{\tau, *}\epsilon_X^{-1}\mathcal{F}'$ であり、証明が完了する。
\end{proof}

\begin{lemma}
\label{sites-cohomology-lemma-cohomological-descent-general}
状況 \ref{sites-cohomology-situation-compare} のもとで、任意の $X \in
\mathcal{C}$ に対して圏
$\mathcal{A}_X \subset \textit{Ab}(\mathcal{C}_\tau/X)$
は弱 Serre 部分圏であり、関手
$$
R\epsilon_{X, *} :
D^+_{\mathcal{A}_X}(\mathcal{C}_\tau/X)
\longrightarrow
D^+_{\mathcal{A}'_X}(\mathcal{C}_{\tau'}/X)
$$
は圏同値で、その擬逆は $\epsilon_X^{-1}$ で与えられる。
\end{lemma}

\begin{proof}
Homology, Lemma \ref{homology-lemma-characterize-weak-serre-subcategory}
に挙げられた条件を $\mathcal{A}_X$ について確認する必要がある。
$\varphi : \mathcal{F} \to \mathcal{G}$ が $\mathcal{A}_X$ の射ならば、
$\epsilon_{X, *}\varphi :
\epsilon_{X, *}\mathcal{F} \to \epsilon_{X, *}\mathcal{G}$
は $\mathcal{A}'_X$ の射である。$\mathcal{A}'_X$ は
$\textit{Ab}(\mathcal{C}_{\tau'}/X)$ の弱 Serre 部分圏なので、
$\Ker(\epsilon_{X, *}\varphi)$ と $\Coker(\epsilon_{X, *}\varphi)$ は
$\mathcal{A}'_X$ の対象である。$\epsilon_X^{-1}$ を適用すると完全列
$$
0 \to
\epsilon_X^{-1}\Ker(\epsilon_{X, *}\varphi) \to
\mathcal{F} \to \mathcal{G} \to
\epsilon_X^{-1}\Coker(\epsilon_{X, *}\varphi) \to 0
$$
を得るので、$\Ker(\varphi)$ と $\Coker(\varphi)$ は $\mathcal{A}_X$ に属する。
最後に、$\textit{Ab}(\mathcal{C}_\tau/X)$ における短完全列
$$
0 \to \mathcal{F}_1 \to \mathcal{F}_2 \to \mathcal{F}_3 \to 0
$$
を考え、$\mathcal{F}_1$ と $\mathcal{F}_3$ は $\mathcal{A}_X$ に属するとする。
$\epsilon_{X, *}$ を適用すると完全列
$$
0 \to 
\epsilon_{X, *}\mathcal{F}_1 \to
\epsilon_{X, *}\mathcal{F}_2 \to
\epsilon_{X, *}\mathcal{F}_3 \to
R^1\epsilon_{X, *}\mathcal{F}_1 = 0
$$
を得る；最後の消滅は補題 \ref{sites-cohomology-lemma-V-C-all-n-general} による。
$\mathcal{A}'_X$ は $\textit{Ab}(\mathcal{C}_{\tau'}/X)$ の弱 Serre 部分圏なので、
$\epsilon_{X, *}\mathcal{F}_2$ は $\mathcal{A}'_X$ に属する。
$\epsilon_X$ によって引き戻すと、$\mathcal{F}_2$ は
$\mathcal{A}_X$ に属すると結論できる。

\medskip\noindent
以上により $\mathcal{A}_X$ は $\textit{Ab}(\mathcal{C}_\tau/X)$ の
弱 Serre 部分圏であり、圏 $D^+_{\mathcal{A}_X}(\mathcal{C}_\tau/X)$
を考えることができる。$\epsilon_X^{-1} : \mathcal{A}'_X \to \mathcal{A}_X$
は圏同値であり、
$\mathcal{F}' \to R\epsilon_{X, *}\epsilon_X^{-1}\mathcal{F}'$
は任意の $\mathcal{F}' \in \mathcal{A}'_X$ に対して同型であることに注意する。
実際、補題 \ref{sites-cohomology-lemma-V-C-all-n-general} によりすべての $n$ について
$(V_n)$ が成り立つ。したがって補題 \ref{sites-cohomology-lemma-equivalence-bounded} により結論を得る。
\end{proof}

\begin{lemma}
\label{sites-cohomology-lemma-compare-cohomology-general}
状況 \ref{sites-cohomology-situation-compare} のもとで $X \in \mathcal{C}$ とする。
\begin{enumerate}
\item $\mathcal{F}' \in \mathcal{A}'_X$ に対して
$H^n_{\tau'}(X, \mathcal{F}') = H^n_\tau(X, \epsilon_X^{-1}\mathcal{F}')$。
\item $K' \in D^+_{\mathcal{A}'_X}(\mathcal{C}_{\tau'}/X)$ に対して
$H^n_{\tau'}(X, K') = H^n_\tau(X, \epsilon_X^{-1}K')$。
\end{enumerate}
\end{lemma}

\begin{proof}
補題 \ref{sites-cohomology-lemma-V-C-all-n-general} と
注意 \ref{sites-cohomology-remark-before-Leray} から従う。
\end{proof}
























\section{ハウスドルフかつ局所準コンパクトな空間上のコホモロジー}
\label{sites-cohomology-section-cohomology-LC}

\noindent
文献ではしばしば「局所コンパクト」と言うところを、本書では引き続き
「ハウスドルフかつ局所準コンパクト」と言うことにする。
$\textit{LC}$ を、対象がハウスドルフかつ局所準コンパクトな位相空間で、
射が連続写像である圏とする。

\begin{lemma}
\label{sites-cohomology-lemma-LC-basic}
圏 $\textit{LC}$ はファイバー積と終対象をもち、したがって任意の有限極限をもつ。
$\textit{LC}$ の射 $X \to Z$ と $Y \to Z$ が与えられ、
$X$ と $Y$ が準コンパクトならば、$X \times_Z Y$ は準コンパクトである。
\end{lemma}

\begin{proof}
終対象は一点空間である。$\textit{LC}$ の射 $X \to Z$ と
$Y \to Z$ が与えられたとき、ファイバー積 $X \times_Z Y$ は
$X \times Y$ の部分空間である。Topology, Section
\ref{topology-section-Hausdorff} により $X \times Y$ はハウスドルフなので、
$X \times_Z Y$ もハウスドルフである。

\medskip\noindent
$X$ と $Y$ が準コンパクトならば、Topology, Theorem
\ref{topology-theorem-tychonov} により $X \times Y$ は準コンパクトである。
$X \times_Z Y$ は $X \times Y$ の閉部分集合なので
（Topology, Lemma \ref{topology-lemma-fibre-product-closed}）、
Topology, Lemma \ref{topology-lemma-closed-in-quasi-compact} により
$X \times_Z Y$ は準コンパクトである。

\medskip\noindent
最後に一般の場合へ戻る。$x \in X$ かつ $y \in Y$ ならば、
準コンパクト近傍 $x \in E \subset X$ と $y \in F \subset Y$ を選べる。
上の結果により $E \times_Z F$ は $(x, y)$ の準コンパクト近傍である。
したがって Topology, Lemma
\ref{topology-lemma-locally-quasi-compact-Hausdorff} により、
$X \times_Z Y$ は $\textit{LC}$ の対象である。
\end{proof}

\noindent
$\textit{LC}$ には通常のものより強い位相を入れることができる。

\begin{definition}
\label{sites-cohomology-definition-covering-LC}
$\{f_i : X_i \to X\}$ を、圏 $\textit{LC}$ において終域を固定した射の族とする。
任意の $x \in X$ に対して $i_1, \ldots, i_n \in I$ と準コンパクト部分集合
$E_j \subset X_{i_j}$ が存在し、$\bigcup f_{i_j}(E_j)$ が $x$ の近傍となるとき、
この族を {\it qc 被覆}\footnote{これは標準的な記法ではない。
スキームの fpqc 被覆を想起させるためにこの記法を選んだ。}という。
\end{definition}

\noindent
$\textit{LC}$ の対象の開被覆 $X = \bigcup U_i$ は、$X$ が局所準コンパクトなので
qc 被覆 $\{U_i \to X\}$ を与えることに注意する。
まず、お決まりの補題から始める。

\begin{lemma}
\label{sites-cohomology-lemma-qc}
$X$ をハウスドルフかつ局所準コンパクトな空間、すなわち
$\textit{LC}$ の対象とする。
\begin{enumerate}
\item $X' \to X$ が $\textit{LC}$ における同型ならば、
$\{X' \to X\}$ は qc 被覆である。
\item $\{f_i : X_i \to X\}_{i\in I}$ が qc 被覆で、各 $i$ に対して
$\{g_{ij} : X_{ij} \to X_i\}_{j\in J_i}$ が qc 被覆ならば、
$\{X_{ij} \to X\}_{i \in I, j\in J_i}$ は qc 被覆である。
\item $\{X_i \to X\}_{i\in I}$ が qc 被覆で、
$X' \to X$ が $\textit{LC}$ の射ならば、
$\{X' \times_X X_i \to X'\}_{i\in I}$ は qc 被覆である。
\end{enumerate}
\end{lemma}

\begin{proof}
（1）は、開被覆が qc 被覆であるという上の注意から従う。

\medskip\noindent
（2）の証明。$x \in X$ とする。$i_1, \ldots, i_n \in I$ と準コンパクトな
$E_a \subset X_{i_a}$ を、$\bigcup f_{i_a}(E_a)$ が $x$ の近傍となるように選ぶ。
各 $e \in E_a$ に対し、有限部分集合 $J_e \subset J_{i_a}$ と準コンパクトな
$F_{e, j} \subset X_{ij}$（$j \in J_e$）で、$\bigcup g_{ij}(F_{e, j})$ が
$e$ の近傍となるものを選べる。$E_a$ は準コンパクトなので、有限個の
$e_1, \ldots, e_{m_a}$ を選んで
$$
E_a \subset
\bigcup\nolimits_{k = 1, \ldots, m_a}
\bigcup\nolimits_{j \in J_{e_k}} g_{ij}(F_{e_k, j})
$$
とできる。このとき
$$
\bigcup\nolimits_{a = 1, \ldots, n}
\bigcup\nolimits_{k = 1, \ldots, m_a}
\bigcup\nolimits_{j \in J_{e_k}} f_i(g_{ij}(F_{e_k, j}))
$$
は $x$ の近傍である。

\medskip\noindent
（3）の証明。点 $x' \in X'$ をとり、その像を $x \in X$ とする。
$i_1, \ldots, i_n \in I$ と準コンパクト部分集合
$E_j \subset X_{i_j}$ を、$\bigcup f_{i_j}(E_j)$ が $x$ の近傍となるように選ぶ。
$x'$ の準コンパクト近傍 $F \subset X'$ を、その像が $x$ の準コンパクト近傍
$\bigcup f_{i_j}(E_j)$ に含まれるように選ぶ。このとき
$F \times_X E_j \subset X' \times_X X_{i_j}$ は準コンパクト部分集合であり、
$F$ は写像 $\coprod F \times_X E_j \to F$ の像である。
したがって基底変換は qc 被覆であり、証明が完了する。
\end{proof}

\noindent
$\textit{LC}$ のすべての対象はハウスドルフなので、$\textit{LC}$ の任意の射
$f : X \to Y$ は位相空間の分離的な連続写像である。したがって $f$ が
位相空間の固有写像であることと、$f$ が普遍閉であることは同値である。
Topology, Section \ref{topology-section-proper} の議論を参照せよ。

\begin{lemma}
\label{sites-cohomology-lemma-proper-surjective-is-qc-covering}
$f : X \to Y$ を $\textit{LC}$ の射とする。
$f$ が固有かつ全射ならば、$\{f : X \to Y\}$ は qc 被覆である。
\end{lemma}

\begin{proof}
点 $y \in Y$ をとる。各 $x \in X_y$ に対して準コンパクト近傍
$E_x \subset X$ を選び、開集合 $x \in U_x \subset E_x$ を選ぶ。
$f$ は固有なのでファイバー $X_y$ は準コンパクトであり、
$x_1, \ldots, x_n \in X_y$ を
$X_y \subset U_{x_1} \cup \ldots \cup U_{x_n}$ となるように選べる。
$f(E_{x_1}) \cup \ldots \cup f(E_{x_n})$ が $y$ の近傍であることを示す。
実際、$f$ は閉写像なので
（Topology, Theorem \ref{topology-theorem-characterize-proper}）、
$Z = f(X \setminus U_{x_1} \cup \ldots \cup U_{x_n})$ は
$y$ を含まない $Y$ の閉部分集合である。$f$ は全射なので、所望のとおり
$Y \setminus Z$ は $f(E_{x_1}) \cup \ldots \cup f(E_{x_n})$ に含まれる。
\end{proof}

\noindent
いくつかの集合論的問題を除けば、補題 \ref{sites-cohomology-lemma-qc} は、
$\textit{LC}$ が qc 被覆の族とともにサイトをなすことを示している。
このサイトを、集合論的問題を避けるため適切に修正したうえで
$\textit{LC}_{qc}$ と表す。

\begin{remark}[集合論的問題]
\label{sites-cohomology-remark-set-theoretic-LC}
圏 $\textit{LC}$ の対象は真の類をなすので、これは「大きな」圏である。
同様に、被覆も真の類をなす。位相空間 $X$ の {\it サイズ}を、
$X$ の点の集合の濃度として定める。たとえば Sets, Equation
(\ref{sets-equation-bound}) のように、濃度上の関数 $Bound$ を選ぶ。
さらに、$S_0$ を $\textit{LC}$ の対象からなる初期集合、たとえば
$S_0 = \{(\mathbf{R}, \text{ユークリッド位相})\}$ とする。
Sets, Lemma \ref{sets-lemma-construct-category} とまったく同様に、
$\textit{LC}_\alpha = \textit{LC} \cap V_\alpha$
が $S_0$ を含み、$\textit{LC}$ に存在するすべての可算極限と可算余極限のもとで
閉じるような極限順序数 $\alpha$ を選べる。さらに、
$X \in \textit{LC}_\alpha$、$Y \in \textit{LC}$ で
$\text{size}(Y) \leq Bound(\text{size}(X))$ ならば、$Y$ は
$\textit{LC}_\alpha$ の対象と同型である。次に Sets, Lemma
\ref{sets-lemma-coverings-site} を適用して、$\textit{LC}_\alpha$ 上の
qc 被覆からなる集合 $\text{Cov}$ を、$\textit{LC}_\alpha$ の各 qc 被覆が
この集合の被覆と組合せ論的に同値となるように選ぶ。
こうしてサイト $(\textit{LC}_\alpha, \text{Cov})$ を得て、これを
$\textit{LC}_{qc}$ と表す。
\end{remark}

\noindent
注意 \ref{sites-cohomology-remark-set-theoretic-LC} のサイト $\textit{LC}_{qc}$ には
第二の位相がある。すなわち、対象 $X$ に対し、各 $X_i \to X$ が
開埋め込みであるような $\textit{LC}_{qc}$ の被覆 $\{X_i \to X\}$ を
すべて考える。このサイトを $\textit{LC}_{Zar}$ と表す。恒等関手
$\textit{LC}_{Zar} \to \textit{LC}_{qc}$ は連続であり、サイトの射
$$
\epsilon : \textit{LC}_{qc} \longrightarrow \textit{LC}_{Zar}
$$
を定める。Section \ref{sites-cohomology-section-compare} を参照せよ。
ハウスドルフかつ局所準コンパクトな位相空間 $X$、より正確には
$X \in \Ob(\textit{LC}_{qc})$ に対し、誘導される射を
$$
\epsilon_X : \textit{LC}_{qc}/X \longrightarrow \textit{LC}_{Zar}/X
$$
と表す（Sites, Lemma \ref{sites-lemma-localize-morphism} を参照）。
$X_{Zar}$ を、対象が $X$ の開集合であるサイトとする；Sites, Example
\ref{sites-example-site-topological} を参照せよ。サイトの射
$$
\pi_X : \textit{LC}_{Zar}/X \longrightarrow X_{Zar}
$$
があり、これは連続関手
$X_{Zar} \to \textit{LC}_{Zar}/X$, $U \mapsto U$
によって与えられる。実際、$X_{Zar}$ はファイバー積と終対象をもち、
上の関手はそれらと可換なので、Sites, Proposition
\ref{sites-proposition-get-morphism} を適用できる。等式
$\Sh(X_{Zar}) = \Sh(X)$ を用いて、しばしば $\pi$ をトポスの射
$$
\pi_X : \Sh(\textit{LC}_{Zar}/X) \longrightarrow \Sh(X)
$$
とみなす。

\begin{lemma}
\label{sites-cohomology-lemma-describe-pullback-pi}
$X$ を $\textit{LC}_{qc}$ の対象とし、$\mathcal{F}$ を $X$ 上の層とする。
対応
$$
\textit{LC}_{qc}/X \longrightarrow \textit{Sets},\quad
(f : Y \to X) \longmapsto \Gamma(Y, f^{-1}\mathcal{F})
$$
は層であり、したがって特に $\textit{LC}_{Zar}/X$ 上の層でもある。
この層は、$\textit{LC}_{Zar}/X$ 上では $\pi_X^{-1}\mathcal{F}$ に等しく、
$\textit{LC}_{qc}/X$ 上では $\epsilon_X^{-1}\pi_X^{-1}\mathcal{F}$ に等しい。
\end{lemma}

\begin{proof}
補題の式で与えられる前層を $\mathcal{G}$ と表す。もちろん、式中の
引き戻し $f^{-1}$ は位相空間上の層の通常の引き戻しを意味する。
定義から直ちに、$\mathcal{G}$ は Zar 位相に関する層である。

\medskip\noindent
$Y \to X$ を $\textit{LC}_{qc}$ の射とし、
$\mathcal{V} = \{g_i : Y_i \to Y\}_{i \in I}$ を qc 被覆とする。
$\mathcal{G}$ が qc 位相に関する層であることを示すには、写像
$\mathcal{G}(Y) \to H^0(\mathcal{V}, \mathcal{G})$
が同型であることを示せば十分である；Sites, Section
\ref{sites-section-sheafification} を参照せよ。まず、この写像は単射である。
実際、qc 被覆は全射であり、切片の等しさは茎で検出できる
（Sheaves, Lemmas \ref{sheaves-lemma-sheaf-subset-stalks} および
\ref{sheaves-lemma-stalk-pullback-presheaf} を用いよ）。したがって
$\mathcal{G}$ は $\textit{LC}_{qc}$ 上の分離前層であるから、任意の元
$(s_i) \in H^0(\mathcal{V}, \mathcal{G})$
が、$\mathcal{V}$ を細分で置き換えた後に $\mathcal{G}(Y)$ の像に入ることを
示せば十分である（Sites, Theorem \ref{sites-theorem-plus}）。

\medskip\noindent
$Y_{i, Zar}$ 上の層と $Y_i$ 上の層を同一視すると、
$\mathcal{G}|_{Y_{i, Zar}}$ は連続写像 $g_i : Y_i \to Y$ による
$f^{-1}\mathcal{F}$ の引き戻しである。したがって開被覆
$Y_i = \bigcup V_{ij}$ を、各 $j$ に対して開集合 $W_{ij} \subset Y$ と
切片 $t_{ij} \in \mathcal{G}(W_{ij})$ が存在し、$V_{ij}$ が $W_{ij}$ に写り、
$s|_{V_{ij}}$ が $t_{ij}$ の引き戻しとなるように選べる。言い換えると、
被覆 $\{Y_i \to Y\}$ を細分した後、開集合 $W_i \subset Y$ があって
$Y_i \to Y$ が $W_i$ を経由し、$W_i$ 上の $\mathcal{G}$ の切片 $t_i$ が
与えられた切片 $s_i$ に制限されると仮定してよい。さらに、$y \in Y$ が
$Y_i \to Y$ と $Y_j \to Y$ の双方の像に属するならば、茎
$f^{-1}\mathcal{F}_y$ における像 $t_{i, y}$ と $t_{j, y}$ は一致する
（$s_i$ と $s_j$ が $Y_i \times_Y Y_j$ 上で一致するためである）。
したがって各 $y \in Y$ に対し、$f^{-1}\mathcal{F}_y$ の良定義な元 $t_y$ があり、
$y$ が $Y_i \to Y$ の像に属するとき $t_{i,y}$ と一致する。
元 $(t_y)$ が $Y$ 上の $f^{-1}\mathcal{F}$ の大域切片から来ることを示せば、
補題の証明が完了する。

\medskip\noindent
これを $Y$ 上で局所的に示せば十分である；Sheaves, Section
\ref{sheaves-section-sheafification} を参照せよ。$y_0 \in Y$ とする。
$i_1, \ldots, i_n \in I$ と準コンパクト部分集合 $E_j \subset Y_{i_j}$ を、
$\bigcup g_{i_j}(E_j)$ が $y_0$ の近傍となるように選ぶ。
$V \subset Y$ を、$\bigcup g_{i_j}(E_j)$ と
$W_{i_1} \cap \ldots \cap W_{i_n}$ に含まれる $y_0$ の開近傍とする。
$t_{i_1, y_0} = \ldots = t_{i_n, y_0}$ なので、$V$ を縮小すれば、
$f^{-1}\mathcal{F}$ の切片 $t_{i_j}|_V$（$j = 1, \ldots, n$）は
一致すると仮定できる。$V \subset \bigcup g_{i_j}(E_j)$ だから、
$(t_y)_{y \in V}$ はこの切片から来る。

\medskip\noindent
$\mathcal{G}$ が $\textit{LC}_{qc}$ 上で
$\epsilon_X^{-1}\pi_X^{-1}\mathcal{F}$ に、また $\textit{LC}_{Zar}$ 上で
$\pi_X^{-1}\mathcal{F}$ に等しいことがまだ残っている。いずれの場合も、
引き戻しは前層
$$
(f : Y \to X)
\longmapsto
\colim_{f(Y) \subset U \subset X} \mathcal{F}(U)
$$
をとってから層化することで定義される。Zar 位相で層化するとちょうど
層 $\mathcal{G}$ が得られ、$\mathcal{G}$ が qc 層であることから、
qc 位相でも同様に成り立つ。
\end{proof}

\noindent
$X \in \Ob(\textit{LC}_{Zar})$ とし、$\mathcal{H}$ を
$\textit{LC}_{Zar}/X$ 上のアーベル層とする。このとき、
$\textit{LC}_{Zar}/X$ の対象 $U$ 上での $\mathcal{H}$ のコホモロジーを
$H^n_{Zar}(U, \mathcal{H})$ と書く。

\begin{lemma}
\label{sites-cohomology-lemma-collect-true-things-Zar}
$X$ を $\textit{LC}_{Zar}$ の対象とする。このとき、
\begin{enumerate}
\item $\mathcal{F} \in \textit{Ab}(X)$ に対して
$H^n_{Zar}(X, \pi_X^{-1}\mathcal{F}) = H^n(X, \mathcal{F})$。
\item $\pi_{X, *} : \textit{Ab}(\textit{LC}_{Zar}/X) \to \textit{Ab}(X)$
は完全である。
\item 随伴の単位 $\text{id} \to \pi_{X, *} \circ \pi_X^{-1}$ は同型である。
\item $K \in D(X)$ に対して標準写像
$K \to R\pi_{X, *} \pi_X^{-1}K$ は同型である。
\end{enumerate}
$f : X \to Y$ を $\textit{LC}_{Zar}$ の射とする。このとき、
\begin{enumerate}
\item[(5)] トポスの可換図式
$$
\xymatrix{
\Sh(\textit{LC}_{Zar}/X) \ar[r]_{f_{Zar}} \ar[d]_{\pi_X} &
\Sh(\textit{LC}_{Zar}/Y) \ar[d]^{\pi_Y} \\
\Sh(X_{Zar}) \ar[r]^f &
\Sh(Y_{Zar})
}
$$
がある。
\item[(6)] $L \in D^+(Y)$ に対して
$H^n_{Zar}(X, \pi_Y^{-1}L) = H^n(X, f^{-1}L)$。
\item[(7)] $f$ が固有ならば、次が成り立つ。
\begin{enumerate}
\item $\pi_Y^{-1} \circ f_* = f_{Zar, *} \circ \pi_X^{-1}$ は
$\Sh(X) \to \Sh(\textit{LC}_{Zar}/Y)$ という関手の同型である。
\item $\pi_Y^{-1} \circ Rf_* = Rf_{Zar, *} \circ \pi_X^{-1}$ は
$D^+(X) \to D^+(\textit{LC}_{Zar}/Y)$ という関手の同型である。
\end{enumerate}
\end{enumerate}
\end{lemma}

\begin{proof}
（1）の証明。等式
$H^n_{Zar}(X, \pi_X^{-1}\mathcal{F}) = H^n(X, \mathcal{F})$ は、
$\textit{LC}_{Zar}$ における $X$ の被覆が $X$ の開被覆と同じであるという
自明な観察から従う一般的事実である。詳しい証明を見たい読者は、関手
$X_{Zar} \to \textit{LC}_{Zar}$ に補題 \ref{sites-cohomology-lemma-cohomology-bigger-site}
を適用すればよい。

\medskip\noindent
（2）の証明。$\pi_X$ の定義に用いた関手
$X_{Zar} \to \textit{LC}_{Zar}/X$ に Sites, Lemma
\ref{sites-lemma-bigger-site} を適用すると、あるトポスの射
$\tau_X : \Sh(X_{Zar}) \to \Sh(\textit{LC}_{Zar})$ に対して
$\pi_{X, *} = \tau_X^{-1}$ となるので、主張が従う。

\medskip\noindent
（3）の証明。Sites, Lemma \ref{sites-lemma-bigger-site} により
$\tau_X^{-1} \circ \pi_X^{-1}$ は恒等関手なので、主張が従う。
あるいは、補題 \ref{sites-cohomology-lemma-describe-pullback-pi} における
$\pi_X^{-1}$ の明示的記述からも導ける。

\medskip\noindent
（4）の証明。$K$ を表すアーベル層の複体に（3）を適用する。

\medskip\noindent
（5）の証明。トポスの射 $f_{Zar}$ は Sites, Lemma
\ref{sites-lemma-relocalize} の適用から得られ、この場合には Sites, Lemma
\ref{sites-lemma-relocalize-given-fibre-products} により連続関手
$Z/Y \mapsto Z \times_Y X/X$ から得られる。対応する連続関手が正しく合成されるので
図式は可換である（Sites, Lemma
\ref{sites-lemma-composition-morphisms-sites} を参照）。

\medskip\noindent
（6）の証明。$\mathcal{G} \in \textit{Ab}(Y)$ に対して
$H^n_{Zar}(X, \pi_Y^{-1}\mathcal{G}) =
H^n_{Zar}(X, f_{Zar}^{-1}\pi_Y^{-1}\mathcal{G})$
である；補題 \ref{sites-cohomology-lemma-cohomology-of-open} を参照せよ。
（5）の図式の可換性により、これは
$H^n_{Zar}(X, \pi_X^{-1}f^{-1}\mathcal{G})$ に等しい。
したがって $L$ が次数 $0$ の一つの層からなる場合は（1）から結論を得る。
一般の場合は、$L$ を下に有界なアーベル層の複体で表せば従う。

\medskip\noindent
（7a）の証明。$\mathcal{F}$ を $X$ 上の層とする。
$g : Z \to Y$ を $\textit{LC}_{Zar}/Y$ の対象とし、ファイバー積
$$
\xymatrix{
Z' \ar[r]_{f'} \ar[d]_{g'} & Z \ar[d]^g \\
X \ar[r]^f & Y
}
$$
を考える。このとき
$$
(f_{Zar, *}\pi_X^{-1}\mathcal{F})(Z/Y) =
(\pi_X^{-1}\mathcal{F})(Z'/X)  =
\Gamma(Z', (g')^{-1}\mathcal{F})  =
\Gamma(Z, f'_*(g')^{-1}\mathcal{F})
$$
であり、第二の等号は補題 \ref{sites-cohomology-lemma-describe-pullback-pi} による。一方、
$$
(\pi_Y^{-1}f_*\mathcal{F})(Z/Y) = \Gamma(Z, g^{-1}f_*\mathcal{F})
$$
も補題 \ref{sites-cohomology-lemma-describe-pullback-pi} による。したがって集合の層に対する
固有基底変換（Cohomology, Lemma
\ref{cohomology-lemma-proper-base-change-sheaves-of-sets}）により、
二つの集合は標準的に同型である。この同型は制限写像と両立し、同型
$\pi_Y^{-1}f_*\mathcal{F} = f_{Zar, *}\pi_X^{-1}\mathcal{F}$
を定める。したがって関手の同型
$\pi_Y^{-1} \circ f_* = f_{Zar, *} \circ \pi_X^{-1}$ を得る。

\medskip\noindent
（7b）の証明。$K \in D^+(X)$ とする。補題
\ref{sites-cohomology-lemma-unbounded-describe-higher-direct-images} により、
$Rf_{Zar, *}\pi_X^{-1}K$ の第 $n$ コホモロジー層は前層
$$
(g : Z \to Y) \longmapsto H^n_{Zar}(Z', \pi_X^{-1}K)
$$
に付随する層である。記法は上と同じである。次の等式に注意する。
\begin{align*}
H^n_{Zar}(Z', \pi_X^{-1}K)
& =
H^n(Z', (g')^{-1}K) \\
& =
H^n(Z, Rf'_*(g')^{-1}K) \\
& =
H^n(Z, g^{-1}Rf_*K) \\
& =
H^n_{Zar}(Z, \pi_Y^{-1}Rf_*K)
\end{align*}
第一の等号は $K$ と $g' : Z' \to X$ に（6）を適用したものである。
第二の等号は $f' : Z' \to Z$ に対する Leray
（Cohomology, Lemma \ref{cohomology-lemma-before-Leray}）である。
第三の等号は固有基底変換定理（Cohomology, Theorem
\ref{cohomology-theorem-proper-base-change}）である。
第四の等号は $g : Z \to Y$ と $Rf_*K$ に（6）を適用したものである。
したがって $Rf_{Zar, *}\pi_X^{-1}K$ と $\pi_Y^{-1}Rf_*K$ は
同じコホモロジー層をもつ。標準基底変換写像
$\pi_Y^{-1}Rf_*K \to Rf_{Zar, *}\pi_X^{-1}K$ がこの同型を誘導することの
確認は省略する。
\end{proof}

\noindent
補題 \ref{sites-cohomology-lemma-describe-pullback-pi} の状況において、$\epsilon$ と $\pi$ の合成および
等式 $\Sh(X) = \Sh(X_{Zar})$ はトポスの射
$$
a_X : \Sh(\textit{LC}_{qc}/X) \longrightarrow \Sh(X)
$$
を定める。

\begin{lemma}
\label{sites-cohomology-lemma-push-pull-LC}
$f : X \to Y$ を $\textit{LC}_{qc}$ の射とする。このときトポスの可換図式
$$
\vcenter{
\xymatrix{
\Sh(\textit{LC}_{qc}/X) \ar[r]_{f_{qc}} \ar[d]_{\epsilon_X} &
\Sh(\textit{LC}_{qc}/Y) \ar[d]^{\epsilon_Y} \\
\Sh(\textit{LC}_{Zar}/X) \ar[r]^{f_{Zar}} &
\Sh(\textit{LC}_{Zar}/Y)
}
}
\quad\text{および}\quad
\vcenter{
\xymatrix{
\Sh(\textit{LC}_{qc}/X) \ar[r]_{f_{qc}} \ar[d]_{a_X} &
\Sh(\textit{LC}_{qc}/Y) \ar[d]^{a_Y} \\
\Sh(X) \ar[r]^f &
\Sh(Y)
}
}
$$
があり、$a_X = \pi_X \circ \epsilon_X$, $a_Y = \pi_X \circ \epsilon_X$ である。
$f$ が固有ならば、$a_Y^{-1} \circ f_* = f_{qc, *} \circ a_X^{-1}$ である。
\end{lemma}

\begin{proof}
トポスの射 $f_{qc}$ は Sites, Lemma \ref{sites-lemma-relocalize} のものであり、
この場合には連続関手 $Z/Y \mapsto Z \times_Y X/X$ から得られる；Sites, Lemma
\ref{sites-lemma-relocalize-given-fibre-products} を参照せよ。
対応する連続関手が正しく合成されるので左の図式は可換である
（Sites, Lemma \ref{sites-lemma-composition-morphisms-sites} を参照）。
右の図式が可換であることは、左の図式の可換性と補題
\ref{sites-cohomology-lemma-collect-true-things-Zar} の（5）から従う。

\medskip\noindent
最後の主張の証明。補題 \ref{sites-cohomology-lemma-collect-true-things-Zar} の（7a）の証明を
繰り返してもよいが、ここではそれから主張を導く。
$\epsilon_{Y, *}$ は台となる前層上で恒等関手なので、同型を反映する。
補題 \ref{sites-cohomology-lemma-describe-pullback-pi} の記述により
$\epsilon_{Y, *} \circ a_Y^{-1} = \pi_Y^{-1}$ であり、$X$ についても同様である。
標準写像 $a_Y^{-1}f_*\mathcal{F} \to f_{qc, *}a_X^{-1}\mathcal{F}$ が
同型であることを示すには、写像
$$
\pi_Y^{-1}f_*\mathcal{F} =
\epsilon_{Y, *}a_Y^{-1}f_*\mathcal{F} \to
\epsilon_{Y, *}f_{qc, *}a_X^{-1}\mathcal{F} =
f_{Zar, *}\epsilon_{X, *}a_X^{-1}\mathcal{F} =
f_{Zar, *}\pi_X^{-1}\mathcal{F}
$$
が同型であることを示せば十分である。これは補題
\ref{sites-cohomology-lemma-collect-true-things-Zar} の（7a）である。
\end{proof}

\begin{lemma}
\label{sites-cohomology-lemma-compare-qc-zar}
比較射 $\epsilon : \textit{LC}_{qc} \to \textit{LC}_{Zar}$ を考える。
$\mathcal{P}$ を位相空間の固有写像の類とする。
$X \in \textit{LC}_{Zar}$ に対して、
$\mathcal{A}'_X \subset \textit{Ab}(\textit{LC}_{Zar}/X)$
を、$\mathcal{F} \in \textit{Ab}(X)$ に対する形 $\pi_X^{-1}\mathcal{F}$ の
層からなる充満部分圏とする。このとき、
(\ref{sites-cohomology-item-base-change-P}),
(\ref{sites-cohomology-item-restriction-A}),
(\ref{sites-cohomology-item-A-sheaf}),
(\ref{sites-cohomology-item-A-and-P}), および
(\ref{sites-cohomology-item-refine-tau-by-P})
は状況 \ref{sites-cohomology-situation-compare} において成り立つ。
\end{lemma}

\begin{proof}
まず Homology, Lemma \ref{homology-lemma-characterize-weak-serre-subcategory} の
条件（1）、（2）、（3）、（4）を確認して、
$\mathcal{A}'_X \subset \textit{Ab}(\textit{LC}_{Zar}/X)$ が
弱 Serre 部分圏であることを示す。補題 \ref{sites-cohomology-lemma-collect-true-things-Zar} の
（3）により $\pi_X^{-1}$ は完全かつ充満忠実なので、（1）、（2）、（3）は直ちに従う。
短完全列
$0 \to \pi_X^{-1}\mathcal{F} \to \mathcal{G} \to \pi_X^{-1}\mathcal{F}' \to 0$
が $\textit{Ab}(\textit{LC}_{Zar}/X)$ にあるならば、補題
\ref{sites-cohomology-lemma-collect-true-things-Zar} の（2）により
$0 \to \mathcal{F} \to \pi_{X, *}\mathcal{G} \to \mathcal{F}' \to 0$
は完全である。したがって
$\mathcal{G} = \pi_X^{-1}\pi_{X, *}\mathcal{G}$ は $\mathcal{A}'_X$ に属し、
最後の条件も確認できた。

\medskip\noindent
性質（\ref{sites-cohomology-item-base-change-P}）は、補題 \ref{sites-cohomology-lemma-LC-basic} と、
固有写像の基底変換が固有であるという事実から従う
（Topology, Theorem \ref{topology-theorem-characterize-proper} および
Lemma \ref{topology-lemma-base-change-separated} を参照）。

\medskip\noindent
性質（\ref{sites-cohomology-item-restriction-A}）は補題
\ref{sites-cohomology-lemma-collect-true-things-Zar} の可換図式（5）から従う。

\medskip\noindent
性質（\ref{sites-cohomology-item-A-sheaf}）は補題 \ref{sites-cohomology-lemma-describe-pullback-pi} である。

\medskip\noindent
性質（\ref{sites-cohomology-item-A-and-P}）は補題 \ref{sites-cohomology-lemma-collect-true-things-Zar} の
（7）（b）である。

\medskip\noindent
（\ref{sites-cohomology-item-refine-tau-by-P}）の証明。qc 被覆 $\{U_i \to U\}$ が
与えられたとする。$u \in U$ に対し、$i_1, \ldots, i_m \in I$ と
準コンパクト部分集合 $E_j \subset U_{i_j}$ を、
$\bigcup f_{i_j}(E_j)$ が $u$ の近傍となるように選ぶ。
$Y = \coprod_{j = 1, \ldots, m} E_j \to U$ は、ハウスドルフ準コンパクト空間の間の
連続写像として固有であることに注意する
（Topology, Lemma \ref{topology-lemma-closed-map}）。
$\bigcup f_{i_j}(E_j)$ に含まれる開近傍 $u \in V$ を選ぶ。
このとき $Y \times_U V \to V$ は全射な固有射であり、補題
\ref{sites-cohomology-lemma-proper-surjective-is-qc-covering} により $qc$ 被覆である。
これをすべての $u \in U$ に対して行えるので、
（\ref{sites-cohomology-item-refine-tau-by-P}）が成り立つ。
\end{proof}

\begin{lemma}
\label{sites-cohomology-lemma-V-C-all-n}
記法は上のとおりとする。
\begin{enumerate}
\item $X \in \Ob(\textit{LC}_{qc})$ と $X$ 上のアーベル層 $\mathcal{F}$ に対して
$\epsilon_{X, *}a_X^{-1}\mathcal{F} = \pi_X^{-1}\mathcal{F}$ であり、
$i > 0$ に対して $R^i\epsilon_{X, *}(a_X^{-1}\mathcal{F}) = 0$ である。
\item $\textit{LC}_{qc}$ の固有射 $f : X \to Y$ と $X$ 上のアーベル層
$\mathcal{F}$ に対して、すべての $i$ について
$a_Y^{-1}(R^if_*\mathcal{F}) = R^if_{qc, *}(a_X^{-1}\mathcal{F})$
である。
\item $X \in \Ob(\textit{LC}_{qc})$ と $K \in D^+(X)$ に対して、写像
$\pi_X^{-1}K \to R\epsilon_{X, *}(a_X^{-1}K)$ は同型である。
\item $\textit{LC}_{qc}$ の固有射 $f : X \to Y$ と $K \in D^+(X)$ に対して
$a_Y^{-1}(Rf_*K) = Rf_{qc, *}(a_X^{-1}K)$ である。
\end{enumerate}
\end{lemma}

\begin{proof}
補題 \ref{sites-cohomology-lemma-compare-qc-zar} により、Section
\ref{sites-cohomology-section-compare-general} のすべての補題を現在の設定に適用できる。
その結果を読み替えるため、補題 \ref{sites-cohomology-lemma-A} の圏 $\mathcal{A}_X$ は
$a_X^{-1} : \textit{Ab}(X) \to \textit{Ab}(\textit{LC}_{qc}/X)$ の
本質像であることに注意する。

\medskip\noindent
（1）はすべての $n$ に対する $(V_n)$ と同値であり、補題
\ref{sites-cohomology-lemma-V-C-all-n-general} により成り立つ。

\medskip\noindent
（2）は補題 \ref{sites-cohomology-lemma-V-implies-C-general} の結論に
$\epsilon_Y^{-1}$ を適用すれば従う。

\medskip\noindent
（3）は補題 \ref{sites-cohomology-lemma-V-C-all-n-general} の（1）から従う。実際、
$\pi_X^{-1}K$ は $D^+_{\mathcal{A}'_X}(\textit{LC}_{Zar}/X)$ に属し、
$a_X^{-1} = \epsilon_X^{-1} \circ a_X^{-1}$ である。

\medskip\noindent
（4）も同じ理由で補題 \ref{sites-cohomology-lemma-V-C-all-n-general} の（2）から従う。
\end{proof}

\begin{lemma}
\label{sites-cohomology-lemma-cohomological-descent-LC}
$X$ を $\textit{LC}_{qc}$ の対象とする。$K \in D^+(X)$ に対して写像
$$
K \longrightarrow Ra_{X, *}a_X^{-1}K
$$
は同型である。ここで $a_X : \Sh(\textit{LC}_{qc}/X) \to \Sh(X)$ は上のものである。
\end{lemma}

\begin{proof}
まず、$K$ が一つのアーベル層で与えられる場合に帰着する。
実際、$K$ を下に有界な複体 $\mathcal{F}^\bullet$ で表す。層の場合から
$\mathcal{F}^n = a_{X, *} a_X^{-1} \mathcal{F}^n$
であり、$q > 0$ に対して層 $R^qa_{X, *}a_X^{-1}\mathcal{F}^n$ は零である。
$a_X^{-1}\mathcal{F}^\bullet$ と関手 $a_{X,*}$ に Leray の非輪状性補題
（Derived Categories, Lemma \ref{derived-lemma-leray-acyclicity}）を適用して
結論を得る。以下では $K = \mathcal{F}$ と仮定する。

\medskip\noindent
補題 \ref{sites-cohomology-lemma-describe-pullback-pi} により
$a_{X, *}a_X^{-1}\mathcal{F} = \mathcal{F}$ である。したがって $q > 0$ に対して
$R^qa_{X, *}a_X^{-1}\mathcal{F} = 0$ を示せば十分である。
このために $a_X = \epsilon_X \circ \pi_X$ と Leray スペクトル系列の補題
\ref{sites-cohomology-lemma-relative-Leray} を用いる。補題 \ref{sites-cohomology-lemma-V-C-all-n} により、
$i > 0$ に対して $R^i\epsilon_{X, *}(a_X^{-1}\mathcal{F}) = 0$ であり、
$\epsilon_{X, *}a_X^{-1}\mathcal{F} = \pi_X^{-1}\mathcal{F}$ である。
補題 \ref{sites-cohomology-lemma-collect-true-things-Zar} により、$j > 0$ に対して
$R^j\pi_{X, *}(\pi_X^{-1}\mathcal{F}) = 0$ である。これで証明が完了する。
\end{proof}

\begin{lemma}
\label{sites-cohomology-lemma-compare-cohomology-LC}
$X \in \Ob(\textit{LC}_{qc})$ とし、
$a_X : \Sh(\textit{LC}_{qc}/X) \to \Sh(X)$ を上のものとする。このとき、
\begin{enumerate}
\item $X$ 上のアーベル層 $\mathcal{F}$ に対して
$H^n(X, \mathcal{F}) = H^n_{qc}(X, a_X^{-1}\mathcal{F})$。
\item $K \in D^+(X)$ に対して $H^n(X, K) = H^n_{qc}(X, a_X^{-1}K)$。
\end{enumerate}
たとえば $A$ がアーベル群ならば、
$H^n(X, \underline{A}) = H^n_{qc}(X, \underline{A})$ である。
\end{lemma}

\begin{proof}
補題 \ref{sites-cohomology-lemma-cohomological-descent-LC} と
注意 \ref{sites-cohomology-remark-before-Leray} から従う。
\end{proof}











\section{Ext に対するスペクトル系列}
\label{sites-cohomology-section-spectral-sequence-ext}

\noindent
この節では、Ext 関手を考える際に現れるさまざまなスペクトル系列をまとめる。
環付きサイト $(\mathcal{C}, \mathcal{O})$ 上の加群の複体の任意の対
$\mathcal{G}^\bullet, \mathcal{F}^\bullet$ に対して、一般的な規約
Derived Categories, Section \ref{derived-section-ext} に従い、
$$
\Ext^n_\mathcal{O}(\mathcal{G}^\bullet, \mathcal{F}^\bullet)
=
\Hom_{D(\mathcal{O})}(\mathcal{G}^\bullet, \mathcal{F}^\bullet[n])
$$
と表す。

\begin{example}
\label{sites-cohomology-example-hom-complex-into-sheaf}
$(\mathcal{C}, \mathcal{O})$ を環付きサイトとする。
$\mathcal{K}^\bullet$ を上に有界な $\mathcal{O}$-加群の複体とし、
$\mathcal{F}$ を $\mathcal{O}$-加群とする。このとき第 $E_2$ 頁が
$$
E_2^{i, j} =
\Ext_\mathcal{O}^i(H^{-j}(\mathcal{K}^\bullet), \mathcal{F})
\Rightarrow
\Ext_\mathcal{O}^{i + j}(\mathcal{K}^\bullet, \mathcal{F})
$$
であるスペクトル系列と、第 $E_1$ 頁が
$$
E_1^{i, j} =
\Ext_\mathcal{O}^j(\mathcal{K}^{-i}, \mathcal{F})
\Rightarrow
\Ext_\mathcal{O}^{i + j}(\mathcal{K}^\bullet, \mathcal{F}).
$$
であるもう一つのスペクトル系列がある。これらを構成するには、入射分解
$\mathcal{F} \to \mathcal{I}^\bullet$ を選び、二重複体
$\Hom_\mathcal{O}(\mathcal{K}^\bullet, \mathcal{I}^\bullet)$ から生じる
二つのスペクトル系列を考える；Homology, Section
\ref{homology-section-double-complex} を参照せよ。
\end{example}








\section{カップ積}
\label{sites-cohomology-section-cup-product}

\noindent
$(\mathcal{C}, \mathcal{O})$ を環付きサイトとし、$K, M$ を
$D(\mathcal{O})$ の対象とする。$A = \Gamma(\mathcal{C}, \mathcal{O})$ とおく。
この設定における（大域）カップ積は $D(A)$ における写像
$$
R\Gamma(\mathcal{C}, K) \otimes_A^\mathbf{L}
R\Gamma(\mathcal{C}, M)
\longrightarrow
R\Gamma(\mathcal{C}, K \otimes_\mathcal{O}^\mathbf{L} M)
$$
である。これを、注意 \ref{sites-cohomology-remark-cup-product} における環付きトポスの射
$(\Sh(\mathcal{C}), \mathcal{O}) \to (\Sh(pt), A)$
に対する相対カップ積として定義する。

\medskip\noindent
相対カップ積の自然な両立性を定式化し、証明しよう。すなわち、環付きトポスの射
$f : (\Sh(\mathcal{C}), \mathcal{O}_\mathcal{C}) \to
(\Sh(\mathcal{D}), \mathcal{O}_\mathcal{D})$
があるとする。$\mathcal{K}^\bullet$ と $\mathcal{M}^\bullet$ を
$\mathcal{O}_\mathcal{C}$-加群の複体とする。素朴なカップ積
$$
\text{Tot}(
f_*\mathcal{K}^\bullet
\otimes_{\mathcal{O}_\mathcal{D}}
f_*\mathcal{M}^\bullet)
\longrightarrow
f_*\text{Tot}(\mathcal{K}^\bullet
\otimes_{\mathcal{O}_\mathcal{C}}
\mathcal{M}^\bullet)
$$
がある。これは相対カップ積と次のように関係する。

\begin{lemma}
\label{sites-cohomology-lemma-cup-compatible-with-naive}
上の状況において、次の図式は可換である。
$$
\xymatrix{
f_*\mathcal{K}^\bullet
\otimes_{\mathcal{O}_\mathcal{D}}^\mathbf{L}
f_*\mathcal{M}^\bullet \ar[r] \ar[d]
&
Rf_*\mathcal{K}^\bullet
\otimes_{\mathcal{O}_\mathcal{D}}^\mathbf{L}
Rf_*\mathcal{M}^\bullet \ar[d]^{\text{注意 \ref{sites-cohomology-remark-cup-product}}} \\
\text{Tot}(
f_*\mathcal{K}^\bullet
\otimes_{\mathcal{O}_\mathcal{D}}
f_*\mathcal{M}^\bullet) \ar[d]_{\text{素朴なカップ積}} &
Rf_*(\mathcal{K}^\bullet
\otimes_{\mathcal{O}_\mathcal{C}}^\mathbf{L}
\mathcal{M}^\bullet) \ar[d] \\
f_*\text{Tot}(\mathcal{K}^\bullet
\otimes_{\mathcal{O}_\mathcal{C}}
\mathcal{M}^\bullet) \ar[r] &
Rf_*\text{Tot}(\mathcal{K}^\bullet
\otimes_{\mathcal{O}_\mathcal{C}}
\mathcal{M}^\bullet)
}
$$
\end{lemma}

\begin{proof}
注意 \ref{sites-cohomology-remark-cup-product} の構成により、図式を時計回りにたどって得る写像
$$
f_*\mathcal{K}^\bullet
\otimes_{\mathcal{O}_\mathcal{D}}^\mathbf{L}
f_*\mathcal{M}^\bullet 
\longrightarrow
Rf_*\text{Tot}(\mathcal{K}^\bullet
\otimes_{\mathcal{O}_\mathcal{C}}
\mathcal{M}^\bullet)
$$
は、写像
\begin{align*}
Lf^*(f_*\mathcal{K}^\bullet
\otimes_{\mathcal{O}_\mathcal{D}}^\mathbf{L}
f_*\mathcal{M}^\bullet)
& =
Lf^*f_*\mathcal{K}^\bullet
\otimes_{\mathcal{O}_\mathcal{D}}^\mathbf{L}
Lf^*f_*\mathcal{M}^\bullet \\
& \to
Lf^*Rf_*\mathcal{K}^\bullet
\otimes_{\mathcal{O}_\mathcal{D}}^\mathbf{L}
Lf^*Rf_*\mathcal{M}^\bullet \\
& \to
\mathcal{K}^\bullet
\otimes_{\mathcal{O}_\mathcal{D}}^\mathbf{L}
\mathcal{M}^\bullet \\
& \to
\text{Tot}(\mathcal{K}^\bullet
\otimes_{\mathcal{O}_\mathcal{C}}
\mathcal{M}^\bullet)
\end{align*}
に随伴する。補題 \ref{sites-cohomology-lemma-adjoints-push-pull-compatibility} により、これはまた
\begin{align*}
Lf^*(f_*\mathcal{K}^\bullet
\otimes_{\mathcal{O}_\mathcal{D}}^\mathbf{L}
f_*\mathcal{M}^\bullet)
& =
Lf^*f_*\mathcal{K}^\bullet
\otimes_{\mathcal{O}_\mathcal{D}}^\mathbf{L}
Lf^*f_*\mathcal{M}^\bullet \\
& \to
f^*f_*\mathcal{K}^\bullet
\otimes_{\mathcal{O}_\mathcal{D}}^\mathbf{L}
f^*f_*\mathcal{M}^\bullet \\
& \to
\mathcal{K}^\bullet
\otimes_{\mathcal{O}_\mathcal{D}}^\mathbf{L}
\mathcal{M}^\bullet \\
& \to
\text{Tot}(\mathcal{K}^\bullet
\otimes_{\mathcal{O}_\mathcal{C}}
\mathcal{M}^\bullet)
\end{align*}
に等しい。反時計回りにたどると、写像
\begin{align*}
Lf^*(f_*\mathcal{K}^\bullet
\otimes_{\mathcal{O}_\mathcal{D}}^\mathbf{L}
f_*\mathcal{M}^\bullet)
& \to
Lf^*\text{Tot}(
f_*\mathcal{K}^\bullet
\otimes_{\mathcal{O}_\mathcal{D}}
f_*\mathcal{M}^\bullet) \\
& \to
Lf^*f_*\text{Tot}(\mathcal{K}^\bullet
\otimes_{\mathcal{O}_\mathcal{C}}
\mathcal{M}^\bullet) \\
& \to
Lf^*Rf_*\text{Tot}(\mathcal{K}^\bullet
\otimes_{\mathcal{O}_\mathcal{C}}
\mathcal{M}^\bullet) \\
& \to
\text{Tot}(\mathcal{K}^\bullet
\otimes_{\mathcal{O}_\mathcal{C}}
\mathcal{M}^\bullet)
\end{align*}
に随伴する写像を得る。補題 \ref{sites-cohomology-lemma-adjoints-push-pull-compatibility} により、
これはまた
\begin{align*}
Lf^*(f_*\mathcal{K}^\bullet
\otimes_{\mathcal{O}_\mathcal{D}}^\mathbf{L}
f_*\mathcal{M}^\bullet)
& \to
Lf^*\text{Tot}(
f_*\mathcal{K}^\bullet
\otimes_{\mathcal{O}_\mathcal{D}}
f_*\mathcal{M}^\bullet) \\
& \to
Lf^*f_*\text{Tot}(\mathcal{K}^\bullet
\otimes_{\mathcal{O}_\mathcal{C}}
\mathcal{M}^\bullet) \\
& \to
f^*f_*\text{Tot}(\mathcal{K}^\bullet
\otimes_{\mathcal{O}_\mathcal{C}}
\mathcal{M}^\bullet) \\
& \to
\text{Tot}(\mathcal{K}^\bullet
\otimes_{\mathcal{O}_\mathcal{C}}
\mathcal{M}^\bullet)
\end{align*}
に等しい。ここで次の図式を考えれば証明が完了する。
$$
\xymatrix{
Lf^*(f_*\mathcal{K}^\bullet
\otimes_{\mathcal{O}_\mathcal{D}}^\mathbf{L}
f_*\mathcal{M}^\bullet) \ar[d] \ar[rr] & &
Lf^*f_*\mathcal{K}^\bullet \otimes_{\mathcal{O}_\mathcal{C}}^\mathbf{L}
Lf^*f_*\mathcal{M}^\bullet \ar[d] \\
Lf^*\text{Tot}(
f_*\mathcal{K}^\bullet
\otimes_{\mathcal{O}_\mathcal{D}}
f_*\mathcal{M}^\bullet) \ar[d]_{naive} \ar[r] &
f^*\text{Tot}(
f_*\mathcal{K}^\bullet
\otimes_{\mathcal{O}_\mathcal{D}}
f_*\mathcal{M}^\bullet) \ar[ldd]^{naive} \ar[dd] &
f^*f_*\mathcal{K}^\bullet \otimes_{\mathcal{O}_\mathcal{C}}^\mathbf{L}
f^*f_*\mathcal{M}^\bullet \ar[dd] \ar[ldd] \\
Lf^*f_*\text{Tot}(\mathcal{K}^\bullet
\otimes_{\mathcal{O}_\mathcal{C}}
\mathcal{M}^\bullet) \ar[d] \\
f^*f_*\text{Tot}(\mathcal{K}^\bullet \otimes_{\mathcal{O}_\mathcal{C}}
\mathcal{M}^\bullet) \ar[rd] &
\text{Tot}(f^*f_*\mathcal{K}^\bullet \otimes_{\mathcal{O}_\mathcal{C}}
f^*f_*\mathcal{M}^\bullet) \ar[d] &
\mathcal{K}^\bullet \otimes_{\mathcal{O}_\mathcal{C}}^\mathbf{L}
\mathcal{M}^\bullet \ar[ld] \\
& \text{Tot}(\mathcal{K}^\bullet
\otimes_{\mathcal{O}_\mathcal{C}}
\mathcal{M}^\bullet)
}
$$
この図式のすべての多角形は可換である。上の多角形は補題
\ref{sites-cohomology-lemma-tensor-pull-compatibility} により可換である。
二つの素朴なカップ積を含む正方形は、$Lf^* \to f^*$ が加群の複体について
関手的なので可換である。同様に、二つの写像
$\mathcal{A}^\bullet \otimes^\mathbf{L} \mathcal{B}^\bullet \to
\text{Tot}(\mathcal{A}^\bullet \otimes \mathcal{B}^\bullet)$
を含む正方形も可換である。最後に、残る正方形の可換性は複体のレベルで成り立ち、
$f^*$ と $f_*$ の随伴性による素朴なカップ積の定義とみなせる。
図式の外周をたどる二つの写像が上で与えた二つの写像なので、証明が完了する。
\end{proof}

\begin{lemma}
\label{sites-cohomology-lemma-cup-product-associative}
$f : (\Sh(\mathcal{C}), \mathcal{O}) \to (\Sh(\mathcal{C}'), \mathcal{O}')$
を環付きトポスの射とする。注意 \ref{sites-cohomology-remark-cup-product} の相対カップ積は、
図式
$$
\xymatrix{
Rf_*K \otimes_{\mathcal{O}'}^\mathbf{L}
Rf_*L \otimes_{\mathcal{O}'}^\mathbf{L}
Rf_*M \ar[r] \ar[d] &
Rf_*(K \otimes_\mathcal{O}^\mathbf{L} L)
\otimes_{\mathcal{O}'}^\mathbf{L} Rf_*M \ar[d] \\
Rf_*K \otimes_{\mathcal{O}'}^\mathbf{L}
Rf_*(L \otimes_\mathcal{O}^\mathbf{L} M) \ar[r] &
Rf_*(K \otimes_\mathcal{O}^\mathbf{L} 
L \otimes_\mathcal{O}^\mathbf{L} M)
}
$$
がすべての $K, L, M \in D(\mathcal{O})$ に対して $D(\mathcal{O}')$ で
可換となるという意味で結合的である。
\end{lemma}

\begin{proof}
いずれの側をたどっても、$D(\mathcal{O})$ における明らかな写像
\begin{align*}
Lf^*(Rf_*K \otimes_{\mathcal{O}'}^\mathbf{L}
Rf_*L \otimes_{\mathcal{O}'}^\mathbf{L}
Rf_*M) & =
Lf^*(Rf_*K) \otimes_\mathcal{O}^\mathbf{L}
Lf^*(Rf_*L) \otimes_\mathcal{O}^\mathbf{L}
Lf^*(Rf_*M) \\
& \to
K \otimes_\mathcal{O}^\mathbf{L} 
L \otimes_\mathcal{O}^\mathbf{L} M
\end{align*}
に随伴する写像を得る。
\end{proof}

\begin{lemma}
\label{sites-cohomology-lemma-cup-product-commutative}
$f : (\Sh(\mathcal{C}), \mathcal{O}) \to (\Sh(\mathcal{C}'), \mathcal{O}')$
を環付きトポスの射とする。注意 \ref{sites-cohomology-remark-cup-product} の相対カップ積は、
図式
$$
\xymatrix{
Rf_*K \otimes_{\mathcal{O}'}^\mathbf{L} Rf_*L \ar[r] \ar[d]_\psi &
Rf_*(K \otimes_\mathcal{O}^\mathbf{L} L) \ar[d]^{Rf_*\psi} \\
Rf_*L \otimes_{\mathcal{O}'}^\mathbf{L} Rf_*K \ar[r] &
Rf_*(L \otimes_\mathcal{O}^\mathbf{L} K)
}
$$
がすべての $K, L \in D(\mathcal{O})$ に対して $D(\mathcal{O}')$ で
可換となるという意味で可換である。ここで $\psi$ は導来圏上の可換性制約である
（補題 \ref{sites-cohomology-lemma-symmetric-monoidal-derived}）。
\end{lemma}

\begin{proof}
省略する。
\end{proof}

\begin{lemma}
\label{sites-cohomology-lemma-compose-cup-product}
$f : (\Sh(\mathcal{C}), \mathcal{O}) \to (\Sh(\mathcal{C}'), \mathcal{O}')$
および $f' : (\Sh(\mathcal{C}'), \mathcal{O}') \to
(\Sh(\mathcal{C}''), \mathcal{O}'')$ を環付きトポスの射とする。
注意 \ref{sites-cohomology-remark-cup-product} の相対カップ積は、図式
$$
\xymatrix{
R(f' \circ f)_*K \otimes_{\mathcal{O}''}^\mathbf{L} R(f' \circ f)_*L
\ar@{=}[rr] \ar[d] & &
Rf'_*Rf_*K \otimes_{\mathcal{O}''}^\mathbf{L} Rf'_*Rf_*L \ar[d] \\
R(f' \circ f)_*(K \otimes_\mathcal{O}^\mathbf{L} L) \ar@{=}[r] &
Rf'_*Rf_*(K \otimes_\mathcal{O}^\mathbf{L} L) &
Rf'_*(Rf_*K \otimes_{\mathcal{O}'}^\mathbf{L}  Rf_*L) \ar[l]
}
$$
がすべての $K, L \in D(\mathcal{O})$ に対して $D(\mathcal{O}'')$ で
可換となるという意味で合成と両立する。
\end{lemma}

\begin{proof}
これは、いずれの向きに図式をたどっても、$D(\mathcal{O})$ における写像
\begin{align*}
& L(f' \circ f)^*\left(R(f' \circ f)_*K
\otimes_{\mathcal{O}''}^\mathbf{L}
R(f' \circ f)_*L\right) \\
& =
L(f' \circ f)^*R(f' \circ f)_*K
\otimes_\mathcal{O}^\mathbf{L}
L(f' \circ f)^*R(f' \circ f)_*L) \\
& \to
K \otimes_\mathcal{O}^\mathbf{L} L
\end{align*}
に随伴する写像を得るためである。これを見るには、余単位の合成
$$
L(f' \circ f)^*R(f' \circ f)_* =
Lf^* L(f')^* Rf'_* Rf_*  \to
Lf^* Rf_* \to \text{id}
$$
が $L(f' \circ f)^*$ と $R(f' \circ f)_*$ に対する余単位であることを用いる。
Categories, Lemma \ref{categories-lemma-compose-counits} を参照せよ。
\end{proof}

\begin{lemma}
\label{sites-cohomology-lemma-base-change-cup-product}
環付きトポスの可換正方形
$$
\xymatrix{
(\Sh(\mathcal{C}'), \mathcal{O}_{\mathcal{C}'}) \ar[r]_{g'} \ar[d]_{f'} &
(\Sh(\mathcal{C}, \mathcal{O}_{\mathcal{C}}) \ar[d]^f \\
(\Sh(\mathcal{D}'), \mathcal{O}_{\mathcal{D}'}) \ar[r]^g &
(\Sh(\mathcal{D}), \mathcal{O}_{\mathcal{D}}) 
}
$$
を考える。$K, L \in D(\mathcal{O}_{\mathcal{C}})$ とする。
相対カップ積は、図式
$$
\xymatrix{
Lg^*(Rf_*K \otimes_{\mathcal{O}_{\mathcal{D}}}^\mathbf{L} Rf_*L)
\ar[r] \ar@{=}[d] &
Lg^*(Rf_*(K \otimes_{\mathcal{O}_{\mathcal{C}}}^\mathbf{L} L)) \ar[d] \\
Lg^*Rf_*K \otimes_{\mathcal{O}_{\mathcal{D}'}}^\mathbf{L} Lg^*Rf_*L \ar[d] &
R(f')_*L(g')^*(K \otimes_{\mathcal{O}_{\mathcal{C}}}^\mathbf{L} L) \ar@{=}[d] \\
R(f')_*(L(g')^*K \otimes_{\mathcal{O}_{\mathcal{D}'}} R(f')_*(L(g')^*L \ar[r] &
R(f')_*(L(g')^*K \otimes_{\mathcal{O}_{\mathcal{C}'}}^\mathbf{L} L(g')^*L)
}
$$
が $D(\mathcal{O}_{\mathcal{D}'})$ で可換となるという意味でこの正方形と両立する。
横の矢印は相対カップ積（注意 \ref{sites-cohomology-remark-cup-product}）により与えられ、
縦の矢印は基底変換写像（注意 \ref{sites-cohomology-remark-base-change}）と補題
\ref{sites-cohomology-lemma-pullback-tensor-product} により与えられる。
\end{lemma}

\begin{proof}
省略する。
\end{proof}






\section{Hom 複体}
\label{sites-cohomology-section-hom-complexes}

\noindent
$(\mathcal{C}, \mathcal{O})$ を環付きサイトとする。
$\mathcal{L}^\bullet$ と $\mathcal{M}^\bullet$ を二つの $\mathcal{O}$-加群の
複体とする。$\mathcal{O}$-加群の複体
$\SheafHom^\bullet(\mathcal{L}^\bullet, \mathcal{M}^\bullet)$ を構成する。
すなわち、各 $n$ に対して
$$
\SheafHom^n(\mathcal{L}^\bullet, \mathcal{M}^\bullet) =
\prod\nolimits_{n = p + q}
\SheafHom_\mathcal{O}(\mathcal{L}^{-q}, \mathcal{M}^p)
$$
とおく。$\SheafHom^n$ は、$\mathcal{L}^\bullet$ から
$\mathcal{M}^\bullet$ への次数 $n$ の斉次な $\mathcal{O}$-線形写像
（両者を次数付き $\mathcal{O}$-加群とみなす）すべてからなる
$\mathcal{O}$-加群の層と考えるとよい。この用語のもとで、微分を規則
$$
\text{d}(f) =
\text{d}_\mathcal{M} \circ f - (-1)^n f \circ \text{d}_\mathcal{L}
$$
により定義する。ここで
$f \in \SheafHom^n_\mathcal{O}(\mathcal{L}^\bullet, \mathcal{M}^\bullet)$ である。
$\text{d}^2 = 0$ の確認は省略する。この構成は Differential Graded Algebra,
Example \ref{dga-example-category-complexes} の特別な場合である。
構成から直ちに、すべての $n \in \mathbf{Z}$ と各
$U \in \Ob(\mathcal{C})$ に対して
\begin{equation}
\label{sites-cohomology-equation-cohomology-hom-complex}
H^n(\Gamma(U, \SheafHom^\bullet(\mathcal{L}^\bullet, \mathcal{M}^\bullet))) =
\Hom_{K(\mathcal{O}_U)}(\mathcal{L}^\bullet|_U, \mathcal{M}^\bullet[n]|_U)
\end{equation}
を得る。同様に、大域切片の複体について
\begin{equation}
\label{sites-cohomology-equation-global-cohomology-hom-complex}
H^n(\Gamma(\mathcal{C},
\SheafHom^\bullet(\mathcal{L}^\bullet, \mathcal{M}^\bullet))) =
\Hom_{K(\mathcal{O})}(\mathcal{L}^\bullet, \mathcal{M}^\bullet[n])
\end{equation}
を得る。

\begin{lemma}
\label{sites-cohomology-lemma-compose}
$(\mathcal{C}, \mathcal{O})$ を環付きサイトとする。
$\mathcal{O}$-加群の複体
$\mathcal{K}^\bullet, \mathcal{L}^\bullet, \mathcal{M}^\bullet$
が与えられたとき、同型
$$
\SheafHom^\bullet(\mathcal{K}^\bullet,
\SheafHom^\bullet(\mathcal{L}^\bullet, \mathcal{M}^\bullet))
=
\SheafHom^\bullet(\text{Tot}(\mathcal{K}^\bullet \otimes_\mathcal{O}
\mathcal{L}^\bullet), \mathcal{M}^\bullet)
$$
があり、これは $\mathcal{K}^\bullet, \mathcal{L}^\bullet,
\mathcal{M}^\bullet$ について関手的な $\mathcal{O}$-加群の複体の同型である。
\end{lemma}

\begin{proof}
省略する。ヒント：More on Algebra, Lemma
\ref{more-algebra-lemma-compose} とまったく同じ方法で証明される。
\end{proof}

\begin{lemma}
\label{sites-cohomology-lemma-composition}
$(\mathcal{C}, \mathcal{O})$ を環付きサイトとする。$\mathcal{O}$-加群の複体
$\mathcal{K}^\bullet, \mathcal{L}^\bullet, \mathcal{M}^\bullet$
が与えられたとき、標準射
$$
\text{Tot}\left(
\SheafHom^\bullet(\mathcal{L}^\bullet, \mathcal{M}^\bullet)
\otimes_\mathcal{O}
\SheafHom^\bullet(\mathcal{K}^\bullet, \mathcal{L}^\bullet)
\right)
\longrightarrow
\SheafHom^\bullet(\mathcal{K}^\bullet, \mathcal{M}^\bullet)
$$
が $\mathcal{O}$-加群の複体の間にある。
\end{lemma}

\begin{proof}
省略する。ヒント：More on Algebra, Lemma
\ref{more-algebra-lemma-composition} とまったく同じ方法で証明される。
\end{proof}

\begin{lemma}
\label{sites-cohomology-lemma-diagonal-better}
$(\mathcal{C}, \mathcal{O})$ を環付きサイトとする。$\mathcal{O}$-加群の複体
$\mathcal{K}^\bullet, \mathcal{L}^\bullet, \mathcal{M}^\bullet$
が与えられたとき、標準射
$$
\text{Tot}\left(
\mathcal{K}^\bullet \otimes_\mathcal{O}
\SheafHom^\bullet(\mathcal{M}^\bullet, \mathcal{L}^\bullet)
\right)
\longrightarrow
\SheafHom^\bullet(\mathcal{M}^\bullet,
\text{Tot}(\mathcal{K}^\bullet \otimes_\mathcal{O} \mathcal{L}^\bullet))
$$
があり、これは三つの複体すべてについて関手的な
$\mathcal{O}$-加群の複体の射である。
\end{lemma}

\begin{proof}
省略する。ヒント：More on Algebra, Lemma
\ref{more-algebra-lemma-diagonal-better} とまったく同じ方法で証明される。
\end{proof}

\begin{lemma}
\label{sites-cohomology-lemma-diagonal}
$(\mathcal{C}, \mathcal{O})$ を環付きサイトとする。$\mathcal{O}$-加群の複体
$\mathcal{K}^\bullet, \mathcal{L}^\bullet, \mathcal{M}^\bullet$
が与えられたとき、標準射
$$
\mathcal{K}^\bullet
\longrightarrow
\SheafHom^\bullet(\mathcal{L}^\bullet,
\text{Tot}(\mathcal{K}^\bullet \otimes_\mathcal{O} \mathcal{L}^\bullet))
$$
があり、これは両方の複体について関手的な $\mathcal{O}$-加群の複体の射である。
\end{lemma}

\begin{proof}
省略する。ヒント：More on Algebra, Lemma
\ref{more-algebra-lemma-diagonal} とまったく同じ方法で証明される。
\end{proof}

\begin{lemma}
\label{sites-cohomology-lemma-evaluate-and-more}
$(\mathcal{C}, \mathcal{O})$ を環付きサイトとする。$\mathcal{O}$-加群の複体
$\mathcal{K}^\bullet, \mathcal{L}^\bullet, \mathcal{M}^\bullet$
が与えられたとき、標準射
$$
\text{Tot}(\SheafHom^\bullet(\mathcal{L}^\bullet,
\mathcal{M}^\bullet) \otimes_\mathcal{O} \mathcal{K}^\bullet)
\longrightarrow
\SheafHom^\bullet(\SheafHom^\bullet(\mathcal{K}^\bullet,
\mathcal{L}^\bullet), \mathcal{M}^\bullet)
$$
があり、これは三つの複体すべてについて関手的な
$\mathcal{O}$-加群の複体の射である。
\end{lemma}

\begin{proof}
省略する。ヒント：More on Algebra, Lemma
\ref{more-algebra-lemma-evaluate-and-more} とまったく同じ方法で証明される。
\end{proof}

\begin{lemma}
\label{sites-cohomology-lemma-RHom-into-K-injective}
$(\mathcal{C}, \mathcal{O})$ を環付きサイトとし、$L$ と $M$ を
$D(\mathcal{O})$ の対象とする。$\mathcal{I}^\bullet$ を $M$ を表す
$\mathcal{O}$-加群の K-入射複体とし、$\mathcal{L}^\bullet$ を $L$ を表す
$\mathcal{O}$-加群の複体とする。このとき、すべての
$U \in \Ob(\mathcal{C})$ に対して
$$
H^0(\Gamma(U, \SheafHom^\bullet(\mathcal{L}^\bullet, \mathcal{I}^\bullet))) =
\Hom_{D(\mathcal{O}_U)}(L|_U, M|_U)
$$
である。同様に、
$H^0(\Gamma(\mathcal{C},
\SheafHom^\bullet(\mathcal{L}^\bullet, \mathcal{I}^\bullet))) =
\Hom_{D(\mathcal{O})}(L, M)$.
\end{lemma}

\begin{proof}
次の等式が成り立つ。
\begin{align*}
H^0(\Gamma(U, \SheafHom^\bullet(\mathcal{L}^\bullet, \mathcal{I}^\bullet)))
& =
\Hom_{K(\mathcal{O}_U)}(\mathcal{L}^\bullet|_U, \mathcal{I}^\bullet|_U) \\
& =
\Hom_{D(\mathcal{O}_U)}(L|_U, M|_U)
\end{align*}
第一の等号は（\ref{sites-cohomology-equation-cohomology-hom-complex}）である。
第二の等号は、補題 \ref{sites-cohomology-lemma-restrict-K-injective-to-open} により
$\mathcal{I}^\bullet|_U$ が K-入射なので成り立つ。最後の等式の証明も同様だが、
（\ref{sites-cohomology-equation-global-cohomology-hom-complex}）を用いる。
\end{proof}

\begin{lemma}
\label{sites-cohomology-lemma-RHom-well-defined}
$(\mathcal{C}, \mathcal{O})$ を環付きサイトとする。
$(\mathcal{I}')^\bullet \to \mathcal{I}^\bullet$ を
$\mathcal{O}$-加群の K-入射複体の間の擬同型とし、
$(\mathcal{L}')^\bullet \to \mathcal{L}^\bullet$ を
$\mathcal{O}$-加群の複体の間の擬同型とする。このとき
$$
\SheafHom^\bullet(\mathcal{L}^\bullet, (\mathcal{I}')^\bullet)
\longrightarrow
\SheafHom^\bullet((\mathcal{L}')^\bullet, \mathcal{I}^\bullet)
$$
は擬同型である。
\end{lemma}

\begin{proof}
$M$ を $\mathcal{I}^\bullet$ と $(\mathcal{I}')^\bullet$ が表す
$D(\mathcal{O})$ の対象とし、$L$ を $\mathcal{L}^\bullet$ と
$(\mathcal{L}')^\bullet$ が表す $D(\mathcal{O})$ の対象とする。
補題 \ref{sites-cohomology-lemma-RHom-into-K-injective} により、層
$$
H^0(\SheafHom^\bullet(\mathcal{L}^\bullet, (\mathcal{I}')^\bullet))
\quad\text{and}\quad
H^0(\SheafHom^\bullet((\mathcal{L}')^\bullet, \mathcal{I}^\bullet))
$$
はいずれも前層
$$
U \longmapsto \Hom_{D(\mathcal{O}_U)}(L|_U, M|_U)
$$
に付随する層に等しい。したがってこの写像は擬同型である。
\end{proof}

\begin{lemma}
\label{sites-cohomology-lemma-RHom-from-K-flat-into-K-injective}
$(\mathcal{C}, \mathcal{O})$ を環付きサイトとする。$\mathcal{I}^\bullet$ を
$\mathcal{O}$-加群の K-入射複体とし、$\mathcal{L}^\bullet$ を
$\mathcal{O}$-加群の K-平坦複体とする。このとき
$\SheafHom^\bullet(\mathcal{L}^\bullet, \mathcal{I}^\bullet)$ は
$\mathcal{O}$-加群の K-入射複体である。
\end{lemma}

\begin{proof}
実際、$\mathcal{K}^\bullet$ が $\mathcal{O}$-加群の非輪状複体ならば、
\begin{align*}
\Hom_{K(\mathcal{O})}(\mathcal{K}^\bullet,
\SheafHom^\bullet(\mathcal{L}^\bullet, \mathcal{I}^\bullet))
& =
H^0(\Gamma(\mathcal{C},
\SheafHom^\bullet(\mathcal{K}^\bullet,
\SheafHom^\bullet(\mathcal{L}^\bullet, \mathcal{I}^\bullet)))) \\
& =
H^0(\Gamma(\mathcal{C}, \SheafHom^\bullet(\text{Tot}(
\mathcal{K}^\bullet \otimes_\mathcal{O} \mathcal{L}^\bullet),
\mathcal{I}^\bullet))) \\
& =
\Hom_{K(\mathcal{O})}(
\text{Tot}(\mathcal{K}^\bullet \otimes_\mathcal{O} \mathcal{L}^\bullet),
\mathcal{I}^\bullet) \\
& =
0
\end{align*}
第一の等号は（\ref{sites-cohomology-equation-global-cohomology-hom-complex}）による。
第二の等号は補題 \ref{sites-cohomology-lemma-compose} による。
第三の等号は（\ref{sites-cohomology-equation-global-cohomology-hom-complex}）による。
最後の等号は、
$\text{Tot}(\mathcal{K}^\bullet \otimes_\mathcal{O} \mathcal{L}^\bullet)$
が $\mathcal{L}^\bullet$ の K-平坦性
（定義 \ref{sites-cohomology-definition-K-flat}）により非輪状であり、
$\mathcal{I}^\bullet$ が K-入射だからである。
\end{proof}








\section{導来圏における内部 Hom}
\label{sites-cohomology-section-internal-hom}

\noindent
$(\mathcal{C}, \mathcal{O})$ を環付きサイトとする。$L, M$ を
$D(\mathcal{O})$ の対象とする。$D(\mathcal{O})$ の任意の第三の対象
$K$ に対して標準的全単射
\begin{equation}
\label{sites-cohomology-equation-internal-hom}
\Hom_{D(\mathcal{O})}(K, R\SheafHom(L, M))
=
\Hom_{D(\mathcal{O})}(K \otimes_\mathcal{O}^\mathbf{L} L, M)
\end{equation}
が存在するような $D(\mathcal{O})$ の対象 $R\SheafHom(L, M)$ を構成したい。
米田の補題（Categories, Lemma \ref{categories-lemma-yoneda}）により、
この式は $R\SheafHom(L, M)$ を一意的同型を除いて定めることに注意する。

\medskip\noindent
そのような対象を構成するため、$M$ を表す $\mathcal{O}$-加群の
K-入射複体 $\mathcal{I}^\bullet$ と、$L$ を表す任意の
$\mathcal{O}$-加群の複体 $\mathcal{L}^\bullet$ を選ぶ。そして
$$
R\SheafHom(L, M) = \SheafHom^\bullet(\mathcal{L}^\bullet, \mathcal{I}^\bullet)
$$
と置く。右辺は節 \ref{sites-cohomology-section-hom-complexes} で構成した
$\mathcal{O}$-加群の複体である。これは補題
\ref{sites-cohomology-lemma-RHom-well-defined} により良定義である。こうして関手
$$
D(\mathcal{O})^{opp} \times D(\mathcal{O}) \longrightarrow D(\mathcal{O}),
\quad
(K, L) \longmapsto R\SheafHom(K, L)
$$
を得る。（\ref{sites-cohomology-equation-internal-hom}）を証明する準備として、
$R\SheafHom(K, L)$ のコホモロジー群を計算する。

\begin{lemma}
\label{sites-cohomology-lemma-section-RHom-over-U}
$(\mathcal{C}, \mathcal{O})$ を環付きサイトとする。$L, M$ を
$D(\mathcal{O})$ の対象とする。$\mathcal{C}$ の任意の対象 $U$ に対して
$$
H^0(U, R\SheafHom(L, M)) =
\Hom_{D(\mathcal{O}_U)}(L|_U, M|_U)
$$
が成り立ち、また $H^0(\mathcal{C}, R\SheafHom(L, M)) =
\Hom_{D(\mathcal{O})}(L, M)$ が成り立つ。
\end{lemma}

\begin{proof}
K-入射な $\mathcal{O}$-加群の複体 $\mathcal{I}^\bullet$ で $M$ を表し、
K-平坦な複体 $\mathcal{L}^\bullet$ で $L$ を表すように選ぶ。このとき
$\SheafHom^\bullet(\mathcal{L}^\bullet, \mathcal{I}^\bullet)$
は補題 \ref{sites-cohomology-lemma-RHom-from-K-flat-into-K-injective} により K-入射である。
したがって $U$ 上のコホモロジーは単に $U$ 上の切断を取ることで計算でき、
結論は補題 \ref{sites-cohomology-lemma-RHom-into-K-injective} から従う。
\end{proof}

\begin{lemma}
\label{sites-cohomology-lemma-internal-hom}
$(\mathcal{C}, \mathcal{O})$ を環付きサイトとする。$K, L, M$ を
$D(\mathcal{O})$ の対象とする。上で述べた構成により標準同型
$$
R\SheafHom(K, R\SheafHom(L, M)) =
R\SheafHom(K \otimes_\mathcal{O}^\mathbf{L} L, M)
$$
が $D(\mathcal{O})$ に存在し、$K, L, M$ に関して関手的である。
$H^0(\mathcal{C}, -)$ を取ると（\ref{sites-cohomology-equation-internal-hom}）が得られる。
\end{lemma}

\begin{proof}
$M$ を表す K-入射複体 $\mathcal{I}^\bullet$ と、$L$ を表す
K-平坦な $\mathcal{O}$-加群の複体 $\mathcal{L}^\bullet$ を選ぶ。
任意の $\mathcal{O}$-加群の複体 $\mathcal{K}^\bullet$ に対して
$$
\SheafHom^\bullet(\mathcal{K}^\bullet,
\SheafHom^\bullet(\mathcal{L}^\bullet, \mathcal{I}^\bullet))
=
\SheafHom^\bullet(
\text{Tot}(\mathcal{K}^\bullet \otimes_\mathcal{O} \mathcal{L}^\bullet),
\mathcal{I}^\bullet)
$$
が補題 \ref{sites-cohomology-lemma-compose} により成り立つ。左辺は
$R\SheafHom(K, R\SheafHom(L, M))$ を表し（補題
\ref{sites-cohomology-lemma-RHom-from-K-flat-into-K-injective} を用いる）、右辺は
$R\SheafHom(K \otimes_\mathcal{O}^\mathbf{L} L, M)$ を表すことに注意する。
これで補題の表示式が証明された。大域切断を取り、補題
\ref{sites-cohomology-lemma-section-RHom-over-U} を用いると
（\ref{sites-cohomology-equation-internal-hom}）を得る。
\end{proof}

\begin{lemma}
\label{sites-cohomology-lemma-restriction-RHom-to-U}
$(\mathcal{C}, \mathcal{O})$ を環付きサイトとする。$K, L$ を
$D(\mathcal{O})$ の対象とする。$R\SheafHom(K, L)$ の構成は制限と可換である。
すなわち、$\mathcal{C}$ の任意の対象 $U$ に対して
$R\SheafHom(K|_U, L|_U) = R\SheafHom(K, L)|_U$.
\end{lemma}

\begin{proof}
これは構成と補題 \ref{sites-cohomology-lemma-restrict-K-injective-to-open} から明らかである。
\end{proof}

\begin{lemma}
\label{sites-cohomology-lemma-RHom-triangulated}
$(\mathcal{C}, \mathcal{O})$ を環付きサイトとする。二変数関手
$R\SheafHom(- , -)$ は、どちらの変数についても特別三角形を
特別三角形へ移す。
\end{lemma}

\begin{proof}
これは、対応
$$
(\mathcal{L}^\bullet, \mathcal{M}^\bullet) \longmapsto
\SheafHom^\bullet(\mathcal{L}^\bullet, \mathcal{M}^\bullet)
$$
が、どちらの変数についても複体の各次数で分裂する短完全列を
各次数で分裂する短完全列へ移すことから従う。詳細は省略する。
\end{proof}

\begin{lemma}
\label{sites-cohomology-lemma-internal-hom-evaluate}
$(\mathcal{C}, \mathcal{O})$ を環付きサイトとする。$K, L, M$ を
$D(\mathcal{O})$ の対象とする。標準射
$$
R\SheafHom(L, M) \otimes_\mathcal{O}^\mathbf{L} K
\longrightarrow
R\SheafHom(R\SheafHom(K, L), M)
$$
が $D(\mathcal{O})$ に存在し、$K, L, M$ に関して関手的である。
\end{lemma}

\begin{proof}
$M$ を表す K-入射複体 $\mathcal{I}^\bullet$、$L$ を表す
K-入射複体 $\mathcal{J}^\bullet$、および $K$ を表す
K-平坦複体 $\mathcal{K}^\bullet$ を選ぶ。この写像は、補題
\ref{sites-cohomology-lemma-evaluate-and-more} の写像
$$
\text{Tot}(\SheafHom^\bullet(\mathcal{J}^\bullet,
\mathcal{I}^\bullet) \otimes_\mathcal{O} \mathcal{K}^\bullet)
\longrightarrow
\SheafHom^\bullet(\SheafHom^\bullet(\mathcal{K}^\bullet,
\mathcal{J}^\bullet), \mathcal{I}^\bullet)
$$
を用いて定義される。上記の複体の選び方により、左辺は
$R\SheafHom(L, M) \otimes_\mathcal{O}^\mathbf{L} K$ を表し、右辺は
$R\SheafHom(R\SheafHom(K, L), M)$ を表す。これが
$D(\mathcal{O})$ の三つの対象すべてに関して関手的であることの証明は省略する。
\end{proof}

\begin{lemma}
\label{sites-cohomology-lemma-internal-hom-composition}
\begin{slogan}
RSheafHom 上の合成。
\end{slogan}
$(\mathcal{C}, \mathcal{O})$ を環付きサイトとする。
$D(\mathcal{O})$ の $K, L, M$ が与えられたとき、標準射
$$
R\SheafHom(L, M) \otimes_\mathcal{O}^\mathbf{L} R\SheafHom(K, L)
\longrightarrow R\SheafHom(K, M)
$$
が $D(\mathcal{O})$ に存在する。
\end{lemma}

\begin{proof}
$M$ を表す K-入射複体 $\mathcal{I}^\bullet$、$L$ を表す
K-入射複体 $\mathcal{J}^\bullet$、および $K$ を表す任意の
$\mathcal{O}$-加群の複体 $\mathcal{K}^\bullet$ を選ぶ。補題
\ref{sites-cohomology-lemma-composition} により複体の写像
$$
\text{Tot}\left(
\SheafHom^\bullet(\mathcal{J}^\bullet, \mathcal{I}^\bullet)
\otimes_\mathcal{O}
\SheafHom^\bullet(\mathcal{K}^\bullet, \mathcal{J}^\bullet)
\right)
\longrightarrow
\SheafHom^\bullet(\mathcal{K}^\bullet, \mathcal{I}^\bullet)
$$
が存在する。$\mathcal{O}$-加群の複体
$\SheafHom^\bullet(\mathcal{J}^\bullet, \mathcal{I}^\bullet)$、
$\SheafHom^\bullet(\mathcal{K}^\bullet, \mathcal{J}^\bullet)$、および
$\SheafHom^\bullet(\mathcal{K}^\bullet, \mathcal{I}^\bullet)$
は、それぞれ $R\SheafHom(L, M)$、$R\SheafHom(K, L)$、および
$R\SheafHom(K, M)$ を表す。K-平坦複体 $\mathcal{H}^\bullet$ と擬同型
$\mathcal{H}^\bullet \to
\SheafHom^\bullet(\mathcal{K}^\bullet, \mathcal{J}^\bullet)$,
を選ぶと、写像
$$
\text{Tot}\left(
\SheafHom^\bullet(\mathcal{J}^\bullet, \mathcal{I}^\bullet)
\otimes_\mathcal{O} \mathcal{H}^\bullet
\right)
\longrightarrow
\text{Tot}\left(
\SheafHom^\bullet(\mathcal{J}^\bullet, \mathcal{I}^\bullet)
\otimes_\mathcal{O}
\SheafHom^\bullet(\mathcal{K}^\bullet, \mathcal{J}^\bullet)
\right)
$$
が存在し、その始域は
$R\SheafHom(L, M) \otimes_\mathcal{O}^\mathbf{L} R\SheafHom(K, L)$
を表す。表示された二つの矢を合成すると求める写像が得られる。
構成が関手的であることの証明は省略する。
\end{proof}

\begin{lemma}
\label{sites-cohomology-lemma-internal-hom-diagonal-better}
$(\mathcal{C}, \mathcal{O})$ を環付きサイトとする。
$D(\mathcal{O})$ の $K, L, M$ が与えられたとき、標準射
$$
K \otimes_\mathcal{O}^\mathbf{L} R\SheafHom(M, L)
\longrightarrow
R\SheafHom(M, K \otimes_\mathcal{O}^\mathbf{L} L)
$$
が $D(\mathcal{O})$ に存在し、$K, L, M$ に関して関手的である。
\end{lemma}

\begin{proof}
$K$ を表す K-平坦複体 $\mathcal{K}^\bullet$ と、$L$ を表す
K-入射複体 $\mathcal{I}^\bullet$ を選び、$M$ を表す任意の
$\mathcal{O}$-加群の複体 $\mathcal{M}^\bullet$ を選ぶ。擬同型
$\text{Tot}(\mathcal{K}^\bullet \otimes_{\mathcal{O}_X} \mathcal{I}^\bullet)
\to \mathcal{J}^\bullet$
を選ぶ。ただし $\mathcal{J}^\bullet$ は K-入射とする。このとき写像
$$
\text{Tot}\left(
\mathcal{K}^\bullet \otimes_\mathcal{O}
\SheafHom^\bullet(\mathcal{M}^\bullet, \mathcal{I}^\bullet)
\right)
\to
\SheafHom^\bullet(\mathcal{M}^\bullet,
\text{Tot}(\mathcal{K}^\bullet \otimes_\mathcal{O} \mathcal{I}^\bullet))
\to
\SheafHom^\bullet(\mathcal{M}^\bullet, \mathcal{J}^\bullet)
$$
を用いる。ここで第一の写像は補題 \ref{sites-cohomology-lemma-diagonal-better} の写像である。
\end{proof}

\begin{lemma}
\label{sites-cohomology-lemma-internal-hom-diagonal}
$(\mathcal{C}, \mathcal{O})$ を環付きサイトとする。
$D(\mathcal{O})$ の $K, L$ が与えられたとき、標準射
$$
K \longrightarrow R\SheafHom(L, K \otimes_\mathcal{O}^\mathbf{L} L)
$$
が $D(\mathcal{O})$ に存在し、$K$ と $L$ の双方に関して関手的である。
\end{lemma}

\begin{proof}
$K$ を表す K-平坦複体 $\mathcal{K}^\bullet$ と、$L$ を表す任意の
$\mathcal{O}$-加群の複体 $\mathcal{L}^\bullet$ を選ぶ。K-入射複体
$\mathcal{J}^\bullet$ と擬同型
$\text{Tot}(\mathcal{K}^\bullet \otimes_\mathcal{O} \mathcal{L}^\bullet)
\to \mathcal{J}^\bullet$ を選ぶ。このとき
$$
\mathcal{K}^\bullet \to
\SheafHom^\bullet(\mathcal{L}^\bullet,
\text{Tot}(\mathcal{K}^\bullet \otimes_\mathcal{O} \mathcal{L}^\bullet))
\to
\SheafHom^\bullet(\mathcal{L}^\bullet, \mathcal{J}^\bullet)
$$
を用いる。ここで第一の写像は補題 \ref{sites-cohomology-lemma-diagonal} に由来する。
\end{proof}

\begin{lemma}
\label{sites-cohomology-lemma-dual}
$(\mathcal{C}, \mathcal{O})$ を環付きサイトとする。$L$ を
$D(\mathcal{O})$ の対象とする。$L^\vee = R\SheafHom(L, \mathcal{O})$
と置く。$D(\mathcal{O})$ の $M$ に対して標準射
\begin{equation}
\label{sites-cohomology-equation-eval}
M \otimes^\mathbf{L}_\mathcal{O} L^\vee \longrightarrow R\SheafHom(L, M)
\end{equation}
が存在し、これは標準射
$$
H^0(\mathcal{C}, M \otimes_\mathcal{O}^\mathbf{L} L^\vee)
\longrightarrow
\Hom_{D(\mathcal{O})}(L, M)
$$
を誘導する。この射は $D(\mathcal{O})$ の $M$ に関して関手的である。
\end{lemma}

\begin{proof}
写像（\ref{sites-cohomology-equation-eval}）は、同一視
$M = R\SheafHom(\mathcal{O}, M)$ を用いた補題
\ref{sites-cohomology-lemma-internal-hom-composition} の特別な場合である。
\end{proof}

\begin{remark}
\label{sites-cohomology-remark-projection-formula-for-internal-hom}
$f : (\Sh(\mathcal{C}), \mathcal{O}_\mathcal{C}) \to
(\Sh(\mathcal{D}), \mathcal{O}_\mathcal{D})$ を環付きトポスの射とする。
$K, L$ を $D(\mathcal{O}_\mathcal{C})$ の対象とする。標準射
$$
Rf_*R\SheafHom(L, K) \longrightarrow R\SheafHom(Rf_*L, Rf_*K)
$$
が存在すると主張する。実際、（\ref{sites-cohomology-equation-internal-hom}）により、これは写像
$Rf_*R\SheafHom(L, K) \otimes_{\mathcal{O}_\mathcal{D}}^\mathbf{L} Rf_*L
\to Rf_*K$ と同じことである。このために合成
$$
Rf_*R\SheafHom(L, K) \otimes_{\mathcal{O}_\mathcal{D}}^\mathbf{L} Rf_*L \to
Rf_*(R\SheafHom(L, K) \otimes_{\mathcal{O}_\mathcal{C}}^\mathbf{L} L) \to
Rf_*K
$$
を用いることができる。ここで第一の矢は相対カップ積
（注 \ref{sites-cohomology-remark-cup-product}）であり、第二の矢は標準射
$R\SheafHom(L, K) \otimes_{\mathcal{O}_\mathcal{C}}^\mathbf{L} L \to K$
に $Rf_*$ を適用したものである。この標準射は補題
\ref{sites-cohomology-lemma-internal-hom-composition}（一つの箇所を
$\mathcal{O}_\mathcal{C}$ とする）から得られる。
\end{remark}

\begin{remark}
\label{sites-cohomology-remark-prepare-fancy-base-change}
$h : (\Sh(\mathcal{C}), \mathcal{O}) \to (\Sh(\mathcal{C}'), \mathcal{O}')$
を環付きトポスの射とする。$K, L$ を $D(\mathcal{O}')$ の対象とする。
標準射
$$
Lh^*R\SheafHom(K, L) \longrightarrow R\SheafHom(Lh^*K, Lh^*L)
$$
が $D(\mathcal{O})$ に存在すると主張する。実際、
（\ref{sites-cohomology-equation-internal-hom}）、すなわち補題
\ref{sites-cohomology-lemma-internal-hom} で証明した式により、
そのような写像は写像
$$
Lh^*R\SheafHom(K, L) \otimes^\mathbf{L} Lh^*K \longrightarrow Lh^*L
$$
と同じことである。この矢の始域は補題
\ref{sites-cohomology-lemma-pullback-tensor-product} により
$Lh^*(\SheafHom(K, L) \otimes^\mathbf{L} K)$ であるから、標準射
$$
R\SheafHom(K, L) \otimes^\mathbf{L} K \longrightarrow L.
$$
を構成すれば十分である。そのため、（\ref{sites-cohomology-equation-internal-hom}）を介して
$$
\text{id} :
R\SheafHom(K, L)
\longrightarrow
R\SheafHom(K, L)
$$
に対応する矢を取る。
\end{remark}

\begin{remark}
\label{sites-cohomology-remark-fancy-base-change}
環付きトポスの可換図式
$$
\xymatrix{
(\Sh(\mathcal{C}'), \mathcal{O}_{\mathcal{C}'})
\ar[r]_h \ar[d]_{f'} &
(\Sh(\mathcal{C}), \mathcal{O}_\mathcal{C}) \ar[d]^f \\
(\Sh(\mathcal{D}'), \mathcal{O}_{\mathcal{D}'})
\ar[r]^g &
(\Sh(\mathcal{D}), \mathcal{O}_\mathcal{D})
}
$$
が与えられているとする。$K, L$ を $D(\mathcal{O}_\mathcal{C})$ の
対象とする。標準的な基底変換写像
$$
Lg^*Rf_*R\SheafHom(K, L)
\longrightarrow
R(f')_*R\SheafHom(Lh^*K, Lh^*L)
$$
が $D(\mathcal{O}_{\mathcal{D}'})$ に存在すると主張する。実際、合成
\begin{align*}
L(f')^*Lg^*Rf_*R\SheafHom(K, L)
& =
Lh^*Lf^*Rf_*R\SheafHom(K, L) \\
& \to
Lh^*R\SheafHom(K, L) \\
& \to
R\SheafHom(Lh^*K, Lh^*L)
\end{align*}
に随伴する写像を取る。ここで第一の矢は随伴写像
$Lf^*Rf_* \to \text{id}$ を用い、第二の矢は注
\ref{sites-cohomology-remark-prepare-fancy-base-change} で構成した標準射である。
\end{remark}




\section{大域導来 Hom}
\label{sites-cohomology-section-global-RHom}

\noindent
$(\Sh(\mathcal{C}), \mathcal{O})$ を環付きトポスとする。
$K, L \in D(\mathcal{O})$ とする。導来圏における内部 Hom の構成を用いると、
良定義な対象
$$
R\Hom_\mathcal{O}(K, L) = R\Gamma(\mathcal{C}, R\SheafHom(K, L))
$$
を $D(\Gamma(\mathcal{C}, \mathcal{O}))$ に得る。補題
\ref{sites-cohomology-lemma-section-RHom-over-U} により
$$
H^0(R\Hom_\mathcal{O}(K, L)) = \Hom_{D(\mathcal{O})}(K, L)
$$
および
$$
H^p(R\Hom_\mathcal{O}(K, L)) = \Ext_{D(\mathcal{O})}^p(K, L)
$$
が成り立つ。$f : (\mathcal{C}', \mathcal{O}') \to
(\mathcal{C}, \mathcal{O})$ が環付きトポスの射ならば、標準射
$$
R\Hom_\mathcal{O}(K, L) \longrightarrow R\Hom_{\mathcal{O}'}(Lf^*K, Lf^*L)
$$
が $D(\Gamma(\mathcal{O}))$ に存在する。これは注
\ref{sites-cohomology-remark-prepare-fancy-base-change} で定義した写像の大域切断を取って得られる。









\section{導来関手 $Lg_!$}
\label{sites-cohomology-section-derived-lower-shriek}

\noindent
本節では、$Lg^*$ と $Rg_*$ に加えて導来関手 $Lg_!$ も存在するような
環付きトポスの射 $g$ を調べる。

\begin{lemma}
\label{sites-cohomology-lemma-pullback-injective-pre-limp}
サイトの連続かつ余連続な関手
$u : \mathcal{C} \to \mathcal{D}$ を取る。$g : \Sh(\mathcal{C}) \to
\Sh(\mathcal{D})$ を対応するトポスの射とする。
$\mathcal{O}_\mathcal{D}$ を環の層、$\mathcal{I}$ を入射的
$\mathcal{O}_\mathcal{D}$-加群とする。このとき
$H^p(U, g^{-1}\mathcal{I}) = 0$ がすべての $p > 0$ と
$U \in \Ob(\mathcal{C})$ に対して成り立つ。
\end{lemma}

\begin{proof}
任意の $\mathcal{C}$ の被覆
$\mathcal{U} = \{U_i \to U\}$ に対して、すべての高次
{\v C}ech コホモロジー群
$\check H^p(\mathcal{U}, g^{-1}\mathcal{I})$
の消滅を証明できれば、補題 \ref{sites-cohomology-lemma-cech-vanish-collection} から
補題の消滅が従う。$u$ は連続なので、
$u(\mathcal{U}) = \{u(U_i) \to u(U)\}$ は $\mathcal{D}$ の被覆であり、
$u(U_{i_0} \times_U \ldots \times_U U_{i_n}) =
u(U_{i_0}) \times_{u(U)} \ldots \times_{u(U)} u(U_{i_n})$.
したがって
$$
\check H^p(\mathcal{U}, g^{-1}\mathcal{I}) =
\check H^p(u(\mathcal{U}), \mathcal{I})
$$
が成り立つ。これは Sites, Lemma \ref{sites-lemma-when-shriek} により
$g^{-1} = u^p$ だからである。$\mathcal{I}$ は入射的
$\mathcal{O}_\mathcal{D}$-加群なので、これらの {\v C}ech
コホモロジー群は消滅する。補題
\ref{sites-cohomology-lemma-injective-module-trivial-cech} を参照せよ。
\end{proof}

\begin{lemma}
\label{sites-cohomology-lemma-existence-derived-lower-shriek}
サイトの連続かつ余連続な関手
$u : \mathcal{C} \to \mathcal{D}$ を取る。$g : \Sh(\mathcal{C}) \to
\Sh(\mathcal{D})$ を対応するトポスの射とする。
$\mathcal{O}_\mathcal{D}$ を環の層とし、
$\mathcal{O}_\mathcal{C} = g^{-1}\mathcal{O}_\mathcal{D}$.
と置く。関手 $g_! : \textit{Mod}(\mathcal{O}_\mathcal{C}) \to
\textit{Mod}(\mathcal{O}_\mathcal{D})$
（Modules on Sites, Lemma
\ref{sites-modules-lemma-lower-shriek-modules} 参照）は左導来関手
$$
Lg_! : D(\mathcal{O}_\mathcal{C}) \longrightarrow D(\mathcal{O}_\mathcal{D})
$$
をもち、これは $g^*$ の左随伴である。さらに $U \in \Ob(\mathcal{C})$
に対して
$$
Lg_!(j_{U!}\mathcal{O}_U) =
g_!j_{U!}\mathcal{O}_U =
j_{u(U)!} \mathcal{O}_{u(U)}.
$$
が成り立つ。ここで $j_{U!}$ と $j_{u(U)!}$ は、局所化射
$j_U : \mathcal{C}/U \to \mathcal{C}$ および
$j_{u(U)} : \mathcal{D}/u(U) \to \mathcal{D}$ に付随する零による拡張である。
\end{lemma}

\begin{proof}
$Lg_!$ の構成には Derived Categories, Proposition
\ref{derived-proposition-left-derived-exists} を用いる。そのため、この命題の仮定
(1)、(2)、(3)、(4)、(5) を検証しなければならない。まず $g_!$ は左随伴なので
右完全であり、すべての余極限と可換する。したがって (5) が成り立つ。
環付きサイト上の加群の圏はグロタンディーク・アーベル圏なので、条件 (3) と (4)
も成り立つ。$\mathcal{P} \subset
\Ob(\textit{Mod}(\mathcal{O}_\mathcal{C}))$ を、$j_{U!}\mathcal{O}_U$
の形の加群の直和である $\mathcal{O}_\mathcal{C}$-加群全体とする。
$g_!j_{U!}\mathcal{O}_U = j_{u(U)!} \mathcal{O}_{u(U)}$ であることに注意する。
Modules on Sites, Lemma
\ref{sites-modules-lemma-lower-shriek-modules} の証明を参照せよ。
任意の $\mathcal{O}_\mathcal{C}$-加群は $\mathcal{P}$ の対象の商である。
Modules on Sites, Lemma
\ref{sites-modules-lemma-module-quotient-flat} を参照せよ。よって (1) が成り立つ。
最後に (2) を証明する。
$\mathcal{K}^\bullet$ を、すべての $n$ に対して
$\mathcal{K}^n \in \mathcal{P}$ を満たす、上に有界な非輪状な
$\mathcal{O}_\mathcal{C}$-加群の複体とする。
$g_!\mathcal{K}^\bullet$ が完全であることを示さなければならない。
そのためには、任意の入射的 $\mathcal{O}_\mathcal{D}$-加群
$\mathcal{I}$ に対して
$$
\Hom_{D(\mathcal{O}_\mathcal{D})}(
g_!\mathcal{K}^\bullet, \mathcal{I}[n]) = 0
$$
がすべての $n \in \mathbf{Z}$ について成り立つことを示せば十分である。
$\mathcal{I}$ は入射的なので
\begin{align*}
\Hom_{D(\mathcal{O}_\mathcal{D})}(
g_!\mathcal{K}^\bullet, \mathcal{I}[n])
& =
\Hom_{K(\mathcal{O}_\mathcal{D})}(
g_!\mathcal{K}^\bullet, \mathcal{I}[n]) \\
& =
H^n(\Hom_{\mathcal{O}_\mathcal{D}}(
g_!\mathcal{K}^\bullet, \mathcal{I})) \\
& =
H^n(\Hom_{\mathcal{O}_\mathcal{C}}(
\mathcal{K}^\bullet, g^{-1}\mathcal{I}))
\end{align*}
が成り立つ。最後の等号は $g_!$ と $g^{-1}$ の随伴性による。

\medskip\noindent
$g^{-1}\mathcal{I}$ が入射的 $\mathcal{O}_\mathcal{C}$-加群ならば、
この群の消滅は明らかであろう。しかし $g_!$ は一般には完全でないので、
$g^{-1}\mathcal{I}$ は必ずしも入射的 $\mathcal{O}_\mathcal{C}$-加群ではない。
一方、
$$
\Ext^p_{\mathcal{O}_\mathcal{C}}(
j_{U!}\mathcal{O}_U, g^{-1}\mathcal{I}) =
H^p(U, g^{-1}\mathcal{I}) = 0 \text{ for }p \geq 1
$$
である。ここで第一の等号は
$\Hom_{\mathcal{O}_\mathcal{C}}(j_{U!}\mathcal{O}_U, \mathcal{H}) =
\mathcal{H}(U)$ に導来関手を取ることから従い、$p > 0$ および
$U \in \Ob(\mathcal{C})$ に対する
$H^p(U, g^{-1}\mathcal{I})$ の消滅は補題
\ref{sites-cohomology-lemma-pullback-injective-pre-limp} から従う。各 $\mathcal{K}^{-q}$ は
$j_{U!}\mathcal{O}_U$ の形の加群の直和なので、
$$
\Ext^p_{\mathcal{O}_\mathcal{C}}(\mathcal{K}^{-q}, g^{-1}\mathcal{I}) = 0
\text{ for }p \geq 1\text{ and all }q
$$
が分かる。スペクトル系列（例
\ref{sites-cohomology-example-hom-complex-into-sheaf} 参照）
$$
E_1^{p, q} = \Ext^q_{\mathcal{O}_\mathcal{C}}(
\mathcal{K}^{-p}, g^{-1}\mathcal{I})
\Rightarrow
\Ext^{p + q}_{\mathcal{O}_\mathcal{C}}(
\mathcal{K}^\bullet, g^{-1}\mathcal{I}) = 0.
$$
を用いる。$\mathcal{K}^\bullet$ は非輪状なので導来圏で零となり、したがって
Ext 群も消滅するから、このスペクトル系列は零へ収束することに注意する。
上で証明した高次 Ext の消滅により、$E_1$ 頁で零でない項は
$E_1^{p, 0} = \Hom_{\mathcal{O}_\mathcal{C}}(
\mathcal{K}^{-p}, g^{-1}\mathcal{I})$ のみである。したがって複体
$\Hom_{\mathcal{O}_\mathcal{C}}(
\mathcal{K}^\bullet, g^{-1}\mathcal{I})$
は、望んだとおり非輪状である。

\medskip\noindent
以上により左導来関手 $Lg_!$ が存在する。これは
$g^{-1} = g^* = Rg^* = Lg^*$ の左随伴である。すなわち
\begin{equation}
\label{sites-cohomology-equation-to-prove}
\Hom_{D(\mathcal{O}_\mathcal{C})}(K, g^*L) =
\Hom_{D(\mathcal{O}_\mathcal{D})}(Lg_!K, L)
\end{equation}
が Derived Categories, Lemma
\ref{derived-lemma-derived-adjoint-functors} により成り立つ。これで証明は完了する。
\end{proof}

\begin{remark}
\label{sites-cohomology-remark-when-derived-shriek-equal}
注意！$u : \mathcal{C} \to \mathcal{D}$、$g$、$\mathcal{O}_\mathcal{D}$、
$\mathcal{O}_\mathcal{C}$ を補題
\ref{sites-cohomology-lemma-existence-derived-lower-shriek} のものとする。一般には図式
$$
\xymatrix{
D(\mathcal{O}_\mathcal{C}) \ar[r]_{Lg_!} \ar[d]_{forget} &
D(\mathcal{O}_\mathcal{D}) \ar[d]^{forget} \\
D(\mathcal{C}) \ar[r]^{Lg^{Ab}_!} &
D(\mathcal{D})
}
$$
は可換ではない。ここで関手 $Lg_!^{Ab}$ は、$\mathcal{C}$ と
$\mathcal{D}$ の双方で定数層 $\mathbf{Z}$ を構造層として用いて補題
\ref{sites-cohomology-lemma-existence-derived-lower-shriek} で構成したものである。一般には
$g_! = g_!^{Ab}$ でさえない（Modules on Sites, Remark
\ref{sites-modules-remark-when-shriek-equal} 参照）。しかしこの現象は
{\bf $g_! = g_!^{Ab}$ であっても起こりうる}！実際、補題
\ref{sites-cohomology-lemma-existence-derived-lower-shriek} の証明における $Lg_!$ の構成から、
$Lg_!$ が $Lg_!^{\textit{Ab}}$ と一致するための必要十分条件は、標準射
$$
Lg^{Ab}_!j_{U!}\mathcal{O}_U \longrightarrow j_{u(U)!}\mathcal{O}_{u(U)}
$$
が $\mathcal{C}$ のすべての対象 $U$ に対して $D(\mathcal{D})$ の同型である
ことだと分かる。一般に言えるのは、自然変換
$$
Lg_!^{Ab} \circ forget \longrightarrow forget \circ Lg_!
$$
が存在することだけである。
\end{remark}

\begin{lemma}
\label{sites-cohomology-lemma-pullback-injective-limp}
サイトの連続かつ余連続な関手
$u : \mathcal{C} \to \mathcal{D}$ を取る。$g : \Sh(\mathcal{C}) \to
\Sh(\mathcal{D})$ を対応するトポスの射とする。
$\mathcal{O}_\mathcal{D}$ を環の層、$\mathcal{I}$ を入射的
$\mathcal{O}_\mathcal{D}$-加群とする。もし
$g_!^{Sh} : \Sh(\mathcal{C}) \to \Sh(\mathcal{D})$
がファイバー積と可換するならば\footnote{$\mathcal{C}$ が有限連結極限をもち、
$u$ がそれらと可換するときに成り立つ。Sites, Lemma
\ref{sites-lemma-preserve-equalizers} を参照せよ。}、
$g^{-1}\mathcal{I}$ は完全非輪状である。
\end{lemma}

\begin{proof}
補題 \ref{sites-cohomology-lemma-characterize-limp} の判定条件を用いる。条件 (1) は補題
\ref{sites-cohomology-lemma-pullback-injective-pre-limp} により成り立つ。
$K' \to K$ を $\mathcal{C}$ 上の集合の層の全射とする。
$g_!^{Sh}$ は左随伴なので、$g_!^{Sh}K' \to g_!^{Sh}K$ は全射である。
次に注意する：
\begin{align*}
H^0(K' \times_K \ldots \times_K K', g^{-1}\mathcal{I}) 
& =
H^0(g_!^{Sh}(K' \times_K \ldots \times_K K'), \mathcal{I}) \\
& =
H^0(g_!^{Sh}K' \times_{g_!^{Sh}K} \ldots \times_{g_!^{Sh}K} g_!^{Sh}K',
\mathcal{I})
\end{align*}
これは $g_!^{Sh}$ に関する仮定による。$\mathcal{I}$ は入射的加群なので、
補題 \ref{sites-cohomology-lemma-direct-image-injective-sheaf} を恒等射に適用すれば完全非輪状である。
したがって補題 \ref{sites-cohomology-lemma-characterize-limp} の逆を用いると、複体
$$
0 \to H^0(K, g^{-1}\mathcal{I}) \to H^0(K', g^{-1}\mathcal{I}) \to
H^0(K' \times_K K', g^{-1}\mathcal{I}) \to \ldots
$$
は望んだとおり完全である。
\end{proof}

\begin{lemma}
\label{sites-cohomology-lemma-pullback-same-cohomology}
サイトの連続かつ余連続な関手
$u : \mathcal{C} \to \mathcal{D}$ を取る。$g : \Sh(\mathcal{C}) \to
\Sh(\mathcal{D})$ を対応するトポスの射とする。
$U \in \Ob(\mathcal{C})$ とする。
\begin{enumerate}
\item $D(\mathcal{D})$ の $M$ に対して
$R\Gamma(U, g^{-1}M) = R\Gamma(u(U), M)$.
\item $\mathcal{O}_\mathcal{D}$ が環の層であり、
$\mathcal{O}_\mathcal{C} = g^{-1}\mathcal{O}_\mathcal{D}$
であるなら、$D(\mathcal{O}_\mathcal{D})$ の $M$ に対して
$R\Gamma(U, g^*M) = R\Gamma(u(U), M)$.
\end{enumerate}
\end{lemma}

\begin{proof}
下に有界な場合、(1) と (2) は $K$ を下に有界な入射的対象の複体で表し、
補題 \ref{sites-cohomology-lemma-pullback-injective-pre-limp} と Leray の非輪状性補題を
用いれば分かる。非有界の場合、まず (1) が (2) の特別な場合であることに
注意する。(2) については
$$
R\Gamma(U, g^*M) =
R\Hom_{\mathcal{O}_\mathcal{C}}(j_{U!}\mathcal{O}_U, g^*M) =
R\Hom_{\mathcal{O}_\mathcal{D}}(j_{u(U)!}\mathcal{O}_{u(U)}, M) =
R\Gamma(u(U), M)
$$
を用いることができる。中央の等号は補題
\ref{sites-cohomology-lemma-existence-derived-lower-shriek} の帰結である。
\end{proof}

\begin{lemma}
\label{sites-cohomology-lemma-special-square-cocontinuous}
環付きトポスの可換図式
$$
\xymatrix{
(\Sh(\mathcal{C}'), \mathcal{O}_{\mathcal{C}'})
\ar[r]_{(g', (g')^\sharp)} \ar[d]_{(f', (f')^\sharp)} &
(\Sh(\mathcal{C}), \mathcal{O}_\mathcal{C}) \ar[d]^{(f, f^\sharp)} \\
(\Sh(\mathcal{D}'), \mathcal{O}_{\mathcal{D}'}) \ar[r]^{(g, g^\sharp)} &
(\Sh(\mathcal{D}), \mathcal{O}_\mathcal{D})
}
$$
が与えられているとする。次を仮定する。
\begin{enumerate}
\item $f$、$f'$、$g$、$g'$ は、Sites, Lemma
\ref{sites-lemma-cocontinuous-morphism-topoi} における余連続関手
$u$、$u'$、$v$、$v'$ にそれぞれ対応する。
\item $v \circ u' = u \circ v'$,
\item $v$ と $v'$ は連続かつ余連続である。
\item $\mathcal{D}'$ の任意の対象 $V'$ に対して、$v$ が与える関手
${}^{u'}_{V'}\mathcal{I} \to {}^{\ \ \ u}_{v(V')}\mathcal{I}$
は共終である。
\item $g^{-1}\mathcal{O}_{\mathcal{D}} = \mathcal{O}_{\mathcal{D}'}$
かつ $(g')^{-1}\mathcal{O}_{\mathcal{C}} = \mathcal{O}_{\mathcal{C}'}$ である。
\item $g'_! : \textit{Ab}(\mathcal{C}') \to \textit{Ab}(\mathcal{C})$
は完全である\footnote{$\mathcal{C}'$ にファイバー積と等化子が存在し、
$v'$ がそれらと可換するときに成り立つ。Modules on Sites, Lemma
\ref{sites-modules-lemma-exactness-lower-shriek} を参照せよ。}。
\end{enumerate}
このとき、関手 $D(\mathcal{O}_\mathcal{C}) \to
D(\mathcal{O}_{\mathcal{D}'})$ として
$Rf'_* \circ (g')^* = g^* \circ Rf_*$ である。
\end{lemma}

\begin{proof}
条件 (5) により $g^* = Lg^* = g^{-1}$ かつ
$(g')^* = L(g')^* = (g')^{-1}$ である。補題
\ref{sites-cohomology-lemma-modules-abelian-unbounded} により、アーベル層の導来圏
$D(\mathcal{C})$ 上で結果を証明すれば十分である。
$K \in D(\mathcal{C})$ を取り、$K$ を表す $\mathcal{C}$ 上の
アーベル層の K-入射複体 $\mathcal{I}^\bullet$ を取る。Derived Categories,
Lemma \ref{derived-lemma-adjoint-preserve-K-injectives} と仮定 (6) により、
$(g')^{-1}\mathcal{I}^\bullet$ は $\mathcal{C}'$ 上のアーベル層の
K-入射複体である。Modules on Sites, Lemma
\ref{sites-modules-lemma-special-square-cocontinuous} により
$f'_*(g')^{-1}\mathcal{I}^\bullet = g^{-1}f_*\mathcal{I}^\bullet$ を得る。
$f_*\mathcal{I}^\bullet$ は $Rf_*K$ を表し、
$f'_*(g')^{-1}\mathcal{I}^\bullet$ は $Rf'_*(g')^{-1}K$ を表すので、結論が従う。
\end{proof}

\begin{lemma}
\label{sites-cohomology-lemma-special-square-continuous}
環付きトポスの可換図式
$$
\xymatrix{
(\Sh(\mathcal{C}'), \mathcal{O}_{\mathcal{C}'})
\ar[r]_{(g', (g')^\sharp)} \ar[d]_{(f', (f')^\sharp)} &
(\Sh(\mathcal{C}), \mathcal{O}_\mathcal{C}) \ar[d]^{(f, f^\sharp)} \\
(\Sh(\mathcal{D}'), \mathcal{O}_{\mathcal{D}'}) \ar[r]^{(g, g^\sharp)} &
(\Sh(\mathcal{D}), \mathcal{O}_\mathcal{D})
}
$$
と関手の図式
$$
\xymatrix{
\mathcal{C}' \ar[r]_{v'} &
\mathcal{C} \\
\mathcal{D}' \ar[r]^v \ar[u]^{u'} &
\mathcal{D} \ar[u]_u
}
$$
が与えられているとする。Sites, Sections
\ref{sites-section-morphism-sites} および
\ref{sites-section-cocontinuous-morphism-topoi} の記法の下で、次を仮定する。
\begin{enumerate}
\item $u$ と $u'$ は連続であり、射 $f$ と $f'$ をそれぞれ与える。
\item $v$ と $v'$ は余連続であり、射 $g$ と $g'$ をそれぞれ与える。
\item $u \circ v = v' \circ u'$,
\item $v$ と $v'$ は連続かつ余連続である。
\item $g^{-1}\mathcal{O}_{\mathcal{D}} = \mathcal{O}_{\mathcal{D}'}$
かつ $(g')^{-1}\mathcal{O}_{\mathcal{C}} = \mathcal{O}_{\mathcal{C}'}$ である。
\end{enumerate}
このとき、関手 $D^+(\mathcal{O}_\mathcal{C}) \to
D^+(\mathcal{O}_{\mathcal{D}'})$ として
$Rf'_* \circ (g')^* = g^* \circ Rf_*$ である。さらに
\begin{enumerate}
\item[(6)] $g'_! : \textit{Ab}(\mathcal{C}') \to \textit{Ab}(\mathcal{C})$
が完全である\footnote{$\mathcal{C}'$ にファイバー積と等化子が存在し、
$v'$ がそれらと可換するときに成り立つ。Modules on Sites, Lemma
\ref{sites-modules-lemma-exactness-lower-shriek} を参照せよ。}、
\end{enumerate}
ならば、関手 $D(\mathcal{O}_\mathcal{C}) \to
D(\mathcal{O}_{\mathcal{D}'})$ として
$Rf'_* \circ (g')^* = g^* \circ Rf_*$ である。
\end{lemma}

\begin{proof}
条件 (5) により $g^* = Lg^* = g^{-1}$ かつ
$(g')^* = L(g')^* = (g')^{-1}$ である。補題
\ref{sites-cohomology-lemma-modules-abelian-unbounded} により、アーベル層の導来圏
$D^+(\mathcal{C})$ または $D(\mathcal{C})$ 上で結果を証明すれば十分である。

\medskip\noindent
$K \in D^+(\mathcal{C})$ を取る。$K$ を表す、$\mathcal{C}$ 上の
入射的アーベル層からなる下に有界な複体 $\mathcal{I}^\bullet$ を取る。
補題 \ref{sites-cohomology-lemma-pullback-injective-pre-limp} により、任意の $p > 0$、
任意の $q$、任意の $U' \in \Ob(\mathcal{C}')$ に対して
$H^p(U', (g')^{-1}\mathcal{I}^q) = 0$ である。補題
\ref{sites-cohomology-lemma-higher-direct-images} により、
$R^pf'_*(g')^{-1}\mathcal{I}^q$ は前層
$V' \mapsto H^p(u'(V'), (g')^{-1}\mathcal{I}^q)$ に付随する層であることを
想起する。したがって $(g')^{-1}\mathcal{I}^q$ は関手 $f'_*$ に関して
右非輪状である。Leray の非輪状性補題（Derived Categories, Lemma
\ref{derived-lemma-leray-acyclicity}）により、
$f'_*(g')^*\mathcal{I}^\bullet$ は $Rf'_*(g')^{-1}K$ を表す。
Modules on Sites, Lemma
\ref{sites-modules-lemma-special-square-continuous} により
$f'_*(g')^{-1}\mathcal{I}^\bullet = g^{-1}f_*\mathcal{I}^\bullet$ を得る。
$g^{-1}f_*\mathcal{I}^\bullet$ は $g^{-1}Rf_*K$ を表すので、結論が従う。

\medskip\noindent
$K \in D(\mathcal{C})$ を取る。$K$ を表す $\mathcal{C}$ 上の
アーベル層の K-入射複体 $\mathcal{I}^\bullet$ を取る。Derived Categories,
Lemma \ref{derived-lemma-adjoint-preserve-K-injectives} と仮定 (6) により、
$(g')^{-1}\mathcal{I}^\bullet$ は $\mathcal{C}'$ 上のアーベル層の
K-入射複体である。Modules on Sites, Lemma
\ref{sites-modules-lemma-special-square-continuous} により
$f'_*(g')^{-1}\mathcal{I}^\bullet = g^{-1}f_*\mathcal{I}^\bullet$ を得る。
$f_*\mathcal{I}^\bullet$ は $Rf_*K$ を表し、
$f'_*(g')^{-1}\mathcal{I}^\bullet$ は $Rf'_*(g')^{-1}K$ を表すので、結論が従う。
\end{proof}








\section{ファイバー圏に対する導来関手 $Lg_!$}
\label{sites-cohomology-section-derived-lower-shriek-fibred}

\noindent
本節では、節 \ref{sites-cohomology-section-derived-lower-shriek} で論じた状況のいくつかの
特別な場合を詳しく調べる。加群に対する関手 $g_!$ とアーベル群の層に対する
関手 $g_!$ が一致することを確認する。初読では本節を飛ばすことを勧める。

\begin{situation}
\label{sites-cohomology-situation-fibred-category}
ここで $(\mathcal{D}, \mathcal{O}_\mathcal{D})$ は環付きサイトであり、
$p : \mathcal{C} \to \mathcal{D}$ はファイバー圏である。
$\mathcal{C}$ には $\mathcal{D}$ から誘導される位相を入れる
（Stacks, Section \ref{stacks-section-topology}）。$p$ に付随するトポスの射を
$\pi : \Sh(\mathcal{C}) \to \Sh(\mathcal{D})$ と書く
（Stacks, Lemma \ref{stacks-lemma-topology-inherited-functorial}）。
$\mathcal{O}_\mathcal{C} = \pi^{-1}\mathcal{O}_\mathcal{D}$ と置くと、環付きトポスの射
$$
\pi :
(\Sh(\mathcal{C}), \mathcal{O}_\mathcal{C})
\longrightarrow
(\Sh(\mathcal{D}), \mathcal{O}_\mathcal{D})
$$
を得る。
\end{situation}

\begin{lemma}
\label{sites-cohomology-lemma-fibred-category-with-object}
状況 \ref{sites-cohomology-situation-fibred-category} の仮定と記法を用いる。
$U \in \Ob(\mathcal{C})$ に対し、誘導されるトポスの射
$$
\pi_U : \Sh(\mathcal{C}/U) \longrightarrow \Sh(\mathcal{D}/p(U))
$$
を考える。このときトポスの射
$$
\sigma : \Sh(\mathcal{D}/p(U)) \to \Sh(\mathcal{C}/U)
$$
が存在し、$\pi_U \circ \sigma = \text{id}$ かつ
$\sigma^{-1} = \pi_{U, *}$ である。
\end{lemma}

\begin{proof}
$\pi_U$ は $\pi$ を局所化へ制限したものであることに注意する。Sites, Lemma
\ref{sites-lemma-localize-cocontinuous} を参照せよ。
$\mathcal{D}/p(U)$ の対象 $V \to p(U)$ に対し、$p$ がファイバー圏であるため
存在する、$\mathcal{D}$ 上の $\mathcal{C}$ の強デカルト射を
$V \times_{p(U)} U \to U$ と書く。関手
$$
v : \mathcal{D}/p(U) \to \mathcal{C}/U,\quad
V/p(U) \mapsto V \times_{p(U)} U/U
$$
は $\mathcal{C}$ の位相の定義により連続である。さらに、強デカルト射の定義により
$p$ の右随伴である。したがって Sites, Section
\ref{sites-section-cocontinuous-adjoint} で論じた状況にあり、層
$\pi_{U, *}\mathcal{F}$ は $V \mapsto
\mathcal{F}(V \times_{p(U)} U)$ に等しいことが分かる（特に Sites, Lemma
\ref{sites-lemma-have-functor-other-way-morphism} を参照せよ）。

\medskip\noindent
ここではさらに、関手 $v$ は余連続でもある（$\mathcal{C}$ の被覆に現れる射は
すべて強デカルトだからである）。したがって $v$ は補題に記した射 $\sigma$ を
定める。等式 $\sigma^{-1} = \pi_{U, *}$ は定義から直ちに従う。
$\pi_U^{-1}\mathcal{G}$ は規則
$U'/U \mapsto \mathcal{G}(p(U')/p(U))$ で与えられるので、
$\sigma^{-1} \circ \pi_U^{-1} = \text{id}$ が従い、これにより
$\pi_U \circ \sigma = \text{id}$ が証明される。
\end{proof}

\begin{situation}
\label{sites-cohomology-situation-morphism-fibred-categories}
$(\mathcal{D}, \mathcal{O}_\mathcal{D})$ を環付きサイトとする。
$u : \mathcal{C}' \to \mathcal{C}$ を $\mathcal{D}$ 上のファイバー圏の
$1$-射とする（Categories, Definition
\ref{categories-definition-fibred-categories-over-C}）。
$\mathcal{C}$ と $\mathcal{C}'$ にそれぞれの誘導位相を入れ
（Stacks, Definition \ref{stacks-definition-topology-inherited}）、
$\pi : \Sh(\mathcal{C}) \to \Sh(\mathcal{D})$、
$\pi' : \Sh(\mathcal{C}') \to \Sh(\mathcal{D})$、および
$g : \Sh(\mathcal{C}') \to \Sh(\mathcal{C})$
を対応するトポスの射とする（Stacks, Lemma
\ref{stacks-lemma-topology-inherited-functorial}）。
$\mathcal{O}_\mathcal{C} = \pi^{-1}\mathcal{O}_\mathcal{D}$ および
$\mathcal{O}_{\mathcal{C}'} = (\pi')^{-1}\mathcal{O}_\mathcal{D}$ と置く。
$g^{-1}\mathcal{O}_\mathcal{C} = \mathcal{O}_{\mathcal{C}'}$ であるから、
$$
\xymatrix{
(\Sh(\mathcal{C}'), \mathcal{O}_{\mathcal{C}'}) \ar[rd]_{\pi'} \ar[rr]_g & &
(\Sh(\mathcal{C}), \mathcal{O}_\mathcal{C}) \ar[ld]^\pi \\
& (\Sh(\mathcal{D}), \mathcal{O}_\mathcal{D})
}
$$
は環付きトポスの射の可換図式である。
\end{situation}

\begin{lemma}
\label{sites-cohomology-lemma-morphism-fibred-categories-with-object}
状況 \ref{sites-cohomology-situation-morphism-fibred-categories} の仮定と記法を用いる。
$U' \in \Ob(\mathcal{C}')$ に対し $U = u(U')$、$V = p'(U')$ と置き、
誘導される環付きトポスの射
$$
\xymatrix{
(\Sh(\mathcal{C}'/U'), \mathcal{O}_{U'}) \ar[rd]_{\pi'_{U'}} \ar[rr]_{g'} & &
(\Sh(\mathcal{C}), \mathcal{O}_U) \ar[ld]^{\pi_U} \\
& (\Sh(\mathcal{D}/V), \mathcal{O}_V)
}
$$
を考える。このときトポスの射
$$
\sigma' : \Sh(\mathcal{D}/V) \to \Sh(\mathcal{C}'/U'),
$$
が存在する。$\sigma = g' \circ \sigma'$ と置けば、
$\pi'_{U'} \circ \sigma' = \text{id}$、$\pi_U \circ \sigma = \text{id}$、
$(\sigma')^{-1} = \pi'_{U', *}$、および $\sigma^{-1} = \pi_{U, *}$ である。
\end{lemma}

\begin{proof}
$v' : \mathcal{D}/V \to \mathcal{C}'/U'$ を、$p' : \mathcal{C}' \to
\mathcal{D}'$ と対象 $U'$ から出発して補題
\ref{sites-cohomology-lemma-fibred-category-with-object} の証明で構成した関手とする。
$u$ は $\mathcal{D}$ 上のファイバー圏の $1$-射なので、強デカルト射を
強デカルト射へ移す。したがって関手 $v = u \circ v'$ は、
$p : \mathcal{C} \to \mathcal{D}$ と $U$ に関する補題
\ref{sites-cohomology-lemma-fibred-category-with-object} の証明の関手である。よって本補題は
その補題から従う。
\end{proof}

\begin{lemma}
\label{sites-cohomology-lemma-properties-lower-shriek-fibred-category}
状況 \ref{sites-cohomology-situation-morphism-fibred-categories} の仮定と記法を用いる。
\begin{enumerate}
\item 加群およびアーベル層上の $g^* = g^{-1}$ は、それぞれ左随伴
$g_! : \textit{Mod}(\mathcal{O}_{\mathcal{C}'}) \to
\textit{Mod}(\mathcal{O}_\mathcal{C})$ および
$g_!^{\textit{Ab}} : \textit{Ab}(\mathcal{C}') \to \textit{Ab}(\mathcal{C})$
をもつ。
\item 図式
$$
\xymatrix{
\textit{Mod}(\mathcal{O}_{\mathcal{C}'}) \ar[d] \ar[r]_{g_!} &
\textit{Mod}(\mathcal{O}_\mathcal{C}) \ar[d] \\
\textit{Ab}(\mathcal{C}') \ar[r]^{g_!^{\textit{Ab}}} &
\textit{Ab}(\mathcal{C})
}
$$
は可換である。
\item 加群およびアーベル層の導来圏上の $g^* = g^{-1}$ は、それぞれ左随伴
$Lg_! : D(\mathcal{O}_{\mathcal{C}'}) \to D(\mathcal{O}_\mathcal{C})$
および
$Lg_!^{\textit{Ab}} : D(\mathcal{C}') \to D(\mathcal{C})$
をもつ。
\item 図式
$$
\xymatrix{
D(\mathcal{O}_{\mathcal{C}'}) \ar[d] \ar[r]_{Lg_!} &
D(\mathcal{O}_\mathcal{C}) \ar[d] \\
D(\mathcal{C}') \ar[r]^{Lg_!^{\textit{Ab}}} &
D(\mathcal{C})
}
$$
は可換である。
\end{enumerate}
\end{lemma}

\begin{proof}
Stacks, Lemma \ref{stacks-lemma-topology-inherited-functorial} により、関手
$u$ は連続かつ余連続である。したがって関手 $g_!$、$g_!^{\textit{Ab}}$、
$Lg_!$、$Lg_!^{\textit{Ab}}$ の存在については Modules on Sites, Sections
\ref{sites-modules-section-exactness-lower-shriek}、
\ref{sites-modules-section-lower-shriek-modules} および本章の節
\ref{sites-cohomology-section-derived-lower-shriek} を参照すればよい。

\medskip\noindent
(2) を証明するには、標準射
$$
g_!^{\textit{Ab}}j_{U'!}\mathcal{O}_{U'} \to j_{u(U')!}\mathcal{O}_{u(U')}
$$
が $\mathcal{C}'$ のすべての対象 $U'$ に対して同型であることを示せば十分である。
Modules on Sites, Remark
\ref{sites-modules-remark-when-shriek-equal} を参照せよ。同様に (4) を
証明するには、標準射
$$
Lg_!^{\textit{Ab}}j_{U'!}\mathcal{O}_{U'} \to j_{u(U')!}\mathcal{O}_{u(U')}
$$
が $\mathcal{C}'$ のすべての対象 $U'$ に対して $D(\mathcal{C})$ の同型である
ことを示せば十分である。注 \ref{sites-cohomology-remark-when-derived-shriek-equal} を参照せよ。
これは前の公式も含意するので、このことを示す。

\medskip\noindent
局所化射 $j$ に対して関手 $j_!$ と $j_!^{\textit{Ab}}$ が一致すること
（Modules on Sites, Remark
\ref{sites-modules-remark-localize-shriek-equal} 参照）、および $j_!$ が完全で
あること（Modules on Sites, Lemma
\ref{sites-modules-lemma-extension-by-zero-exact}）を用いる。補題
\ref{sites-cohomology-lemma-morphism-fibred-categories-with-object} の記法を採用する。
$Lg_!^{\textit{Ab}} \circ j_{U'!} = j_{U!} \circ L(g')^{\textit{Ab}}_!$ なので
（Sites, Lemma \ref{sites-lemma-localize-cocontinuous} の可換性と随伴関手の
一意性による）、$L(g')^{\textit{Ab}}_!\mathcal{O}_{U'} = \mathcal{O}_U$
を証明すれば十分である。補題
\ref{sites-cohomology-lemma-morphism-fibred-categories-with-object} の結果を用いると、
$D(\mathcal{C}/u(U'))$ の任意の対象 $E$ に対して次の等式列を得る。
\begin{align*}
\Hom_{D(\mathcal{C}/U)}(L(g')_!^{\textit{Ab}}\mathcal{O}_{U'}, E)
& =
\Hom_{D(\mathcal{C}'/U')}(\mathcal{O}_{U'}, (g')^{-1}E) \\
& =
\Hom_{D(\mathcal{C}'/U')}((\pi'_{U'})^{-1}\mathcal{O}_V, (g')^{-1}E) \\
& =
\Hom_{D(\mathcal{D}/V)}(\mathcal{O}_V, R\pi'_{U', *}(g')^{-1}E) \\
& =
\Hom_{D(\mathcal{D}/V)}(\mathcal{O}_V, (\sigma')^{-1}(g')^{-1}E) \\
& =
\Hom_{D(\mathcal{D}/V)}(\mathcal{O}_V, \sigma^{-1}E) \\
& =
\Hom_{D(\mathcal{D}/V)}(\mathcal{O}_V, \pi_{U, *}E) \\
& =
\Hom_{D(\mathcal{C}/U)}(\pi_U^{-1}\mathcal{O}_V, E) \\
& =
\Hom_{D(\mathcal{C}/U)}(\mathcal{O}_U, E)
\end{align*}
米田の補題により結論が従う。
\end{proof}

\begin{remark}
\label{sites-cohomology-remark-fibred-category}
状況 \ref{sites-cohomology-situation-fibred-category} の仮定と記法を用いる。
$\mathcal{C}' = \mathcal{D}$ と置き、$u$ を $\mathcal{C}$ の構造関手とすると、
状況 \ref{sites-cohomology-situation-morphism-fibred-categories} の場合が得られることに注意する。
したがって補題 \ref{sites-cohomology-lemma-properties-lower-shriek-fibred-category} により、
$forget \circ \pi_! = \pi_!^{\textit{Ab}} \circ forget$ および
$forget \circ L\pi_! = L\pi_!^{\textit{Ab}} \circ forget$ を満たす関手
$\pi_!$、$\pi_!^{\textit{Ab}}$、$L\pi_!$、$L\pi_!^{\textit{Ab}}$ が存在する。
\end{remark}

\begin{remark}
\label{sites-cohomology-remark-morphism-fibred-categories}
状況 \ref{sites-cohomology-situation-morphism-fibred-categories} の仮定と記法を用いる。
$\mathcal{F}$ を $\mathcal{C}$ 上のアーベル層、$\mathcal{F}'$ を
$\mathcal{C}'$ 上のアーベル層とし、$t : \mathcal{F}' \to
g^{-1}\mathcal{F}$ を写像とする。このとき標準射
$$
L\pi'_!(\mathcal{F}') \longrightarrow L\pi_!(\mathcal{F})
$$
を得る。これには、$t$ に随伴する写像
$g_!\mathcal{F}' \to \mathcal{F}$、写像
$Lg_!(\mathcal{F}') \to g_!\mathcal{F}'$、および等式
$L\pi'_! = L\pi_! \circ Lg_!$ を用いる。
\end{remark}

\begin{lemma}
\label{sites-cohomology-lemma-compute-pi-shriek}
状況 \ref{sites-cohomology-situation-fibred-category} の仮定と記法を用いる。
$\mathcal{F} \in \textit{Ab}(\mathcal{C})$ に対し、層
$\pi_!\mathcal{F}$ は前層
$$
V \longmapsto \colim_{\mathcal{C}_V^{opp}} \mathcal{F}|_{\mathcal{C}_V}
$$
に付随する層である。制限写像は証明中に記す。
\end{lemma}

\begin{proof}
補題の規則を $\mathcal{H}$ と書く。$\mathcal{D}$ の射
$h : V' \to V$ に対し、ファイバー圏の引き戻し関手
$h^* : \mathcal{C}_V \to \mathcal{C}_{V'}$ が存在する
（Categories, Definition
\ref{categories-definition-pullback-functor-fibred-category}）。さらに
$U \in \Ob(\mathcal{C}_V)$ に対し、$h$ を覆う強デカルト射
$h^*U \to U$ が存在する。これらの強デカルト射に沿う制限は関手の変換
$$
\mathcal{F}|_{\mathcal{C}_V}
\longrightarrow
\mathcal{F}|_{\mathcal{C}_{V'}} \circ h^*.
$$
を定める。したがって余極限の間の写像
$\mathcal{H}(V) \to \mathcal{H}(V')$ が得られる。Categories, Lemma
\ref{categories-lemma-functorial-colimit} を参照せよ。

\medskip\noindent
補題を証明するため、
$$
\Mor_{\textit{PSh}(\mathcal{D})}(\mathcal{H}, \mathcal{G}) =
\Mor_{\Sh(\mathcal{C})}(\mathcal{F}, \pi^{-1}\mathcal{G})
$$
が $\mathcal{C}$ 上の任意の層 $\mathcal{G}$ に対して成り立つことを示す。
左辺の元は、$\mathcal{C}$ のすべての $U$ に対する写像
$\mathcal{F}(U) \to \mathcal{G}(p(U))$ の整合的な族である。
$\mathcal{C}$ の位相の選び方により
$\pi^{-1}\mathcal{G}(U) = \mathcal{G}(p(U))$ なので、右辺についても
同じ記述が成り立ち、結論を得る。
\end{proof}





\section{圏上のホモロジー}
\label{sites-cohomology-section-homology}

\noindent
一点上の圏の場合、左導来関手 $L\pi_!$ をホモロジー関手と呼ぶ。

\begin{example}[一点上の圏]
\label{sites-cohomology-example-category-to-point}
$\mathcal{C}$ を圏とし、カオス位相を入れる（Sites, Example
\ref{sites-example-indiscrete}）。したがって $\mathcal{C}$ 上では前層と層が
一致する。$*$ を一つの対象と一つの射だけをもつ圏とすると、関手
$p : \mathcal{C} \to *$ は余連続かつ連続である。
$\pi : \Sh(\mathcal{C}) \to \Sh(*)$ を対応するトポスの射とする。
$B$ を環とする。$*$ には環の層 $B$ を入れ、$\mathcal{C}$ には
$\mathcal{O}_\mathcal{C} = \pi^{-1}B$ を入れる。後者を
$\underline{B}$ と書く。このようにして
$$
\pi : (\Sh(\mathcal{C}), \underline{B}) \to (\Sh(*), B)
$$
は状況 \ref{sites-cohomology-situation-fibred-category} の例となる。注
\ref{sites-cohomology-remark-fibred-category} により、加群上の $\pi_!$ とアーベル層上の
$\pi_!$ を区別する必要はない。補題 \ref{sites-cohomology-lemma-compute-pi-shriek} により
$\pi_!\mathcal{F} = \colim_{\mathcal{C}^{opp}} \mathcal{F}$.
である。したがって $L_n\pi_!$ は余極限を取る関手の第 $n$ 左導来関手である。
以下では
$$
H_n(\mathcal{C}, \mathcal{F}) = L_n\pi_!(\mathcal{F})
$$
と書き、これを $\mathcal{C}$ 上の {\it $\mathcal{F}$ の第 $n$ ホモロジー群}
と呼ぶ。
\end{example}

\begin{example}[ホモロジーの計算]
\label{sites-cohomology-example-left-derived-colimits}
例 \ref{sites-cohomology-example-category-to-point} では関手 $H_n(\mathcal{C}, -)$ を
次のように計算できる。$\mathcal{F} \in
\Ob(\textit{Ab}(\mathcal{C}))$ とし、鎖複体
$$
K_\bullet(\mathcal{F}) :
\ \ldots \to
\bigoplus\nolimits_{U_2 \to U_1 \to U_0} \mathcal{F}(U_0)
\to
\bigoplus\nolimits_{U_1 \to U_0} \mathcal{F}(U_0)
\to
\bigoplus\nolimits_{U_0} \mathcal{F}(U_0)
$$
を考える。境界写像は
$$
(U_2 \to U_1 \to U_0, s)
\longmapsto
(U_1 \to U_0, s) - (U_2 \to U_0, s) + (U_2 \to U_1, s|_{U_1})
$$
で与えられ、他の次数でも同様である。構成により
$$
H_0(\mathcal{C}, \mathcal{F}) =
\colim_{\mathcal{C}^{opp}} \mathcal{F} =
H_0(K_\bullet(\mathcal{F})),
$$
である。Categories, Lemma
\ref{categories-lemma-colimits-coproducts-coequalizers} を参照せよ。
$K_\bullet(\mathcal{F})$ の構成は $\mathcal{F}$ に関して関手的であり、
$\textit{Ab}(\mathcal{C})$ の短完全列を複体の短完全列へ移す。
したがって関手列 $\mathcal{F} \mapsto H_n(K_\bullet(\mathcal{F}))$ は
$\delta$-関手をなす。Homology, Definition
\ref{homology-definition-cohomological-delta-functor} および Lemma
\ref{homology-lemma-long-exact-sequence-cochain} を参照せよ。
$\mathcal{F} = j_{U!}\mathbf{Z}_U$ のとき、複体
$K_\bullet(\mathcal{F})$ は、各項が
$$
X_n = \coprod\nolimits_{U_n \to \ldots \to U_1 \to U_0}
\Mor_\mathcal{C}(U_0, U)
$$
である単体的集合 $X_\bullet$ 上の自由 $\mathbf{Z}$-加群に付随する複体である。
この単体的集合は一元集合 $\{*\}$ 上の定値単体的集合とホモトピー同値である。
実際、写像 $X_\bullet \to \{*\}$ は明らかであり、写像
$\{*\} \to X_n$ は $*$ を
$(U \to \ldots \to U, \text{id}_U)$ へ送ることで与えられる。また二つの写像
$X_\bullet \to X_\bullet$ の間のホモトピーを定める写像
$$
h_{n, i} : X_n \longrightarrow X_n
$$
は（Simplicial, Lemma \ref{simplicial-lemma-relations-homotopy}）、規則
$$
h_{n, i} :
(U_n \to \ldots \to U_0, f)
\longmapsto
(U_n \to \ldots \to U_i \to U \to \ldots \to U, \text{id})
$$
で $i > 0$ に対して与えられ、$h_{n, 0} = \text{id}$ である。検証は省略する。
これにより $K_\bullet(j_{U!}\mathbf{Z}_U)$ は負次数で自明なコホモロジーを
もつ（Simplicial, Remark \ref{simplicial-remark-homotopy-better} の関手性と
Simplicial, Lemma \ref{simplicial-lemma-homotopy-s-N} の結果による）。
したがって、例えば Homology, Lemma
\ref{homology-lemma-efface-implies-universal} と Derived Categories,
Lemma \ref{derived-lemma-right-derived-delta-functor} の双対により、
$K_\bullet(\mathcal{F})$ は $H_0(\mathcal{C}, -)$ の左導来関手
$H_n(\mathcal{C}, -)$ を計算する。
\end{example}

\begin{example}
\label{sites-cohomology-example-morphism-categories}
$u : \mathcal{C}' \to \mathcal{C}$ を関手とする。例
\ref{sites-cohomology-example-category-to-point} と同様に $\mathcal{C}'$ と
$\mathcal{C}$ にカオス位相を入れる。$*$ を一つの対象と一つの射だけをもつ圏と
すると、関手 $u$、$\mathcal{C}' \to *$、$\mathcal{C} \to *$ は
余連続かつ連続である。
$g : \Sh(\mathcal{C}') \to \Sh(\mathcal{C})$,
$\pi' : \Sh(\mathcal{C}') \to \Sh(*)$, and
$\pi : \Sh(\mathcal{C}) \to \Sh(*)$,
を対応するトポスの射とする。$B$ を環とする。$*$ には環の層 $B$ を入れ、
$\mathcal{C}'$ と $\mathcal{C}$ には定値層 $\underline{B}$ を入れる。
このようにして
$$
\xymatrix{
(\Sh(\mathcal{C}'), \underline{B}) \ar[rd]_{\pi'} \ar[rr]_g & &
(\Sh(\mathcal{C}), \underline{B}) \ar[ld]^\pi \\
& (\Sh(*), B)
}
$$
は状況 \ref{sites-cohomology-situation-morphism-fibred-categories} の例となる。したがって補題
\ref{sites-cohomology-lemma-properties-lower-shriek-fibred-category} を $g$ に適用でき、
加群上の $g_!$ とアーベル層上の $g_!$ を区別する必要はない。特に注
\ref{sites-cohomology-remark-morphism-fibred-categories} は標準射
$$
H_n(\mathcal{C}', \mathcal{F}')
\longrightarrow
H_n(\mathcal{C}, \mathcal{F})
$$
を与える。ただし $\mathcal{F} \in \textit{Ab}(\mathcal{C})$、
$\mathcal{F}' \in \textit{Ab}(\mathcal{C}')$ とし、写像
$t : \mathcal{F}' \to g^{-1}\mathcal{F}$ が与えられているものとする。
例 \ref{sites-cohomology-example-left-derived-colimits} のホモロジー計算で見ると、これらの写像は
複体の写像
$$
K_\bullet(\mathcal{F}') \longrightarrow K_\bullet(\mathcal{F})
$$
に由来し、その規則は
$$
(U'_n \to \ldots \to U'_0, s') \longmapsto
(u(U'_n) \to \ldots \to u(U'_0), t(s'))
$$
である。記法の意味は明らかであろう。
\end{example}

\begin{remark}
\label{sites-cohomology-remark-map-evaluation-to-derived}
例 \ref{sites-cohomology-example-category-to-point} の記法と仮定を用いる。
$\mathcal{F}^\bullet$ を $\mathcal{C}$ 上のアーベル層の有界複体とする。
$\mathcal{C}$ の任意の対象 $U$ に対して標準射
$$
\mathcal{F}^\bullet(U) \longrightarrow L\pi_!(\mathcal{F}^\bullet)
$$
が $D(\textit{Ab})$ に存在する。$\mathcal{F}^\bullet$ が
$\underline{B}$-加群の複体ならば、この写像は $D(B)$ に属する。これを証明する
ため、射影的対象の複体 $\mathcal{P}^\bullet$ からの擬同型
$\mathcal{P}^\bullet \to \mathcal{F}^\bullet$ を取り、
$L\pi_!(\mathcal{F}^\bullet)$ を計算することに注意する。しかし位相がカオスなので、
$\mathcal{P}^\bullet(U) \to \mathcal{F}^\bullet(U)$ は擬同型であり、
$D(\textit{Ab})$、それぞれ $D(B)$ で逆にできる。これと、$\pi_!$ を余極限として
計算することから得られる標準射
$\mathcal{P}^\bullet(U) \to \pi_!(\mathcal{P}^\bullet)$ を合成すれば、
求める矢を得る。
\end{remark}

\begin{lemma}
\label{sites-cohomology-lemma-initial-final}
例 \ref{sites-cohomology-example-category-to-point} の記法と仮定を用いる。
$\mathcal{C}$ が始対象または終対象をもつならば、$D(\textit{Ab})$、それぞれ
$D(B)$ 上で $L\pi_! \circ \pi^{-1} = \text{id}$ である。
\end{lemma}

\begin{proof}
もし $\mathcal{C}$ が始対象をもつなら、$\pi_!$ はこの対象での評価により計算され、
主張は明らかである。もし $\mathcal{C}$ が終対象をもつなら、$R\pi_*$ はこの対象での
評価により計算される。したがって $D(\textit{Ab})$、それぞれ $D(B)$ 上で
$R\pi_* \circ \pi^{-1} \cong \text{id}$ である。これにより
$\pi^{-1} : D(\textit{Ab}) \to D(\mathcal{C})$,
それぞれ $\pi^{-1} : D(B) \to D(\underline{B})$ は充満忠実である。
Categories, Lemma \ref{categories-lemma-adjoint-fully-faithful} を参照せよ。
すると同じ補題により、望んだとおり
$L\pi_! \circ \pi^{-1} = \text{id}$ が従う。
\end{proof}

\begin{lemma}
\label{sites-cohomology-lemma-change-of-rings}
例 \ref{sites-cohomology-example-category-to-point} の記法と仮定を用いる。
$B \to B'$ を環準同型とする。環付きトポスの可換図式
$$
\xymatrix{
(\Sh(\mathcal{C}), \underline{B}) \ar[d]_\pi &
(\Sh(\mathcal{C}), \underline{B'}) \ar[d]^{\pi'} \ar[l]^h \\
(*, B) & (*, B') \ar[l]_f
}
$$
このとき $L\pi_! \circ Lh^* = Lf^* \circ L\pi'_!$ である。
\end{lemma}

\begin{proof}
どちらの関手も明らかな関手 $D(B') \to D(\underline{B})$ の右随伴である。
\end{proof}

\begin{lemma}
\label{sites-cohomology-lemma-compute-by-cosimplicial-resolution}
例 \ref{sites-cohomology-example-category-to-point} の記法と仮定を用いる。
$U_\bullet$ を $\mathcal{C}$ の余単体的対象とし、任意の
$U \in \Ob(\mathcal{C})$ に対して単体的集合
$\Mor_\mathcal{C}(U_\bullet, U)$
が一元集合上の定値単体的集合とホモトピー同値であるとする。このとき
$$
L\pi_!(\mathcal{F}) = \mathcal{F}(U_\bullet)
$$
が $D(\textit{Ab})$、それぞれ $D(B)$ において成り立ち、
$\textit{Ab}(\mathcal{C})$、それぞれ $\textit{Mod}(\underline{B})$ の
$\mathcal{F}$ に関して関手的である。
\end{lemma}

\begin{proof}
補題 \ref{sites-cohomology-lemma-properties-lower-shriek-fibred-category} により $L\pi_!$ は
加群とアーベル層で一致するので、$\mathcal{F}$ がアーベル層の場合を証明すれば
十分である。$U \in \Ob(\mathcal{C})$ に対して、アーベル層
$j_{U!}\mathbf{Z}_U$ は $\textit{Ab}(\mathcal{C})$ の射影的対象である。実際、
$\Hom(j_{U!}\mathbf{Z}_U, \mathcal{F}) = \mathcal{F}(U)$
であり、位相がカオスなので切断を取る関手は完全である。Modules on Sites,
Lemma \ref{sites-modules-lemma-module-quotient-flat} により、任意のアーベル層は
$j_{U!}\mathbf{Z}_U$ の直和の商である。したがって分解
$$
\ldots \to \mathcal{G}^{-1} \to \mathcal{G}^0 \to \mathcal{F} \to 0
$$
を、その各項が上記の形の層の直和となるように選び、
$L\pi_!(\mathcal{F}) = \pi_!(\mathcal{G}^\bullet)$ として計算できる。
二重複体 $A^{\bullet, \bullet} = \mathcal{G}^\bullet(U_\bullet)$ を考える。
写像 $\mathcal{G}^0 \to \mathcal{F}$ は複体の写像
$A^{0, \bullet} \to \mathcal{F}(U_\bullet)$ を与える。
$\pi_!$ は $\mathcal{C}^{opp}$ 上の余極限を取ることで計算されるので
（補題 \ref{sites-cohomology-lemma-compute-pi-shriek}）、二つの合成
$\mathcal{G}^m(U_1) \to \mathcal{G}^m(U_0) \to \pi_!\mathcal{G}^m$
は等しい。したがって複体の標準射
$$
\text{Tot}(A^{\bullet, \bullet})
\longrightarrow
\pi_!(\mathcal{G}^\bullet) = L\pi_!(\mathcal{F})
$$
を得る。補題を証明するには、複体
$$
\ldots \to \mathcal{G}^m(U_1) \to \mathcal{G}^m(U_0) \to
\pi_!\mathcal{G}^m \to 0
$$
が完全であることを示せば十分である。Homology, Lemma
\ref{homology-lemma-double-complex-gives-resolution} を参照せよ。
層 $\mathcal{G}^m$ は層 $j_{U!}\mathbf{Z}_U$ の直和なので、
$\mathcal{G} = j_{U!}\mathbf{Z}_U$ に帰着する。複体
$j_{U!}\mathbf{Z}_U(U_\bullet)$ は、単体的集合
$\Mor_\mathcal{C}(U_\bullet, U)$ 上の自由 $\mathbf{Z}$-加群に付随する
アーベル群の複体である。この単体的集合は一元集合とホモトピー同値と仮定した。
したがって
$$
j_{U!}\mathbf{Z}_U(U_\bullet) \to \mathbf{Z}
$$
はアーベル群のホモトピー同値、したがって擬同型である（Simplicial, Remark
\ref{simplicial-remark-homotopy-better} および Lemma
\ref{simplicial-lemma-homotopy-s-N}）。補題
\ref{sites-cohomology-lemma-properties-lower-shriek-fibred-category} の証明で示したように
$\pi_!j_{U!}\mathbf{Z}_U = \mathbf{Z}$ なので、これで証明は完了する。
\end{proof}

\begin{lemma}
\label{sites-cohomology-lemma-get-it-now}
例 \ref{sites-cohomology-example-morphism-categories} の記法と仮定を用いる。
$\mathcal{C}'$ の余単体的対象 $U'_\bullet$ が存在し、補題
\ref{sites-cohomology-lemma-compute-by-cosimplicial-resolution} を $\mathcal{C}'$ の
$U'_\bullet$ と $\mathcal{C}$ の $u(U'_\bullet)$ の双方に適用できるとする。
このとき
$L\pi'_! \circ g^{-1} = L\pi_!$
が関手 $D(\mathcal{C}) \to D(\textit{Ab})$ として、またそれぞれ
$D(\mathcal{C}, \underline{B}) \to D(B)$ の関手として成り立つ。
\end{lemma}

\begin{proof}
補題 \ref{sites-cohomology-lemma-compute-by-cosimplicial-resolution} と、$g^{-1}$ が $u$ との
前合成で与えられることから直ちに従う。
\end{proof}

\begin{lemma}
\label{sites-cohomology-lemma-product-categories}
$\mathcal{C}_i$（$i = 1, 2$）を圏とする。
$u_i : \mathcal{C}_1 \times \mathcal{C}_2 \to \mathcal{C}_i$ を射影関手とする。
$B$ を環とし、
$g_i : (\Sh(\mathcal{C}_1 \times \mathcal{C}_2), \underline{B}) \to
(\Sh(\mathcal{C}_i), \underline{B})$ を対応する環付きトポスの射とする。
例 \ref{sites-cohomology-example-morphism-categories} を参照せよ。
$K_i \in D(\mathcal{C}_i, B)$ に対して
$$
L(\pi_1 \times \pi_2)_!(
g_1^{-1}K_1 \otimes_{\underline{B}}^\mathbf{L} g_2^{-1}K_2)
=
L\pi_{1, !}(K_1) \otimes_B^\mathbf{L} L\pi_{2, !}(K_2)
$$
が $D(B)$ において成り立つ。記法の意味は明らかであろう。
\end{lemma}

\begin{proof}
両辺は余極限と可換するので、$U \in \Ob(\mathcal{C}_1)$ および
$V \in \Ob(\mathcal{C}_2)$ に対する
$K_1 = j_{U!}\underline{B}_U$ と $K_2 = j_{V!}\underline{B}_V$ の場合を
証明すれば十分である。補題 \ref{sites-cohomology-lemma-existence-derived-lower-shriek} における
$L\pi_!$ の構成を参照せよ。この場合
$$
g_1^{-1}K_1 \otimes_{\underline{B}}^\mathbf{L} g_2^{-1}K_2 =
g_1^{-1}K_1 \otimes_{\underline{B}} g_2^{-1}K_2 =
j_{(U, V)!}\underline{B}_{(U, V)}
$$
検証は省略する。したがって、補題の公式の左辺と右辺はいずれも $B$ に評価される
ので結果が従う。補題 \ref{sites-cohomology-lemma-existence-derived-lower-shriek} における
$L\pi_!$ の構成を参照せよ。
\end{proof}

\begin{lemma}
\label{sites-cohomology-lemma-eilenberg-zilber}
例 \ref{sites-cohomology-example-category-to-point} の記法と仮定を用いる。
補題 \ref{sites-cohomology-lemma-compute-by-cosimplicial-resolution} を適用できる
$\mathcal{C}$ の余単体的対象 $U_\bullet$ が存在するとする。このとき
$$
L\pi_!(K_1 \otimes^\mathbf{L}_{\underline{B}} K_2) =
L\pi_!(K_1) \otimes^\mathbf{L}_B L\pi_!(K_2)
$$
がすべての $K_i \in D(\underline{B})$ に対して成り立つ。
\end{lemma}

\begin{proof}
圏と関手の図式
$$
\xymatrix{
& & \mathcal{C} \\
\mathcal{C} \ar[r]^-u &
\mathcal{C} \times \mathcal{C} \ar[rd]^{u_2} \ar[ru]_{u_1} \\
& & \mathcal{C}
}
$$
を考える。ここで $u$ は対角関手、$u_i$ は射影関手である。これは環付きトポスの
射 $g$、$g_1$、$g_2$ を与える。$\mathcal{C}$ の任意の対象
$(U_1, U_2)$ に対して
$$
\Mor_{\mathcal{C} \times \mathcal{C}}(u(U_\bullet), (U_1, U_2)) =
\Mor_\mathcal{C}(U_\bullet, U_1) \times \Mor_\mathcal{C}(U_\bullet, U_2)
$$
であり、これは Simplicial, Lemma
\ref{simplicial-lemma-products-homotopy} により一点とホモトピー同値である。
したがって補題 \ref{sites-cohomology-lemma-get-it-now} により、
$D(\mathcal{C} \times \mathcal{C}, B)$ の任意の $K$ に対して
$L\pi_!(g^{-1}K) = L(\pi \times \pi)_!(K)$ である。
$K = g_1^{-1}K_1 \otimes_B^\mathbf{L} g_2^{-1}K_2$ と取る。
$g^{-1} = g^* = Lg^*$ は導来テンソル積と可換するので（補題
\ref{sites-cohomology-lemma-pullback-tensor-product}）、
$g^{-1}K = K_1 \otimes^\mathbf{L}_{\underline{B}} K_2$ である。
最後に補題 \ref{sites-cohomology-lemma-product-categories} を適用する。
\end{proof}

\begin{remark}[単体的加群]
\label{sites-cohomology-remark-simplicial-modules}
$\mathcal{C} = \Delta$ とし、$B$ を任意の環とする。これは例
\ref{sites-cohomology-example-category-to-point} の特別な場合であり、補題
\ref{sites-cohomology-lemma-compute-by-cosimplicial-resolution} の仮定が成り立つ。実際、$U_\bullet$ を恒等関手が
与える $\Delta$ の余単体的対象とする。条件を検証するには、
$[m] \in \Ob(\Delta)$ に対して単体的集合
$\Delta[m] : n \mapsto \Mor_\Delta([n], [m])$ が一点とホモトピー同値であることを
示さなければならない。これは Simplicial, Example
\ref{simplicial-example-simplex-contractible} で説明されている。

\medskip\noindent
この状況では圏 $\textit{Mod}(\underline{B})$ は単体的 $B$-加群の圏にほかならず、
関手 $L\pi_!$ は単体的 $B$-加群 $M_\bullet$ を、それに付随する
$B$-加群の複体 $s(M_\bullet)$ へ送る。したがって上の結果は単体的加群に関する
結果として読み替えられる。例えば補題 \ref{sites-cohomology-lemma-eilenberg-zilber} の特別な場合は
次である。$M_\bullet$、$M'_\bullet$ が平坦な単体的 $B$-加群ならば、複体
$s(M_\bullet \otimes_B M'_\bullet)$ は、二重複体
$s(M_\bullet) \otimes_B s(M'_\bullet)$ に付随する全複体と擬同型である。
（ヒント：平坦性を用いて導来テンソル積を通常のテンソル積へ置き換える。）
これは \cite{Eilenberg-Zilber} にある Eilenberg--Zilber の定理の特別な場合である。
\end{remark}

\begin{lemma}
\label{sites-cohomology-lemma-O-homology-qis}
$\mathcal{C}$ をカオス位相を入れた圏とする。
$\mathcal{O} \to \mathcal{O}'$ を $\mathcal{C}$ 上の環の層の写像とする。
次を仮定する。
\begin{enumerate}
\item 補題 \ref{sites-cohomology-lemma-compute-by-cosimplicial-resolution} におけるような
$\mathcal{C}$ の余単体的対象 $U_\bullet$ が存在する。
\item $L\pi_!\mathcal{O} \to L\pi_!\mathcal{O}'$ は同型である。
\end{enumerate}
$D(\mathcal{O})$ の $K$ に対して
$$
L\pi_!(K) = L\pi_!(K \otimes_\mathcal{O}^\mathbf{L} \mathcal{O}')
$$
が $D(\textit{Ab})$ において成り立つ。
\end{lemma}

\begin{proof}
注意：この証明で $L\pi_!$ は、アーベル層上の $\pi_!$ の左導来関手を表す。
$L\pi_!$ は余極限と可換するので、$\mathcal{O}$-加群の上に有界な複体について
証明すれば十分である（Derived Categories, Proposition
\ref{derived-proposition-left-derived-exists} の議論と比較するか、単に上に有界な
複体に限定せよ）。そのような任意の複体は、各項が
$U \in \Ob(\mathcal{C})$ に対する $j_{U!}\mathcal{O}_U$ の直和である
上に有界な複体と擬同型である。Modules on Sites, Lemma
\ref{sites-modules-lemma-module-quotient-flat} を参照せよ。したがって補題を
$j_{U!}\mathcal{O}_U$ について証明すれば十分である。仮定により
$$
S_\bullet = \Mor_\mathcal{C}(U_\bullet, U)
$$
は一元集合上の定値単体的集合とホモトピー同値な単体的集合である。
$P_n = \mathcal{O}(U_n)$ および $P'_n = \mathcal{O}'(U_n)$ と置く。
単体的アーベル群
$$
X_\bullet : n \longmapsto \bigoplus\nolimits_{s \in S_n} P_n
$$
に付随する複体は、補題 \ref{sites-cohomology-lemma-compute-by-cosimplicial-resolution} により
$L\pi_!(j_{U!}\mathcal{O}_U)$ を計算する。
$j_{U!}\mathcal{O}_U$ は平坦な $\mathcal{O}$-加群なので
$j_{U!}\mathcal{O}_U \otimes^\mathbf{L}_\mathcal{O} \mathcal{O}' =
j_{U!}\mathcal{O}'_U$ であり、これの $L\pi_!$ は単体的アーベル群
$$
X'_\bullet : n \longmapsto \bigoplus\nolimits_{s \in S_n} P'_n
$$
に付随する複体によって計算される。
単体的集合 $T_\bullet$ に、各項が
$\bigoplus_{t \in T_n} P_n$ である単体的アーベル群を対応させる規則は関手なので、
$X_\bullet \to P_\bullet$ は単体的アーベル群のホモトピー同値である。
同様に、単体的集合 $T_\bullet$ に、各項が
$\bigoplus_{t \in T_n} P'_n$ である単体的アーベル群を対応させる規則も関手である。
したがって $X'_\bullet \to P'_\bullet$ は単体的アーベル群のホモトピー同値である。
仮定により $P_\bullet \to P'_\bullet$ は擬同型である（実際、補題
\ref{sites-cohomology-lemma-compute-by-cosimplicial-resolution} により、$P_\bullet$ と
$P'_\bullet$ はそれぞれ $L\pi_!\mathcal{O}$ と $L\pi_!\mathcal{O}'$ を計算する）。
ゆえに望みどおり $X_\bullet$ と $X'_\bullet$ は擬同型である。
\end{proof}

\begin{remark}
\label{sites-cohomology-remark-O-homology-B-homology-general}
例 \ref{sites-cohomology-example-category-to-point} におけるような $\mathcal{C}$ と $B$ を取る。
補題 \ref{sites-cohomology-lemma-compute-by-cosimplicial-resolution} におけるような余単体的対象が
存在すると仮定する。$\mathcal{O} \to \underline{B}$ を、同型
$L\pi_!\mathcal{O} \to L\pi_!\underline{B}$ を誘導する
$\mathcal{C}$ 上の環の層の写像とする。このとき三角圏の完全関手
$$
L\pi_! : D(\mathcal{O}) \longrightarrow D(B)
$$
を得る。実際、$D(\mathcal{O})$ の任意の対象 $K$ に対して、補題
\ref{sites-cohomology-lemma-O-homology-qis} により
$L\pi^{\textit{Ab}}_!(K) =
L\pi^{\textit{Ab}}_!(K \otimes_{\mathcal{O}}^\mathbf{L} \underline{B})$
である。したがって、表示された関手を
$- \otimes^\mathbf{L}_\mathcal{O} \underline{B}$ と関手
$L\pi_! : D(\underline{B}) \to D(B)$ の合成として定義できる。
言い換えると、補題による $L\pi_!(K)$ と
$L\pi_!(K \otimes_\mathcal{O}^\mathbf{L} \underline{B})$ の
（標準的かつ関手的な）同一視から、$L\pi_!(K)$ 上の $B$-加群構造を得る。
\end{remark}







\section{導来関手 $L\pi_!$ と $Lg_!$ の計算}
\label{sites-cohomology-section-calculate}

\noindent
本節では節 \ref{sites-cohomology-section-homology} の結果を適用し、状況
\ref{sites-cohomology-situation-fibred-category} における $L\pi_!$ と、状況
\ref{sites-cohomology-situation-morphism-fibred-categories} における $Lg_!$ を計算する。

\begin{lemma}
\label{sites-cohomology-lemma-compute-left-derived-pi-shriek-pre}
状況 \ref{sites-cohomology-situation-fibred-category} の仮定と記法を用いる。
$\textit{PAb}(\mathcal{C})$ の $\mathcal{F}$ と $n \geq 0$ に対し、
前層
$$
V \longmapsto H_n(\mathcal{C}_V, \mathcal{F}|_{\mathcal{C}_V})
$$
に付随する $\mathcal{D}$ 上のアーベル層を $L_n(\mathcal{F})$ とする。
制限写像は証明中で指定する。このとき
$L_n(\mathcal{F}) = L_n(\mathcal{F}^\#)$ である。
\end{lemma}

\begin{proof}
$\mathcal{D}$ の射 $h : V' \to V$ に対して、ファイバー圏の引き戻し関手
$h^* : \mathcal{C}_V \to \mathcal{C}_{V'}$ が存在する（Categories, Definition
\ref{categories-definition-pullback-functor-fibred-category}）。
さらに $U \in \Ob(\mathcal{C}_V)$ に対して、$h$ を覆う強デカルト射
$h^*U \to U$ が存在する。これらの強デカルト射に沿う制限は関手の変換
$$
\mathcal{F}|_{\mathcal{C}_V}
\longrightarrow
\mathcal{F}|_{\mathcal{C}_{V'}} \circ h^*.
$$
を定める。例 \ref{sites-cohomology-example-morphism-categories} により、望む制限写像
$$
H_n(\mathcal{C}_V, \mathcal{F}|_{\mathcal{C}_V})
\longrightarrow
H_n(\mathcal{C}_{V'}, \mathcal{F}|_{\mathcal{C}_{V'}})
$$
を得る。この前層を $L_{n, p}(\mathcal{F})$ と書く。したがって
$L_n(\mathcal{F}) = L_{n, p}(\mathcal{F})^\#$ である。
標準写像 $\gamma : \mathcal{F} \to \mathcal{F}^+$
（Sites, Theorem \ref{sites-theorem-plus}）は標準写像
$L_{n, p}(\mathcal{F}) \to L_{n, p}(\mathcal{F}^+)$ を定める。
この写像が層化後に同型となることを証明しなければならない。

\medskip\noindent
例 \ref{sites-cohomology-example-left-derived-colimits} で与えたホモロジーの計算を用いる。
$\mathcal{F}$ のファイバー圏 $\mathcal{C}_V$ への制限に付随する複体を
$K_\bullet(\mathcal{F}|_{\mathcal{C}_V})$ と書く。上の考察により、複体の前層
$K_\bullet(\mathcal{F})$
$$
V \longmapsto K_\bullet(\mathcal{F}|_{\mathcal{C}_V})
$$
を得る。そのコホモロジー前層は $L_{n, p}(\mathcal{F})$ である。
したがって
$$
K_\bullet(\mathcal{F}) \longrightarrow K_\bullet(\mathcal{F}^+)
$$
が層化後に同型となることを示せば十分である。

\medskip\noindent
単射性。$V$ を $\mathcal{D}$ の対象とし、$\xi \in K_n(\mathcal{F})(V)$ を
$K_n(\mathcal{F}^+)(V)$ で零に写る元とする。
$K_n(\mathcal{F})(V_j)$ において $\xi|_{V_j}$ が零となるような被覆
$\{V_j \to V\}$ が存在することを示さなければならない。
$$
\xi = \sum (U_{i, n + 1} \to \ldots \to U_{i, 0}, \sigma_i)
$$
と書く。ここで $\sigma_i \in \mathcal{F}(U_{i, 0})$ である。
$\mathcal{C}_V$ の各射列 $U_n \to \ldots \to U_0$ が高々一度現れるように整える。
複体 $K_\bullet$ の定義における和は直和なので、これが
$K_\bullet(\mathcal{F}^+)(V)$ で零に写るためには、すべての $\sigma_i$ が
$\mathcal{F}^+(U_{i, 0})$ で零に写らなければならない。
$\mathcal{F}^+$ の構成により、$\sigma_i|_{U_{i, 0, j}}$ が零となる被覆
$\{U_{i, 0, j} \to U_{i, 0}\}$ が存在する。
$\mathcal{C}$ 上の位相の構成により、$U_{i, 0, j} \to U_{i, 0}$ は、ある射
$V_{i, j} \to V$ の引き戻しとして書け（Categories, Definition
\ref{categories-definition-pullback-functor-fibred-category}）、しかも各
$\{V_{i, j} \to V\}$ は被覆である。これらすべての被覆
$\{V_{i, j} \to V\}$ を細分する被覆 $\{V_j \to V\}$ を選ぶ。
すると $\xi|_{V_j} = 0$ は明らかである。

\medskip\noindent
全射性。証明は省略する。ヒント：単射性の証明と同様に論ぜよ。
\end{proof}

\begin{lemma}
\label{sites-cohomology-lemma-compute-left-derived-pi-shriek}
状況 \ref{sites-cohomology-situation-fibred-category} の仮定と記法を用いる。
$\textit{Ab}(\mathcal{C})$ の $\mathcal{F}$ と $n \geq 0$ に対して、層
$L_n\pi_!(\mathcal{F})$ は補題
\ref{sites-cohomology-lemma-compute-left-derived-pi-shriek-pre} で構成した層
$L_n(\mathcal{F})$ に等しい。
\end{lemma}

\begin{proof}
$\textit{PAb}(\mathcal{C}) \to \textit{Ab}(\mathcal{C})$ の関手列
$\mathcal{F} \mapsto L_n(\mathcal{F})$ を考える。
各 $V \in \Ob(\mathcal{D})$ に対して関手列
$H_n(\mathcal{C}_V, - )$ は $\delta$-関手をなすので、関手
$\mathcal{F} \mapsto L_n(\mathcal{F})$ もそうである。
目標は、これらが普遍 $\delta$-関手をなすことを示すことである。
そのために、これらの関手が消えるいくつかのアーベル前層を構成する。

\medskip\noindent
$U' \in \Ob(\mathcal{C})$ に対してアーベル前層
$\mathcal{F}_{U'} = j_{U'!}^{\textit{PAb}}\mathbf{Z}_{U'}$
を考える（Modules on Sites, Remark
\ref{sites-modules-remark-localize-presheaves}）。次を思い出そう。
$$
\mathcal{F}_{U'}(U) =
\bigoplus\nolimits_{\Mor_\mathcal{C}(U, U')} \mathbf{Z}
$$
$U$ が $\mathcal{D})$ の $V = p(U)$ 上にあり、$U'$ が $V' = p(U')$ 上にあるとする。
このとき任意の射 $a : U \to U'$ は
$U \to h^*U' \to U'$ と一意的に分解する。ここで
$h = p(a) : V \to V'$ である（Categories, Definition
\ref{categories-definition-pullback-functor-fibred-category} 参照）。
したがって
$$
\mathcal{F}_{U'}|_{\mathcal{C}_V}
=
\bigoplus\nolimits_{h \in \Mor_\mathcal{D}(V, V')}
j_{h^*U'!}\mathbf{Z}_{h^*U'}
$$
である。ここで
$j_{h^*U'} : \Sh(\mathcal{C}_V/h^*U') \to \Sh(\mathcal{C}_V)$
は局所化射である。層 $j_{h^*U'!}\mathbf{Z}_{h^*U'}$ の高次ホモロジー群は
消える（例 \ref{sites-cohomology-example-left-derived-colimits} 参照）。
したがって、すべての $n > 0$ とすべての $U'$ に対して
$L_n(\mathcal{F}_{U'}) = 0$ である。ゆえに任意のアーベル前層
$\mathcal{F}$ は、すべての $n > 0$ に対して
$L_n(\mathcal{G}) = 0$ となるアーベル前層 $\mathcal{G}$ の商である
（Modules on Sites, Lemma
\ref{sites-modules-lemma-module-quotient-flat}）。
$L_n(\mathcal{F}) = L_n(\mathcal{F}^\#)$ なので、アーベル層についても同じことが
成り立つ。したがって関手列 $L_n(-)$ は $\textit{Ab}(\mathcal{C})$ 上の
普遍デルタ関手である（Homology, Lemma
\ref{homology-lemma-efface-implies-universal}）。
補題 \ref{sites-cohomology-lemma-compute-pi-shriek} により $n = 0$ では
$H^{-n}(L\pi_!(-))$ と一致するので、普遍 $\delta$-関手の一意性
（Homology, Lemma
\ref{homology-lemma-uniqueness-universal-delta-functor}）と
Derived Categories, Lemma
\ref{derived-lemma-right-derived-delta-functor} により結論を得る。
\end{proof}

\begin{lemma}
\label{sites-cohomology-lemma-compute-left-derived-g-shriek}
状況 \ref{sites-cohomology-situation-morphism-fibred-categories} の仮定と記法を用いる。
$\mathcal{C}'$ 上のアーベル層 $\mathcal{F}'$ に対して、層
$L_ng_!(\mathcal{F}')$ は前層
$$
U \longmapsto H_n(\mathcal{I}_U, \mathcal{F}'_U)
$$
に付随する層である。記法と制限写像については証明を参照せよ。
\end{lemma}

\begin{proof}
$p(U) = V$ とする。圏 $\mathcal{I}_U$ は、$\mathcal{C}$ の射
$\varphi : U \to u(U')$ で $p(\varphi) = \text{id}_V$ を満たすもの、すなわち
ファイバー圏 $\mathcal{C}_V$ の射であるものからなる対 $(U', \varphi)$ の圏である。
射 $(U'_1, \varphi_1) \to (U'_2, \varphi_2)$ は、
$\varphi_2 = u(a) \circ \varphi_1$ を満たすファイバー圏
$\mathcal{C}'_V$ の射 $a : U'_1 \to U'_2$ によって与えられる。
前層 $\mathcal{F}'_U$ は $(U', \varphi)$ を $\mathcal{F}'(U')$ へ送る。
以下で制限写像を構成する。

\medskip\noindent
Categories, Lemma
\ref{categories-lemma-ameliorate-morphism-fibred-categories} におけるような
$u$ の分解
$$
\xymatrix{
\mathcal{C}' \ar@<1ex>[r]^{u'} &
\mathcal{C}'' \ar[r]^{u''} \ar@<1ex>[l]^w & \mathcal{C}
}
$$
を選ぶ。このとき $g_! = g''_! \circ g'_!$ であり、導来関手についても同様である。
一方、関手 $g'_!$ は完全である。Modules on Sites, Lemma
\ref{sites-modules-lemma-have-left-adjoint} を参照せよ。
したがって $\mathcal{F}'' = g'_!\mathcal{F}'$ と置けば
$Lg_!(\mathcal{F}') = Lg''_!(\mathcal{F}'')$ を得る。
$h : \Sh(\mathcal{C}'') \to \Sh(\mathcal{C}')$ を $w$ に付随するトポスの射とすると、
$\mathcal{F}'' = h^{-1}\mathcal{F}'$ であることに注意する。Sites, Lemma
\ref{sites-lemma-have-left-adjoint} を参照せよ。
関手 $u''$ は $\mathcal{C}''$ を $\mathcal{C}$ 上のファイバー圏にするので、
$L_ng''_!$ の計算に補題
\ref{sites-cohomology-lemma-compute-left-derived-pi-shriek} を適用できる。
Categories, Lemma
\ref{categories-lemma-ameliorate-morphism-fibred-categories} の証明における
$\mathcal{C}''$ の構成によれば、ファイバー圏 $\mathcal{C}''_U$ は
$\mathcal{I}_U$ に等しいので、結果が従う。さらに
$h^{-1}\mathcal{F}'|_{\mathcal{C}''_U}$ は上で述べた規則によって与えられる
（Stacks, Lemma
\ref{stacks-lemma-topology-inherited-functorial} により $w$ は連続かつ余連続なので、
Sites, Lemma \ref{sites-lemma-when-shriek} を適用できる）。
\end{proof}








\section{単体的加群}
\label{sites-cohomology-section-simplicial-modules}

\noindent
$A_\bullet$ を単体環とする。$A_\bullet$ は、カオス位相を入れた
$\Delta$ 上の層とみなせることを思い出そう。Simplicial, Section
\ref{simplicial-section-simplicial-presheaves} を参照せよ。
すると $A_\bullet$ 上の単体的加群 $M_\bullet$ は、ちょうど $\Delta$ 上の
$A_\bullet$-加群の層である。言い換えると、各 $n \geq 0$ に対して
$A_n$-加群 $M_n$ があり、各写像 $\varphi : [n] \to [m]$ に対して対応する写像
$$
M_\bullet(\varphi) : M_m \longrightarrow M_n
$$
があり、これは $A_\bullet(\varphi)$-線形であって、これらの写像は通常の仕方で
合成される。

\medskip\noindent
$\mathcal{C}$ をサイトとする。$\mathcal{C}$ 上の
{\it 環の単体的層} $\mathcal{A}_\bullet$ とは、$\mathcal{C}$ 上の環の層の圏に
おける単体的対象である。この場合、対応
$U \mapsto \mathcal{A}_\bullet(U)$ は単体環の層であり、実際この二つの概念は
同値である。単体的アーベル層、Lie 代数の単体的層などについても同様である。

\medskip\noindent
しかし、上の単体環の場合と同様に、単体的層には別の見方がある。すなわち射影
$$
p : \Delta \times \mathcal{C} \longrightarrow \mathcal{C}
$$
を考える。これはファイバー圏を定め、その強デカルト射はちょうど
$([n], U) \to ([n], V)$ の形の射である。圏 $\Delta \times \mathcal{C}$ には
$\mathcal{C}$ から誘導される位相を入れる（Stacks, Section
\ref{stacks-section-topology} 参照）。$\Delta \times \mathcal{C}$ における被覆の
簡明な記述（Stacks, Lemma
\ref{stacks-lemma-topology-inherited}）から、$\mathcal{C}$ 上の環の単体的層と
$\Delta \times \mathcal{C}$ 上の環の層が同じものであることが直ちに従う。

\medskip\noindent
単体環上の単体的加群の場合との類比により、環の単体的層上の単体的加群を
次のように定義する。

\begin{definition}
\label{sites-cohomology-definition-simplicial-module}
$\mathcal{C}$ をサイトとし、$\mathcal{A}_\bullet$ を $\mathcal{C}$ 上の
環の単体的層とする。{\it 単体的 $\mathcal{A}_\bullet$-加群}
$\mathcal{F}_\bullet$（{\it $\mathcal{A}_\bullet$-加群の単体的層}ともいう）とは、
$\mathcal{A}_\bullet$ に付随する $\Delta \times \mathcal{C}$ 上の環の層を係数とする
加群の層である。
\end{definition}

\noindent
単体的加群の圏 $\textit{Mod}(\mathcal{A}_\bullet)$ と、それに対応する導来圏
$D(\mathcal{A}_\bullet)$ を得る。環の単体的層の写像
$\mathcal{A}_\bullet \to \mathcal{B}_\bullet$ が与えられると、関手
$$
- \otimes^\mathbf{L}_{\mathcal{A}_\bullet} \mathcal{B}_\bullet :
D(\mathcal{A}_\bullet)
\longrightarrow
D(\mathcal{B}_\bullet)
$$
を得る。さらに、前節までの結果は関手
$$
L\pi_! : D(\mathcal{A}_\bullet) \longrightarrow D(\mathcal{C})
$$
を定める。単体的加群 $\mathcal{F}_\bullet$ に対して、対象
$L\pi_!(\mathcal{F}_\bullet)$ は付随する鎖複体
$s(\mathcal{F}_\bullet)$ によって表される
（Simplicial, Section \ref{simplicial-section-complexes}）。
これは補題 \ref{sites-cohomology-lemma-compute-left-derived-pi-shriek} および
\ref{sites-cohomology-lemma-compute-by-cosimplicial-resolution} から従う。

\begin{lemma}
\label{sites-cohomology-lemma-base-change-by-qis}
$\mathcal{C}$ をサイトとする。$\mathcal{A}_\bullet \to \mathcal{B}_\bullet$ を
$\mathcal{C}$ 上の環の単体的層の準同型とする。
$L\pi_!\mathcal{A}_\bullet \to L\pi_!\mathcal{B}_\bullet$ が
$D(\mathcal{C})$ における同型ならば
$$
L\pi_!(K) =
L\pi_!(K \otimes^\mathbf{L}_{\mathcal{A}_\bullet} \mathcal{B}_\bullet)
$$
がすべての $D(\mathcal{A}_\bullet)$ の $K$ に対して成り立つ。
\end{lemma}

\begin{proof}
$([n], U)$ を $\Delta \times \mathcal{C}$ の対象とする。
$L\pi_!$ は余極限と可換するので、$\mathcal{O}$-加群の上に有界な複体について
証明すれば十分である（Derived Categories, Proposition
\ref{derived-proposition-left-derived-exists} の議論と比較するか、単に上に有界な
複体に限定せよ）。そのような任意の複体は、各項が平坦加群である上に有界な複体と
擬同型である。Modules on Sites, Lemma
\ref{sites-modules-lemma-module-quotient-flat} を参照せよ。
したがって、平坦な $\mathcal{A}_\bullet$-加群 $\mathcal{F}$ について補題を
証明すれば十分である。この場合、導来テンソル積は通常のテンソル積であり、これも
層である。ゆえに補題 \ref{sites-cohomology-lemma-compute-left-derived-pi-shriek} により、等式の
両辺のコホモロジー層を補題
\ref{sites-cohomology-lemma-compute-left-derived-pi-shriek-pre} の手順で計算できる。
したがって、$\mathcal{F}$ のファイバー圏への制限、すなわち
$\Delta \times U$ への制限について結果を証明すれば十分である。
この場合、結果は補題 \ref{sites-cohomology-lemma-O-homology-qis} から従う。
\end{proof}

\begin{remark}
\label{sites-cohomology-remark-homology-augmentation}
$\mathcal{C}$ をサイトとする。
$\epsilon : \mathcal{A}_\bullet \to \mathcal{O}$ を環の層の圏における増大
（Simplicial, Definition
\ref{simplicial-definition-augmentation}）とする。$\epsilon$ が擬同型
$s(\mathcal{A}_\bullet) \to \mathcal{O}$ を誘導すると仮定する。
この場合、三角圏の完全関手
$$
L\pi_! : D(\mathcal{A}_\bullet) \longrightarrow D(\mathcal{O})
$$
を得る。実際、$D(\mathcal{A}_\bullet)$ の任意の対象 $K$ に対して、補題
\ref{sites-cohomology-lemma-base-change-by-qis} により
$L\pi_!(K) = L\pi_!(K \otimes_{\mathcal{A}_\bullet}^\mathbf{L} \mathcal{O})$
である。したがって、表示された関手を
$- \otimes^\mathbf{L}_{\mathcal{A}_\bullet} \mathcal{O}$ と、注意
\ref{sites-cohomology-remark-fibred-category} の関手
$L\pi_! : D(\Delta \times \mathcal{C}, \pi^{-1}\mathcal{O}) \to
D(\mathcal{O})$ の合成として定義できる。言い換えると、補題による
$L\pi_!(K)$ と
$L\pi_!(K \otimes_{\mathcal{A}_\bullet}^\mathbf{L} \mathcal{O})$ の
（標準的かつ関手的な）同一視から、$L\pi_!(K)$ 上の
$\mathcal{O}$-加群構造を得る。
\end{remark}






\section{圏上のコホモロジー}
\label{sites-cohomology-section-cohomology}

\noindent
例 \ref{sites-cohomology-example-category-to-point} の状況では、導来関手 $L\pi_!$ に加えて
関手 $R\pi_*$ もある。$\mathcal{C}$ 上のアーベル層 $\mathcal{F}$ に対して
$H_n(\mathcal{C}, \mathcal{F}) = H^{-n}(L\pi_!\mathcal{F})$ および
$H^n(\mathcal{C}, \mathcal{F}) = H^n(R\pi_*\mathcal{F})$ である。

\begin{example}[コホモロジーの計算]
\label{sites-cohomology-example-right-derived-limits}
例 \ref{sites-cohomology-example-category-to-point} において、関手 $H^n(\mathcal{C}, -)$ は
次のように計算できる。
$\mathcal{F} \in \Ob(\textit{Ab}(\mathcal{C}))$ とし、余鎖複体
$$
K^\bullet(\mathcal{F}) :
\prod\nolimits_{U_0} \mathcal{F}(U_0)
\to
\prod\nolimits_{U_0 \to U_1} \mathcal{F}(U_0)
\to
\prod\nolimits_{U_0 \to U_1 \to U_2} \mathcal{F}(U_0)
\to \ldots
$$
を考える。遷移写像は
$$
(s_{U_0 \to U_1})
\longmapsto
((U_0 \to U_1 \to U_2) \mapsto s_{U_0 \to U_1} - s_{U_0 \to U_2}
+ s_{U_1 \to U_2}|_{U_0})
$$
で与えられ、他の次数でも同様である。構成により
$$
H^0(\mathcal{C}, \mathcal{F}) =
\lim_{\mathcal{C}^{opp}} \mathcal{F} =
H^0(K^\bullet(\mathcal{F})),
$$
である。Categories, Lemma
\ref{categories-lemma-limits-products-equalizers} を参照せよ。
$K^\bullet(\mathcal{F})$ の構成は $\mathcal{F}$ に関して関手的であり、
$\textit{Ab}(\mathcal{C})$ の短完全列を複体の短完全列へ移す。
したがって関手列
$\mathcal{F} \mapsto H^n(K^\bullet(\mathcal{F}))$ は $\delta$-関手をなす。
Homology, Definition
\ref{homology-definition-cohomological-delta-functor} および Lemma
\ref{homology-lemma-long-exact-sequence-cochain} を参照せよ。
$\mathcal{C}$ の対象 $U$ に対し、$p_U^{-1}$ が $U$ での評価に等しい対応する点を
$p_U : \Sh(*) \to \Sh(\mathcal{C})$ と書く。Sites, Example
\ref{sites-example-indiscrete-points} を参照せよ。
$A$ をアーベル群とし、$\mathcal{F} = p_{U, *}A$ と置く。この場合、
複体 $K^\bullet(\mathcal{F})$ の各項は $\text{Map}(X_n, A)$ である。ここで
$$
X_n = \coprod\nolimits_{U_0 \to \ldots \to U_{n - 1} \to U_n}
\Mor_\mathcal{C}(U, U_0)
$$
である。この単体的集合は一元集合 $\{*\}$ 上の定値単体的集合とホモトピー同値で
ある。実際、写像 $X_\bullet \to \{*\}$ は明らかであり、写像
$\{*\} \to X_n$ は $*$ を $(U \to \ldots \to U, \text{id}_U)$ へ送ることで
与えられ、二つの写像 $X_\bullet \to X_\bullet$ の間のホモトピーを定める写像
$$
h_{n, i} : X_n \longrightarrow X_n
$$
（Simplicial, Lemma
\ref{simplicial-lemma-relations-homotopy}）は規則
$$
h_{n, i} :
(U_0 \to \ldots \to U_n, f)
\longmapsto
(U \to \ldots \to U \to U_i \to \ldots \to U_n, \text{id})
$$
によって与えられる。ただし $i > 0$ とし、$h_{n, 0} = \text{id}$ である。
検証は省略する。$\text{Map}(-, A)$ は反変関手なので、
$K^\bullet(p_{U, *}A)$ の正次数のコホモロジーは自明である
（Simplicial, Remark
\ref{simplicial-remark-homotopy-better} の関手性と Simplicial, Lemma
\ref{simplicial-lemma-homotopy-s-Q} の結果による）。
したがって、$\mathcal{F}$ が $p_{U, *}A$ の形の層の積である場合にも
$K^\bullet(\mathcal{F})$ は正次数で非輪状である。
$\mathcal{C}$ 上の任意のアーベル層はそのような積へ埋め込まれるので、
例えば Homology, Lemma
\ref{homology-lemma-efface-implies-universal} および Derived Categories,
Lemma \ref{derived-lemma-right-derived-delta-functor} により、
$K^\bullet(\mathcal{F})$ は $H^0(\mathcal{C}, -)$ の左導来関手
$H^n(\mathcal{C}, -)$ を計算すると結論する。
\end{example}

\begin{example}[Ext の計算]
\label{sites-cohomology-example-computing-exts}
例 \ref{sites-cohomology-example-category-to-point} において、さらに $\mathcal{C}$ 上の環の層
$\mathcal{O}$ が与えられているとする。
$\mathcal{F}$、$\mathcal{G}$ を $\mathcal{O}$-加群とする。次数 $n$ の項が
$$
\prod\nolimits_{U_0 \to U_1 \to \ldots \to U_n}
\Hom_{\mathcal{O}(U_n)}(\mathcal{G}(U_n), \mathcal{F}(U_0))
$$
であり、遷移写像が
$$
(\varphi_{U_0 \to U_1})
\longmapsto
((U_0 \to U_1 \to U_2) \mapsto
\varphi_{U_0 \to U_1} \circ \rho^{U_2}_{U_1}
- \varphi_{U_0 \to U_2}
+ \rho^{U_1}_{U_0} \circ \varphi_{U_1 \to U_2}
$$
で与えられる複体 $K^\bullet(\mathcal{G}, \mathcal{F})$ を考える。
他の次数でも同様である。ここで $\rho$ は制限写像を表す。構成により
$$
\Hom_\mathcal{O}(\mathcal{G}, \mathcal{F}) =
H^0(K^\bullet(\mathcal{G}, \mathcal{F}))
$$
がすべての $\mathcal{O}$-加群の対 $\mathcal{F}, \mathcal{G}$ に対して成り立つ。
対応
$(\mathcal{G}, \mathcal{F}) \mapsto K^\bullet(\mathcal{G}, \mathcal{F})$
は、第一変数の直和を積へ移し、第二変数の積と可換する双関手である。
次を主張する。
$$
\Ext^i_\mathcal{O}(\mathcal{G}, \mathcal{F}) =
H^i(K^\bullet(\mathcal{G}, \mathcal{F}))
$$
$i \geq 0$ に対してこの等式が成り立つための十分条件は、次のいずれかである。
\begin{enumerate}
\item すべての $U \in \Ob(\mathcal{C})$ に対して $\mathcal{G}(U)$ が
射影的 $\mathcal{O}(U)$-加群である。
\item すべての $U \in \Ob(\mathcal{C})$ に対して $\mathcal{F}(U)$ が
入射的 $\mathcal{O}(U)$-加群である。
\end{enumerate}
実際、場合 (1) には関手 $K^\bullet(\mathcal{G}, -)$ は
$\mathcal{O}$-加群の圏からアーベル群の余鎖複体の圏への完全関手である。
したがって例 \ref{sites-cohomology-example-right-derived-limits} と同様に議論すると、
$\mathcal{O}(U)$-加群 $A$ に対して
$\mathcal{F} = p_{U, *}A$ であるとき
$K^\bullet(\mathcal{G}, \mathcal{F})$ が正次数で非輪状であることを示せば十分である。
短完全列
\begin{equation}
\label{sites-cohomology-equation-split}
0 \to \mathcal{G}' \to \bigoplus j_{U_i!}\mathcal{O}_{U_i} \to
\mathcal{G} \to 0
\end{equation}
を選ぶ。Modules on Sites, Lemma
\ref{sites-modules-lemma-module-quotient-flat} を参照せよ。
(1) は中央と右の層について成り立つので、$\mathcal{G}'$ についても成り立ち、
(\ref{sites-cohomology-equation-split}) を $\mathcal{C}$ の対象上で評価すると加群の分裂完全列を得る。
任意の $\mathcal{F}$ に対して複体の短完全列
$$
0 \to
K^\bullet(\mathcal{G}, \mathcal{F}) \to
\prod K^\bullet(j_{U_i!}\mathcal{O}_{U_i}, \mathcal{F}) \to
K^\bullet(\mathcal{G}', \mathcal{F}) \to 0
$$
を得る。特に $\mathcal{F} = p_{U, *}A$ と取れる。
$H^0$ 上では
$$
0 \to \Hom(\mathcal{G}, p_{U, *}A) \to
\Hom(\prod j_{U_i!}\mathcal{O}_{U_i}, p_{U, *}A) \to
\Hom(\mathcal{G}', p_{U, *}A) \to 0
$$
を得る。これは
$\Hom(\mathcal{H}, p_{U, *}A) = \Hom_{\mathcal{O}(U)}(\mathcal{H}(U), A)$
であり、(\ref{sites-cohomology-equation-split}) の $U$ 上の切断の列が分裂完全なので完全である。
したがって次元移動を用いると、すべての
$U, U' \in \Ob(\mathcal{C})$ に対して
$K^\bullet(j_{U'!}\mathcal{O}_{U'}, p_{U, *}A)$ が正次数で非輪状であることを
証明すれば十分だと分かる。この場合
$K^n(j_{U'!}\mathcal{O}_{U'}, p_{U, *}A)$ は
$$
\prod\nolimits_{U \to U_0 \to U_1 \to \ldots \to U_n \to U'} A
$$
に等しい。言い換えると、
$K^\bullet(j_{U'!}\mathcal{O}_{U'}, p_{U, *}A)$ は各項が
$\text{Map}(X_\bullet, A)$ である複体である。ここで
$$
X_n = \coprod\nolimits_{U_0 \to \ldots \to U_{n - 1} \to U_n}
\Mor_\mathcal{C}(U, U_0) \times \Mor_\mathcal{C}(U_n, U')
$$
である。この単体的集合は一元集合 $\{*\}$ 上の定値単体的集合とホモトピー同値で
ある。これは例 \ref{sites-cohomology-example-right-derived-limits} の対応する主張とまったく同じ
方法で証明できる。これで主張の証明は完了する。

\medskip\noindent
場合 (2) の議論も同様である（ただし双対である）。
\end{example}






\section{圏上の加群}
\label{sites-cohomology-section-modules-cohomology}

\noindent
本節の内容は、スキームの圏上の亜群スタックにおける準連接加群の導来圏の
一変種を定義するために用いられる。Sheaves on Stacks, Section
\ref{stacks-sheaves-section-QC} を参照せよ。

\medskip\noindent
$\mathcal{C}$ を圏とし、カオス位相を入れたサイトとみなす。
例 \ref{sites-cohomology-example-computing-exts} と同様に、$\mathcal{O}$ を
$\mathcal{C}$ 上の環の層とする。言い換えると、$\mathcal{O}$ は圏
$\mathcal{C}$ 上の環の前層である。Categories, Definition
\ref{categories-definition-presheaf} を参照せよ。

\begin{definition}
\label{sites-cohomology-definition-cartesian}
上の状況で、$D(\mathcal{O}) = D(\mathcal{C}, \mathcal{O})$ の対象 $K$ で、
$\mathcal{C}$ のすべての $U \to V$ に対して標準写像
$$
R\Gamma(V, K) \otimes_{\mathcal{O}(V)}^\mathbf{L} \mathcal{O}(U)
\longrightarrow
R\Gamma(U, K)
$$
が $D(\mathcal{O}(U))$ における同型となるものからなる充満部分圏を
$\mathit{QC}(\mathcal{C}, \mathcal{O})$、または単に
{\it $\mathit{QC}(\mathcal{O})$} と書く。
\end{definition}

\begin{lemma}
\label{sites-cohomology-lemma-cartesian-good-sub}
上の状況で、部分圏 $\mathit{QC}(\mathcal{O})$ は $D(\mathcal{O})$ の
厳密充満で飽和な三角部分圏であり、任意の直和で保たれる。
\end{lemma}

\begin{proof}
$U$ を $\mathcal{C}$ の対象とする。$\mathcal{C}$ 上の位相はカオスなので、
関手 $\mathcal{F} \mapsto \mathcal{F}(U)$ は完全であり、直和と可換する。
したがって完全関手 $K \mapsto R\Gamma(U, K)$ は、$K$ を任意の
$\mathcal{O}$-加群の複体 $\mathcal{F}^\bullet$ で表し
$\mathcal{F}^\bullet(U)$ を取ることで計算される。ゆえに
$R\Gamma(U, -)$ は直和と可換する。Injectives, Lemma
\ref{injectives-lemma-derived-products} を参照せよ。
同様に、$\mathcal{C}$ の射 $U \to V$ が与えられると、導来テンソル積関手
$- \otimes_{\mathcal{O}(V)}^\mathbf{L} \mathcal{O}(U) :
D(\mathcal{O}(V)) \to D(\mathcal{O}(U))$
は完全であり、直和と可換する。これらの観察から補題は直ちに従う。
詳細は省略する。
\end{proof}

\begin{lemma}
\label{sites-cohomology-lemma-QC-bounded-above}
上の状況で、$M$ を $\mathit{QC}(\mathcal{O})$ の対象とし、
$b \in \mathbf{Z}$ はすべての $i > b$ に対して $H^i(M) = 0$ を満たすとする。
このとき $H^b(M)$ は Modules on Sites, Definition
\ref{sites-modules-definition-site-local} の意味で
$(\mathcal{C}, \mathcal{O})$ 上の準連接加群である。
\end{lemma}

\begin{proof}
Modules on Sites, Lemma
\ref{sites-modules-lemma-chaotic-quasi-coherent} により、
$\mathcal{C}$ のすべての射 $U \to V$ に対して写像
$$
H^p(M)(V) \otimes_{\mathcal{O}(V)} \mathcal{O}(U) \to H^b(M)(U)
$$
が同型であることを示せば十分である。仮定により写像
$$
R\Gamma(V, M) \otimes_{\mathcal{O}(V)}^\mathbf{L} \mathcal{O}(U)
\to R\Gamma(U, M)
$$
は同型である。したがって、例えば Tor スペクトル系列により結果が従う。
詳細は省略する。
\end{proof}

\begin{lemma}
\label{sites-cohomology-lemma-QC-final-object}
上の状況で、$\mathcal{C}$ が終対象 $X$ をもつとする。
$R = \mathcal{O}(X)$ と置き、
$f : (\mathcal{C}, \mathcal{O}) \to (pt, R)$ を明らかなサイトの射とする。
このとき $Lf^*$ と $Rf_*$ により
$\mathit{QC}(\mathcal{O}) = D(R)$ である。
\end{lemma}

\begin{proof}
省略する。
\end{proof}

\begin{lemma}
\label{sites-cohomology-lemma-QC-hom-out-of}
上の状況で、$K$ を $\mathit{QC}(\mathcal{O})$ の対象とし、
$M$ を $D(\mathcal{O})$ の任意の対象とする。
$\mathcal{C}$ のすべての対象 $U$ に対して
$$
\Hom_{D(\mathcal{O}_U)}(K|_U, M|_U) =
R\Hom_{\mathcal{O}(U)}(R\Gamma(U, K), R\Gamma(U, M))
$$
である。
\end{lemma}

\begin{proof}
$\mathcal{C}$ を $\mathcal{C}/U$ で置き換えてよい。したがって
$U = X$ が $\mathcal{C}$ の終対象であると仮定してよい。
補題 \ref{sites-cohomology-lemma-QC-final-object} により $K = Lf^*P$ である。ここで
$P = R\Gamma(U, K) = R\Gamma(X, K) = Rf_*K$ である。
$Lf^*$ は $Rf_*(-) = R\Gamma(U, -)$ の左随伴なので、結果が従う。
\end{proof}

\noindent
$(\mathcal{C}, \mathcal{O})$ を上のとおりとする。
$\mathcal{O}$-加群の複体 $\mathcal{F}^\bullet$ の
{\it 大きさ} $|\mathcal{F}^\bullet|$ を
$$
|\mathcal{F}^\bullet| =
\left|
\coprod\nolimits_{i \in \mathbf{Z},\ U \in \Ob(\mathcal{C})} \mathcal{F}^i(U)
\right|
$$
で定義する。$D(\mathcal{O})$ の対象 $K$ の{\it 大きさ} $|K|$ を濃度
$$
|K| = \min
\left\{
\left|
\mathcal{F}^\bullet
\right|
\text{ where }\mathcal{F}^\bullet\text{ represents }K
\right\}
$$
で定義する。濃度の性質により、この最小値は存在する。

\begin{lemma}
\label{sites-cohomology-lemma-build-from}
上の状況で、次の性質をもつ濃度 $\kappa$ が存在する。
$\mathcal{O}$-加群の複体 $\mathcal{F}^\bullet$ と部分集合
$\Omega^i_U \subset \mathcal{F}^i(U)$ が与えられると、
$\Omega^i_U \subset \mathcal{H}^i(U)$ および
$|\mathcal{H}^\bullet| \leq \max(\kappa, |\bigcup \Omega^i_U|)$ を満たす
部分複体 $\mathcal{H}^\bullet \subset \mathcal{F}^\bullet$ が存在する。
\end{lemma}

\begin{proof}
$\mathcal{H}^i(U)$ を、任意の射 $f : U \to V$ に沿う制限による
$\Omega^i_V$ の像と $\text{d}(\Omega^{i - 1}_U)$ で生成される
$\mathcal{F}^i(U)$ の $\mathcal{O}(U)$-部分加群とする。
$\mathcal{H}^i(U)$ の濃度は、$\aleph_0$、$\mathcal{O}(U)$ の濃度、
$\text{Arrows}(\mathcal{C})$ の濃度、および
$|\bigcup \Omega^i_U|$ の最大値で抑えられる。詳細は省略する。
\end{proof}

\begin{lemma}
\label{sites-cohomology-lemma-qis-small}
上の状況で、次の性質をもつ濃度 $\kappa$ が存在する。
$D(\mathcal{O})$ の対象 $K$ を表す $\mathcal{O}$-加群の複体
$\mathcal{F}^\bullet$ が与えられると、$K$ を表し、かつ
$|\mathcal{H}^\bullet| \leq \max(\kappa, |K|)$ を満たす部分複体
$\mathcal{H}^\bullet \subset \mathcal{F}^\bullet$ が存在する。
\end{lemma}

\begin{proof}
まず、各 $i$ と $U$ に対して部分集合
$\Omega^i_U \subset \Ker(\text{d} :
\mathcal{F}^i(U) \to \mathcal{F}^{i + 1}(U))$
であって
$H^i(K)(U) = H^i(\mathcal{F}^\bullet(U))$.
へ全単射に写るものを選ぶ。$K$ は大きさ $|K|$ の複体で表せるので
$|\Omega^i_U| \leq |K|$ である。補題 \ref{sites-cohomology-lemma-build-from} を適用すると、
$\Omega^i_U$ を含み、大きさが高々 $\max(\kappa, |K|)$ である部分複体
$\mathcal{S}^\bullet \subset \mathcal{F}^\bullet$ を得る。
したがって $H^i(\mathcal{S}^\bullet) \to H^i(\mathcal{F}^\bullet)$ は
層の全射である。

\medskip\noindent
部分複体
$$
\mathcal{S}^\bullet =
\mathcal{S}_0^\bullet \subset
\mathcal{S}_1^\bullet \subset
\mathcal{S}_2^\bullet \subset \ldots \subset \mathcal{F}^\bullet
$$
を帰納的に構成する。これらの大きさは
$\leq \max(\kappa, |K|)$ であり、
$H^i(\mathcal{S}_n^\bullet) \to H^i(\mathcal{F}^\bullet)$ の核は
$H^i(\mathcal{S}_n^\bullet) \to H^i(\mathcal{S}_{n + 1}^\bullet)$ の核と
同じであるようにする。これができれば
$\mathcal{H}^\bullet = \bigcup \mathcal{S}_n^\bullet$ を解として取れる。

\medskip\noindent
与えられた $\mathcal{S}_n^\bullet$ から $\mathcal{S}_{n + 1}^\bullet$ を
構成する。すべての $U$ と $i$ に対して
$\Omega^{i - 1}_U \subset \mathcal{F}^{i - 1}(U)$ を、
$\text{d} : \mathcal{F}^{i - 1}(U) \to \mathcal{F}^i(U)$ が
$\Omega^{i - 1}_U$ を
$$
\mathcal{S}_n^i(U) \cap
\text{Im}(\text{d} : \mathcal{F}^{i - 1}(U)\to \mathcal{F}^i(U))
$$
へ全単射に写す部分集合とする。
$\mathcal{S}_n^i(U)$ はそのように抑えられているので
$|\Omega^i_U| \leq |K|$ であることに注意する。
部分集合
$$
\mathcal{S}^i(U) \cup \Omega^i_U \subset \mathcal{F}^i(U)
$$
に補題 \ref{sites-cohomology-lemma-build-from} を適用して
$\mathcal{S}_{n + 1}^\bullet$ を得る。残りは明らかである。
\end{proof}

\begin{lemma}
\label{sites-cohomology-lemma-cartesian-colimit-kappa}
上の状況で、次の性質をもつ濃度 $\kappa$ が存在する。
\begin{enumerate}
\item $\mathit{QC}(\mathcal{O})$ のすべての非零対象 $K$ に対して、
$|E| \leq \kappa$ を満たす $\mathit{QC}(\mathcal{O})$ の非零射
$E \to K$ が存在する。
\item $\mathit{QC}(\mathcal{O})$ の任意の射
$\alpha : E \to \bigoplus_n K_n$ で $|E| \leq \kappa$ を満たすものに対し、
$|E_n| \leq \kappa$ を満たす $\mathit{QC}(\mathcal{O})$ の射
$E_n \to K_n$ が存在し、$\alpha$ は
$\bigoplus E_n \to \bigoplus K_n$ を経由する。
\end{enumerate}
\end{lemma}

\begin{proof}
$\kappa$ を次の濃度の集合の上界とする。
\begin{enumerate}
\item すべての $U \in \Ob(\mathcal{C})$ に対する
$|\coprod_V j_{U!}\mathcal{O}_U(V)|$。
\item $\mathcal{C}$ のすべての射 $U \to V$ に対して More on Algebra, Lemma
\ref{more-algebra-lemma-enlarge} で得られる濃度
$\kappa(\mathcal{O}(V) \to \mathcal{O}(U))$。
\item 補題 \ref{sites-cohomology-lemma-qis-small} で得られる濃度。
\end{enumerate}
$\mathit{QC}(\mathcal{O})$ の対象を表す任意の複体
$\mathcal{F}^\bullet$ と、$|\mathcal{S}^\bullet| \leq \kappa$ を満たす任意の
部分複体 $\mathcal{S}^\bullet \subset \mathcal{F}^\bullet$ に対し、
$\mathcal{S}^\bullet$ を含む $\mathcal{F}^\bullet$ の部分複体
$\mathcal{H}^\bullet$ であって、$\mathit{QC}(\mathcal{O})$ の対象を表し、
$|\mathcal{H}^\bullet| \leq \kappa$ を満たすものが存在すると主張する。
次の二段落で、この主張から補題が従うことを示す。

\medskip\noindent
(1) のように、$K$ を $\mathit{QC}(\mathcal{O})$ の非零対象とし、
$\mathcal{O}$-加群の複体 $\mathcal{F}^\bullet$ で表す。
するとある $i$ に対して $H^i(\mathcal{F}^\bullet)$ は非零である。
したがって $\mathcal{C}$ の対象 $U$ と、$d(s) = 0$ を満たし
$U$ 上で $H^i(\mathcal{F}^\bullet)$ の非零切断を定める切断
$s \in \mathcal{F}^i(U)$ が存在する。このとき
$s : j_{U!}\mathcal{O}_U[-i] \to \mathcal{F}^\bullet$ の像は、
$|\mathcal{S}^\bullet| \leq \kappa$ を満たす部分複体
$\mathcal{S}^\bullet \subset \mathcal{F}^\bullet$ である。
主張を適用すると、$|\mathcal{H}^\bullet| \leq \kappa$ を満たす
$\mathit{QC}(\mathcal{O})$ の非零射
$\mathcal{H}^\bullet \to \mathcal{F}^\bullet$ を得る。よって (1) が成り立つ。

\medskip\noindent
$\alpha : E \to \bigoplus K_n$ を (2) のとおりとする。
$K_n$ を表す任意の複体 $\mathcal{K}_n^\bullet$ を選ぶ。
すると $\bigoplus \mathcal{K}_n^\bullet$ は $\bigoplus K_n$ を表す。
導来圏の構成により、$E$ を複体 $\mathcal{E}^\bullet$ で表し、
$\alpha$ を複体の射
$a : \mathcal{E}^\bullet \to \bigoplus \mathcal{K}_n^\bullet$ で
表すことができる。補題 \ref{sites-cohomology-lemma-qis-small} と上の $\kappa$ の選び方により
$|\mathcal{E}^\bullet| \leq \kappa$ と仮定してよい。
主張により、$a_n : \mathcal{E}^\bullet \to \mathcal{K}_n^\bullet$ の像を含み、
$|E_n| \leq \kappa$ を満たす $\mathit{QC}(\mathcal{O})$ の対象 $E_n$ を表す
部分複体 $\mathcal{E}_n^\bullet \subset \mathcal{K}_n^\bullet$ を得る。
これは望んだものである。

\medskip\noindent
主張を証明する。$\mathcal{F}^\bullet$ を $\mathit{QC}(\mathcal{O})$ の
対象を表す複体とし、$\mathcal{S}^\bullet \subset \mathcal{F}^\bullet$ を
大きさ $\leq \kappa$ の部分複体とする。大きさ $\leq \kappa$ の部分複体
$$
\mathcal{S}^\bullet = \mathcal{S}_0^\bullet \subset
\mathcal{S}_1^\bullet \subset
\mathcal{S}_2^\bullet \subset \ldots \subset \mathcal{F}^\bullet
$$
を帰納的に構成し、$\mathcal{C}$ のすべての射 $f : U \to V$ と
すべての $i \in \mathbf{Z}$ に対して次を満たすようにする。
\begin{enumerate}
\item 射
$H^i(\mathcal{S}_n^\bullet(V)
\otimes_{\mathcal{O}(V)}^\mathbf{L} \mathcal{O}(U)) \to
H^i(\mathcal{S}_n^\bullet(U))$
の核は
$H^i(\mathcal{S}_{n + 1}^\bullet(V)
\otimes_{\mathcal{O}(V)}^\mathbf{L} \mathcal{O}(U))$ で零に写る。
\item 射
$H^i(\mathcal{S}_n^\bullet(U)) \to H^i(\mathcal{S}_{n + 1}^\bullet(U))$
の像は
$H^i(\mathcal{S}_{n + 1}^\bullet(V)
\otimes_{\mathcal{O}(V)}^\mathbf{L} \mathcal{O}(U)) \to
H^i(\mathcal{S}_{n + 1}^\bullet(U))$ の像に含まれる。
\end{enumerate}
これができれば
$\mathcal{H}^\bullet = \bigcup \mathcal{S}_n^\bullet$.
とおける。実際、導来テンソル積と環上の加群複体のコホモロジーを取る操作は
フィルター余極限と可換なので、条件 (1) と (2) を合わせると
$$
\mathcal{H}^\bullet(V) \otimes_{\mathcal{O}(V)}^\mathbf{L} \mathcal{O}(U)
\longrightarrow
\mathcal{H}^\bullet(U)
$$
がすべての次数でコホモロジー上の同型となり、したがって
$\mathcal{C}$ のすべての $f : U \to V$ に対して
$D(\mathcal{O}(U))$ における同型となることが保証される。
よって $\mathcal{H}^\bullet$ は望みどおり
$\mathit{QC}(\mathcal{O})$ の対象を表す。

\medskip\noindent
与えられた $\mathcal{S}_n$ から $\mathcal{S}_{n + 1}$ を構成する。
$\mathcal{C}$ のすべての射 $f : U \to V$ に対して可換図式
$$
\xymatrix{
\mathcal{S}_n^\bullet(V)
\ar[r] \ar[d] &
\mathcal{S}_n^\bullet(U) \ar[d] \\
\mathcal{F}^\bullet(V)
\ar[r] &
\mathcal{F}^\bullet(U)
}
$$
を考える。これは環準同型 $\mathcal{O}(V) \to \mathcal{O}(U)$ に対する
More on Algebra, Lemma \ref{more-algebra-lemma-enlarge} の図式である。
すなわち下段は
$$
\mathcal{F}^\bullet(V) \otimes_{\mathcal{O}(V)}^\mathbf{L} \mathcal{O}(U)
\longrightarrow
\mathcal{F}^\bullet(U)
$$
という $D(\mathcal{O}(U))$ における同型を誘導する。したがって部分複体
$$
\mathcal{S}_n^\bullet(V) \subset M^\bullet_f \subset \mathcal{F}^\bullet(V)
\quad\text{and}\quad
\mathcal{S}_n^\bullet(U) \subset N^\bullet_f \subset \mathcal{F}^\bullet(U)
$$
を More on Algebra, Lemma \ref{more-algebra-lemma-enlarge} のとおりに選べる。
特に $|N^i_f|, |M^i_f| \leq \kappa$ である。次に部分集合
$$
\mathcal{S}_n^i(U) \amalg \coprod\nolimits_{f : U \to V} N^i_f
\amalg \coprod\nolimits_{g : W \to U} M^i_g
\subset
\mathcal{F}^i(U)
$$
を用いて補題 \ref{sites-cohomology-lemma-build-from} を適用し、部分複体
$$
\mathcal{S}_n^\bullet \subset
\mathcal{S}_{n + 1}^\bullet \subset
\mathcal{F}^\bullet
$$
を得る。これはそれらの部分集合を含み、
$|\mathcal{S}_{n + 1}^\bullet| \leq \kappa$ を満たす。
More on Algebra, Lemma \ref{more-algebra-lemma-enlarge} の構成により、
$\mathcal{S}_n^\bullet(V) \subset M^\bullet_f$ および
$\mathcal{S}_n^\bullet(U) \subset N^\bullet_f$ に対して対応する主張が
成り立つので、条件 (1) と (2) が成り立つ。これで証明は完了する。
\end{proof}

\begin{proposition}
\label{sites-cohomology-proposition-cartesian-brown}
$\mathcal{C}$ を圏とし、カオス位相を入れたサイトとみなす。
$\mathcal{O}$ を $\mathcal{C}$ 上の環の層とする。
$\mathit{QC}(\mathcal{O})$ を定義 \ref{sites-cohomology-definition-cartesian} のとおりとすると、
\begin{enumerate}
\item $\mathit{QC}(\mathcal{O})$ は $D(\mathcal{O})$ の厳密充満で飽和な
三角部分圏であり、任意の直和で保たれる。
\item 直和を直積に移す任意の反変コホモロジー関手
$H : \mathit{QC}(\mathcal{O}) \to \textit{Ab}$
は表現可能である。
\item 直和を直和に移す三角圏の任意の完全関手
$F : \mathit{QC}(\mathcal{O}) \to \mathcal{D}$ は完全な右随伴をもつ。
\item 包含関手 $\mathit{QC}(\mathcal{O}) \to D(\mathcal{O})$ は
完全な右随伴をもつ。
\end{enumerate}
\end{proposition}

\begin{proof}
(1) は補題 \ref{sites-cohomology-lemma-cartesian-good-sub} である。
(2) は補題 \ref{sites-cohomology-lemma-cartesian-colimit-kappa} と
Derived Categories, Lemma \ref{derived-lemma-brown-bis} から従う。
(3) は補題 \ref{sites-cohomology-lemma-cartesian-colimit-kappa} と
Derived Categories, Proposition \ref{derived-proposition-brown-bis} から従う。
(4) は (3) の特別な場合である。
\end{proof}

\noindent
$u : \mathcal{C}' \to \mathcal{C}$ を圏の間の関手とする。
$\mathcal{C}$ と $\mathcal{C}'$ をカオス位相を入れたサイトとみなすと、
$u$ は連続かつ余連続な関手である。したがってトポスの射
$g : \Sh(\mathcal{C}') \to \Sh(\mathcal{C})$ を得る。Sites, Lemma
\ref{sites-lemma-cocontinuous-morphism-topoi} を参照せよ。
さらに、$\mathcal{C}$ 上の環の層 $\mathcal{O}$、
$\mathcal{C}'$ 上の環の層 $\mathcal{O}'$、および写像
$g^\sharp : g^{-1}\mathcal{O} \to \mathcal{O}'$.
が与えられているとする。対応する環付きトポスの射を単に
$g : (\Sh(\mathcal{C}'), \mathcal{O}') \to (\Sh(\mathcal{C}), \mathcal{O})$
と書く。Modules on Sites, Section \ref{sites-modules-section-ringed-topoi}
を参照せよ。

\begin{lemma}
\label{sites-cohomology-lemma-cartesian-functorial}
$g : (\Sh(\mathcal{C}'), \mathcal{O}') \to (\Sh(\mathcal{C}), \mathcal{O})$
を上のとおりとする。このとき関手 $Lg^* : D(\mathcal{O}) \to D(\mathcal{O}')$
は $\mathit{QC}(\mathcal{O})$ を $\mathit{QC}(\mathcal{O}')$ の中へ写す。
\end{lemma}

\begin{proof}
$U' \in \Ob(\mathcal{C}')$ とし、$\mathcal{C}$ におけるその像を
$U = u(U')$ とする。$pt$ で一つの対象と一つの射だけをもつ圏を表す。
表示された環付きトポスを $(\Sh(pt), \mathcal{O}'(U'))$ および
$(\Sh(pt), \mathcal{O}(U))$ と書く。もちろん、これらの環付きトポス上の
加群の導来圏をそれぞれ $D(\mathcal{O}'(U'))$ および
$D(\mathcal{O}(U))$ と同一視する。このとき環付きトポスの可換図式
$$
\xymatrix{
(\Sh(pt), \mathcal{O}'(U')) \ar[rr]_{U'} \ar[d] & &
(\Sh(\mathcal{C}'), \mathcal{O}') \ar[d]^g \\
(\Sh(pt), \mathcal{O}(U)) \ar[rr]^U & &
(\Sh(\mathcal{C}), \mathcal{O})
}
$$
を得る。下の水平射による引き戻しは $D(\mathcal{O})$ の $K$ を
$R\Gamma(U, K)$ に写す。左の垂直射による引き戻しは $M$ を
$M \otimes_{\mathcal{O}(U)}^\mathbf{L} \mathcal{O}'(U')$ に写す。
図式をどちらの向きに回っても同じ結果を得る
（補題 \ref{sites-cohomology-lemma-derived-pullback-composition}）ので、
$$
R\Gamma(U', Lg^*K) =
R\Gamma(U, K) \otimes_{\mathcal{O}(U)}^\mathbf{L} \mathcal{O}'(U')
$$
を得る。最後に、$f' : U' \to V'$ を $\mathcal{C}'$ の射とし、
$\mathcal{C}$ におけるその像を
$f = u(f') : U = u(U') \to V = u(V')$ と書く。
$K$ が $\mathit{QC}(\mathcal{O})$ に属するならば
\begin{align*}
R\Gamma(V', Lg^*K) \otimes_{\mathcal{O}'(V')}^\mathbf{L} \mathcal{O}'(U')
& =
R\Gamma(V, K) \otimes_{\mathcal{O}(V)}^\mathbf{L} \mathcal{O}'(V')
\otimes_{\mathcal{O}'(V')}^\mathbf{L} \mathcal{O}'(U') \\
& =
R\Gamma(V, K) \otimes_{\mathcal{O}(V)}^\mathbf{L} \mathcal{O}'(U') \\
& =
R\Gamma(V, K) \otimes_{\mathcal{O}(V)}^\mathbf{L} \mathcal{O}(U)
\otimes_{\mathcal{O}(U)}^\mathbf{L} \mathcal{O}'(U') \\
& =
R\Gamma(U, K) \otimes_{\mathcal{O}(U)}^\mathbf{L} \mathcal{O}'(U') \\
& =
R\Gamma(U', Lg^*K)
\end{align*}
となり、これは望んだものである。ここでは上の観察を $U'$ と $V'$ の
両方に用いた。
\end{proof}

\begin{lemma}
\label{sites-cohomology-lemma-cartesian-flat-quasi-coherent}
$\mathcal{C}$ を圏とし、カオス位相を入れたサイトとみなす。
$\mathcal{O}$ を $\mathcal{C}$ 上の環の層とする。
$\mathcal{C}$ のすべての $U \to V$ に対して制限写像
$\mathcal{O}(V) \to \mathcal{O}(U)$ が平坦な環準同型であると仮定する。
このとき $\mathit{QC}(\mathcal{O})$ は、コホモロジー層が準連接である複体からなる
部分圏 $D_\QCoh(\mathcal{O}) \subset D(\mathcal{O})$ と一致する。
\end{lemma}

\begin{proof}
仮定のもとで $\QCoh(\mathcal{O}) \subset \textit{Mod}(\mathcal{O})$ は
弱 Serre 部分圏であることを思い出す。Modules on Sites, Lemma
\ref{sites-modules-lemma-chaotic-flat-quasi-coherent} を参照せよ。
したがって $D(\mathcal{O})$ の充満部分圏
$$
D_\QCoh(\mathcal{O}) = D_{\QCoh(\mathcal{O})}(\textit{Mod}(\mathcal{O}))
$$
を取ることができる。Derived Categories, Section
\ref{derived-section-triangulated-sub} を参照せよ。
（厳密にいえば、補題の証明にはこれは必要ない。）

\medskip\noindent
$M$ を $\mathit{QC}(\mathcal{O})$ の対象とする。$\mathcal{C}$ のすべての射
$U \to V$ に対して制限写像
$\mathcal{O}(V) \to \mathcal{O}(U)$ は平坦なので、
\begin{align*}
H^i(M)(U)
& =
H^i(R\Gamma(U, M)) \\
& =
H^i(R\Gamma(V, M) \otimes_{\mathcal{O}(V)}^\mathbf{L} \mathcal{O}(U)) \\
& =
H^i(R\Gamma(V, M)) \otimes_{\mathcal{O}(V)} \mathcal{O}(U) \\
& =
H^i(M)(V) \otimes_{\mathcal{O}(V)} \mathcal{O}(U)
\end{align*}
となる。ゆえに Modules on Sites, Lemma
\ref{sites-modules-lemma-chaotic-quasi-coherent} により
$H^i(M)$ は準連接である。上の最初と最後の等式は、
$\mathcal{C}$ 上の位相がカオスであるため、$\mathcal{C}$ の対象上の切断を
取る関手が完全であることから従う。

\medskip\noindent
逆に $M$ が $D_\QCoh(\mathcal{O})$ の対象ならば、Modules on Sites, Lemma
\ref{sites-modules-lemma-chaotic-quasi-coherent} により、写像
$R\Gamma(V, M) \to R\Gamma(U, M)$ は同型
$H^i(M)(U) \to H^i(M)(V) \otimes_{\mathcal{O}(V)} \mathcal{O}(U)$
を誘導する。したがって $\mathcal{O}(V) \to \mathcal{O}(U)$ の平坦性により
$R\Gamma(V, K) \otimes_{\mathcal{O}(V)}^\mathbf{L} \mathcal{O}(U)
\to R\Gamma(U, K)$ は $D(\mathcal{O}(U))$ における同型であり、
$M$ は $\mathit{QC}(\mathcal{O})$ に属すると結論する。
\end{proof}

\begin{lemma}
\label{sites-cohomology-lemma-cartesion-plus-topology}
$\epsilon : (\mathcal{C}_\tau, \mathcal{O}_\tau) \to
(\mathcal{C}_{\tau'}, \mathcal{O}_{\tau'})$ を
節 \ref{sites-cohomology-section-compare} のとおりとする。次を仮定する。
\begin{enumerate}
\item $\tau'$ は圏 $\mathcal{C}$ 上のカオス位相である。
\item すべての $U \in \Ob(\mathcal{C})$ と、すべての
$\mathcal{O}(U)$-加群の K-平坦複体 $M^\bullet$ に対して写像
$$
M^\bullet \longrightarrow
R\Gamma((\mathcal{C}/U)_\tau,
(M^\bullet \otimes_{\mathcal{O}(U)} \mathcal{O}_U)^\#)
$$
は擬同型である（説明については証明を参照せよ）。
\end{enumerate}
このとき $\epsilon^*$ と $R\epsilon_*$ は、$\mathit{QC}(\mathcal{O})$ と、
$R\epsilon_*K$ が $\mathit{QC}(\mathcal{O})$ に属するような対象 $K$ からなる
$D(\mathcal{C}_\tau, \mathcal{O}_\tau)$ の充満部分圏との間に、
互いに擬逆な同値を定める\footnote{これは、$\mathcal{C}$ のすべての
$U \to V$ に対して
$R\Gamma(V, K) \otimes_{\mathcal{O}(V)}^\mathbf{L} \mathcal{O}(U)
\to R\Gamma(U, K)$ が同型であることを意味する。}。
\end{lemma}

\begin{proof}
節 \ref{sites-cohomology-section-compare} で行った観察を、以後断りなく用いる。
$R\epsilon_*$ は充満忠実であり
$\epsilon^* \circ R\epsilon_* = \text{id}$ なので、補題を証明するには
$\mathit{QC}(\mathcal{O})$ の $M$ に対して
$R\epsilon_*(\epsilon^*M) = M$ であることを示せば十分である。
条件 (2) は、まさにこれを示すために必要な条件である。
実際、$M$ を表す、各項が平坦な $\mathcal{O}$-加群の K-平坦複体
$\mathcal{M}^\bullet$ を選ぶ。すると $\epsilon^*M$ は
$\mathcal{M}^\bullet$ の $\tau$-層化
$(\mathcal{M}^\bullet)^\#$ で表される。
$U \in \Ob(\mathcal{C})$ とする。Leray により
$$
R\Gamma(U, R\epsilon_*(\epsilon^*M)) =
R\Gamma((\mathcal{C}/U)_\tau, (\mathcal{M}^\bullet)^\#|_{\mathcal{C}/U}) =
R\Gamma((\mathcal{C}/U)_\tau, (\mathcal{M}^\bullet|_{\mathcal{C}/U})^\#)
$$
を得る。最後の等式は、層化が $\mathcal{C}/U$ への制限と可換するためである。
いつものように、$\mathcal{O}$ の $\mathcal{C}/U$ への制限を
$\mathcal{O}_U$ と書く。写像
$$
\mathcal{M}^\bullet(U) \otimes_{\mathcal{O}(U)} \mathcal{O}_U
\longrightarrow
\mathcal{M}^\bullet|_{\mathcal{C}/U}
$$
を考える。これは（$\tau'$-位相における）$\mathcal{O}_U$-加群の複体の射である。
$\mathcal{M}^\bullet$ の選び方により、複体 $\mathcal{M}^\bullet(U)$ は
$\mathcal{O}(U)$-加群の K-平坦複体である。補題
\ref{sites-cohomology-lemma-pullback-K-flat} を参照し、$U$ の $\mathcal{C}$ への包含が
環付きトポスの射
$(\Sh(pt), \mathcal{O}(U)) \to (\Sh(\mathcal{C}_{\tau'}), \mathcal{O})$.
を定めることを用いよ。$M$ は $\mathit{QC}(\mathcal{O})$ に属するので、
表示された射は擬同型である。層化は完全だから、層化した後も同じことが
成り立つ。したがって
$$
R\Gamma(U, R\epsilon_*(\epsilon^*M)) =
R\Gamma((\mathcal{C}/U)_\tau,
(M^\bullet \otimes_{\mathcal{O}(U)} \mathcal{O}_U)^\#)
$$
となり、仮定 (2) によれば、これは望みどおり
$R\Gamma(U, M) = \mathcal{M}^\bullet(U)$ に等しい。
\end{proof}

\begin{lemma}
\label{sites-cohomology-lemma-QC-hom-out-of-plus-topology}
記法と仮定を補題 \ref{sites-cohomology-lemma-cartesion-plus-topology} のとおりとする。
$K$ を $\mathit{QC}(\mathcal{O})$ の対象とし、$M$ を
$D(\mathcal{O}_\tau)$ の任意の対象とする。
$\mathcal{C}$ のすべての対象 $U$ に対して
$$
\Hom_{D((\mathcal{O}_U)_\tau)}(\epsilon^*K|_U, M|_U) =
R\Hom_{\mathcal{O}(U)}(R\Gamma(U, K), R\Gamma(U, M))
$$
が成り立つ。
\end{lemma}

\begin{proof}
随伴により
$$
\Hom_{D((\mathcal{O}_U)_\tau)}(\epsilon^*K|_U, M|_U) =
\Hom_{D((\mathcal{O}_U)_{\tau'})}(K|_U, R\epsilon_*M|_U)
$$
である。したがって補題 \ref{sites-cohomology-lemma-QC-hom-out-of} と、Leray による等式
$$
R\Gamma(U, M) = R\Gamma(U, R\epsilon_*M)
$$
から結果が従う。
\end{proof}







\section{狭義完全複体}
\label{sites-cohomology-section-strictly-perfect}

\noindent
本節は Cohomology, Section
\ref{cohomology-section-strictly-perfect} の類似である。

\begin{definition}
\label{sites-cohomology-definition-strictly-perfect}
$(\mathcal{C}, \mathcal{O})$ を環付きサイトとする。
$\mathcal{E}^\bullet$ を $\mathcal{O}$-加群の複体とする。
有限個を除くすべての $i$ に対して $\mathcal{E}^i$ が零であり、
すべての $i$ に対して $\mathcal{E}^i$ が有限自由
$\mathcal{O}$-加群の直和因子であるとき、
$\mathcal{E}^\bullet$ は {\it 狭義完全} であるという。
\end{definition}

\noindent
$U$ を $\mathcal{C}$ の対象とする。「$\mathcal{E}^\bullet$ を
$\mathcal{O}_U$-加群の狭義完全複体とする」としばしばいうが、これは
$\mathcal{E}^\bullet$ が環付きサイト
$(\mathcal{C}/U, \mathcal{O}_U)$ 上の加群の狭義完全複体であることを意味する。
Modules on Sites, Definition
\ref{sites-modules-definition-localize-ringed-site} を参照せよ。

\begin{lemma}
\label{sites-cohomology-lemma-cone}
狭義完全複体の射の錐は狭義完全である。
\end{lemma}

\begin{proof}
定義から直ちに従う。
\end{proof}

\begin{lemma}
\label{sites-cohomology-lemma-tensor}
二つの狭義完全複体のテンソル積に付随する全複体は狭義完全である。
\end{lemma}

\begin{proof}
省略する。
\end{proof}

\begin{lemma}
\label{sites-cohomology-lemma-strictly-perfect-pullback}
$(f, f^\sharp) : (\mathcal{C}, \mathcal{O}_\mathcal{C}) \to
(\mathcal{D}, \mathcal{O}_\mathcal{D})$
を環付きトポスの射とする。$\mathcal{F}^\bullet$ が
$\mathcal{O}_\mathcal{D}$-加群の狭義完全複体ならば、
$f^*\mathcal{F}^\bullet$ は $\mathcal{O}_\mathcal{C}$-加群の
狭義完全複体である。
\end{lemma}

\begin{proof}
Modules on Sites, Lemma \ref{sites-modules-lemma-global-pullback} で、
有限自由加群の引き戻しは有限自由であることを見た。関手 $f^*$ は
加法的なので直和因子を保つ。したがって補題が従う。
\end{proof}

\begin{lemma}
\label{sites-cohomology-lemma-local-lift-map}
$(\mathcal{C}, \mathcal{O})$ を環付きサイトとし、$U$ を
$\mathcal{C}$ の対象とする。$\mathcal{O}_U$-加群の実線図式
$$
\xymatrix{
\mathcal{E} \ar@{..>}[dr] \ar[r] & \mathcal{F} \\
& \mathcal{G} \ar[u]_p
}
$$
が与えられ、$\mathcal{E}$ が有限自由 $\mathcal{O}_U$-加群の直和因子であり、
$p$ が全射であるとする。このとき、各 $U_i$ 上で図式を可換にする
点線の矢印が存在するような被覆 $\{U_i \to U\}$ が存在する。
\end{lemma}

\begin{proof}
ある $n$ に対して $\mathcal{E} = \mathcal{O}_U^{\oplus n}$ と仮定してよい。
この場合、点線の矢印を見つけることは、基底元の像を
$\Gamma(U, \mathcal{F})$ に持ち上げることと同値である。
層の全射の特徴づけ（Sites, Section
\ref{sites-section-sheaves-injective}）により、これは局所的に可能である。
\end{proof}

\begin{lemma}
\label{sites-cohomology-lemma-local-homotopy}
$(\mathcal{C}, \mathcal{O})$ を環付きサイトとし、$U$ を
$\mathcal{C}$ の対象とする。
\begin{enumerate}
\item $\alpha : \mathcal{E}^\bullet \to \mathcal{F}^\bullet$ を
$\mathcal{O}_U$-加群の複体の射とし、$\mathcal{E}^\bullet$ は狭義完全、
$\mathcal{F}^\bullet$ は非輪状であるとする。このとき、各
$\alpha|_{U_i}$ が零とホモトピックになるような被覆
$\{U_i \to U\}$ が存在する。
\item $\alpha : \mathcal{E}^\bullet \to \mathcal{F}^\bullet$ を
$\mathcal{O}_U$-加群の複体の射とし、$\mathcal{E}^\bullet$ は狭義完全、
$i < a$ に対して $\mathcal{E}^i = 0$、$i \geq a$ に対して
$H^i(\mathcal{F}^\bullet) = 0$ であるとする。このとき、各
$\alpha|_{U_i}$ が零とホモトピックになるような被覆
$\{U_i \to U\}$ が存在する。
\end{enumerate}
\end{lemma}

\begin{proof}
第1の主張は第2の主張から従うので、(2) のみを証明する。
複体 $\mathcal{E}^\bullet$ の長さに関する帰納法で証明する。
$\mathcal{E}$ が有限自由 $\mathcal{O}$-加群の直和因子であり、
$n \geq a$ が整数であって
$\mathcal{E}^\bullet \cong \mathcal{E}[-n]$ である場合、結論は補題
\ref{sites-cohomology-lemma-local-lift-map} と、仮定した
$H^n(\mathcal{F}^\bullet)$ の消滅により
$\mathcal{F}^{n - 1} \to \Ker(\mathcal{F}^n \to \mathcal{F}^{n + 1})$
が全射であることから従う。
$i \in [a, b]$ を除いて $\mathcal{E}^i$ が零である場合、複体の分裂完全列
$$
0 \to \mathcal{E}^b[-b] \to \mathcal{E}^\bullet \to
\sigma_{\leq b - 1}\mathcal{E}^\bullet \to 0
$$
を得、これは $K(\mathcal{O}_U)$ に完全三角を定める。したがって、
三角圏の公理により完全列
$$
\Hom_{K(\mathcal{O}_U)}(
\sigma_{\leq b - 1}\mathcal{E}^\bullet, \mathcal{F}^\bullet)
\to
\Hom_{K(\mathcal{O}_U)}(\mathcal{E}^\bullet, \mathcal{F}^\bullet)
\to
\Hom_{K(\mathcal{O}_U)}(\mathcal{E}^b[-b], \mathcal{F}^\bullet)
$$
を得る。上の議論により、合成
$\mathcal{E}^b[-b] \to \mathcal{F}^\bullet$ は $U$ のある被覆の各要素上で
零とホモトピックである。したがって、考えている写像は上の表示された
完全列の左辺の元から来ると仮定してよい。この元は帰納法の仮定により
$U$ のある被覆の各要素上で零である。
\end{proof}

\begin{lemma}
\label{sites-cohomology-lemma-lift-through-quasi-isomorphism}
$(\mathcal{C}, \mathcal{O})$ を環付きサイトとし、$U$ を
$\mathcal{C}$ の対象とする。$\mathcal{O}_U$-加群の複体の実線図式
$$
\xymatrix{
\mathcal{E}^\bullet \ar@{..>}[dr] \ar[r]_\alpha & \mathcal{F}^\bullet \\
& \mathcal{G}^\bullet \ar[u]_f
}
$$
が与えられ、$\mathcal{E}^\bullet$ は狭義完全で、$j < a$ に対して
$\mathcal{E}^j = 0$、$j > a$ に対して $H^j(f)$ は同型、
$j = a$ に対しては全射であるとする。このとき被覆
$\{U_i \to U\}$ が存在し、各 $i$ に対して、ホモトピーを除いて図式を
可換にする $U_i$ 上の点線の矢印が存在する。
\end{lemma}

\begin{proof}
写像 $f$ に関する仮定から、錐 $C(f)^\bullet$ の次数 $\geq a$ の
コホモロジー層は消滅する。したがって補題
\ref{sites-cohomology-lemma-local-homotopy} により、合成
$\mathcal{E}^\bullet \to \mathcal{F}^\bullet \to C(f)^\bullet$
が各 $U_i$ 上で零とホモトピックになるような被覆
$\{U_i \to U\}$ が存在する。
$$
\mathcal{G}^\bullet \to \mathcal{F}^\bullet \to C(f)^\bullet \to
\mathcal{G}^\bullet[1]
$$
は $K(\mathcal{O}_{U_i})$ において完全三角へ制限されるので、
$\alpha|_{U_i}$ はホモトピーを除いて写像
$\alpha_i : \mathcal{E}^\bullet|_{U_i} \to \mathcal{G}^\bullet|_{U_i}$
へ持ち上がる。これが望む結論である。
\end{proof}

\begin{lemma}
\label{sites-cohomology-lemma-local-actual}
$(\mathcal{C}, \mathcal{O})$ を環付きサイトとし、$U$ を
$\mathcal{C}$ の対象とする。$\mathcal{E}^\bullet$、$\mathcal{F}^\bullet$ を
$\mathcal{O}_U$-加群の複体とし、$\mathcal{E}^\bullet$ は狭義完全であるとする。
\begin{enumerate}
\item 任意の元
$\alpha \in \Hom_{D(\mathcal{O}_U)}(\mathcal{E}^\bullet, \mathcal{F}^\bullet)$
に対して、$\alpha|_{U_i}$ が複体の射
$\alpha_i : \mathcal{E}^\bullet|_{U_i} \to \mathcal{F}^\bullet|_{U_i}$
で与えられるような被覆 $\{U_i \to U\}$ が存在する。
\item 複体の射 $\alpha : \mathcal{E}^\bullet \to \mathcal{F}^\bullet$ が与えられ、
その群
$\Hom_{D(\mathcal{O}_U)}(\mathcal{E}^\bullet, \mathcal{F}^\bullet)$
における像が零であるとする。このとき、$\alpha|_{U_i}$ が零と
ホモトピックになるような被覆 $\{U_i \to U\}$ が存在する。
\end{enumerate}
\end{lemma}

\begin{proof}
(1) を証明する。導来圏の構成により、$\alpha = f^{-1}\beta$ となる
擬同型 $f : \mathcal{F}^\bullet \to \mathcal{G}^\bullet$ と複体の射
$\beta : \mathcal{E}^\bullet \to \mathcal{G}^\bullet$ をとることができる。
したがって結論は補題 \ref{sites-cohomology-lemma-lift-through-quasi-isomorphism} から従う。
(2) の証明は省略する。
\end{proof}

\begin{lemma}
\label{sites-cohomology-lemma-Rhom-strictly-perfect}
$(\mathcal{C}, \mathcal{O})$ を環付きサイトとする。
$\mathcal{E}^\bullet$、$\mathcal{F}^\bullet$ を $\mathcal{O}$-加群の複体とし、
$\mathcal{E}^\bullet$ は狭義完全であるとする。このとき内部 Hom
$R\SheafHom(\mathcal{E}^\bullet, \mathcal{F}^\bullet)$ は、各項が
$$
\mathcal{H}^n =
\bigoplus\nolimits_{n = p + q}
\SheafHom_\mathcal{O}(\mathcal{E}^{-q}, \mathcal{F}^p)
$$
であり、微分が節 \ref{sites-cohomology-section-internal-hom} で記述したものである
複体 $\mathcal{H}^\bullet$ によって表される。
\end{lemma}

\begin{proof}
K-入射複体への擬同型
$\mathcal{F}^\bullet \to \mathcal{I}^\bullet$ を選ぶ。
$(\mathcal{H}')^\bullet$ を、各項が
$$
(\mathcal{H}')^n =
\prod\nolimits_{n = p + q}
\SheafHom_\mathcal{O}(\mathcal{L}^{-q}, \mathcal{I}^p)
$$
である複体とする。節 \ref{sites-cohomology-section-internal-hom} の構成により、
これは $R\SheafHom(\mathcal{E}^\bullet, \mathcal{F}^\bullet)$ を表す。
写像
$$
\mathcal{H}^\bullet \longrightarrow (\mathcal{H}')^\bullet
$$
が擬同型であることを示せば十分である。$\mathcal{C}$ の対象 $U$ に対して、
直接確かめれば
$$
H^0(\mathcal{H}^\bullet(U)) =
\Hom_{K(\mathcal{O}_U)}(\mathcal{E}^\bullet|_U, \mathcal{I}^\bullet|_U)
\to
H^0((\mathcal{H}')^\bullet(U)) =
\Hom_{D(\mathcal{O}_U)}(\mathcal{E}^\bullet|_U, \mathcal{I}^\bullet|_U)
$$
を得る。補題 \ref{sites-cohomology-lemma-local-actual} により、
$U \mapsto H^0(\mathcal{H}^\bullet(U))$ の層化は
$U \mapsto H^0((\mathcal{H}')^\bullet(U))$ の層化に等しい。
他のコホモロジー層についても同様に論じられる。したがって
$\mathcal{H}^\bullet$ は $(\mathcal{H}')^\bullet$ と擬同型であり、
補題が証明された。
\end{proof}

\begin{lemma}
\label{sites-cohomology-lemma-Rhom-complex-of-direct-summands-finite-free}
$(\mathcal{C}, \mathcal{O})$ を環付きサイトとする。
$\mathcal{E}^\bullet$、$\mathcal{F}^\bullet$ を
$\mathcal{O}$-加群の複体とし、次を仮定する。
\begin{enumerate}
\item $n \ll 0$ に対して $\mathcal{F}^n = 0$。
\item $n \gg 0$ に対して $\mathcal{E}^n = 0$。
\item $\mathcal{E}^n$ は有限自由 $\mathcal{O}$-加群の直和因子と同型である。
\end{enumerate}
このとき内部 Hom
$R\SheafHom(\mathcal{E}^\bullet, \mathcal{F}^\bullet)$ は、各項が
$$
\mathcal{H}^n =
\bigoplus\nolimits_{n = p + q}
\SheafHom_\mathcal{O}(\mathcal{E}^{-q}, \mathcal{F}^p)
$$
であり、微分が節 \ref{sites-cohomology-section-internal-hom} で記述したものである
複体 $\mathcal{H}^\bullet$ によって表される。
\end{lemma}

\begin{proof}
$\mathcal{I}^\bullet$ が入射加群からなる下に有界な複体となるように、
擬同型 $\mathcal{F}^\bullet \to \mathcal{I}^\bullet$ を選ぶ。
$\mathcal{I}^\bullet$ は K-入射的であることに注意する
（Derived Categories, Lemma
\ref{derived-lemma-bounded-below-injectives-K-injective}）。
したがって節 \ref{sites-cohomology-section-internal-hom} の構成から、
$R\SheafHom(\mathcal{E}^\bullet, \mathcal{F}^\bullet)$ は、各項が
$$
(\mathcal{H}')^n =
\prod\nolimits_{n = p + q}
\SheafHom_\mathcal{O}(\mathcal{E}^{-q}, \mathcal{I}^p) =
\bigoplus\nolimits_{n = p + q}
\SheafHom_\mathcal{O}(\mathcal{E}^{-q}, \mathcal{I}^p)
$$
である複体 $(\mathcal{H}')^\bullet$ によって表される
（非零の項は有限個しかないので等号が成り立つ）。
$\mathcal{H}^\bullet$ は、各項が
$\SheafHom_\mathcal{O}(\mathcal{E}^{-q}, \mathcal{F}^p)$
である二重複体に付随する全複体であり、
$(\mathcal{H}')^\bullet$ についても同様である。
自然な写像 $(\mathcal{H}')^\bullet \to \mathcal{H}^\bullet$ は
二重複体の写像から来る。したがって、この写像が擬同型であることを示すには、
二重複体のスペクトル系列
（Homology, Lemma \ref{homology-lemma-first-quadrant-ss}）
$$
{}'E_1^{p, q} =
H^p(\SheafHom_\mathcal{O}(\mathcal{E}^{-q}, \mathcal{F}^\bullet))
$$
を用いればよい。これは $H^{p + q}(\mathcal{H}^\bullet)$ に収束し、
$(\mathcal{H}')^\bullet$ についても同様である。補題の証明を終えるには、
$\mathcal{E}$ が有限自由 $\mathcal{O}$-加群の直和因子であるとき、
$\mathcal{F}^\bullet \to \mathcal{I}^\bullet$ がコホモロジー層上で同型
$$
H^p(\SheafHom_\mathcal{O}(\mathcal{E}, \mathcal{F}^\bullet))
\longrightarrow
H^p(\SheafHom_\mathcal{O}(\mathcal{E}, \mathcal{I}^\bullet))
$$
を誘導することを示せば十分である。$\mathcal{E}$ が有限自由である場合、
これは明らかなので結論が従う。
\end{proof}









\section{擬コヒーレント加群}
\label{sites-cohomology-section-pseudo-coherent}

\noindent
本節では擬コヒーレント複体を論じる。

\begin{definition}
\label{sites-cohomology-definition-pseudo-coherent}
$(\mathcal{C}, \mathcal{O})$ を環付きサイトとし、$\mathcal{E}^\bullet$ を
$\mathcal{O}$-加群の複体とする。$m \in \mathbf{Z}$ とする。
\begin{enumerate}
\item $\mathcal{C}$ のすべての対象 $U$ に対して被覆 $\{U_i \to U\}$ と、
各 $i$ に対して複体の射
$\alpha_i : \mathcal{E}_i^\bullet \to \mathcal{E}^\bullet|_{U_i}$
で、$\mathcal{E}_i^\bullet$ は $\mathcal{O}_{U_i}$-加群の狭義完全複体であり、
$j > m$ に対して $H^j(\alpha_i)$ が同型、$H^m(\alpha_i)$ が全射となる
ものが存在するとき、$\mathcal{E}^\bullet$ は
{\it $m$-擬コヒーレント} であるという。
\item すべての $m$ に対して $m$-擬コヒーレントであるとき、
$\mathcal{E}^\bullet$ は {\it 擬コヒーレント} であるという。
\item $D(\mathcal{O})$ の対象 $E$ が {\it $m$-擬コヒーレント}
（それぞれ {\it 擬コヒーレント}）であるとは、$E$ が
$\mathcal{O}$-加群の $m$-擬コヒーレント（それぞれ擬コヒーレント）複体で
表されることをいう。
\end{enumerate}
\end{definition}

\noindent
$\mathcal{C}$ が準コンパクトな終対象 $X$ をもつならば
（例えば $X$ のすべての被覆が有限被覆で細分できるならば）、
$D(\mathcal{O})$ の $m$-擬コヒーレント対象は $D^-(\mathcal{O})$ に属する。
しかし一般にはそうとは限らない。

\begin{lemma}
\label{sites-cohomology-lemma-pseudo-coherent-independent-representative}
$(\mathcal{C}, \mathcal{O})$ を環付きサイトとし、$E$ を
$D(\mathcal{O})$ の対象とする。
\begin{enumerate}
\item $\mathcal{C}$ が終対象 $X$ をもち、被覆 $\{U_i \to X\}$、
$\mathcal{O}_{U_i}$-加群の狭義完全複体 $\mathcal{E}_i^\bullet$、および
$D(\mathcal{O}_{U_i})$ における写像
$\alpha_i : \mathcal{E}_i^\bullet \to E|_{U_i}$ で、$j > m$ に対して
$H^j(\alpha_i)$ が同型、$H^m(\alpha_i)$ が全射となるものが存在すれば、
$E$ は $m$-擬コヒーレントである。
\item $E$ が $m$-擬コヒーレントならば、$E$ を表す任意の
$\mathcal{O}$-加群の複体は $m$-擬コヒーレントである。
\item $\mathcal{C}$ のすべての対象 $U$ に対して、$E|_{U_i}$ が
$m$-擬コヒーレントとなる被覆 $\{U_i \to U\}$ が存在するならば、
$E$ は $m$-擬コヒーレントである。
\end{enumerate}
\end{lemma}

\begin{proof}
任意の $E$ を表す複体 $\mathcal{F}^\bullet$ をとり、$X$、
$\{U_i \to X\}$、および $\alpha_i : \mathcal{E}_i \to E|_{U_i}$ を
(1) のとおりとする。$\mathcal{F}^\bullet$ が複体として
$m$-擬コヒーレントであることを示す。これにより、$\mathcal{C}$ が終対象を
もつ場合の (1) と (2) が証明される。補題 \ref{sites-cohomology-lemma-local-actual} により、
被覆 $\{U_i \to X\}$ を細分すれば、写像 $\alpha_i$ を複体の射
$\alpha_i : \mathcal{E}_i^\bullet \to \mathcal{F}^\bullet|_{U_i}$.
で表せる。仮定により、$j > m$ に対して $H^j(\alpha_i)$ は同型であり、
$H^m(\alpha_i)$ は全射である。したがって $\mathcal{F}^\bullet$ は
$m$-擬コヒーレントである。

\medskip\noindent
(2) を証明する。上の議論から、$\mathcal{C}$ のすべての対象 $U$ に対して
$\mathcal{F}^\bullet|_U$ は $\mathcal{O}_U$-加群の複体として
$m$-擬コヒーレントである。定義から形式的に、$\mathcal{F}^\bullet$ は
$m$-擬コヒーレントである。

\medskip\noindent
(3) を証明する。定義と Sites, Definition
\ref{sites-definition-site} の (2) から従う。
\end{proof}

\begin{lemma}
\label{sites-cohomology-lemma-pseudo-coherent-pullback}
$(f, f^\sharp) : (\mathcal{C}, \mathcal{O}_\mathcal{C}) \to
(\mathcal{D}, \mathcal{O}_\mathcal{D})$
を環付きサイトの射とする。$E$ を $D(\mathcal{O}_\mathcal{C})$ の対象とする。
$E$ が $m$-擬コヒーレントならば、$Lf^*E$ は $m$-擬コヒーレントである。
\end{lemma}

\begin{proof}
$f$ が関手 $u : \mathcal{D} \to \mathcal{C}$ によって与えられるとする。
$U$ を $\mathcal{C}$ の対象とする。Sites, Lemma
\ref{sites-lemma-morphism-of-sites-covering} により、被覆 $\{U_i \to U\}$ と、
各 $i$ に対して、ある $\mathcal{D}$ の対象 $V_i$ への射
$U_i \to u(V_i)$ を見つけられる。補題
\ref{sites-cohomology-lemma-pseudo-coherent-independent-representative} により、
$Lf^*E|_{U_i}$ が $m$-擬コヒーレントであることを示せば十分である。
このためには $Lf^*E|_{u(V_i)}$ が $m$-擬コヒーレントであることを
示せばよい。実際 $Lf^*E|_{U_i}$ は、Modules on Sites, Lemma
\ref{sites-modules-lemma-relocalize} を通じて
$Lf^*E|_{u(V_i)}$ を $\mathcal{C}/U_i$ へ制限したものである。
Modules on Sites, Lemma
\ref{sites-modules-lemma-localize-morphism-ringed-sites} の可換図式により、
環付きサイトの射
$(\mathcal{C}/u(V_i), \mathcal{O}_{u(V_i)}) \to
(\mathcal{D}/V_i, \mathcal{O}_{V_i})$ に対して補題を証明すれば十分である。
したがって $\mathcal{D}$ は終対象 $Y$ をもち、$X = u(Y)$ は
$\mathcal{C}$ の終対象であると仮定してよい。

\medskip\noindent
各 $i$ に対して $\mathcal{O}_{V_i}$-加群の狭義完全複体
$\mathcal{F}_i^\bullet$ と、$D(\mathcal{O}_{V_i})$ における射
$\alpha_i : \mathcal{F}_i^\bullet \to E|_{V_i}$ で、$j > m$ に対して
$H^j(\alpha_i)$ が同型、$H^m(\alpha_i)$ が全射となるものが存在するような
被覆 $\{V_i \to Y\}$ をとる。上と同じ議論により、
$(\mathcal{C}/u(V_i), \mathcal{O}_{u(V_i)}) \to
(\mathcal{D}/V_i, \mathcal{O}_{V_i})$ に対して結論を証明すれば十分である。
ゆえに、$\mathcal{O}_\mathcal{D}$-加群の狭義完全複体
$\mathcal{F}^\bullet$ と $D(\mathcal{O}_\mathcal{D})$ における射
$\alpha : \mathcal{F}^\bullet \to E$ で、$j > m$ に対して
$H^j(\alpha)$ が同型、$H^m(\alpha)$ が全射となるものが存在すると
仮定してよい。この場合、完全三角
$$
\mathcal{F}^\bullet \to E \to C \to \mathcal{F}^\bullet[1]
$$
を選ぶ。$\alpha$ に関する仮定は、すべての $j \geq m$ に対して
コホモロジー層 $H^j(C)$ が零であることにほかならない。
$Lf^*$ を適用すると完全三角
$$
Lf^*\mathcal{F}^\bullet \to Lf^*E \to Lf^*C \to Lf^*\mathcal{F}^\bullet[1]
$$
を得る。左導来関手としての $Lf^*$ の構成から、$j \geq m$ に対して
$H^j(Lf^*C) = 0$ である（Derived Categories, Lemma
\ref{derived-lemma-negative-vanishing} の双対）。したがって
$j > m$ に対して $H^j(Lf^*\alpha)$ は同型であり、
$H^m(Lf^*\alpha)$ は全射である。他方、$\mathcal{F}^\bullet$ は
平坦な $\mathcal{O}_\mathcal{D}$-加群からなる上に有界な複体なので、
$Lf^*\mathcal{F}^\bullet = f^*\mathcal{F}^\bullet$ である。
補題 \ref{sites-cohomology-lemma-strictly-perfect-pullback} を適用して結論を得る。
\end{proof}

\begin{lemma}
\label{sites-cohomology-lemma-cone-pseudo-coherent}
$(\mathcal{C}, \mathcal{O})$ を環付きサイトとし、$m \in \mathbf{Z}$ とする。
$(K, L, M, f, g, h)$ を $D(\mathcal{O})$ の完全三角とする。
\begin{enumerate}
\item $K$ が $(m + 1)$-擬コヒーレントで、$L$ が
$m$-擬コヒーレントならば、$M$ は $m$-擬コヒーレントである。
\item $K$ と $M$ が $m$-擬コヒーレントならば、$L$ は
$m$-擬コヒーレントである。
\item $L$ が $(m + 1)$-擬コヒーレントで、$M$ が
$m$-擬コヒーレントならば、$K$ は $(m + 1)$-擬コヒーレントである。
\end{enumerate}
\end{lemma}

\begin{proof}
(1) を証明する。$U$ を $\mathcal{C}$ の対象とする。被覆
$\{U_i \to U\}$ と $D(\mathcal{O}_{U_i})$ における写像
$\alpha_i : \mathcal{K}_i^\bullet \to K|_{U_i}$ で、
$\mathcal{K}_i^\bullet$ は狭義完全、$j > m + 1$ に対して
$H^j(\alpha_i)$ は同型、$j = m + 1$ に対して全射となるものを選ぶ。
$\mathcal{K}_i^\bullet$ を
$\sigma_{\geq m + 1}\mathcal{K}_i^\bullet$
で置き換えられるので、$j < m + 1$ に対して
$\mathcal{K}_i^j = 0$ と仮定してよい。被覆を細分した後、
$D(\mathcal{O}_{U_i})$ における写像
$\beta_i : \mathcal{L}_i^\bullet \to L|_{U_i}$ で、
$\mathcal{L}_i^\bullet$ は狭義完全、$j > m$ に対して
$H^j(\beta)$ は同型、$j = m$ に対して全射となるものを選べる。
補題 \ref{sites-cohomology-lemma-lift-through-quasi-isomorphism} により、被覆をさらに
細分すれば、複体の写像
$\gamma_i : \mathcal{K}^\bullet \to \mathcal{L}^\bullet$
で、図式
$$
\xymatrix{
K|_{U_i} \ar[r] & L|_{U_i} \\
\mathcal{K}_i^\bullet \ar[u]^{\alpha_i} \ar[r]^{\gamma_i} &
\mathcal{L}_i^\bullet \ar[u]_{\beta_i}
}
$$
が $D(\mathcal{O}_{U_i})$ で可換になるものを見つけられる
（このためには写像 $\alpha_i$、$\beta_i$ および
$K|_{U_i} \to L|_{U_i}$ を実際の複体の写像で表す必要がある。
細部は省略する）。錐 $C(\gamma_i)^\bullet$ は狭義完全である
（補題 \ref{sites-cohomology-lemma-cone}）。図式の可換性から、完全三角の射
$$
(\mathcal{K}_i^\bullet, \mathcal{L}_i^\bullet, C(\gamma_i)^\bullet)
\longrightarrow
(K|_{U_i}, L|_{U_i}, M|_{U_i}).
$$
が存在する。誘導されるコホモロジー長完全列の写像と
Homology, Lemmas \ref{homology-lemma-four-lemma} および
\ref{homology-lemma-five-lemma} から、
$C(\gamma_i)^\bullet \to M|_{U_i}$ は次数 $> m$ のコホモロジー上で
同型を、次数 $m$ で全射を誘導する。したがって補題
\ref{sites-cohomology-lemma-pseudo-coherent-independent-representative} により、
$M$ は $m$-擬コヒーレントである。

\medskip\noindent
主張 (2) と (3) は、完全三角を回転して (1) を適用すれば従う。
\end{proof}

\begin{lemma}
\label{sites-cohomology-lemma-tensor-pseudo-coherent}
$(\mathcal{C}, \mathcal{O})$ を環付きサイトとし、$K,L$ を
$D(\mathcal{O})$ の対象とする。
\begin{enumerate}
\item $K$ が $n$-擬コヒーレントで、$i > a$ に対して
$H^i(K) = 0$ であり、$L$ が $m$-擬コヒーレントで、$j > b$ に対して
$H^j(L) = 0$ ならば、$t = \max(m + a, n + b)$ とおくと
$K \otimes_\mathcal{O}^\mathbf{L} L$ は $t$-擬コヒーレントである。
\item $K$ と $L$ が擬コヒーレントならば、
$K \otimes_\mathcal{O}^\mathbf{L} L$ は擬コヒーレントである。
\end{enumerate}
\end{lemma}

\begin{proof}
(1) を証明する。$U$ を $\mathcal{C}$ の対象とする。$U$ を被覆の各要素で
置き換え、$\mathcal{C}$ を局所化 $\mathcal{C}/U$ で置き換えることにより、
狭義完全複体 $\mathcal{K}^\bullet$、$\mathcal{L}^\bullet$ と写像
$\alpha : \mathcal{K}^\bullet \to K$、$\beta : \mathcal{L}^\bullet \to L$
が存在し、$i > n$ に対して $H^i(\alpha)$ が同型、$i = n$ に対して
全射、また $i > m$ に対して $H^i(\beta)$ が同型、$i = m$ に対して
全射であると仮定してよい。このとき写像
$$
\alpha \otimes^\mathbf{L} \beta :
\text{Tot}(\mathcal{K}^\bullet \otimes_\mathcal{O} \mathcal{L}^\bullet)
\to K \otimes_\mathcal{O}^\mathbf{L} L
$$
は $i > t$ に対して次数 $i$ のコホモロジー層上で同型を、
$i = t$ に対して全射を誘導する。これは Tor のスペクトル系列から従う
（細部は省略する）。

\medskip\noindent
(2) を証明する。$U$ を $\mathcal{C}$ の対象とする。まず $U$ を被覆の各要素で
置き換え、$\mathcal{C}$ を局所化 $\mathcal{C}/U$ で置き換えることにより、
$K$ と $L$ が上に有界な場合へ帰着できる。このとき主張は (1) から
直ちに従う。
\end{proof}

\begin{lemma}
\label{sites-cohomology-lemma-summands-pseudo-coherent}
$(\mathcal{C}, \mathcal{O})$ を環付きサイトとし、$m \in \mathbf{Z}$ とする。
$D(\mathcal{O})$ において $K \oplus L$ が $m$-擬コヒーレント
（それぞれ擬コヒーレント）ならば、$K$ と $L$ もそうである。
\end{lemma}

\begin{proof}
$K \oplus L$ が $m$-擬コヒーレントであると仮定する。
$U$ を $\mathcal{C}$ の対象とする。$U$ を被覆の各要素で置き換えれば
$K \oplus L \in D^-(\mathcal{O}_U)$、したがって
$L \in D^-(\mathcal{O}_U)$ と仮定してよい。完全三角
$$
(K \oplus L, K \oplus L, L \oplus L[1]) =
(K, K, 0) \oplus (L, L, L \oplus L[1])
$$
が存在することに注意する。Derived Categories, Lemma
\ref{derived-lemma-direct-sum-triangles} を参照せよ。
補題 \ref{sites-cohomology-lemma-cone-pseudo-coherent} により、
$L \oplus L[1]$ は $m$-擬コヒーレントである。
したがって $L[1] \oplus L[2]$ も $m$-擬コヒーレントであり、
帰納的に $L[n] \oplus L[n + 1]$ は $m$-擬コヒーレントである。
$L$ は上に有界なので、十分大きな $n$ に対して $L[n]$ は
$m$-擬コヒーレントである。したがって完全三角
$$
(L[n], L[n] \oplus L[n - 1], L[n - 1])
$$
を用いて逆向きに進めると、
$L[n - 1], L[n - 2], \ldots, L$ が $m$-擬コヒーレントであるとわかる。
\end{proof}

\begin{lemma}
\label{sites-cohomology-lemma-finite-cohomology}
$(\mathcal{C}, \mathcal{O})$ を環付きサイトとする。
$K$ を $D(\mathcal{O})$ の対象とする。$m \in \mathbf{Z}$ とする。
\begin{enumerate}
\item $K$ が $m$-擬コヒーレントで、$i > m$ に対して
$H^i(K) = 0$ ならば、$H^m(K)$ は有限型 $\mathcal{O}$-加群である。
\item $K$ が $m$-擬コヒーレントで、$i > m + 1$ に対して
$H^i(K) = 0$ ならば、$H^{m + 1}(K)$ は有限表示
$\mathcal{O}$-加群である。
\end{enumerate}
\end{lemma}

\begin{proof}
(1) を証明する。$U$ を $\mathcal{C}$ の対象とする。$H^m(K)$ が、
$U$ のある被覆の各要素上で有限個の切断により生成されることを示さなければ
ならない（Modules on Sites, Definition
\ref{sites-modules-definition-site-local} を参照せよ）。
したがって証明中、被覆 $\{U_i \to U\}$ を（有限回）選び、
$\mathcal{C}$ を $\mathcal{C}/U_i$ で、$U$ を $U_i$ で置き換えてよい。
特に定義により、狭義完全複体 $\mathcal{E}^\bullet$ と写像
$\alpha : \mathcal{E}^\bullet \to K$ で、次数 $> m$ のコホモロジー上で
同型を、次数 $m$ で全射を誘導するものが存在すると仮定してよい。
$\mathcal{E}^\bullet$ に対して結論を証明すれば十分である。
$\mathcal{E}^n \not = 0$ となる最大の整数を $n$ とする。
$n = m$ ならば、$H^m(\mathcal{E}^\bullet)$ は $\mathcal{E}^n$ の商であり、
結論は明らかである。$n > m$ ならば、$H^n(E^\bullet) = 0$ なので
$\mathcal{E}^{n - 1} \to \mathcal{E}^n$ は全射である。補題
\ref{sites-cohomology-lemma-local-lift-map} により、$U$ を被覆の各要素で置き換えれば、
この全射の切断を見つけて
$\mathcal{E}^{n - 1} = \mathcal{E}' \oplus \mathcal{E}^n$.
と書ける。したがって、次数 $n - 1$ の項が $\mathcal{E}'$、
次数 $n$ の項が $0$ であることを除いて $\mathcal{E}^\bullet$ と同じ複体
$(\mathcal{E}')^\bullet$ に対して結論を証明すれば十分である。
$n$ に関する帰納法で結論を得る。

\medskip\noindent
(2) を証明する。$U$ を $\mathcal{C}$ の対象とする。(1) の証明と同様に、
$U$ 上局所的に作業してよい。したがって、狭義完全複体
$\mathcal{E}^\bullet$ と写像 $\alpha : \mathcal{E}^\bullet \to K$ で、
次数 $> m$ のコホモロジー上で同型を、次数 $m$ で全射を誘導するものが
存在すると仮定してよい。(1) の証明と同様に、$i > m + 1$ に対して
$\mathcal{E}^i = 0$ である場合へ帰着できる。このとき
$H^{m + 1}(K) \cong H^{m + 1}(\mathcal{E}^\bullet) =
\Coker(\mathcal{E}^m \to \mathcal{E}^{m + 1})$
であり、これは有限表示である。
\end{proof}








\section{Tor 次元}
\label{sites-cohomology-section-tor}

\noindent
本節では、平坦加群による分解をより詳しく調べる。

\begin{definition}
\label{sites-cohomology-definition-tor-amplitude}
$(\mathcal{C}, \mathcal{O})$ を環付きサイトとし、$E$ を
$D(\mathcal{O})$ の対象とする。$a \leq b$ なる
$a, b \in \mathbf{Z}$ とする。
\begin{enumerate}
\item すべての $\mathcal{O}$-加群 $\mathcal{F}$ と
すべての $i \not \in [a, b]$ に対して
$H^i(E \otimes_\mathcal{O}^\mathbf{L} \mathcal{F}) = 0$
であるとき、$E$ は {\it $[a,b]$ に Tor 振幅をもつ} という。
\item ある $a,b$ に対して $[a,b]$ に Tor 振幅をもつとき、
$E$ は {\it 有限 Tor 次元をもつ} という。
\item $\mathcal{C}$ の任意の対象 $U$ に対して、すべての $i$ について
$E|_{U_i}$ が有限 Tor 次元をもつような被覆 $\{U_i \to U\}$ が存在するとき、
$E$ は {\it 局所的に有限 Tor 次元をもつ} という。
\end{enumerate}
$D(\mathcal{O})$ の対象とみなした $\mathcal{F}[0]$ が
$[-d,0]$ に Tor 振幅をもつとき、$\mathcal{O}$-加群
$\mathcal{F}$ は {\it Tor 次元 $\leq d$} をもつという。
\end{definition}

\noindent
定義における $E$ が有限 Tor 次元をもつならば、$E$ は
$D^b(\mathcal{O})$ の対象であることに注意する。これは上の定義で
$\mathcal{F} = \mathcal{O}$ とすればわかる。

\begin{lemma}
\label{sites-cohomology-lemma-last-one-flat}
$(\mathcal{C}, \mathcal{O})$ を環付きサイトとする。
$\mathcal{E}^\bullet$ を、$[a,b]$ に Tor 振幅をもつ、平坦な
$\mathcal{O}$-加群からなる上に有界な複体とする。
このとき $\Coker(d_{\mathcal{E}^\bullet}^{a - 1})$ は
平坦な $\mathcal{O}$-加群である。
\end{lemma}

\begin{proof}
複体 $\mathcal{E}^\bullet$ は平坦加群からなる上に有界な複体なので、
$\mathcal{O}$-加群 $\mathcal{F}$ に対して
$\mathcal{E}^\bullet \otimes_\mathcal{O} \mathcal{F} =
\mathcal{E}^\bullet \otimes_\mathcal{O}^{\mathbf{L}} \mathcal{F}$
である。したがって、すべての $\mathcal{O}$-加群 $\mathcal{F}$ に対して列
$$
\mathcal{E}^{a - 2} \otimes_\mathcal{O} \mathcal{F} \to
\mathcal{E}^{a - 1} \otimes_\mathcal{O} \mathcal{F} \to
\mathcal{E}^a \otimes_\mathcal{O} \mathcal{F}
$$
は中央で完全である。また
$\mathcal{E}^{a - 2} \to \mathcal{E}^{a - 1} \to \mathcal{E}^a \to
\Coker(d^{a - 1}) \to 0$
は平坦分解なので、すべての $\mathcal{O}$-加群 $\mathcal{F}$ に対して
$\text{Tor}_1^\mathcal{O}(\Coker(d^{a - 1}), \mathcal{F}) = 0$
であることが従う。ゆえに $\Coker(d^{a - 1})$ は平坦である。
補題 \ref{sites-cohomology-lemma-flat-tor-zero} を参照せよ。
\end{proof}

\begin{lemma}
\label{sites-cohomology-lemma-tor-amplitude}
$(\mathcal{C}, \mathcal{O})$ を環付きサイトとし、$E$ を
$D(\mathcal{O})$ の対象とする。$a \leq b$ なる
$a, b \in \mathbf{Z}$ とする。次は同値である。
\begin{enumerate}
\item $E$ は $[a,b]$ に Tor 振幅をもつ。
\item $E$ は平坦な $\mathcal{O}$-加群からなる複体
$\mathcal{E}^\bullet$ で表され、$i \not \in [a,b]$ に対して
$\mathcal{E}^i = 0$ である。
\end{enumerate}
\end{lemma}

\begin{proof}
(2) が成り立つならば
$E \otimes_\mathcal{O}^\mathbf{L} \mathcal{F} =
\mathcal{E}^\bullet \otimes_\mathcal{O} \mathcal{F}$
と計算でき、(1) は明らかである。

\medskip\noindent
(1) を仮定する。節 \ref{sites-cohomology-section-flat} により、$E$ を平坦な
$\mathcal{O}$-加群からなる上に有界な複体 $\mathcal{K}^\bullet$ で表せる。
$\mathcal{K}^n \not = 0$ となる最大の整数を $n$ とする。
$n > b$ ならば、$H^n(\mathcal{K}^\bullet) = 0$ なので
$\mathcal{K}^{n - 1} \to \mathcal{K}^n$ は全射である。
$\mathcal{K}^n$ は平坦なので
$\Ker(\mathcal{K}^{n - 1} \to \mathcal{K}^n)$ は平坦である
（Modules on Sites, Lemma \ref{sites-modules-lemma-flat-ses}）。
したがって $\mathcal{K}^\bullet$ を
$\tau_{\leq n - 1}\mathcal{K}^\bullet$ で置き換えてよい。ゆえに $n$ に
関する帰納法により、$K^\bullet$ は平坦な $\mathcal{O}$-加群からなる複体で、
$i > b$ に対して $\mathcal{K}^i = 0$ である場合へ帰着する。

\medskip\noindent
$\mathcal{E}^\bullet = \tau_{\geq a}\mathcal{K}^\bullet$ とおく。
$\mathcal{E}^a$ が平坦であることを除けばすべて明らかであり、この平坦性は
補題 \ref{sites-cohomology-lemma-last-one-flat} と定義から直ちに従う。
\end{proof}

\begin{lemma}
\label{sites-cohomology-lemma-bounded-below-tor-amplitude}
$(\mathcal{C}, \mathcal{O})$ を環付きサイトとし、$E$ を
$D(\mathcal{O})$ の対象とする。$a \in \mathbf{Z}$ とする。次は同値である。
\begin{enumerate}
\item $E$ は $[a, \infty]$ に Tor 振幅をもつ。
\item $E$ は平坦な $\mathcal{O}$-加群からなる K-平坦複体
$\mathcal{E}^\bullet$ で表すことができ、$i \not \in [a, \infty]$ に対して
$\mathcal{E}^i = 0$ である。
\end{enumerate}
さらに、任意の環付きサイトの射による引き戻しが各項の平坦な K-平坦複体と
なるように $\mathcal{E}^\bullet$ を選べる。
\end{lemma}

\begin{proof}
含意 (2) $\Rightarrow$ (1) は直ちに従う。(1) を仮定する。
まず、$E$ を表す各項の平坦な K-平坦複体 $\mathcal{K}^\bullet$ を選ぶ。
補題 \ref{sites-cohomology-lemma-K-flat-resolution} を参照せよ。
任意の $\mathcal{O}$-加群 $\mathcal{M}$ に対して
$$
\mathcal{K}^{n - 1} \otimes_\mathcal{O} \mathcal{M} \to
\mathcal{K}^n \otimes_\mathcal{O} \mathcal{M} \to
\mathcal{K}^{n + 1} \otimes_\mathcal{O} \mathcal{M}
$$
のコホモロジーは $H^n(E \otimes_\mathcal{O}^\mathbf{L} \mathcal{M})$ を
計算する。これは $n < a$ に対して常に零である。したがって複体
$\ldots \to \mathcal{K}^{a - 1} \to \mathcal{K}^a \to \mathcal{K}^{a + 1}$
に補題 \ref{sites-cohomology-lemma-last-one-flat} を適用すると、
$\mathcal{N} = \Coker(\mathcal{K}^{a - 1} \to \mathcal{K}^a)$ は
平坦な $\mathcal{O}$-加群であると結論できる。
$$
\mathcal{E}^\bullet = \tau_{\geq a}\mathcal{K}^\bullet =
(\ldots \to 0 \to \mathcal{N} \to \mathcal{K}^{a + 1} \to \ldots )
$$
とおく。$\mathcal{K}^\bullet \to \mathcal{E}^\bullet$ の核
$\mathcal{L}^\bullet$ は複体
$$
\mathcal{L}^\bullet = (\ldots \to \mathcal{K}^{a - 1} \to
\mathcal{I} \to 0 \to \ldots)
$$
である。ここで $\mathcal{I} \subset \mathcal{K}^a$ は
$\mathcal{K}^{a - 1} \to \mathcal{K}^a$ の像である。
短完全列
$0 \to \mathcal{I} \to \mathcal{K}^a \to \mathcal{N} \to 0$
があるので、$\mathcal{I}$ は平坦な $\mathcal{O}$-加群である。
したがって $\mathcal{L}^\bullet$ は平坦加群からなる上に有界な複体であり、
補題 \ref{sites-cohomology-lemma-bounded-flat-K-flat} により K-平坦である。
補題 \ref{sites-cohomology-lemma-K-flat-two-out-of-three-ses} により、
$\mathcal{E}^\bullet$ も K-平坦である。

\medskip\noindent
最後の主張を証明する。
$f : (\mathcal{C}', \mathcal{O}') \to (\mathcal{C}, \mathcal{O})$
を環付きサイトの射とする。補題 \ref{sites-cohomology-lemma-pullback-K-flat} により、
複体 $f^*\mathcal{K}^\bullet$ は各項が平坦な K-平坦複体である。
複体 $f^*\mathcal{L}^\bullet$ は平坦な $\mathcal{O}'$-加群からなる
上に有界な複体なので K-平坦である。$\mathcal{O}'$-加群の複体の短完全列
$$
0 \to f^*\mathcal{L}^\bullet \to f^*\mathcal{K}^\bullet \to
f^*\mathcal{E}^\bullet \to 0
$$
を得る。これは平坦加群の短完全列
$0 \to \mathcal{I} \to \mathcal{K}^a \to \mathcal{N} \to 0$
の引き戻しが短完全列だからである。補題
\ref{sites-cohomology-lemma-K-flat-two-out-of-three-ses} により、複体
$f^*\mathcal{E}^\bullet$ は K-平坦であり、証明は完了する。
\end{proof}

\begin{lemma}
\label{sites-cohomology-lemma-tor-amplitude-pullback}
$(f, f^\sharp) : (\mathcal{C}, \mathcal{O}_\mathcal{C}) \to
(\mathcal{D}, \mathcal{O}_\mathcal{D})$
を環付きサイトの射とする。$E$ を $D(\mathcal{O}_\mathcal{D})$ の対象とする。
$E$ が $[a,b]$ に Tor 振幅をもつならば、$Lf^*E$ も
$[a,b]$ に Tor 振幅をもつ。
\end{lemma}

\begin{proof}
$E$ が $[a,b]$ に Tor 振幅をもつと仮定する。補題
\ref{sites-cohomology-lemma-tor-amplitude} により、$E$ を平坦な $\mathcal{O}$-加群からなる
$\mathcal{E}^\bullet$ の複体で表し、$i \not \in [a,b]$ に対して
$\mathcal{E}^i = 0$ とできる。このとき $Lf^*E$ は
$f^*\mathcal{E}^\bullet$ で表される。Modules on Sites, Lemma
\ref{sites-modules-lemma-pullback-flat} により加群
$f^*\mathcal{E}^i$ は平坦である。したがって補題
\ref{sites-cohomology-lemma-tor-amplitude} により、$Lf^*E$ は $[a,b]$ に
Tor 振幅をもつと結論する。
\end{proof}

\begin{lemma}
\label{sites-cohomology-lemma-cone-tor-amplitude}
$(\mathcal{C}, \mathcal{O})$ を環付きサイトとする。
$(K,L,M,f,g,h)$ を $D(\mathcal{O})$ の完全三角とし、
$a,b\in\mathbf{Z}$ とする。
\begin{enumerate}
\item $K$ が $[a+1,b+1]$ に Tor 振幅をもち、$L$ が $[a,b]$ に
Tor 振幅をもつならば、$M$ は $[a,b]$ に Tor 振幅をもつ。
\item $K$ と $M$ が $[a,b]$ に Tor 振幅をもつならば、
$L$ は $[a,b]$ に Tor 振幅をもつ。
\item $L$ が $[a+1,b+1]$ に Tor 振幅をもち、$M$ が $[a,b]$ に
Tor 振幅をもつならば、$K$ は $[a+1,b+1]$ に Tor 振幅をもつ。
\end{enumerate}
\end{lemma}

\begin{proof}
省略する。ヒント：完全三角に付随するコホモロジー長完全列と、
$- \otimes_\mathcal{O}^{\mathbf{L}} \mathcal{F}$
が完全三角を保つことから直ちに従う。最も容易なのは (2) であり、
他は平行移動によりこれから従う。
\end{proof}

\begin{lemma}
\label{sites-cohomology-lemma-tensor-tor-amplitude}
$(\mathcal{C}, \mathcal{O})$ を環付きサイトとし、$K,L$ を
$D(\mathcal{O})$ の対象とする。$K$ が $[a,b]$ に、
$L$ が $[c,d]$ に Tor 振幅をもつならば、
$K \otimes_\mathcal{O}^\mathbf{L} L$ は
$[a+c,b+d]$ に Tor 振幅をもつ。
\end{lemma}

\begin{proof}
省略する。ヒント：Tor のスペクトル系列を用いる。
\end{proof}

\begin{lemma}
\label{sites-cohomology-lemma-summands-tor-amplitude}
$(\mathcal{C}, \mathcal{O})$ を環付きサイトとし、$a,b\in\mathbf{Z}$ とする。
$K,L$ を $D(\mathcal{O})$ の対象とする。$K\oplus L$ が
$[a,b]$ に Tor 振幅をもつならば、$K$ と $L$ もそうである。
\end{lemma}

\begin{proof}
Tor 関手が加法的であることから明らかである。
\end{proof}

\begin{lemma}
\label{sites-cohomology-lemma-bounded}
$(\mathcal{C}, \mathcal{O})$ を環付きサイトとする。
$\mathcal{I} \subset \mathcal{O}$ をイデアルの層とし、
$K$ を $D(\mathcal{O})$ の対象とする。
\begin{enumerate}
\item $K \otimes_\mathcal{O}^\mathbf{L} \mathcal{O}/\mathcal{I}$ が
上に有界ならば、すべての $n$ に対して
$K \otimes_\mathcal{O}^\mathbf{L} \mathcal{O}/\mathcal{I}^n$
は一様に上に有界である。
\item $K \otimes_\mathcal{O}^\mathbf{L} \mathcal{O}/\mathcal{I}$ が
$D(\mathcal{O}/\mathcal{I})$ の対象として $[a,b]$ に Tor 振幅をもつならば、
$K \otimes_\mathcal{O}^\mathbf{L} \mathcal{O}/\mathcal{I}^n$ は
すべての $n$ に対して、$D(\mathcal{O}/\mathcal{I}^n)$ の対象として
$[a,b]$ に Tor 振幅をもつ。
\end{enumerate}
\end{lemma}

\begin{proof}
(1) を証明する。
$K \otimes_\mathcal{O}^\mathbf{L} \mathcal{O}/\mathcal{I}$
が上に有界であると仮定し、例えば
$H^i(K \otimes_\mathcal{O}^\mathbf{L} \mathcal{O}/\mathcal{I}) = 0$
が $i > b$ に対して成り立つとする。完全三角
$$
K \otimes_\mathcal{O}^\mathbf{L}
\mathcal{I}^n/\mathcal{I}^{n + 1} \to
K \otimes_\mathcal{O}^\mathbf{L}
\mathcal{O}/\mathcal{I}^{n + 1} \to
K \otimes_\mathcal{O}^\mathbf{L}
\mathcal{O}/\mathcal{I}^n \to
K \otimes_\mathcal{O}^\mathbf{L}
\mathcal{I}^n/\mathcal{I}^{n + 1}[1]
$$
があり、また
$$
K \otimes_\mathcal{O}^\mathbf{L}
\mathcal{I}^n/\mathcal{I}^{n + 1} =
\left(
K \otimes_\mathcal{O}^\mathbf{L}
\mathcal{O}/\mathcal{I}\right)
\otimes_{\mathcal{O}/\mathcal{I}}^\mathbf{L}
\mathcal{I}^n/\mathcal{I}^{n + 1}
$$
であることに注意する。帰納法により、すべての $n$ に対して
$H^i(K \otimes_\mathcal{O}^\mathbf{L} \mathcal{O}/\mathcal{I}^n) = 0$
が $i > b$ に対して成り立つと結論する。

\medskip\noindent
(2) を証明する。$K \otimes_\mathcal{O}^\mathbf{L} \mathcal{O}/\mathcal{I}$ が
$D(\mathcal{O}/\mathcal{I})$ の対象として $[a,b]$ に Tor 振幅をもつと仮定する。
$\mathcal{F}$ を $\mathcal{O}/\mathcal{I}^n$-加群の層とする。このとき有限フィルトレーション
$$
0 \subset \mathcal{I}^{n - 1}\mathcal{F} \subset \ldots
\subset \mathcal{I}\mathcal{F} \subset \mathcal{F}
$$
をもち、その逐次商は $\mathcal{O}/\mathcal{I}$-加群の層である。
したがって $K \otimes_\mathcal{O}^\mathbf{L} \mathcal{O}/\mathcal{I}^n$ が
$[a,b]$ に Tor 振幅をもつことを証明するには、すべての
$\mathcal{O}/\mathcal{I}$-加群 $\mathcal{G}$ に対して
$H^i(K \otimes_\mathcal{O}^\mathbf{L} \mathcal{O}/\mathcal{I}^n
\otimes_{\mathcal{O}/\mathcal{I}^n}^\mathbf{L} \mathcal{G})$
が $i \not \in [a,b]$ に対して零であることを示せば十分である。
$$
\left(K \otimes_\mathcal{O}^\mathbf{L} \mathcal{O}/\mathcal{I}^n\right)
\otimes_{\mathcal{O}/\mathcal{I}^n}^\mathbf{L} \mathcal{G}
=
\left(K \otimes_\mathcal{O}^\mathbf{L} \mathcal{O}/\mathcal{I}\right)
\otimes_{\mathcal{O}/\mathcal{I}}^\mathbf{L} \mathcal{G}
$$
がすべての $\mathcal{O}/\mathcal{I}$-加群の層 $\mathcal{G}$ に対して
成り立つので、結果が従う。
\end{proof}

\begin{lemma}
\label{sites-cohomology-lemma-tor-amplitude-stalk}
$(\mathcal{C}, \mathcal{O})$ を環付きサイトとし、$E$ を
$D(\mathcal{O})$ の対象とする。$a,b \in \mathbf{Z}$ とする。
\begin{enumerate}
\item $E$ が $[a,b]$ に Tor 振幅をもつならば、サイト $\mathcal{C}$ の
すべての点 $p$ に対して、$D(\mathcal{O}_p)$ の対象 $E_p$ は
$[a,b]$ に Tor 振幅をもつ。
\item $\mathcal{C}$ が十分な点をもつならば、逆も成り立つ。
\end{enumerate}
\end{lemma}

\begin{proof}
(1) を証明する。$p$ における茎をとることは環付きサイトの射
$(p, \mathcal{O}_p) \to (\mathcal{C}, \mathcal{O})$ による
引き戻しと同じなので、補題 \ref{sites-cohomology-lemma-tor-amplitude-pullback} を
適用できる。

\medskip\noindent
(2) を証明する。$\mathcal{C}$ が十分な点をもつならば、
$H^i(E \otimes_\mathcal{O}^\mathbf{L} \mathcal{F})$
の消滅は茎で調べられる。Modules on Sites, Lemma
\ref{sites-modules-lemma-check-exactness-stalks} を参照せよ。
$H^i(E \otimes_\mathcal{O}^\mathbf{L} \mathcal{F})_p =
H^i(E_p \otimes_{\mathcal{O}_p}^\mathbf{L} \mathcal{F}_p)$ なので結論を得る。
\end{proof}










\section{完全複体}
\label{sites-cohomology-section-perfect}

\noindent
本節では環付きサイト上の完全複体の性質を論じる。

\begin{definition}
\label{sites-cohomology-definition-perfect}
$(\mathcal{C}, \mathcal{O})$ を環付きサイトとする。
$\mathcal{E}^\bullet$ を $\mathcal{O}$-加群の複体とする。
$\mathcal{C}$ の任意の対象 $U$ に対して被覆 $\{U_i \to U\}$ が存在し、
各 $i$ に対して複体の射
$\mathcal{E}_i^\bullet \to \mathcal{E}^\bullet|_{U_i}$
であって擬同型であり、かつ $\mathcal{E}_i^\bullet$ が狭義完全であるものが
存在するとき、$\mathcal{E}^\bullet$ は {\it 完全} であるという。
$D(\mathcal{O})$ の対象 $E$ が {\it 完全} であるとは、
$E$ が $\mathcal{O}$-加群の完全複体によって表されることをいう。
\end{definition}

\noindent
$\Sh(\mathcal{C})$ が準コンパクトならば
（Sites, Section \ref{sites-section-quasi-compact}）、
$D(\mathcal{O})$ の完全対象は $D^b(\mathcal{O})$ に属する。
しかし、それ以外の場合には必ずしもそうではない。

\begin{lemma}
\label{sites-cohomology-lemma-perfect-independent-representative}
$(\mathcal{C}, \mathcal{O})$ を環付きサイトとする。
$E$ を $D(\mathcal{O})$ の対象とする。
\begin{enumerate}
\item $\mathcal{C}$ が終対象 $X$ をもち、被覆 $\{U_i \to X\}$、
$\mathcal{O}_{U_i}$-加群の狭義完全複体 $\mathcal{E}_i^\bullet$、および
$D(\mathcal{O}_{U_i})$ における同型
 $\alpha_i : \mathcal{E}_i^\bullet \to E|_{U_i}$ が存在するならば、
$E$ は完全である。
\item $E$ が完全ならば、$E$ を表す任意の複体は完全である。
\end{enumerate}
\end{lemma}

\begin{proof}
Lemma \ref{sites-cohomology-lemma-pseudo-coherent-independent-representative} の証明と同じである。
\end{proof}

\begin{lemma}
\label{sites-cohomology-lemma-perfect-precise}
$(\mathcal{C}, \mathcal{O})$ を環付きサイトとする。
$E$ を $D(\mathcal{O})$ の対象とする。
$a \leq b$ を整数とする。$E$ が $[a, b]$ に Tor 振幅をもち、
かつ $(a - 1)$-擬コヒーレントならば、$E$ は完全である。
\end{lemma}

\begin{proof}
$U$ を $\mathcal{C}$ の対象とする。$U$ をある被覆の各要素で、
$\mathcal{C}$ を局所化 $\mathcal{C}/U$ で置き換えることにより、
狭義完全複体 $\mathcal{E}^\bullet$ と写像
$\alpha : \mathcal{E}^\bullet \to E$ で、$i \geq a$ に対して
$H^i(\alpha)$ が同型となるものが存在すると仮定してよい。
$\mathcal{E}^\bullet$ を $\sigma_{\geq a - 1}\mathcal{E}^\bullet$ で
置き換えておく。完全三角
$$
\mathcal{E}^\bullet \to E \to C \to \mathcal{E}^\bullet[1]
$$
を選ぶ。$E$ と $\mathcal{E}^\bullet$ のコホモロジー層の消滅、および
$\alpha$ に関する仮定から、
$\mathcal{K} = \Ker(\mathcal{E}^{a - 1} \to \mathcal{E}^a)$ とおけば
$C \cong \mathcal{K}[2 - a]$ を得る。
$\mathcal{F}$ を $\mathcal{O}$-加群とする。
$- \otimes_\mathcal{O}^\mathbf{L} \mathcal{F}$ を適用すると、
$E$ が $[a, b]$ に Tor 振幅をもつという仮定から、
$\mathcal{K} \otimes_\mathcal{O} \mathcal{F} \to
\mathcal{E}^{a - 1} \otimes_\mathcal{O} \mathcal{F}$ の像は
$\Ker(\mathcal{E}^{a - 1} \otimes_\mathcal{O} \mathcal{F}
\to \mathcal{E}^a \otimes_\mathcal{O} \mathcal{F})$ であることが従う。
ここで $\mathcal{E}' = \Coker(\mathcal{E}^{a - 1} \to \mathcal{E}^a)$
とおけば、$\text{Tor}_1^\mathcal{O}(\mathcal{E}', \mathcal{F}) = 0$
が従う。したがって $\mathcal{E}'$ は平坦である
（Lemma \ref{sites-cohomology-lemma-flat-tor-zero}）。ゆえに Modules on Sites, Lemma
\ref{sites-modules-lemma-flat-locally-finite-presentation} により、
$\mathcal{E}'|_{U_i}$ が有限自由加群の直和因子となるような被覆
$\{U_i \to U\}$ が存在する。従って複体
$$
\mathcal{E}'|_{U_i} \to \mathcal{E}^{a + 1}|_{U_i} \to \ldots \to
\mathcal{E}^b|_{U_i}
$$
は $E|_{U_i}$ と擬同型であり、$E$ は完全である。
\end{proof}

\begin{lemma}
\label{sites-cohomology-lemma-perfect}
$(\mathcal{C}, \mathcal{O})$ を環付きサイトとする。
$E$ を $D(\mathcal{O})$ の対象とする。次は同値である。
\begin{enumerate}
\item $E$ は完全である。
\item $E$ は擬コヒーレントであり、局所的に有限 Tor 次元をもつ。
\end{enumerate}
\end{lemma}

\begin{proof}
(1) を仮定する。$U$ を $\mathcal{C}$ の対象とする。
定義により、各 $E|_{U_i}$ が狭義完全複体によって表されるような
被覆 $\{U_i \to U\}$ が存在する。
従って Lemma \ref{sites-cohomology-lemma-pseudo-coherent-independent-representative} により、
$E$ は擬コヒーレント（すなわち、すべての $m$ に対して
$m$-擬コヒーレント）である。さらに、有限自由加群の直和因子は平坦なので、
Lemma \ref{sites-cohomology-lemma-tor-amplitude} により $E|_{U_i}$ は有限 Tor 次元をもつ。
従って (2) が成り立つ。

\medskip\noindent
(2) を仮定する。$U$ を $\mathcal{C}$ の対象とする。
$U$ をある被覆の各要素で置き換えることにより、$E|_U$ が
$[a, b]$ に Tor 振幅をもつような整数 $a \leq b$ が存在すると仮定してよい。
$E|_U$ はすべての $m$ に対して $m$-擬コヒーレントなので、
Lemma \ref{sites-cohomology-lemma-perfect-precise} により結論が従う。
\end{proof}

\begin{lemma}
\label{sites-cohomology-lemma-perfect-pullback}
$(f, f^\sharp) : (\mathcal{C}, \mathcal{O}_\mathcal{C}) \to
(\mathcal{D}, \mathcal{O}_\mathcal{D})$
を環付きサイトの射とする。
$E$ を $D(\mathcal{O}_\mathcal{D})$ の対象とする。
$E$ が $D(\mathcal{O}_\mathcal{D})$ において完全ならば、
$Lf^*E$ は $D(\mathcal{O}_\mathcal{C})$ において完全である。
\end{lemma}

\begin{proof}
これは Lemma \ref{sites-cohomology-lemma-perfect}、
\ref{sites-cohomology-lemma-tor-amplitude-pullback}、および
\ref{sites-cohomology-lemma-pseudo-coherent-pullback} から従う。
\end{proof}

\begin{lemma}
\label{sites-cohomology-lemma-two-out-of-three-perfect}
$(\mathcal{C}, \mathcal{O})$ を環付きサイトとする。
$(K, L, M, f, g, h)$ を $D(\mathcal{O})$ の完全三角とする。
$K, L, M$ のうち二つが完全ならば、残りの一つも完全である。
\end{lemma}

\begin{proof}
第一の証明：Lemmas \ref{sites-cohomology-lemma-perfect}、
\ref{sites-cohomology-lemma-cone-pseudo-coherent}、および
\ref{sites-cohomology-lemma-cone-tor-amplitude} を組み合わせればよい。
第二の証明（概略）：$K$ と $L$ が完全であるとする。
$U$ を $\mathcal{C}$ の対象とする。$U$ をある被覆の各要素で
置き換えることにより、$K|_U$ と $L|_U$ がそれぞれ狭義完全複体
$\mathcal{K}^\bullet$ と $\mathcal{L}^\bullet$ によって表されると
仮定してよい。さらに $U$ をある被覆の各要素で置き換えることにより、
写像 $K|_U \to L|_U$ が複体の射
$\alpha : \mathcal{K}^\bullet \to \mathcal{L}^\bullet$ によって
与えられると仮定してよい（Lemma \ref{sites-cohomology-lemma-local-actual} 参照）。
このとき $M|_U$ は $\alpha$ の錐と同型であり、これは
Lemma \ref{sites-cohomology-lemma-cone} により狭義完全である。
\end{proof}

\begin{lemma}
\label{sites-cohomology-lemma-tensor-perfect}
$(\mathcal{C}, \mathcal{O})$ を環付きサイトとする。
$K, L$ が $D(\mathcal{O})$ の完全対象ならば、
$K \otimes_\mathcal{O}^\mathbf{L} L$ も完全である。
\end{lemma}

\begin{proof}
Lemmas \ref{sites-cohomology-lemma-perfect}、\ref{sites-cohomology-lemma-tensor-pseudo-coherent}、および
\ref{sites-cohomology-lemma-tensor-tor-amplitude} から従う。
\end{proof}

\begin{lemma}
\label{sites-cohomology-lemma-summands-perfect}
$(\mathcal{C}, \mathcal{O})$ を環付きサイトとする。
$K \oplus L$ が $D(\mathcal{O})$ の完全対象ならば、
$K$ と $L$ も完全である。
\end{lemma}

\begin{proof}
Lemmas \ref{sites-cohomology-lemma-perfect}、\ref{sites-cohomology-lemma-summands-pseudo-coherent}、および
\ref{sites-cohomology-lemma-summands-tor-amplitude} から従う。
\end{proof}














\section{双対}
\label{sites-cohomology-section-duals}

\noindent
本節では、複体の圏および導来圏の双対可能対象を特徴づける。
特に、$D(\mathcal{O})$ の対象が双対をもつための必要十分条件が
完全であることだと分かる（Example \ref{sites-cohomology-example-dual-derived} と
Lemma \ref{sites-cohomology-lemma-left-dual-derived} から従う）。

\begin{lemma}
\label{sites-cohomology-lemma-symmetric-monoidal-cat-complexes}
$(\mathcal{C}, \mathcal{O})$ を環付きサイトとする。
$\mathcal{O}$-加群の複体の圏に、テンソル積を
$\mathcal{F}^\bullet \otimes \mathcal{G}^\bullet =
\text{Tot}(\mathcal{F}^\bullet \otimes_\mathcal{O} \mathcal{G}^\bullet)$
で定めると、対称モノイド圏になる。
\end{lemma}

\begin{proof}
省略する。ヒント：単位 $\mathbf{1}$ として、次数 $0$ の項が
$\mathcal{O}$ で、他の次数では零である複体をとる。このとき明らかな同型
$\text{Tot}(\mathbf{1} \otimes_\mathcal{O} \mathcal{G}^\bullet) =
\mathcal{G}^\bullet$ および
$\text{Tot}(\mathcal{F}^\bullet \otimes_\mathcal{O} \mathbf{1}) =
\mathcal{F}^\bullet$ がある。補題を証明するには種々の図式の可換性を
確かめる必要がある。Categories, Definitions
\ref{categories-definition-monoidal-category} および
\ref{categories-definition-symmetric-monoidal-category} を参照せよ。
いずれの確認も直接的である。
\end{proof}

\begin{example}
\label{sites-cohomology-example-dual}
$(\mathcal{C}, \mathcal{O})$ を環付きサイトとする。
$\mathcal{F}^\bullet$ を $\mathcal{O}$-加群の複体とし、任意の
$U \in \Ob(\mathcal{C})$ に対して、各
$\mathcal{F}^\bullet|_{U_i}$ が狭義完全となるような被覆
$\{U_i \to U\}$ が存在すると仮定する。複体
$$
\mathcal{G}^\bullet = \SheafHom^\bullet(\mathcal{F}^\bullet, \mathcal{O})
$$
を Section \ref{sites-cohomology-section-hom-complexes} のように定める。また
$$
\eta :
\mathcal{O}
\to
\text{Tot}(\mathcal{F}^\bullet \otimes_\mathcal{O} \mathcal{G}^\bullet)
\quad\text{and}\quad
\epsilon :
\text{Tot}(\mathcal{G}^\bullet \otimes_\mathcal{O} \mathcal{F}^\bullet)
\to
\mathcal{O}
$$
を $\eta = \sum \eta_n$ および $\epsilon = \sum \epsilon_n$ で定める。
ここで $\eta_n : \mathcal{O} \to
\mathcal{F}^n \otimes_\mathcal{O} \mathcal{G}^{-n}$
および
$\epsilon_n : \mathcal{G}^{-n} \otimes_\mathcal{O} \mathcal{F}^n
\to \mathcal{O}$ は Modules on Sites, Example
\ref{sites-modules-example-dual} の写像である。このとき Categories,
Definition \ref{categories-definition-dual} の意味で、
$\mathcal{G}^\bullet, \eta, \epsilon$ は $\mathcal{F}^\bullet$ の左双対である。
$(1 \otimes \epsilon) \circ (\eta \otimes 1) = \text{id}_{\mathcal{F}^\bullet}$
および
$(\epsilon \otimes 1) \circ (1 \otimes \eta) =
\text{id}_{\mathcal{G}^\bullet}$ の確認は省略する。
More on Algebra, Lemma \ref{more-algebra-lemma-left-dual-complex} と比較せよ。
\end{example}

\begin{lemma}
\label{sites-cohomology-lemma-left-dual-complex}
$(\mathcal{C}, \mathcal{O})$ を環付きサイトとし、$\mathcal{F}^\bullet$ を
$\mathcal{O}$-加群の複体とする。$\mathcal{F}^\bullet$ が
$\mathcal{O}$-加群の複体のモノイド圏において左双対をもつならば
（Categories, Definition \ref{categories-definition-dual}）、
$\mathcal{C}$ の任意の対象 $U$ に対して、各
$\mathcal{F}^\bullet|_{U_i}$ が狭義完全となるような被覆
$\{U_i \to U\}$ が存在し、その左双対は Example
\ref{sites-cohomology-example-dual} で構成したものである。
\end{lemma}

\begin{proof}
左双対の一意性
（Categories, Remark \ref{categories-remark-left-dual-adjoint}）により、
$\mathcal{F}^\bullet$ が主張どおりであることを示せば最後の主張も従う。
$\mathcal{G}^\bullet, \eta, \epsilon$ を左双対とする。
$\eta = \sum \eta_n$、$\epsilon = \sum \epsilon_n$ と書く。ここで
$\eta_n : \mathcal{O} \to
\mathcal{F}^n \otimes_\mathcal{O} \mathcal{G}^{-n}$
および
$\epsilon_n : \mathcal{G}^{-n} \otimes_\mathcal{O} \mathcal{F}^n
\to \mathcal{O}$ である。
$(1 \otimes \epsilon) \circ (\eta \otimes 1) = \text{id}_{\mathcal{F}^\bullet}$
および
$(\epsilon \otimes 1) \circ (1 \otimes \eta) = \text{id}_{\mathcal{G}^\bullet}$
であるから、Categories, Definition \ref{categories-definition-dual} により直ちに
$(1 \otimes \epsilon_n) \circ (\eta_n \otimes 1) = \text{id}_{\mathcal{F}^n}$
および
$(\epsilon_n \otimes 1) \circ (1 \otimes \eta_n) =
\text{id}_{\mathcal{G}^{-n}}$ を得る。すなわち、
$\mathcal{G}^{-n}$ は $\mathcal{F}^n$ の左双対であり、各
$\mathcal{F}^n$ に Modules on Sites, Lemma
\ref{sites-modules-lemma-left-dual-module} を適用できる。
$U$ を $\mathcal{C}$ の対象とする。各 $i$ に対して有限個の
$\eta_n|_{U_i}$ だけが非零となるような被覆 $\{U_i \to U\}$ が存在する。
従って $U$ を $U_i$ で置き換えることにより、有限個の
$\eta_n|_U$ だけが非零であると仮定してよい。引用した補題から、これは有限個の
$\mathcal{F}^n|_U$ だけが非零であることを意味する。
同じ補題をもう一度用いると、各 $\mathcal{F}^n|_{U_i}$ が有限自由
$\mathcal{O}$-加群の直和因子となるような被覆 $\{U_i \to U\}$ を
見つけることができる。これで証明は完了する。
\end{proof}

\begin{lemma}
\label{sites-cohomology-lemma-dual-perfect-complex}
$(\mathcal{C}, \mathcal{O})$ を環付きサイトとする。
$K$ を $D(\mathcal{O})$ の完全対象とする。このとき
$K^\vee = R\SheafHom(K, \mathcal{O})$ も完全対象であり、
$(K^\vee)^\vee \cong K$ である。$D(\mathcal{O})$ の $M$ に対して、
関手的な同型
$$
M \otimes^\mathbf{L}_\mathcal{O} K^\vee = R\SheafHom_\mathcal{O}(K, M)
$$
および
$$
H^0(\mathcal{C}, M \otimes^\mathbf{L}_\mathcal{O} K^\vee) =
\Hom_{D(\mathcal{O})}(K, M)
$$
が存在する。
\end{lemma}

\begin{proof}
内部 Hom の構成が制限と可換すること
（Lemma \ref{sites-cohomology-lemma-restriction-RHom-to-U}）を以下では断りなく用いる。
$U$ を $\mathcal{C}$ の任意の対象とする。$K^\vee$ が完全であることを
確かめるには、すべての $i$ に対して $K^\vee|_{U_i}$ が完全となるような
被覆 $\{U_i \to U\}$ が存在することを示せば十分である。標準写像
$$
K = R\SheafHom(\mathcal{O}_X, \mathcal{O}_X)
\otimes_{\mathcal{O}_X}^\mathbf{L} K \longrightarrow
R\SheafHom(R\SheafHom(K, \mathcal{O}_X), \mathcal{O}_X) =
(K^\vee)^\vee
$$
がある（Lemma \ref{sites-cohomology-lemma-internal-hom-evaluate} 参照）。すべての $i$ に対して
この写像の $\mathcal{C}/U_i$ への制限が同型となるような被覆
$\{U_i \to U\}$ の存在を示せば十分である。
Lemma \ref{sites-cohomology-lemma-dual} により、最後の主張を示すには写像
(\ref{sites-cohomology-equation-eval})
$$
M \otimes^\mathbf{L}_\mathcal{O} K^\vee
\longrightarrow
R\SheafHom(K, M)
$$
が同型であることを確かめれば十分である。これも上の意味で局所的な問題である。
従って $K$ が狭義完全複体によって表される場合に補題を証明すれば十分である。

\medskip\noindent
$K$ が狭義完全複体 $\mathcal{E}^\bullet$ によって表されると仮定する。
Lemma \ref{sites-cohomology-lemma-Rhom-strictly-perfect} により、$K^\vee$ は
次数 $-n$ の項が
$(\mathcal{E}^n)^\vee =
\SheafHom_\mathcal{O}(\mathcal{E}^n, \mathcal{O})$
である複体によって表される。$\mathcal{E}^n$ は有限自由
$\mathcal{O}$-加群の直和因子なので、$(\mathcal{E}^n)^\vee$ もそうである。
従って $K^\vee$ も狭義完全複体によって表され、完全である。
写像 $K \to (K^\vee)^\vee$ は、符号を除けば同型である評価写像
$\mathcal{E}^n \to ((\mathcal{E}^n)^\vee)^\vee$ によって与えられるので
同型である。(\ref{sites-cohomology-equation-eval}) が同型であることを示すため、
$M$ を K-平坦複体 $\mathcal{F}^\bullet$ で表す。
Lemma \ref{sites-cohomology-lemma-Rhom-strictly-perfect} により、
$R\SheafHom(K, M)$ は各項が
$$
\bigoplus\nolimits_{n = p + q}
\SheafHom_\mathcal{O}(\mathcal{E}^{-q}, \mathcal{F}^p)
$$
である複体によって表される。他方、対象
$M \otimes^\mathbf{L}_\mathcal{O} K^\vee$ は各項が
$$
\bigoplus\nolimits_{n = p + q}
\mathcal{F}^p \otimes_\mathcal{O} (\mathcal{E}^{-q})^\vee
$$
である複体によって表される。従って (\ref{sites-cohomology-equation-eval}) が同型であるという
主張は、標準写像
$$
\mathcal{F}
\otimes_\mathcal{O}
\SheafHom_\mathcal{O}(\mathcal{E}, \mathcal{O})
\longrightarrow
\SheafHom_\mathcal{O}(\mathcal{E}, \mathcal{F})
$$
が、$\mathcal{E}$ が有限自由 $\mathcal{O}$-加群の直和因子で、
$\mathcal{F}$ が任意の $\mathcal{O}$-加群であるとき同型であるという
主張へ帰着する。これは $\mathcal{E}$ が有限自由である場合の対応する主張から
直ちに従う。
\end{proof}

\begin{lemma}
\label{sites-cohomology-lemma-symmetric-monoidal-derived}
$(\mathcal{C}, \mathcal{O})$ を環付きサイトとする。導来圏
$D(\mathcal{O})$ は、テンソル積を導来テンソル積とし、通常の結合制約と
可換制約を備えた対称モノイド圏である（符号規則については
More on Algebra, Section \ref{more-algebra-section-sign-rules} を参照せよ）。
\end{lemma}

\begin{proof}
省略する。Lemma \ref{sites-cohomology-lemma-symmetric-monoidal-cat-complexes} と比較せよ。
\end{proof}

\begin{example}
\label{sites-cohomology-example-dual-derived}
$(\mathcal{C}, \mathcal{O})$ を環付きサイトとし、$K$ を
$D(\mathcal{O})$ の完全対象とする。Lemma
\ref{sites-cohomology-lemma-dual-perfect-complex} のように
$K^\vee = R\SheafHom(K, \mathcal{O})$ とおく。このとき写像
$$
K \otimes_\mathcal{O}^\mathbf{L} K^\vee \longrightarrow R\SheafHom(K, K)
$$
は同型である（同補題による）。写像
$$
\eta :
\mathcal{O}
\longrightarrow
K \otimes_\mathcal{O}^\mathbf{L} K^\vee
$$
を、上の同型のもとで $1$ を $\text{id}_K$ に対応する切断へ送る
写像とする。また
$$
\epsilon : 
K^\vee
\otimes_\mathcal{O}^\mathbf{L} K
\longrightarrow
\mathcal{O}
$$
を評価写像とする（たとえばその構成には Lemma
\ref{sites-cohomology-lemma-internal-hom-composition} を用いられる）。このとき Categories,
Definition \ref{categories-definition-dual} の意味で、
$K^\vee, \eta, \epsilon$ は $K$ の左双対である。
$(1 \otimes \epsilon) \circ (\eta \otimes 1) = \text{id}_K$
および
$(\epsilon \otimes 1) \circ (1 \otimes \eta) =
\text{id}_{K^\vee}$ の確認は省略する。
\end{example}

\begin{lemma}
\label{sites-cohomology-lemma-left-dual-derived}
$(\mathcal{C}, \mathcal{O})$ を環付きサイトとし、$M$ を
$D(\mathcal{O})$ の対象とする。モノイド圏 $D(\mathcal{O})$ において
$M$ が左双対をもつならば
（Categories, Definition \ref{categories-definition-dual}）、
$M$ は完全であり、その左双対は Example
\ref{sites-cohomology-example-dual-derived} で構成したものである。
\end{lemma}

\begin{proof}
$N, \eta, \epsilon$ を左双対とする。$\mathcal{C}$ の任意の対象 $U$ に対して、
制限 $N|_U, \eta|_U, \epsilon|_U$ は $M|_U$ の左双対であることに注意する。

\medskip\noindent
$U$ を $\mathcal{C}$ の対象とする。各 $M|_{U_i}$ が
$D(\mathcal{O}_{U_i})$ の完全対象となるような $\mathcal{C}$ の被覆
$\{U_i \to U\}_{i \in I}$ を見つければ十分である。従って
$\mathcal{C}, \mathcal{O}, M, N, \eta, \epsilon$ をそれぞれ
$\mathcal{C}/U, \mathcal{O}_U, M|_U, N|_U, \eta|_U, \epsilon|_U$
で置き換え、$\mathcal{C}$ が終対象 $X$ をもつと仮定してよい。
さらに証明中、$X$ をその被覆 $\{U_i \to X\}$ の各要素で
（有限回）置き換えてよい。

\medskip\noindent
次の議論を何度か用いる。$M$ を表す $\mathcal{O}$-加群の任意の複体
$\mathcal{M}^\bullet$ を選ぶ。$N$ を表し、各項が平坦
$\mathcal{O}$-加群である K-平坦複体 $\mathcal{N}^\bullet$ を選ぶ。
Lemma \ref{sites-cohomology-lemma-K-flat-resolution} を参照せよ。写像
$$
\eta : \mathcal{O} \to
\text{Tot}(\mathcal{M}^\bullet \otimes_\mathcal{O} \mathcal{N}^\bullet)
$$
$X$ をある被覆の各要素で置き換えることにより、整数 $N$ と、
$i = 1, \ldots, N$ に対する整数 $n_i \in \mathbf{Z}$、ならびに
$\mathcal{M}^{n_i}$ と $\mathcal{N}^{-n_i}$ の切断 $f_i$ と $g_i$ で
$$
\eta(1) = \sum\nolimits_i f_i \otimes g_i
$$
を満たすものを見つけられる。$\mathcal{K}^\bullet \subset
\mathcal{M}^\bullet$ を、$i = 1, \ldots, N$ に対して切断 $f_i$ を含む、
$\mathcal{O}$-加群の任意の部分複体とする。
$\mathcal{N}^n$ の平坦性により
$\text{Tot}(\mathcal{K}^\bullet \otimes_\mathcal{O} \mathcal{N}^\bullet)
\subset
\text{Tot}(\mathcal{M}^\bullet \otimes_\mathcal{O} \mathcal{N}^\bullet)$
なので、$\eta$ は
$$
\tilde \eta :
\mathcal{O} \to
\text{Tot}(\mathcal{K}^\bullet \otimes_\mathcal{O} \mathcal{N}^\bullet)
$$
を経由する。$\mathcal{K}^\bullet$ が表す $D(\mathcal{O})$ の対象を
$K$ とすると、可換図式
$$
\xymatrix{
M \ar[rr]_-{\eta \otimes 1} \ar[rrd]_{\tilde \eta \otimes 1} & &
M \otimes^\mathbf{L} N \otimes^\mathbf{L} M
\ar[r]_-{1 \otimes \epsilon} &
M \\
& &
K \otimes^\mathbf{L} N \otimes^\mathbf{L} M
\ar[u] \ar[r]^-{1 \otimes \epsilon} &
K \ar[u]
}
$$
を得る。上の行の合成は $M$ 上の恒等写像なので、
$D(\mathcal{O})$ において $M$ は $K$ の直和因子である。

\medskip\noindent
上の議論の第一の応用として、すべての $i = 1, \ldots, N$ に対して
$a < n_i < b$ となるように部分複体
$\mathcal{K}^\bullet = \sigma_{\geq a} \tau_{\leq b}\mathcal{M}^\bullet$
を選べる。従って $M$ は $D(\mathcal{O})$ において有界複体の直和因子であり、
$M$ が $D^b(\mathcal{O})$ に属すると仮定してよい
（上の手続きでは $X$ をある被覆の各要素で置き換えていることを思い出そう）。

\medskip\noindent
$M$ は $D^b(\mathcal{O})$ に属するので、$\mathcal{M}^\bullet$ を
平坦加群からなる上に有界な複体に選べる
（Modules, Lemma \ref{modules-lemma-module-quotient-flat} および
Derived Categories, Lemma \ref{derived-lemma-subcategory-left-resolution}）。
すると上の議論で、すべての $i = 1, \ldots, N$ に対して $a < n_i$ と
なるように $\mathcal{K}^\bullet = \sigma_{\geq a}\mathcal{M}^\bullet$ を
選べる。従って $M$ は $D(\mathcal{O})$ において、平坦加群からなる
有界複体の直和因子であると仮定してよい。特に $M$ は有限 Tor 振幅をもつ。

\medskip\noindent
$M$ が $[a, b]$ に Tor 振幅をもつとする。$M$ が
$m$-擬コヒーレントであると仮定し、$X$ をある被覆の各要素で置き換えた後、
$M$ が $(m - 1)$-擬コヒーレントであると仮定できることを示す。
$M$ はいずれにせよ $(b + 1)$-擬コヒーレントなので、Lemma
\ref{sites-cohomology-lemma-perfect-precise} により、これで証明が終わる。
$X$ をある被覆の各要素で置き換えることにより、狭義完全複体
$\mathcal{E}^\bullet$ と $D(\mathcal{O})$ における写像
$\alpha : \mathcal{E}^\bullet \to M$ で、$i > m$ に対して
$H^i(\alpha)$ が同型、$i = m$ に対して全射となるものが存在すると
仮定してよい。また $i < m$ に対して $\mathcal{E}^i = 0$ と仮定する。
完全三角
$$
\mathcal{E}^\bullet \to M \to L \to \mathcal{E}^\bullet[1]
$$
を選ぶ。$i \geq m$ に対して $H^i(L) = 0$ であることに注意する。
従って $L$ は、$i \geq m$ に対して $\mathcal{L}^i = 0$ となる複体
$\mathcal{L}^\bullet$ によって表される。写像
$L \to \mathcal{E}^\bullet[1]$ は複体の射
$\mathcal{L}^\bullet \to \mathcal{E}^\bullet[1]$ によって与えられ、
これは次数 $m - 1$ を除くすべての次数で零であり、次数 $m - 1$ では
写像 $\mathcal{L}^{m - 1} \to \mathcal{E}^m$ を与える。
Derived Categories, Lemma \ref{derived-lemma-negative-exts} を参照せよ。
このとき $M$ は複体
$$
\mathcal{M}^\bullet :
\ldots \to
\mathcal{L}^{m - 2} \to
\mathcal{L}^{m - 1} \to
\mathcal{E}^m \to
\mathcal{E}^{m + 1} \to \ldots
$$
によって表される。この複体に第二段落の議論を適用し、
$i = 1, \ldots, N$ に対する $\mathcal{M}^{n_i}$ の切断 $f_i$ を得る。
$n < m$ に対して、$\mathcal{K}^n \subset \mathcal{L}^n$ を、
$n_i = n$ である切断 $f_i$ と $n_i = n - 1$ である $d(f_i)$ によって
生成される $\mathcal{O}$-部分加群とする。$n \geq m$ に対しては
$\mathcal{K}^n = \mathcal{E}^n$ とおく。明らかに完全三角の射
$$
\xymatrix{
\mathcal{E}^\bullet \ar[r] &
\mathcal{M}^\bullet \ar[r] &
\mathcal{L}^\bullet \ar[r] &
\mathcal{E}^\bullet[1] \\
\mathcal{E}^\bullet \ar[r] \ar[u] &
\mathcal{K}^\bullet \ar[r] \ar[u] &
\sigma_{\leq m - 1}\mathcal{K}^\bullet \ar[r] \ar[u] &
\mathcal{E}^\bullet[1] \ar[u]
}
$$
を得る。すべての射は上で示したものである。
複体 $\mathcal{K}^\bullet$ に対応する $D(\mathcal{O})$ の対象を $K$ とする。
証明の第二段落の議論により、合成 $M \to K \to M$ が $M$ 上の恒等写像に
なるような $D(\mathcal{O})$ の射 $s : M \to K$ を得る。図式
$$
\xymatrix{
\mathcal{E}^\bullet \ar[r] &
\mathcal{K}^\bullet \ar@{=}[r] &
K \\
\mathcal{E}^\bullet \ar[u]^{\text{id}} \ar[r]^i &
\mathcal{M}^\bullet \ar@{=}[r] &
M \ar[u]_s
}
$$
が可換であるかどうかは分からないが、写像 $K \to M$ と合成した後には
可換である。Lemma \ref{sites-cohomology-lemma-local-actual} により、$X$ をある被覆の
各要素で置き換えることにより、$s \circ i$ が複体の射
$\sigma : \mathcal{E}^\bullet \to \mathcal{K}^\bullet$ によって
与えられると仮定してよい。同じ補題により、$\sigma$ と包含
$\mathcal{K}^\bullet \subset \mathcal{M}^\bullet$ の合成が、あるホモトピー
$\{h^i : \mathcal{E}^i \to \mathcal{M}^{i - 1}\}$ によって零と
ホモトピックであると仮定してよい。そこで $\mathcal{K}^{m - 1}$ を
$\mathcal{K}^{m - 1} + \Im(h^m)$ で置き換えると
（この置き換え後も $\mathcal{K}^{m - 1}$ は有限個の大域切断によって
生成されていることに注意する）、$\sigma$ 自身が零とホモトピックになる。
これは可換な実線図式
$$
\xymatrix{
\mathcal{E}^\bullet \ar[r] &
M \ar[r] &
\mathcal{L}^\bullet \ar[r] &
\mathcal{E}^\bullet[1] \\
\mathcal{E}^\bullet \ar[r] \ar[u] &
K \ar[r] \ar[u] &
\sigma_{\leq m - 1}\mathcal{K}^\bullet \ar[r] \ar[u] &
\mathcal{E}^\bullet[1] \ar[u] \\
\mathcal{E}^\bullet \ar[r] \ar[u] &
M \ar[r] \ar[u]^s &
\mathcal{L}^\bullet \ar[r] \ar@{..>}[u] &
\mathcal{E}^\bullet[1] \ar[u]
}
$$
を得ることを意味する。三角圏の公理により、この図式に適合する点線の矢印を得る。
次数 $m - 1$ のコホモロジー層を見ると、
$$
\xymatrix{
H^{m - 1}(M) \ar[r] &
H^{m - 1}(\mathcal{L}^\bullet) \ar[r] &
H^m(\mathcal{E}^\bullet) \\
H^{m - 1}(K) \ar[r] \ar[u] &
H^{m - 1}(\sigma_{\leq m - 1}\mathcal{K}^\bullet) \ar[r] \ar[u] &
H^m(\mathcal{E}^\bullet) \ar[u] \\
H^{m - 1}(M) \ar[r] \ar[u] &
H^{m - 1}(\mathcal{L}^\bullet) \ar[r] \ar[u] &
H^m(\mathcal{E}^\bullet) \ar[u]
}
$$
を得る。左列と右列では縦の合成がいずれも恒等写像なので、中央の縦の合成
$H^{m - 1}(\mathcal{L}^\bullet) \to
H^{m - 1}(\sigma_{\leq m - 1}\mathcal{K}^\bullet) \to
H^{m - 1}(\mathcal{L}^\bullet)$ は全射である。特に
$H^{m - 1}(\sigma_{\leq m - 1}\mathcal{K}^\bullet) \to
H^{m - 1}(\mathcal{L}^\bullet)$ は全射である。
上の完全三角の射がコホモロジー層の長完全列の間に誘導する写像を用いると、
図式追跡により、$i \geq m$ に対して $H^i(K) \to H^i(M)$ が同型、
$i = m - 1$ に対して全射であることが従う。
構成により、ある $r \geq 0$ と全射
$\mathcal{O}^{\oplus r} \to \mathcal{K}^{m - 1}$ を選べる。このとき合成
$$
(\mathcal{O}^{\oplus r} \to \mathcal{E}^m \to
\mathcal{E}^{m + 1} \to \ldots ) \longrightarrow
K \longrightarrow M
$$
は次数 $\geq m$ のコホモロジー層上で同型を、次数 $m - 1$ で全射を
誘導する。これで証明は完了する。
\end{proof}

\begin{lemma}
\label{sites-cohomology-lemma-colim-and-lim-of-duals}
\begin{slogan}
完全対象の系に対する自明な双対性。
\end{slogan}
$(\mathcal{C}, \mathcal{O})$ を環付きサイトとする。
$(K_n)_{n \in \mathbf{N}}$ を $D(\mathcal{O})$ の完全対象の系とする。
$K = \text{hocolim} K_n$ を導来余極限とする
（Derived Categories, Definition \ref{derived-definition-derived-colimit}）。
このとき $D(\mathcal{O})$ の任意の対象 $E$ に対して
$$
R\SheafHom(K, E) = R\lim E \otimes^\mathbf{L}_\mathcal{O} K_n^\vee
$$
が成り立つ。ここで $(K_n^\vee)$ は双対完全複体の逆系である。
\end{lemma}

\begin{proof}
Lemma \ref{sites-cohomology-lemma-dual-perfect-complex} により
$R\lim E \otimes^\mathbf{L}_\mathcal{O} K_n^\vee =
R\lim R\SheafHom(K_n, E)$
であり、これは完全三角
$$
R\lim R\SheafHom(K_n, E) \to
\prod R\SheafHom(K_n, E) \to
\prod R\SheafHom(K_n, E)
$$
に入る。同様に $K$ は完全三角
$\bigoplus K_n \to \bigoplus K_n \to K$ に入るので、
$\prod R\SheafHom(K_n, E) = R\SheafHom(\bigoplus K_n, E)$ を
示せば十分である。これは (\ref{sites-cohomology-equation-internal-hom}) と、導来テンソル積が
直和と可換するという事実の形式的帰結である。
\end{proof}






\section{導来圏の可逆対象}
\label{sites-cohomology-section-invertible-D-or-R}

\noindent
環付きサイトの導来圏における可逆対象を特徴づける
（局所環付きトポスの場合と一般の場合の両方を扱う）。

\begin{lemma}
\label{sites-cohomology-lemma-category-summands-finite-free}
$(\mathcal{C}, \mathcal{O})$ を環付きサイトとする。
$R = \Gamma(\mathcal{C}, \mathcal{O})$ とおく。有限自由
$\mathcal{O}$-加群の直和因子である $\mathcal{O}$-加群の圏は、
有限射影 $R$-加群の圏と同値である。
\end{lemma}

\begin{proof}
有限射影 $R$-加群であることは、有限自由 $R$-加群の直和因子で
あることと同値だと注意する。同値は関手
$\mathcal{E} \mapsto \Gamma(\mathcal{C}, \mathcal{E})$ によって与えられる。
逆関手は次の構成によって与えられる。トポスの射
$f : \Sh(\mathcal{C}) \to \Sh(\text{pt})$ で、$f_*$ が大域切断をとる
関手であり、$f^{-1}$ が集合 $S$、すなわち $\Sh(\text{pt})$ の対象を
値 $S$ の定数層へ送る関手であるものを考える。恒等写像
$R \to f_*\mathcal{O}$ を用いて、環付きトポスの射
$(f, f^\sharp) : (\Sh(\mathcal{C}), \mathcal{O}) \to (\Sh(\text{pt}), R)$
を得る。このとき逆関手は $f^*$ によって与えられる。
\end{proof}

\begin{lemma}
\label{sites-cohomology-lemma-invertible-derived}
$(\mathcal{C}, \mathcal{O})$ を環付きサイトとし、$M$ を
$D(\mathcal{O})$ の対象とする。次は同値である。
\begin{enumerate}
\item $M$ は $D(\mathcal{O})$ において可逆である
（Categories, Definition \ref{categories-definition-invertible} 参照）。
\item 局所有限な\footnote{これは、$\mathcal{C}$ の任意の対象 $U$ に対して
被覆 $\{U_i \to U\}$ が存在し、各 $i$ について層
$\mathcal{O}_n|_{U_i}$ が非零となる $n$ は有限個だけであることを意味する。}
直積分解
$$
\mathcal{O} = \prod\nolimits_{n \in \mathbf{Z}} \mathcal{O}_n
$$
が存在し、各 $n$ に対して可逆 $\mathcal{O}_n$-加群 $\mathcal{H}^n$ が存在して
（Modules on Sites, Definition
\ref{sites-modules-definition-invertible-sheaf}）、
$D(\mathcal{O})$ において $M = \bigoplus \mathcal{H}^n[-n]$ である。
\end{enumerate}
(1) と (2) が成り立つならば、$M$ は $D(\mathcal{O})$ の完全対象である。
$(\mathcal{C}, \mathcal{O})$ が局所環付きサイトならば、これらの条件は
次の条件とも同値である。
\begin{enumerate}
\item[(3)] $\mathcal{C}$ の任意の対象 $U$ に対して被覆
$\{U_i \to U\}$ が存在し、各 $i$ に対して、$M|_{U_i}$ が次数 $n_i$ に
置かれた可逆 $\mathcal{O}_{U_i}$-加群によって表されるような整数
$n_i$ が存在する。
\end{enumerate}
\end{lemma}

\begin{proof}
(2) を仮定する。対象 $R\SheafHom(M, \mathcal{O})$ と合成写像
$$
R\SheafHom(M, \mathcal{O}) \otimes_\mathcal{O}^\mathbf{L} M \to \mathcal{O}
$$
を考える。これが同型であることは局所的に証明してよい。従って
$\mathcal{O} = \prod_{a \leq n \leq b} \mathcal{O}_n$ および
$M = \bigoplus_{a \leq n \leq b} \mathcal{H}^n[-n]$ と仮定してよい。
このとき
$$
R\SheafHom(\mathcal{H}^m, \mathcal{O})
\otimes_\mathcal{O}^\mathbf{L} \mathcal{H}^n
$$
が、$n \not = m$ ならば零、$n = m$ ならば $\mathcal{O}_n$ に等しいことを
示せば十分である。$n \not = m$ の場合は、$\mathcal{O}_n$ と
$\mathcal{O}_m$ が平坦 $\mathcal{O}$-代数であり、
$\mathcal{O}_n \otimes_\mathcal{O} \mathcal{O}_m = 0$ であることから従う。
可逆 $\mathcal{O}$-加群の局所構造
（Modules on Sites, Lemma \ref{sites-modules-lemma-invertible}）を用いて
局所的に作業すれば、$n = m$ の場合の同型も直ちに従う。細部は省略する。
$D(\mathcal{O})$ は対称モノイド圏なので、$M$ は可逆である。

\medskip\noindent
(1) を仮定する。(2) の記述から、$\mathcal{O}_n$ の候補は
$\SheafHom_\mathcal{O}(H^n(M), H^n(M))$ である。
これが環の層の局所有限な族で、かつ
$\mathcal{O} = \prod \mathcal{O}_n$ ならば、直積分解から生じる
$\mathcal{O}$ の冪等元を用いて、直ちに直和分解
$M = \bigoplus H^n(M)[-n]$ を得る。このことから、(2) の証明は
次の意味で局所的に行ってよい。$U$ を $\mathcal{C}$ の対象とする。
上のように $\mathcal{O}_n$ を定めたとき、$\mathcal{C}/U_i$ へ制限した後に
上の主張と (2) の主張が成り立つような $U$ の被覆
$\{U_i \to U\}$ が存在することを示さなければならない。
従って $\mathcal{C}$ が終対象 $X$ をもつと仮定してよく、(2) の証明中、
$X$ をそのある被覆の各要素で有限回置き換えてよい。

\medskip\noindent
$D(\mathcal{O})$ の対象 $N$ と同型
$M \otimes_\mathcal{O}^\mathbf{L} N \cong \mathcal{O}$ を選ぶ。
このとき $N$ はモノイド圏 $D(\mathcal{O})$ における $M$ の左双対であり、
Lemma \ref{sites-cohomology-lemma-left-dual-derived} により $M$ は完全である。
対称性から $N$ も完全である。$X$ をある被覆の各要素で置き換えることにより、
$M$ と $N$ が狭義完全複体 $\mathcal{E}^\bullet$ と
$\mathcal{F}^\bullet$ によって表されると仮定してよい。このとき
$M \otimes_\mathcal{O}^\mathbf{L} N$ は
$\text{Tot}(\mathcal{E}^\bullet \otimes_\mathcal{O} \mathcal{F}^\bullet)$
によって表される。$X$ をさらにそのある被覆の各要素で置き換えることにより、
互いに逆な同型 $\mathcal{O} \to M \otimes_\mathcal{O}^\mathbf{L} N$ と
$M \otimes_\mathcal{O}^\mathbf{L} N \to \mathcal{O}$ が複体の射
$$
\alpha : \mathcal{O} \to
\text{Tot}(\mathcal{E}^\bullet \otimes_\mathcal{O} \mathcal{F}^\bullet)
\quad\text{and}\quad
\beta :
\text{Tot}(\mathcal{E}^\bullet \otimes_\mathcal{O} \mathcal{F}^\bullet)
\to \mathcal{O}
$$
によって与えられると仮定してよい。Lemma \ref{sites-cohomology-lemma-local-actual} を参照せよ。
このとき複体の射として $\beta \circ \alpha = 1$ であり、
$D(\mathcal{O})$ の射として $\alpha \circ \beta = 1$ である。
$X$ をそのある被覆の各要素で置き換えることにより、合成
$\alpha \circ \beta$ は、成分が
$$
\theta^n :
\text{Tot}^n(\mathcal{E}^\bullet \otimes_\mathcal{O} \mathcal{F}^\bullet)
\to
\text{Tot}^{n - 1}(
\mathcal{E}^\bullet \otimes_\mathcal{O} \mathcal{F}^\bullet)
$$
であるホモトピー $\theta$ によって $1$ とホモトピックであると
仮定してよい。これは先ほどと同じ補題による。
$R = \Gamma(\mathcal{C}, \mathcal{O})$ とおく。Lemma
\ref{sites-cohomology-lemma-category-summands-finite-free} により、次を得る。
\begin{enumerate}
\item $M^\bullet = \Gamma(X, \mathcal{E}^\bullet)$ は有限射影
$R$-加群からなる有界複体である。
\item $N^\bullet = \Gamma(X, \mathcal{F}^\bullet)$ は有限射影
$R$-加群からなる有界複体である。
\item $\alpha$ と $\beta$ は複体の射
$a : R \to \text{Tot}(M^\bullet \otimes_R N^\bullet)$ および
$b : \text{Tot}(M^\bullet \otimes_R N^\bullet) \to R$ に対応する。
\item $\theta^n$ は写像
$h^n : \text{Tot}^n(M^\bullet \otimes_R N^\bullet) \to
\text{Tot}^{n - 1}(M^\bullet \otimes_R N^\bullet)$ に対応する。
\item $b \circ a = 1$ および $b \circ a - 1 = dh + hd$。
\end{enumerate}
従って $M^\bullet$ と $N^\bullet$ は $D(R)$ の互いに逆な対象を定める。
More on Algebra, Lemma \ref{more-algebra-lemma-invertible-derived} により、
直積分解 $R = \prod_{a \leq n \leq b} R_n$ と可逆 $R_n$-加群 $H^n$ で、
$M^\bullet \cong \bigoplus_{a \leq n \leq b} H^n[-n]$ を満たすものを得る。
$D(R)$ におけるこの同型は、各 $H^n$ が射影 $R$-加群であるため、複体の射
$$
\bigoplus H^n[-n] \longrightarrow M^\bullet
$$
へ持ち上げられる。これに対応して、再び Lemma
\ref{sites-cohomology-lemma-category-summands-finite-free} を用いると、射
$$
\bigoplus H^n \otimes_R \mathcal{O}[-n] \to \mathcal{E}^\bullet
$$
を得る。これは $D(\mathcal{O})$ における同型である。ここで
$M \otimes_R \mathcal{O}$ は、Lemma
\ref{sites-cohomology-lemma-category-summands-finite-free} の証明で構成した、有限射影
$R$-加群から $\mathcal{O}$-加群への関手を表す。
$\mathcal{O}_n = R_n \otimes_R \mathcal{O}$ とおけば、(2) が従う。

\medskip\noindent
$(\mathcal{C}, \mathcal{O})$ が局所環付きサイトであるとする。
対象 $U$ と有限直積分解
$\mathcal{O}|_U = \prod_{a \leq n \leq b} \mathcal{O}_n|_U$ が与えられたとき、
各 $i$ について $\mathcal{O}_n|_{U_i}$ が非零となる $n$ が高々一つである
ような被覆 $\{U_i \to U\}$ を見つけられる。これは Modules on Sites,
Lemma \ref{sites-modules-lemma-locally-ringed} の (2) と、Modules on Sites,
Definition \ref{sites-modules-definition-locally-ringed} に与えられた
局所環付きサイトの定義から直ちに従う。ここから含意
(2) $\Rightarrow$ (3) は容易に分かる。含意 (3) $\Rightarrow$ (2) は
環付きサイトに何の仮定も置かずに成り立つ。細部は省略する。
\end{proof}










\section{射影公式}
\label{sites-cohomology-section-projection-formula}

\noindent
$f : (\Sh(\mathcal{C}), \mathcal{O}_\mathcal{C}) \to
(\Sh(\mathcal{D}), \mathcal{O}_\mathcal{D})$ を環付きトポスの射とする。
$E \in D(\mathcal{O}_\mathcal{C})$ および
$K \in D(\mathcal{O}_\mathcal{D})$ とする。これ以上の仮定なしに写像
\begin{equation}
\label{sites-cohomology-equation-projection-formula-map}
Rf_*E \otimes^\mathbf{L}_{\mathcal{O}_\mathcal{D}} K
\longrightarrow
Rf_*(E \otimes^\mathbf{L}_{\mathcal{O}_\mathcal{C}} Lf^*K)
\end{equation}
が存在する。すなわち、これは標準写像
$$
Lf^*(Rf_*E \otimes^\mathbf{L}_{\mathcal{O}_\mathcal{D}} K) =
Lf^*Rf_*E \otimes^\mathbf{L}_{\mathcal{O}_\mathcal{C}} Lf^*K
\longrightarrow
E \otimes^\mathbf{L}_{\mathcal{O}_\mathcal{C}} Lf^*K
$$
の随伴である。この標準写像は $Lf^*Rf_*E \to E$ と Lemmas
\ref{sites-cohomology-lemma-pullback-tensor-product} および \ref{sites-cohomology-lemma-adjoint} から得られる。
射影公式のかなり一般的な形は次のとおりである。

\begin{lemma}
\label{sites-cohomology-lemma-projection-formula}
$f : (\Sh(\mathcal{C}), \mathcal{O}_\mathcal{C}) \to
(\Sh(\mathcal{D}), \mathcal{O}_\mathcal{D})$ を環付きトポスの射とする。
$E \in D(\mathcal{O}_\mathcal{C})$ および
$K \in D(\mathcal{O}_\mathcal{D})$ とする。$K$ が完全ならば
$$
Rf_*E \otimes^\mathbf{L}_{\mathcal{O}_\mathcal{D}} K =
Rf_*(E \otimes^\mathbf{L}_{\mathcal{O}_\mathcal{C}} Lf^*K)
$$
が $D(\mathcal{O}_\mathcal{D})$ において成り立つ。
\end{lemma}

\begin{proof}
(\ref{sites-cohomology-equation-projection-formula-map}) が同型であることを確かめるには、
$\mathcal{D}$ 上局所的に作業してよい。すなわち、$\mathcal{D}$ の任意の対象
$V$ に対して、写像の $V_j$ への制限が同型となるような被覆
$\{V_j \to V\}$ を見つければよい。完全対象の定義により、これは $K$ が
$\mathcal{O}_\mathcal{D}$-加群の狭義完全複体によって表されると仮定して
よいことを意味する。全く一般に、$K = K_1 \oplus K_2$ に対して主張が
成り立つための必要十分条件は、$K_1$ と $K_2$ に対して主張が成り立つ
ことである。従って $K$ は有限自由 $\mathcal{O}_\mathcal{D}$-加群からなる
有限複体であると仮定してよい。この場合、愚直切断を用いる簡単な議論により、
$K$ が有限自由 $\mathcal{O}_\mathcal{D}$-加群によって表される場合へ帰着する。
主張は変数 $K$ の有限直和因子をとることに関して不変なので、
$K = \mathcal{O}_\mathcal{D}[n]$ の場合に証明すれば十分であり、
この場合は自明である。
\end{proof}

\begin{remark}
\label{sites-cohomology-remark-compatible-with-diagram}
写像 (\ref{sites-cohomology-equation-projection-formula-map}) は、次の意味で Remark
\ref{sites-cohomology-remark-base-change} の基底変換写像と両立する。すなわち、
$$
\xymatrix{
(\Sh(\mathcal{C}'), \mathcal{O}_{\mathcal{C}'})
\ar[r]_{g'} \ar[d]_{f'} &
(\Sh(\mathcal{C}), \mathcal{O}_\mathcal{C}) \ar[d]^f \\
(\Sh(\mathcal{D}'), \mathcal{O}_{\mathcal{D}'})
\ar[r]^g &
(\Sh(\mathcal{D}), \mathcal{O}_\mathcal{D})
}
$$
が環付きトポスの可換図式であるとする。
$E \in D(\mathcal{O}_\mathcal{C})$ および
$K \in D(\mathcal{O}_\mathcal{D})$ とする。このとき図式
$$
\xymatrix{
Lg^*(Rf_*E \otimes^\mathbf{L}_{\mathcal{O}_\mathcal{D}} K)
\ar[r]_p \ar[d]_t &
Lg^*Rf_*(E \otimes^\mathbf{L}_{\mathcal{O}_\mathcal{C}} Lf^*K)
\ar[d]_b \\
Lg^*Rf_*E \otimes^\mathbf{L}_{\mathcal{O}_{\mathcal{D}'}}
Lg^*K \ar[d]_b &
Rf'_*L(g')^*(E \otimes^\mathbf{L}_{\mathcal{O}_\mathcal{C}}
Lf^*K) \ar[d]_t \\
Rf'_*L(g')^*E \otimes^\mathbf{L}_{\mathcal{O}_{\mathcal{D}'}}
Lg^*K \ar[rd]_p &
Rf'_*(L(g')^*E \otimes^\mathbf{L}_{\mathcal{O}_{\mathcal{D}'}}
L(g')^*Lf^*K) \ar[d]_c \\
& Rf'_*(L(g')^*E \otimes^\mathbf{L}_{\mathcal{O}_{\mathcal{D}'}}
L(f')^*Lg^*K)
}
$$
は可換である。ここで、$t$ と記した矢印は Lemma
\ref{sites-cohomology-lemma-pullback-tensor-product} を、$b$ と記した矢印は Remark
\ref{sites-cohomology-remark-base-change} を、$p$ と記した矢印は
(\ref{sites-cohomology-equation-projection-formula-map}) を適用して得られ、$c$ は
$L(g')^* \circ Lf^* = L(f')^* \circ Lg^*$ から得られる。確認は省略する。
\end{remark}






\section{弱可縮対象}
\label{sites-cohomology-section-w-contractible}

\noindent
サイトの対象 $U$ が {\it 弱可縮} であるとは、集合の層の任意の全射
$\mathcal{F} \to \mathcal{G}$ が全射
$\mathcal{F}(U) \to \mathcal{G}(U)$ を誘導することをいう。
Sites, Definition \ref{sites-definition-w-contractible} を参照せよ。

\begin{lemma}
\label{sites-cohomology-lemma-w-contractible}
$\mathcal{C}$ をサイトとし、$U$ を $\mathcal{C}$ の弱可縮対象とする。このとき
\begin{enumerate}
\item 関手 $\mathcal{F} \mapsto \mathcal{F}(U)$ は完全関手
$\textit{Ab}(\mathcal{C}) \to \textit{Ab}$ である。
\item 任意のアーベル層 $\mathcal{F}$ とすべての $p \geq 1$ に対して
$H^p(U, \mathcal{F}) = 0$ である。
\item 任意の群の層 $\mathcal{G}$ に対して、任意の
$\mathcal{G}$-トーサーは $U$ 上の切断をもつ。
\end{enumerate}
\end{lemma}

\begin{proof}
第一の主張は定義から直ちに従う
（Homology, Section \ref{homology-section-functors} も参照）。
Derived Categories, Lemma \ref{derived-lemma-right-derived-exact-functor} により
高次導来関手は消滅する。$\mathcal{F}$ を $\mathcal{G}$-トーサーとする。
このとき $\mathcal{F} \to *$ は層の全射である。従って (3) も定義から従う。
\end{proof}

\noindent
弱可縮対象を十分にもつことの帰結をここに列挙しておくと便利である。

\begin{proposition}
\label{sites-cohomology-proposition-enough-weakly-contractibles}
$\mathcal{C}$ をサイトとする。$\mathcal{B} \subset \Ob(\mathcal{C})$ とし、
すべての $U \in \mathcal{B}$ は弱可縮であり、$\mathcal{C}$ の任意の対象は
$\mathcal{B}$ の要素による被覆をもつと仮定する。
$\mathcal{O}$ を $\mathcal{C}$ 上の環の層とする。このとき
\begin{enumerate}
\item $\mathcal{O}$-加群の複体
$\mathcal{F}_1 \to \mathcal{F}_2 \to \mathcal{F}_3$ が完全であるための
必要十分条件は、すべての $U \in \mathcal{B}$ に対して
$\mathcal{F}_1(U) \to \mathcal{F}_2(U) \to \mathcal{F}_3(U)$ が
完全であることである。
\item $D(\mathcal{O})$ の任意の対象 $K$ は、その標準切断の導来極限である：
$K = R\lim \tau_{\geq -n} K$。
\item 推移写像が全射である逆系
$\ldots \to \mathcal{F}_3 \to \mathcal{F}_2 \to \mathcal{F}_1$ が
与えられたとき、射影 $\lim \mathcal{F}_n \to \mathcal{F}_1$ は全射である。
\item $\textit{Mod}(\mathcal{O})$ において直積は完全である。
\item $D(\mathcal{O})$ における直積は、任意の代表複体の直積をとることで
計算できる。
\item $(\mathcal{F}_n)$ が $\mathcal{O}$-加群の逆系ならば、
すべての $p > 1$ に対して $R^p\lim \mathcal{F}_n = 0$ であり、
$$
R^1\lim \mathcal{F}_n  =
\Coker(\prod \mathcal{F}_n \to \prod \mathcal{F}_n)
$$
が成り立つ。ここで写像は $(x_n) \mapsto (x_n - f(x_{n + 1}))$ である。
\item $(K_n)$ が $D(\mathcal{O})$ の対象の逆系ならば、短完全列
$$
0 \to R^1\lim H^{p - 1}(K_n) \to H^p(R\lim K_n) \to
\lim H^p(K_n) \to 0
$$
が存在する。
\end{enumerate}
\end{proposition}

\begin{proof}
(1) の証明。列が完全ならば、$\mathcal{C}$ の任意の弱可縮対象で評価すると、
Lemma \ref{sites-cohomology-lemma-w-contractible} により完全列を得る。逆に、すべての
$U \in \mathcal{B}$ に対して
$\mathcal{F}_1(U) \to \mathcal{F}_2(U) \to \mathcal{F}_3(U)$ が
完全であると仮定する。$V$ を $\mathcal{C}$ の対象とし、
$s \in \mathcal{F}_2(V)$ を $\mathcal{F}_2 \to \mathcal{F}_3$ の核の
要素とする。仮定により、$U_i \in \mathcal{B}$ である被覆
$\{U_i \to V\}$ が存在する。このとき $s|_{U_i}$ は切断
$s_i \in \mathcal{F}_1(U_i)$ へ持ち上がる。従って $s$ は像の層
$\Im(\mathcal{F}_1 \to \mathcal{F}_2)$ の切断である。すなわち列
$\mathcal{F}_1 \to \mathcal{F}_2 \to \mathcal{F}_3$ は完全である。

\medskip\noindent
(2) の証明。$d = 0$ とした Lemma \ref{sites-cohomology-lemma-is-limit-dimension} により成り立つ。

\medskip\noindent
(3) の証明。$(\mathcal{F}_n)$ を (2) のような系とし、
$\mathcal{F} = \lim \mathcal{F}_n$ とおく。$U \in \mathcal{B}$ ならば、
すべての推移写像 $\mathcal{F}_{n + 1}(U) \to \mathcal{F}_n(U)$ が
全射なので、$\mathcal{F}(U) = \lim \mathcal{F}_n(U)$ は
$\mathcal{F}_1(U)$ へ全射となる。従って Sites, Definition
\ref{sites-definition-sheaves-injective-surjective} と、任意の対象が
$\mathcal{B}$ の要素による被覆をもつという仮定により、
$\mathcal{F} \to \mathcal{F}_1$ は全射である。

\medskip\noindent
(4) の証明。$\mathcal{O}$-加群の完全列の族
$\mathcal{F}_{i, 1} \to \mathcal{F}_{i, 2} \to \mathcal{F}_{i, 3}$ をとる。
列 $\prod \mathcal{F}_{i, 1} \to \prod \mathcal{F}_{i, 2} \to
\prod \mathcal{F}_{i, 3}$ が完全であることを示したい。(1) の判定法を用いる。
$U \in \mathcal{B}$ とする。このとき
$$
(\prod \mathcal{F}_{i, 1})(U) \to
(\prod \mathcal{F}_{i, 2})(U) \to
(\prod \mathcal{F}_{i, 3})(U)
$$
は
$$
\prod \mathcal{F}_{i, 1}(U) \to
\prod \mathcal{F}_{i, 2}(U) \to
\prod \mathcal{F}_{i, 3}(U)
$$
と同じである。各列
$\mathcal{F}_{i, 1}(U) \to \mathcal{F}_{i, 2}(U) \to \mathcal{F}_{i, 3}(U)$
は (1) により完全である。従って Homology, Lemma
\ref{homology-lemma-product-abelian-groups-exact} により、表示した列は
完全である。再び (1) を用いて結論を得る。

\medskip\noindent
(5) の証明。(4) と、少し一般化した Derived Categories, Lemma
\ref{derived-lemma-products} から従う。

\medskip\noindent
(6) と (7) の証明。導来極限とホモトピー極限、および両者の関係については
Section \ref{sites-cohomology-section-derived-limits} を参照せよ。Derived Categories,
Definition \ref{derived-definition-derived-limit} により完全三角
$$
R\lim K_n \to \prod K_n \to \prod K_n \to R\lim K_n[1]
$$
を得る。コホモロジー層の長完全列をとると
$$
H^{p - 1}(\prod K_n) \to H^{p - 1}(\prod K_n) \to
H^p(R\lim K_n) \to H^p(\prod K_n) \to H^p(\prod K_n)
$$
(4) により直積は完全なので、これは
$$
\prod H^{p - 1}(K_n) \to \prod H^{p - 1}(K_n) \to
H^p(R\lim K_n) \to \prod H^p(K_n) \to \prod H^p(K_n)
$$
となる。まず $(\mathcal{F}_n)$ が (6) のような系である場合の
$K_n = \mathcal{F}_n[0]$ にこれを適用する。(6) が従う。
次に (7) のような $(K_n)$ に適用すると、(7) が従う。
\end{proof}






\section{コンパクト対象}
\label{sites-cohomology-section-compact}

\noindent
本節では、環付きサイト上の加群の導来圏におけるコンパクト対象を調べる。
コンパクト対象は Derived Categories, Definition
\ref{derived-definition-compact-object} で定義されていることを思い出そう。

\begin{lemma}
\label{sites-cohomology-lemma-compact-in-terms-of-generators}
$\mathcal{A}$ を Grothendieck アーベル圏とする。
$S \subset \Ob(\mathcal{A})$ を次を満たす対象の集合とする。
\begin{enumerate}
\item $\mathcal{A}$ の任意の対象は $S$ の要素の直和の商である。
\item 任意の $E \in S$ に対して、関手 $\Hom_\mathcal{A}(E, -)$ は
直和と可換する。
\end{enumerate}
このとき $D(\mathcal{A})$ の任意のコンパクト対象は、$S$ の要素の
有限直和からなる有限複体の $D(\mathcal{A})$ における直和因子である。
\end{lemma}

\begin{proof}
$K \in D(\mathcal{A})$ がコンパクト対象であると仮定する。
$K$ を複体 $K^\bullet$ で表し、写像
$$
K^\bullet
\longrightarrow
\bigoplus\nolimits_{n \geq 0} \tau_{\geq n} K^\bullet
$$
を考える。ここでは標準切断を用いた。Homology, Section
\ref{homology-section-truncations} を参照せよ。各次数で右辺の直和は有限なので、
これは意味をもつ。仮定により、この写像は有限直和を経由する。
従って少なくとも一つの $n$ に対して $K \to \tau_{\geq n} K$ は零である。
すなわち $K$ は $D^{-}(R)$ に属する。

\medskip\noindent
$K$ は、各項が $S$ の対象の直和である上に有界な複体 $K^\bullet$ によって
表せる。Derived Categories, Lemma
\ref{derived-lemma-subcategory-left-resolution} を参照せよ。愚直切断を用いると
$$
K^\bullet = \bigcup\nolimits_{n \leq 0} \sigma_{\geq n}K^\bullet
$$
であることに注意する。Homology, Section
\ref{homology-section-truncations} を参照せよ。従って Derived Categories,
Lemmas \ref{derived-lemma-colim-hocolim} および
\ref{derived-lemma-commutes-with-countable-sums} により、
$1 : K^\bullet \to K^\bullet$ は $D(R)$ において
$\sigma_{\geq n}K^\bullet \to K^\bullet$ を経由する。ゆえに
$1 : K^\bullet \to K^\bullet$ は
$$
K^\bullet \xrightarrow{\varphi} L^\bullet \xrightarrow{\psi} K^\bullet
$$
のように因数分解される。ここで $L^\bullet$ は有界で、各項が $S$ の要素の
直和である複体であり、因数分解は $D(\mathcal{A})$ におけるものとする。
$i \not \in [a, b]$ に対して $L^i$ は零であるとする。
$i < c$ に対して $L^i$ が $S$ の要素の有限直和となるような
$c \leq b + 1$ のうち最大の整数を $c$ とする。
主張：$c < b + 1$ ならば、$c$ が増加するように $L^\bullet$ を修正できる。
帰納法により、この主張から $1_K$ の因数分解
$$
K \xrightarrow{\varphi} L \xrightarrow{\psi} K
$$
が得られる。ここで $L$ は $S$ の要素の有限直和からなる有限複体によって
表され、因数分解は $D(\mathcal{A})$ におけるものとする。
$e = \varphi \circ \psi \in \text{End}_{D(\mathcal{A})}(L)$ は
冪等元であることに注意する。Derived Categories, Lemma
\ref{derived-lemma-projectors-have-images-triangulated} により
$L = \Ker(e) \oplus \Ker(1 - e)$ である。
写像 $\varphi : K \to L$ は $D(R)$ において $\Ker(1 - e)$ との同型を
誘導し、結論を得る。

\medskip\noindent
主張の証明。$L^c = \bigoplus_{\lambda \in \Lambda} E_\lambda$ と書く。
$L^{c - 1}$ は $S$ の要素の有限直和なので、仮定 (2) により有限部分集合
$\Lambda' \subset \Lambda$ で、$L^{c - 1} \to L^c$ が
$\bigoplus_{\lambda \in \Lambda'} E_\lambda \subset L^c$ を経由するものを
見つけられる。複体の写像
$$
\pi :
L^\bullet
\longrightarrow
(\bigoplus\nolimits_{\lambda \in \Lambda \setminus \Lambda'} E_\lambda)[-i]
$$
を、次数 $i$ における $\Lambda \setminus \Lambda'$ に対応する因子への射影で
定める。$K$ に関する仮定により、必要なら $\Lambda'$ をより大きな
有限部分集合で置き換えることで、$D(\mathcal{A})$ において
$\pi \circ \varphi = 0$ と仮定してよい。
$(L')^\bullet \subset L^\bullet$ を $\pi$ の核とする。$\pi$ は全射なので、
複体の短完全列を得、これは $D(\mathcal{A})$ の完全三角を与える
（Derived Categories, Lemma
\ref{derived-lemma-derived-canonical-delta-functor} 参照）。
$\Hom_{D(\mathcal{A})}(K, -)$ はホモロジー的であり
（Derived Categories, Lemma
\ref{derived-lemma-representable-homological} 参照）、かつ
$\pi \circ \varphi = 0$ なので、$(L')^\bullet \to L^\bullet$ との合成が
$\varphi$ となる $D(\mathcal{A})$ の射
$\varphi' : K^\bullet \to (L')^\bullet$ を見つけられる。
$\psi'$ を $\psi$ と $(L')^\bullet \to L^\bullet$ の合成とおけば、
新しい因数分解を得る。$(L')^\bullet$ は次数 $c$ を除いて
$L^\bullet$ と一致し、$(L')^c = \bigoplus_{\lambda \in \Lambda'} E_\lambda$
なので、主張が証明された。
\end{proof}

\begin{lemma}
\label{sites-cohomology-lemma-compact-objects-if-enough-qc}
$(\mathcal{C}, \mathcal{O})$ を環付きサイトとする。$\mathcal{C}$ の任意の対象が
準コンパクト対象による被覆をもつと仮定する。このとき
$D(\mathcal{O})$ の任意のコンパクト対象は、各項が
$j_!\mathcal{O}_U$ の形の $\mathcal{O}$-加群の有限直和であるような
有限複体の $D(\mathcal{O})$ における直和因子である。ここで $U$ は
$\mathcal{C}$ の準コンパクト対象である。
\end{lemma}

\begin{proof}
Lemma \ref{sites-cohomology-lemma-compact-in-terms-of-generators} を適用する。ここで
$S \subset \Ob(\textit{Mod}(\mathcal{O}))$ は、準コンパクトな
$U \in \Ob(\mathcal{C})$ に対する $j_!\mathcal{O}_U$ の形の加群の集合とする。
Modules on Sites, Lemma \ref{sites-modules-lemma-module-quotient-flat} と、
任意の $U$ が準コンパクト対象によって被覆されるという仮定により
仮定 (1) が成り立つ。仮定 (2) は
$$
\Hom_\mathcal{O}(j_!\mathcal{O}_U, \mathcal{F}) = \mathcal{F}(U)
$$
であり、Sites, Lemma \ref{sites-lemma-directed-colimits-sections} により
これが直和と可換することから従う。
\end{proof}

\noindent
上の補題の状況でも、加群 $j_!\mathcal{O}_U$ が必ず
$D(\mathcal{O})$ のコンパクト対象になるとは限らない
（$U$ が $\mathcal{C}$ の準コンパクト対象であってもそうである）。
この問題を扱う二つの補題を述べる。


\begin{lemma}
\label{sites-cohomology-lemma-when-jshriek-lower-compact}
$(\mathcal{C}, \mathcal{O})$ を環付きサイトとし、$U$ を
$\mathcal{C}$ の対象とする。関手
$\mathcal{F} \mapsto H^p(U, \mathcal{F})$ が直和と可換すると仮定する。
このとき $\mathcal{O}$-加群 $j_!\mathcal{O}_U$ は次の意味で
$D^+(\mathcal{O})$ のコンパクト対象である：$D(\mathcal{O})$ における
$M = \bigoplus_{i \in I} M_i$ が下に有界ならば、
$\Hom(j_{U!}\mathcal{O}_U, M) =
\bigoplus_{i \in I} \Hom(j_{U!}\mathcal{O}_U, M_i)$ である。
\end{lemma}

\begin{proof}
Modules on Sites, Equation
(\ref{sites-modules-equation-map-lower-shriek-OU-into-module}) により、
$\Hom(j_{U!}\mathcal{O}_U, -)$ は関手
$\mathcal{F} \mapsto \mathcal{F}(U)$ と同じなので、
$H^p(U, M) = \bigoplus H^p(U, M_i)$ を証明すれば十分である。
$\mathcal{I}_i$（$i \in I$）を入射 $\mathcal{O}$-加群の族とする。
仮定により
$$
H^p(U, \bigoplus\nolimits_{i \in I} \mathcal{I}_i) =
\bigoplus\nolimits_{i \in I} H^p(U, \mathcal{I}_i) = 0
$$
がすべての $p$ に対して成り立つ。$M = \bigoplus M_i$ は下に有界なので、
$n < a$ に対して $H^n(M_i) = 0$ となる $a \in \mathbf{Z}$ が存在する。
従って $M_i$ を表し、$n < a$ に対して $\mathcal{I}_i^n = 0$ となる、
入射 $\mathcal{O}$-加群からなる複体 $\mathcal{I}_i^\bullet$ を選べる。
Derived Categories, Lemma
\ref{derived-lemma-injective-resolutions-exist} を参照せよ。
Injectives, Lemma \ref{injectives-lemma-derived-products} により、直和複体
$\bigoplus \mathcal{I}_i^\bullet$ は $M$ を表す。Leray 非輪状性
（Derived Categories, Lemma \ref{derived-lemma-leray-acyclicity}）により
$$
R\Gamma(U, M) = \Gamma(U, \bigoplus \mathcal{I}_i^\bullet) =
\bigoplus \Gamma(U, \bigoplus \mathcal{I}_i^\bullet) =
\bigoplus R\Gamma(U, M_i)
$$
を得る。これが望む結論である。
\end{proof}

\begin{lemma}
\label{sites-cohomology-lemma-when-jshriek-lower-compact-worked-out}
$(\mathcal{C}, \mathcal{O})$ を、被覆の集合 $\text{Cov}_\mathcal{C}$ をもつ
環付きサイトとする。$\mathcal{B} \subset \Ob(\mathcal{C})$ および
$\text{Cov} \subset \text{Cov}_\mathcal{C}$ を部分集合とする。次を仮定する。
\begin{enumerate}
\item 任意の $\mathcal{U} \in \text{Cov}$ に対して
$\mathcal{U} = \{U_i \to U\}_{i \in I}$ であり、$I$ は有限、
$U, U_i \in \mathcal{B}$、かつすべての
$U_{i_0} \times_U \ldots \times_U U_{i_p}$ は $\mathcal{B}$ に属する。
\item 任意の $U \in \mathcal{B}$ に対して、$\text{Cov}$ に現れる
$U$ の被覆は $U$ の被覆の余終系をなす。
\end{enumerate}
このとき $U \in \mathcal{B}$ に対して、対象 $j_{U!}\mathcal{O}_U$ は
次の意味で $D^+(\mathcal{O})$ のコンパクト対象である：
$D(\mathcal{O})$ における $M = \bigoplus_{i \in I} M_i$ が下に有界ならば、
$\Hom(j_{U!}\mathcal{O}_U, M) =
\bigoplus_{i \in I} \Hom(j_{U!}\mathcal{O}_U, M_i)$ である。
\end{lemma}

\begin{proof}
これは Lemma \ref{sites-cohomology-lemma-when-jshriek-lower-compact} と
Lemma \ref{sites-cohomology-lemma-colim-works-over-collection} から従う。
\end{proof}

\begin{lemma}
\label{sites-cohomology-lemma-when-jshriek-compact}
$(\mathcal{C}, \mathcal{O})$ を環付きサイトとし、$U$ を
$\mathcal{C}$ の対象とする。次を満たす整数 $d$ が存在すれば、
$\mathcal{O}$-加群 $j_!\mathcal{O}_U$ は $D(\mathcal{O})$ の
コンパクト対象である。
\begin{enumerate}
\item すべての $p > d$ に対して $H^p(U, \mathcal{F}) = 0$ である。
\item 関手 $\mathcal{F} \mapsto H^p(U, \mathcal{F})$ は直和と可換する。
\end{enumerate}
\end{lemma}

\begin{proof}
(1) と (2) を仮定する。$D(\mathcal{O})$ の $K$ に対して
$\Hom(j_!\mathcal{O}_U, K) = R\Gamma(U, K)$ であることを思い出そう。
従って $R\Gamma(U, -)$ が直和と可換することを示せばよい。
第一の仮定は、関手 $F = H^0(U, -)$ のコホモロジー次元が有限であることを
意味する。さらに第二の仮定から、入射加群の任意の直和は $F$-非輪状である。
$K_i$ を $D(\mathcal{O})$ の対象の族とする。$K_i$ を表し、各項が入射的である
K-入射的複体 $I_i^\bullet$ を選ぶ。Injectives, Theorem
\ref{injectives-theorem-K-injective-embedding-grothendieck} を参照せよ。
任意の非輪状対象からなる複体に $F$ を適用して $RF$ を計算でき
（Derived Categories, Lemma \ref{derived-lemma-unbounded-right-derived}）、
かつ $\bigoplus K_i$ は $\bigoplus I_i^\bullet$ によって表されるので
（Injectives, Lemma \ref{injectives-lemma-derived-products}）、
$R\Gamma(U, \bigoplus K_i)$ は $\bigoplus H^0(U, I_i^\bullet)$ によって
表される。従って $R\Gamma(U, -)$ は直和と可換する。
\end{proof}

\begin{lemma}
\label{sites-cohomology-lemma-quasi-compact-weakly-contractible-compact}
$(\mathcal{C}, \mathcal{O})$ を環付きサイトとする。$U$ を
$\mathcal{C}$ の準コンパクトかつ弱可縮な対象とする。このとき
$j_!\mathcal{O}_U$ は $D(\mathcal{O})$ のコンパクト対象である。
\end{lemma}

\begin{proof}
Lemmas \ref{sites-cohomology-lemma-when-jshriek-compact} および
\ref{sites-cohomology-lemma-w-contractible} と Modules on Sites, Lemma
\ref{sites-modules-lemma-sections-over-quasi-compact} を組み合わせればよい。
\end{proof}

\begin{lemma}
\label{sites-cohomology-lemma-perfect-is-compact}
$(\mathcal{C}, \mathcal{O})$ を環付きサイトとする。
$\mathcal{C}$ が次の性質をもつと仮定する。
\begin{enumerate}
\item $\mathcal{C}$ は準コンパクトな終対象 $X$ をもつ。
\item $\mathcal{C}$ の任意の準コンパクト対象は、有限で準コンパクト対象から
なる被覆の余終系をもつ。
\item $U$ と $U_i$ が準コンパクトである有限被覆
$\{U_i \to U\}_{i \in I}$ に対して、ファイバー積
$U_i \times_U U_j$ は準コンパクトである。
\end{enumerate}
$K$ を $D(\mathcal{O})$ の完全対象とする。このとき
\begin{enumerate}
\item[(a)] $K$ は次の意味で $D^+(\mathcal{O})$ のコンパクト対象である：
$M = \bigoplus_{i \in I} M_i$ が下に有界ならば、
$\Hom(K, M) = \bigoplus_{i \in I} \Hom(K, M_i)$ である。
\item[(b)] $(\mathcal{C}, \mathcal{O})$ のコホモロジー次元が有限、すなわち
任意の $\mathcal{O}$-加群 $\mathcal{F}$ と $i > d$ に対して
$H^i(X, \mathcal{F}) = 0$ となる $d$ が存在するならば、
$K$ は $D(\mathcal{O})$ のコンパクト対象である。
\end{enumerate}
\end{lemma}

\begin{proof}
$K^\vee$ を $K$ の双対とする。Lemma
\ref{sites-cohomology-lemma-dual-perfect-complex} を参照せよ。このとき
$$
\Hom_{D(\mathcal{O})}(K, M) =
H^0(X, K^\vee \otimes_\mathcal{O}^\mathbf{L} M)
$$
であり、これは $D(\mathcal{O})$ の $M$ に関して関手的である。
$K^\vee \otimes_\mathcal{O}^\mathbf{L} -$ は直和と可換するので、
$R\Gamma(X, -)$ が該当する直和と可換することを示せば十分である。

\medskip\noindent
(a) の証明。上のように言い換えれば、これは $U = X$ とした Lemma
\ref{sites-cohomology-lemma-when-jshriek-lower-compact-worked-out} の特別な場合である。

\medskip\noindent
(b) の証明。$R\Gamma(X, K) = R\Hom(\mathcal{O}, K)$ であり、
Lemma \ref{sites-cohomology-lemma-colim-works-over-collection} により
$H^p(X, -)$ は直和と可換するので、これは Lemma
\ref{sites-cohomology-lemma-when-jshriek-compact} の特別な場合である。
\end{proof}






\section{局所定数コホモロジー層をもつ複体}
\label{sites-cohomology-section-locally-constant}

\noindent
局所定数層は Modules on Sites, Section
\ref{sites-modules-section-locally-constant} で導入されている。
$\mathcal{C}$ をサイトとし、$\Lambda$ を環とする。
$\mathcal{C}$ 上の $\underline{\Lambda}$-加群のなすアーベル圏の導来圏を
$D(\mathcal{C}, \Lambda)$ と書く。

\begin{lemma}
\label{sites-cohomology-lemma-locally-constant}
$\mathcal{C}$ を終対象 $X$ をもつサイトとし、$\Lambda$ を Noether 環とする。
$K \in D^b(\mathcal{C}, \Lambda)$ とし、各 $H^i(K)$ は有限型
$\Lambda$-加群の局所定数層であるとする。このとき、各 $K|_{U_i}$ が
有限型 $\Lambda$-加群の局所定数層からなる複体で表されるような被覆
$\{U_i \to X\}$ が存在する。
\end{lemma}

\begin{proof}
$i \not \in [a, b]$ に対して $H^i(K) = 0$ となる $a \leq b$ をとる。
$b-a$ に関する帰納法により、$M^p$ が有限型 $\Lambda$-加群であり
$p \not \in [a,b]$ に対して $M^p=0$ となる複体
$\underline{M^\bullet}_{U_i}$ で $K|_{U_i}$ を表せるような被覆
$\{U_i \to X\}$ が存在することを示す。$b=a$ の場合は明らかである。
一般の場合、$X$ をある被覆の各要素で置き換えることで、$H^b(K)$ は定数、
例えば $H^b(K)=\underline{M}$ であると仮定してよい。Modules on Sites, Lemma
\ref{sites-modules-lemma-locally-constant-finite-type}
により、$M$ は有限 $\Lambda$-加群である。$M$ の生成元
$x_1,\ldots,x_r$ によって与えられる全射
$\Lambda^{\oplus r} \to M$ を選ぶ。

\medskip\noindent
補題 \ref{sites-cohomology-lemma-kill-cohomology-class-on-covering} のわずかな一般化
（詳細は省略する）により、$x_i \in H^0(X,H^b(K))$ が
$H^b(U_i,K)$ の元へ持ち上がるような被覆 $\{U_i \to X\}$ が存在する。
%% P03-STACKS-E882: frozen source index binding retained; see sidecar.
従って $X$ を各 $U_i$ で置き換えると、射
$\underline{\Lambda^{\oplus r}}[-b] \to K$
であって次数 $b$ のコホモロジー層上で全射を誘導するものが存在する状況に至る。
完全三角
$$
\underline{\Lambda^{\oplus r}}[-b] \to K \to L \to
\underline{\Lambda^{\oplus r}}[-b + 1]
$$
を選ぶ。このとき $L$ のコホモロジー層が非零となりうるのは区間
$[a,b-1]$ のみであり、区間 $[a,b-2]$ では $K$ のコホモロジー層と一致し、
次数 $b-1$ には短完全列
$$
0 \to H^{b - 1}(K) \to H^{b - 1}(L) \to
\underline{\Ker(\Lambda^{\oplus r} \to M)} \to 0
$$
がある。Modules on Sites, Lemma
\ref{sites-modules-lemma-kernel-finite-locally-constant}
により、$H^{b-1}(L)$ は有限型の局所定数層である。
帰納法の仮定から、$M^p$ が有限 $\Lambda$-加群であり
$p \not \in [a,b-1]$ に対して $M^p=0$ となる同型
$\underline{M^\bullet} \to L$ を $D(\mathcal{C}, \underline{\Lambda})$ に得る。
射 $L \to \Lambda^{\oplus r}[-b+1]$ は、局所的に定数である射
$\underline{M^{b-1}} \to \underline{\Lambda^{\oplus r}}$ を与える
（Modules on Sites, Lemma
\ref{sites-modules-lemma-morphism-locally-constant}）。
従って、これは射 $M^{b-1} \to \Lambda^{\oplus r}$ で与えられると仮定してよい。
上の完全三角から、合成
$M^{b-2} \to M^{b-1} \to \Lambda^{\oplus r}$ は零であり、
三角圏の公理から同型
$$
\underline{M^a \to \ldots \to M^{b - 1} \to \Lambda^{\oplus r}}
\longrightarrow K
$$
を $D(\mathcal{C},\Lambda)$ に得る。
\end{proof}

\noindent
$\mathcal{C}$ をサイトとし、$\Lambda$ を環とする。射
$\Sh(\mathcal{C}) \to \Sh(pt)$ を用いると、関手
$D(\Lambda) \to D(\mathcal{C},\Lambda)$、$K \mapsto \underline{K}$ があることが分かる。

\begin{lemma}
\label{sites-cohomology-lemma-map-out-of-locally-constant}
$\mathcal{C}$ を終対象 $X$ をもつサイトとし、$\Lambda$ を環とする。次のデータをとる。
\begin{enumerate}
\item $D(\Lambda)$ の完全対象 $K$、
\item $K$ を表す有限生成射影 $\Lambda$-加群の有限複体 $K^\bullet$、
\item $\Lambda$-加群の層からなる複体 $\mathcal{L}^\bullet$、
\item $D(\mathcal{C},\Lambda)$ における射
$\varphi : \underline{K} \to \mathcal{L}^\bullet$。
\end{enumerate}
このとき、$\varphi|_{U_i}$ を表す複体の射
$\alpha_i : \underline{K}^\bullet|_{U_i} \to \mathcal{L}^\bullet|_{U_i}$
と被覆 $\{U_i \to X\}$ が存在する。
\end{lemma}

\begin{proof}
補題 \ref{sites-cohomology-lemma-local-actual} から直ちに従う。
\end{proof}

\begin{lemma}
\label{sites-cohomology-lemma-locally-constant-map}
$\mathcal{C}$ を終対象 $X$ をもつサイトとし、$\Lambda$ を環とする。
$K,L$ を $D(\Lambda)$ の対象とし、$K$ は完全であるとする。
$\varphi : \underline{K} \to \underline{L}$ を $D(\mathcal{C},\Lambda)$
における射とする。このとき、ある被覆 $\{U_i \to X\}$ が存在し、
各 $i$ に対して $\varphi|_{U_i}$ は $D(\Lambda)$ におけるある射
$\alpha_i : K \to L$ から得られる $\underline{\alpha_i}$ に等しい。
\end{lemma}

\begin{proof}
補題 \ref{sites-cohomology-lemma-map-out-of-locally-constant} と Modules on Sites, Lemma
\ref{sites-modules-lemma-morphism-locally-constant} から従う。
\end{proof}

\begin{lemma}
\label{sites-cohomology-lemma-locally-constant-tensor-product}
$\mathcal{C}$ をサイトとし、$\Lambda$ を Noether 環とする。
$K,L \in D^-(\mathcal{C},\Lambda)$ とする。$K$ と $L$ のコホモロジー層が
有限型 $\Lambda$-加群の局所定数層ならば、
$K \otimes_\Lambda^\mathbf{L} L$
のコホモロジー層も有限型 $\Lambda$-加群の局所定数層である。
\end{lemma}

\begin{proof}
これを補題 \ref{sites-cohomology-lemma-locally-constant} の応用として証明する。
$H^i(K \otimes_\Lambda^\mathbf{L} L)$ は
$H^i(\tau_{\geq i-1}K \otimes_\Lambda^\mathbf{L} \tau_{\geq i-1}L)$
と同じであることに注意する。従って、$K$ と $L$ は有界であると仮定してよい。
補題 \ref{sites-cohomology-lemma-locally-constant} により、$K$ と $L$ は有限型
$\Lambda$-加群の局所定数層からなる複体で表されると仮定してよい。
次いで、これらの複体を有限階数自由 $\Lambda$-加群の上に有界な複体で
置き換えることができる。この場合、結果は明らかである。
\end{proof}

\begin{lemma}
\label{sites-cohomology-lemma-locally-constant-bounded}
$\mathcal{C}$ をサイトとし、$\Lambda$ を Noether 環とする。
$I \subset \Lambda$ をイデアルとし、$K \in D^-(\mathcal{C},\Lambda)$ とする。
$K \otimes_\Lambda^\mathbf{L} \underline{\Lambda/I}$ のコホモロジー層が
有限型 $\Lambda/I$-加群の局所定数層ならば、すべての $n \geq 1$ に対して
$K \otimes_\Lambda^\mathbf{L} \underline{\Lambda/I^n}$
のコホモロジー層は有限型 $\Lambda/I^n$-加群の局所定数層である。
\end{lemma}

\begin{proof}
有限型 $\Lambda$-加群の局所定数層は、$\underline{\Lambda}$-加群全体の
弱 Serre 部分圏をなすことを思い出す。Modules on Sites, Lemma
\ref{sites-modules-lemma-kernel-finite-locally-constant} を参照せよ。
従って、コホモロジー層が有限型 $\Lambda$-加群の局所定数層である複体からなる
$D(\mathcal{C},\Lambda)$ の部分圏は、$D(\mathcal{C},\Lambda)$ の厳密充満飽和三角部分圏をなす。
Derived Categories, Lemma \ref{derived-lemma-cohomology-in-serre-subcategory} を参照せよ。
次に、完全三角
$$
K \otimes_\Lambda^\mathbf{L} \underline{I^n/I^{n + 1}} \to
K \otimes_\Lambda^\mathbf{L} \underline{\Lambda/I^{n + 1}} \to
K \otimes_\Lambda^\mathbf{L} \underline{\Lambda/I^n} \to
K \otimes_\Lambda^\mathbf{L} \underline{I^n/I^{n + 1}}[1]
$$
および同型
$$
K \otimes_\Lambda^\mathbf{L} \underline{I^n/I^{n + 1}}
=
\left(K \otimes_\Lambda^\mathbf{L} \underline{\Lambda/I}\right)
\otimes_{\Lambda/I}^\mathbf{L} \underline{I^n/I^{n + 1}}
$$
を考える。これと補題 \ref{sites-cohomology-lemma-locally-constant-tensor-product} を合わせれば
結果を得る。
\end{proof}
